% 来自2026-10-21统一复核稿；题号、图号和解答保留。
\clearpage
\begingroup
\fancyhead[R]{\small 模块二：静电场与边值问题}
\chapter{模块二：静电场与边值问题}\label{m2:ux6a21ux5757ux4e8cux9759ux7535ux573aux4e0eux8fb9ux503cux95eeux9898}

模块二先讨论由电荷分布求场和势，再转入由边界条件确定电势的边值问题。基本方法包括高斯定律、多极展开、唯一性定理、分离变量法和电像法；选讲部分补充特殊函数与保角变换。

{\def\LTcaptype{none} % do not increment counter
\begin{longtable}[]{@{}
  >{\raggedright\arraybackslash}p{(\linewidth - 4\tabcolsep) * \real{0.3333}}
  >{\raggedright\arraybackslash}p{(\linewidth - 4\tabcolsep) * \real{0.3333}}
  >{\raggedright\arraybackslash}p{(\linewidth - 4\tabcolsep) * \real{0.3333}}@{}}
\toprule\noalign{}
\begin{minipage}[b]{\linewidth}\raggedright
层级
\end{minipage} & \begin{minipage}[b]{\linewidth}\raggedright
内容
\end{minipage} & \begin{minipage}[b]{\linewidth}\raggedright
阅读定位
\end{minipage} \\
\midrule\noalign{}
\endhead
\bottomrule\noalign{}
\endlastfoot
基础必讲 &
静电基本方程、电场线与等势线、电偶极子、低阶多极展开、导体静电平衡、静电能与电容、唯一性定理、直角与球坐标分离变量、柱坐标分离变量、电像法基本模型
& 静电边值问题的常用工具 \\
拔高选讲 &
多极展开数学特性与综合训练、电容系数矩阵与虚功法、局域稳定性分析、球谐展开、特殊函数、保角变换、二维柱坐标通解严密推导、复杂导体边值题库
& 服务于长推导和综合建模 \\
\end{longtable}
}

\par\noindent\begin{minipage}{\linewidth}
\begin{center}
\includegraphics[width=\linewidth,height=.40\textheight,keepaspectratio]{\detokenize{assets/module-02/theory/tikz-v11-module2-method-map.pdf}}
\par\smallskip{\small 图示：模块二由电荷源、边界条件和唯一性组织静电场方法}
\end{center}
\end{minipage}\par

\section{第一单元：静电场、导体与唯一性}\label{m2:ux7b2cux4e00ux5355ux5143ux9759ux7535ux573aux5bfcux4f53ux4e0eux552fux4e00ux6027}

专题八至十六依次讨论源、场和势，以及导体边界、静电能和唯一性定理。

\section{专题八：静电场的基本方程（基础必讲）}\label{m2:ux4e13ux9898ux516bux9759ux7535ux573aux7684ux57faux672cux65b9ux7a0bux57faux7840ux5fc5ux8bb2}

电荷分布给出源项，基本方程约束电场，电势把矢量问题化为标量问题，边界条件再从数学解中确定实际解。电场线、多极展开、唯一性定理、分离变量法和电像法都使用这组关系。

\par\noindent\begin{minipage}{\linewidth}
\begin{center}
\includegraphics[width=\linewidth,height=.40\textheight,keepaspectratio]{\detokenize{assets/module-02/original/tikz-m2-topic8-v1-01-source-field-potential.png}}
\par\smallskip{\small 图示：静电学由电荷源、电场和电势三个层次构成，四条基本关系把它们联系起来}
\end{center}
\end{minipage}\par

\subsection{一、静电问题的前提、对象与记号}\label{m2:ux4e00ux9759ux7535ux95eeux9898ux7684ux524dux63d0ux5bf9ux8c61ux4e0eux8bb0ux53f7}

\subsubsection{1.
静电近似的前提}\label{m2:ux9759ux7535ux8fd1ux4f3cux7684ux524dux63d0}

在所选惯性参考系中，静电场满足：

\[
\frac{\partial \rho}{\partial t}=0,
\qquad
\frac{\partial \mathbf E}{\partial t}=0,
\qquad
\frac{\partial \mathbf B}{\partial t}=0
\]

其中最后一个条件只要求考察区域内没有随时间变化的磁场。由麦克斯韦---法拉第方程

\[
\nabla\times\mathbf E=-\frac{\partial\mathbf B}{\partial t}
\]

可得静电场的局域无旋性：

\[
\boxed{\nabla\times\mathbf E=\mathbf 0}.
\]

必须注意：\textbf{仅仅说``电荷静止''还不充分}。若外界仍有变化磁场，空间中可以同时存在感生电场，此时总电场不再是静电场。

\subsubsection{2.
三个基本对象}\label{m2:ux4e09ux4e2aux57faux672cux5bf9ux8c61}

\begin{itemize}
\tightlist
\item
  \(\rho(\mathbf r)\)：总电荷密度，单位为
  \(\mathrm{C/m^3}\)；点、线、面电荷可用 \(\delta\) 分布统一写成体密度。
\item
  \(\mathbf E(\mathbf r)\)：电场强度，单位为
  \(\mathrm{N/C}=\mathrm{V/m}\)。
\item
  \(\varphi(\mathbf r)\)：电势，单位为
  \(\mathrm V\)；它只确定到一个加法常数。
\end{itemize}

本专题讨论真空中的微观静电方程，因此 \(\rho\)
表示\textbf{总电荷密度}。自由电荷与束缚电荷的区分，要到电介质专题中引入
\(\mathbf D\) 后才有必要。

为避免源点与场点混淆，约定

\[
\mathbf R=\mathbf r-\mathbf r',
\qquad
R=|\mathbf R|,
\]

其中 \(\mathbf r'\) 是源点，\(\mathbf r\) 是场点；不带撇号的 \(\nabla\)
对场点坐标 \(\mathbf r\) 求导。

\subsection{二、从电荷分布到电场与电势}\label{m2:ux4e8cux4eceux7535ux8377ux5206ux5e03ux5230ux7535ux573aux4e0eux7535ux52bf}

\subsubsection{1. 库仑积分}\label{m2:ux5e93ux4ed1ux79efux5206}

位于 \(\mathbf r'\) 的点电荷 \(q\) 在场点 \(\mathbf r\) 产生

\[
\mathbf E(\mathbf r)
=
\frac{1}{4\pi\varepsilon_0}
\frac{q(\mathbf r-\mathbf r')}{|\mathbf r-\mathbf r'|^3}.
\]

利用叠加原理，连续电荷分布产生的电场为

\[
\boxed{
\mathbf E(\mathbf r)
=
\frac{1}{4\pi\varepsilon_0}
\int
\rho(\mathbf r')
\frac{\mathbf r-\mathbf r'}{|\mathbf r-\mathbf r'|^3}
\,dV'
}.
\]

积分覆盖全部源区。若场点落在连续电荷分布内部，积分是否普通收敛要结合
\(\rho\)
的光滑性判断；若含点电荷，则应把相应项单独写出或在分布意义下理解。

\subsubsection{2.
电势积分及其零点}\label{m2:ux7535ux52bfux79efux5206ux53caux5176ux96f6ux70b9}

对于局域电荷分布，若取无穷远为电势零点，可定义

\[
\boxed{
\varphi(\mathbf r)
=
\frac{1}{4\pi\varepsilon_0}
\int
\frac{\rho(\mathbf r')}{|\mathbf r-\mathbf r'|}
\,dV'
}.
\]

更一般地，应写成

\[
\varphi(\mathbf r)
=
\frac{1}{4\pi\varepsilon_0}
\int
\frac{\rho(\mathbf r')}{|\mathbf r-\mathbf r'|}
\,dV'+C.
\]

利用恒等式

\[
\nabla\frac{1}{R}
=
-\frac{\mathbf R}{R^3},
\]

并在允许把微分移入积分号时，有

\[
-\nabla\varphi
=
\frac{1}{4\pi\varepsilon_0}
\int \rho(\mathbf r')
\frac{\mathbf R}{R^3}
\,dV'
=\mathbf E.
\]

因此

\[
\boxed{\mathbf E=-\nabla\varphi}.
\]

电势的加法常数 \(C\)
不影响电场。对于无限长线电荷或无限大平面电荷，\(\varphi(\infty)\)
通常不存在，不能机械地令无穷远电势为零；此时应在有限参考点规定
\(\varphi\)，或只讨论电势差。

\textbf{例 M2-001：匀强电场与平行等势面}

某真空区域内存在匀强静电场 \(\mathbf E=E_0\mathbf e_z\)，\(E_0\)
为常量，电势采用同一任意参考零点。由 \(\mathbf E=-\nabla\varphi\)
求含任意加常数的电势，说明沿 \(+z\)
方向电势如何变化，并描述三维空间中等势面的形状及其与电场方向的关系。

\par\noindent\begin{minipage}{\linewidth}
\textbf{解答}

\begin{center}
\includegraphics[width=\linewidth,height=.40\textheight,keepaspectratio]{\detokenize{assets/module-02/examples/M2-001.pdf}}
\par\smallskip{\small 图 M2-001：匀强电场与平行等势面}
\end{center}
\end{minipage}\par

图 M2-001 的水平截线表示 \(z\) 不变的等势面，场箭头沿其法向。沿 \(z\)
积分 \(-\partial_z\varphi=E_0\)
得线性电势，任意加法常数只改变各面的电势标记。

若 \(\mathbf E=E_0\mathbf e_z\)，则

\[
\varphi(z)=-E_0z+C.
\]

当 \(E_0>0\) 时，沿 \(+z\) 电势下降；\(E_0<0\) 时上升；\(E_0=0\)
时电势处处相同。图 M2-001 的箭头取 \(E_0>0\)。非零场时完整等势面为平面
\(z=\text{常量}\)，其法向平行于电场。

\subsection{三、库仑定律与高斯定律}\label{m2:ux4e09ux5e93ux4ed1ux5b9aux5f8bux4e0eux9ad8ux65afux5b9aux5f8b}

\subsubsection{1.
单个点电荷与有向立体角}\label{m2:ux5355ux4e2aux70b9ux7535ux8377ux4e0eux6709ux5411ux7acbux4f53ux89d2}

设点电荷 \(q\) 位于 \(\mathbf r'\)。闭合曲面上的面元为

\[
d\mathbf S=\hat{\mathbf n}\,dS.
\]

点电荷通过该面元的通量为

\[
\mathbf E\cdot d\mathbf S
=
\frac{q}{4\pi\varepsilon_0}
\frac{\mathbf R\cdot\hat{\mathbf n}}{R^3}
\,dS.
\]

定义面元对源点张开的\textbf{有向立体角}

\[
d\Omega
=
\frac{\hat{\mathbf R}\cdot\hat{\mathbf n}}{R^2}
\,dS,
\]

于是

\[
\mathbf E\cdot d\mathbf S
=
\frac{q}{4\pi\varepsilon_0}d\Omega.
\]

\par\noindent\begin{minipage}{\linewidth}
\begin{center}
\includegraphics[width=\linewidth,height=.40\textheight,keepaspectratio]{\detokenize{assets/module-02/original/tikz-m2-topic8-v2-02-solid-angle-gauss.png}}
\par\smallskip{\small 图示：点电荷对面元的通量等于电荷强度乘以面元张开的有向立体角，闭合曲面的总立体角只计数内部源}
\end{center}
\end{minipage}\par

对简单闭合曲面，有

\[
\oint_{\partial V}d\Omega
=
\begin{cases}
4\pi, & \mathbf r'\in V,\\
0, & \mathbf r'\notin V.
\end{cases}
\]

源在内部时，闭合曲面从源点看覆盖整个方向空间；源在外部时，同一射线对曲面的进入与离开具有相反的有向立体角，贡献成对抵消。因此

\[
\oint_{\partial V}\mathbf E\cdot d\mathbf S
=
\begin{cases}
q/\varepsilon_0, & q\text{ 在 }V\text{ 内},\\
0, & q\text{ 在 }V\text{ 外}.
\end{cases}
\]

这比``电场线穿入与穿出相消''更严格，因为它明确给出了每个面元的带符号权重。

\subsubsection{2.
多个电荷与连续分布}\label{m2:ux591aux4e2aux7535ux8377ux4e0eux8fdeux7eedux5206ux5e03}

对所有点电荷叠加，只有闭合曲面内部的电荷留下贡献：

\[
\boxed{
\oint_{\partial V}\mathbf E\cdot d\mathbf S
=
\frac{Q_V}{\varepsilon_0}
=
\frac{1}{\varepsilon_0}\int_V\rho\,dV
}.
\]

这就是高斯定律的积分形式。它对任意闭合曲面成立，但只有在对称性足以使
\(\mathbf E\cdot d\mathbf S\) 易于积分时，才能直接用来求场。

由散度定理

\[
\oint_{\partial V}\mathbf E\cdot d\mathbf S
=
\int_V\nabla\cdot\mathbf E\,dV,
\]

可得光滑区域中的微分形式

\[
\boxed{\nabla\cdot\mathbf E=\frac{\rho}{\varepsilon_0}}.
\]

含点电荷时应在分布意义下使用

\[
\nabla\cdot\frac{\mathbf r-\mathbf r'}{|\mathbf r-\mathbf r'|^3}
=
4\pi\delta^{(3)}(\mathbf r-\mathbf r').
\]

所以点电荷处不是``把普通发散公式直接代入
\(r=0\)''，而是用任意小闭合面的有限通量来定义集中源强。

\textbf{例 M2-002：同心高斯面只读取内部电荷}

真空中电荷分布以原点为中心球对称，体密度为已知函数
\(\rho(r)\)，不另加破坏球对称性的外场。取半径 \(r>0\)
的同心球面，规定外径向为
\(+\mathbf e_r\)。写出它包围的总电荷积分，利用高斯定律求
\(\mathbf E(r)\)，并说明半径大于 \(r\)
的球对称电荷层是否进入该通量的源项。

\par\noindent\begin{minipage}{\linewidth}
\textbf{解答}

\begin{center}
\includegraphics[width=\linewidth,height=.40\textheight,keepaspectratio]{\detokenize{assets/module-02/examples/M2-002.pdf}}
\par\smallskip{\small 图 M2-002：同心高斯面只读取内部电荷}
\end{center}
\end{minipage}\par

图 M2-002
的同心球面与电荷分布具有相同对称中心，面上径向场分量相同。先对球内密度积分得到
\(Q(r)\)，再把通量写成
\(E_r4\pi r^2\)；球对称性是把场强提出积分号的依据。

设

\[
Q(r)=4\pi\int_0^r\rho(s)s^2\,ds.
\]

取半径为 \(r\) 的同心高斯球面，立即得到

\[
\mathbf E(r)
=
\frac{Q(r)}{4\pi\varepsilon_0r^2}\mathbf e_r.
\]

这里真正使高斯定律能直接解出 \(\mathbf E\)
的，是球对称性，而不是球面本身。

\subsection{四、无旋性、路径无关与电势}\label{m2:ux56dbux65e0ux65cbux6027ux8defux5f84ux65e0ux5173ux4e0eux7535ux52bf}

静电条件下 \(\partial\mathbf B/\partial t=0\)，故

\[
\nabla\times\mathbf E=\mathbf 0.
\]

若曲面 \(S\) 完全位于场的光滑定义域内，则斯托克斯定理给出

\[
\oint_{\partial S}\mathbf E\cdot d\mathbf l
=
\int_S(\nabla\times\mathbf E)\cdot d\mathbf S
=0.
\]

于是从 \(A\) 到 \(B\) 的线积分与路径无关，可以定义

\[
\boxed{
\varphi(B)-\varphi(A)
=
-\int_A^B\mathbf E\cdot d\mathbf l
}.
\]

\par\noindent\begin{minipage}{\linewidth}
\begin{center}
\includegraphics[width=\linewidth,height=.40\textheight,keepaspectratio]{\detokenize{assets/module-02/original/tikz-m2-topic8-v1-03-path-independent-potential-safe-v7.png}}
\par\smallskip{\small 图示：静电场在两点间沿不同路径的线积分相同，电场垂直等势线并指向电势降低方向}
\end{center}
\end{minipage}\par

这里要区分两个层次：

\begin{itemize}
\tightlist
\item
  \(\nabla\times\mathbf E=0\) 是局域结论。
\item
  从局域无旋推出全局单值电势，还需区域单连通，或另行验证所有不可缩回路的环量都为零。
\end{itemize}

由库仑积分得到的真正静电场已经有显式的单值电势积分，因此在排除电荷奇点后的定义域中不会额外产生非零静电环量。拓扑条件主要用于提醒：不能只凭局部微分式，跨过孔洞或奇线草率作全局结论。

\textbf{例 M2-003：不同路径具有相同的端点功}

某静电场中两固定点 \(A,B\) 的电势分别为
\(\varphi_A,\varphi_B\)，两者相对于同一参考零点。一个不会改变原场分布的试探电荷
\(q_0\) 沿任意路径从 \(A\) 移到
\(B\)。求静电场力所做的功，并比较两条不同路径的结果；明确电场力做功与外力准静态做功的符号区别。

\par\noindent\begin{minipage}{\linewidth}
\textbf{解答}

\begin{center}
\includegraphics[width=\linewidth,height=.40\textheight,keepaspectratio]{\detokenize{assets/module-02/examples/M2-003.pdf}}
\par\smallskip{\small 图 M2-003：不同路径具有相同的端点功}
\end{center}
\end{minipage}\par

沿图 M2-003 的绿色或橙色路径由 \(A\) 到 \(B\)，势能改变量均为
\(q_0(\varphi_B-\varphi_A)\)。电场力做功取其相反数，故两条路径都给出
\(q_0(\varphi_A-\varphi_B)\)，不需要逐段计算曲线路径。

若 \(A,B\) 的电势分别为 \(\varphi_A,\varphi_B\)，电场力对电荷 \(q_0\)
做功为

\[
W_{A\to B}
=q_0\int_A^B\mathbf E\cdot d\mathbf l
=q_0(\varphi_A-\varphi_B),
\]

只取决于端点，不取决于实际轨迹。

\subsection{五、泊松方程与拉普拉斯方程}\label{m2:ux4e94ux6ccaux677eux65b9ux7a0bux4e0eux62c9ux666eux62c9ux65afux65b9ux7a0b}

把 \(\mathbf E=-\nabla\varphi\) 代入高斯定律：

\[
\nabla\cdot\mathbf E
=
-\nabla\cdot\nabla\varphi
=
-\nabla^2\varphi
=
\frac{\rho}{\varepsilon_0}.
\]

因此

\[
\boxed{\nabla^2\varphi=-\frac{\rho}{\varepsilon_0}}
\]

称为泊松方程。在\textbf{总电荷密度为零}的区域中，它退化为

\[
\boxed{\nabla^2\varphi=0},
\]

即拉普拉斯方程。

也可从电势积分直接验证。因为

\[
\nabla^2\frac{1}{|\mathbf r-\mathbf r'|}
=
-4\pi\delta^{(3)}(\mathbf r-\mathbf r'),
\]

所以

\[
\nabla^2\varphi(\mathbf r)
=
\frac{1}{4\pi\varepsilon_0}
\int\rho(\mathbf r')
\nabla^2\frac{1}{|\mathbf r-\mathbf r'|}
\,dV'
=
-\frac{\rho(\mathbf r)}{\varepsilon_0}.
\]

\par\noindent\begin{minipage}{\linewidth}
\begin{center}
\includegraphics[width=\linewidth,height=.40\textheight,keepaspectratio]{\detokenize{assets/module-02/theory/unified7-poisson-regions.pdf}}
\par\smallskip{\small 图示：有电荷区域与无电荷区域应分别使用泊松方程和拉普拉斯方程，再用界面条件连接}
\end{center}
\end{minipage}\par

拉普拉斯方程只在无源区域内成立。若点电荷、线电荷、面电荷或体电荷位于区域中，就必须挖去奇点、分区列方程，或把源写入泊松方程；不能把整个空间一律写成
\(\nabla^2\varphi=0\)。

\textbf{例 M2-004：泊松源项决定电势曲率}

真空中某局部区域内体电荷密度为常量
\(\rho_0\)。已知该区域的电势只依赖直角坐标 \(x\)，不依赖
\(y,z\)，尚未指定两端边界电势。由泊松方程求电势的局部一般解，保留待由边界条件确定的积分常数，并说明体电荷密度决定哪一个阶次的电势变化。

\par\noindent\begin{minipage}{\linewidth}
\textbf{解答}

\begin{center}
\includegraphics[width=\linewidth,height=.40\textheight,keepaspectratio]{\detokenize{assets/module-02/examples/M2-004.pdf}}
\par\smallskip{\small 图 M2-004：泊松源项决定电势曲率}
\end{center}
\end{minipage}\par

图 M2-004 用向下弯曲的剖面示意正的
\(\rho_0\)。泊松方程固定二阶导数，积分两次得到二次项与两个常数；图上的高度和整体倾斜尚不能由局部密度单独确定。

若某区内 \(\rho=\rho_0\)，且 \(\varphi\) 只依赖 \(x\)，则

\[
\frac{d^2\varphi}{dx^2}
=
-\frac{\rho_0}{\varepsilon_0}.
\]

因此

\[
\varphi(x)
=
-\frac{\rho_0}{2\varepsilon_0}x^2+Ax+B.
\]

正电荷使电势曲线局部向下弯；常数 \(A,B\) 仍须由边界条件确定。

\subsection{六、面电荷界面的连接条件}\label{m2:ux516dux9762ux7535ux8377ux754cux9762ux7684ux8fdeux63a5ux6761ux4ef6}

设界面法向单位矢量 \(\hat{\mathbf n}\) 从区域 \(1\) 指向区域
\(2\)，界面上存在面电荷密度 \(\sigma\)。令两侧场的极限分别为
\(\mathbf E_1\)、\(\mathbf E_2\)。

\subsubsection{1.
法向分量的跳跃}\label{m2:ux6cd5ux5411ux5206ux91cfux7684ux8df3ux8dc3}

取跨越界面的薄柱面，底面积为
\(\Delta S\)、高度趋于零。侧面通量随高度趋零，两个端面的净通量为

\[
(\mathbf E_2\cdot\hat{\mathbf n}
-\mathbf E_1\cdot\hat{\mathbf n})\Delta S.
\]

柱面包围的电荷为 \(\sigma\Delta S\)，故

\[
\boxed{
\hat{\mathbf n}\cdot(\mathbf E_2-\mathbf E_1)
=
\frac{\sigma}{\varepsilon_0}
}.
\]

\subsubsection{2.
切向分量的连续}\label{m2:ux5207ux5411ux5206ux91cfux7684ux8fdeux7eed}

取跨越界面的薄矩形回路，长边平行于任意切向单位矢量
\(\hat{\mathbf t}\)，短边长度趋于零。静电环量为零，得到

\[
(\mathbf E_2-\mathbf E_1)\cdot\hat{\mathbf t}=0.
\]

因为切向方向任意，所以

\[
\boxed{
\hat{\mathbf n}\times(\mathbf E_2-\mathbf E_1)
=
\mathbf 0
}.
\]

\par\noindent\begin{minipage}{\linewidth}
\begin{center}
\includegraphics[width=\linewidth,height=.40\textheight,keepaspectratio]{\detokenize{assets/module-02/original/tikz-m2-topic8-v1-05-surface-charge-jump-pdf-v11.jpg}}
\par\smallskip{\small 图示：薄柱面给出面电荷两侧法向场的跳跃，薄矩形回路给出切向场的连续}
\end{center}
\end{minipage}\par

\subsubsection{3.
电势的连接条件}\label{m2:ux7535ux52bfux7684ux8fdeux63a5ux6761ux4ef6}

若两侧电场有限，沿一条长度趋于零的法向线段积分，有

\[
\varphi_2-\varphi_1
=
-\lim_{h\to0}\int_{-h}^{h}\mathbf E\cdot d\mathbf l
=0.
\]

因此在不存在理想偶极层时

\[
\boxed{\varphi_2=\varphi_1}.
\]

再由 \(E_n=-\partial\varphi/\partial n\)，法向跳跃条件可写为

\[
\boxed{
\frac{\partial\varphi_2}{\partial n}
-
\frac{\partial\varphi_1}{\partial n}
=
-\frac{\sigma}{\varepsilon_0}
}.
\]

\textbf{例 M2-005：另一侧场为零时的法向跳跃}

在普通面电荷界面附近，区域 1 与区域 2
均为真空，界面实际总面电荷密度为有符号量 \(\sigma\)。统一法向
\(\mathbf n_{12}\) 由区域 1 指向区域 2，并定义
\(E_{in}=\mathbf E_i\cdot\mathbf n_{12}\)。已知紧邻界面的
\(E_{1n}=0\)，求
\(E_{2n}\)。这里只确定法向分量，不假定未给出的切向场为零。

\par\noindent\begin{minipage}{\linewidth}
\textbf{解答}

\begin{center}
\includegraphics[width=\linewidth,height=.40\textheight,keepaspectratio]{\detokenize{assets/module-02/examples/M2-005.pdf}}
\par\smallskip{\small 图 M2-005：另一侧场为零时的法向跳跃}
\end{center}
\end{minipage}\par

图 M2-005
中橙色法向从下侧指向上侧，下侧场为零。薄柱两盖的区域外法向相反，因此净通量读作
\((E_{2n}-E_{1n})\Delta S\)；代入面电荷与零场条件，才得到单侧结果。

若 \(E_{1n}=0\)，界面带面电荷 \(\sigma\)，则

\[
E_{2n}=\frac{\sigma}{\varepsilon_0}.
\]

这个结论不是说每个带电面的单侧场都等于
\(\sigma/\varepsilon_0\)；它还使用了另一侧场为零这一额外条件。

\subsection{七、边界条件、电势零点与相容性}\label{m2:ux4e03ux8fb9ux754cux6761ux4ef6ux7535ux52bfux96f6ux70b9ux4e0eux76f8ux5bb9ux6027}

在给定区域 \(V\) 内求解泊松方程或拉普拉斯方程时，常见边界条件有三类。

\subsubsection{1.
给定电势：狄利克雷条件}\label{m2:ux7ed9ux5b9aux7535ux52bfux72c4ux5229ux514bux96f7ux6761ux4ef6}

\[
\varphi|_{\partial V}=f.
\]

典型情形是给定导体表面电势，或给定无穷远处的渐近电势。

\subsubsection{2.
给定法向导数：诺伊曼条件}\label{m2:ux7ed9ux5b9aux6cd5ux5411ux5bfcux6570ux8bfaux4f0aux66fcux6761ux4ef6}

\[
\left.\frac{\partial\varphi}{\partial n}\right|_{\partial V}=g.
\]

因为 \(E_n=-\partial\varphi/\partial n\)，这等价于给定边界上的法向电场。

\subsubsection{3. 混合条件}\label{m2:ux6df7ux5408ux6761ux4ef6}

在边界的不同部分分别给定电势与法向导数。也有更一般的线性组合条件，但竞赛中最常见的是分段混合条件。

\par\noindent\begin{minipage}{\linewidth}
\begin{center}
\includegraphics[width=\linewidth,height=.40\textheight,keepaspectratio]{\detokenize{assets/module-02/original/tikz-m2-topic8-v1-06-boundary-condition-types.png}}
\par\smallskip{\small 图示：静电边值问题可给定电势、法向导数，或在不同边界段给定不同类型的条件}
\end{center}
\end{minipage}\par

纯诺伊曼问题还必须满足相容性条件。对泊松方程在 \(V\) 内积分：

\[
\oint_{\partial V}
\frac{\partial\varphi}{\partial n}
\,dS
=
\int_V\nabla^2\varphi\,dV
=
-\frac{1}{\varepsilon_0}\int_V\rho\,dV.
\]

所以若 \(g=\partial\varphi/\partial n\)，必须有

\[
\boxed{
\oint_{\partial V}g\,dS
=
-\frac{Q_V}{\varepsilon_0}
}.
\]

满足相容性后，纯诺伊曼问题的电势仍只确定到加法常数，因为 \(\varphi\) 与
\(\varphi+C\)
有相同的法向导数。狄利克雷条件则同时固定了电势零点。关于解的唯一性将在后续专题严格证明。

\textbf{例 M2-006：镜面对称使法向斜率为零}

一个静电边值问题的电荷分布、求解区域及给定边界条件全部关于平面 \(x=0\)
镜面对称，并满足电势解的唯一性条件；两侧使用同一电势参考。在所考察的对称面点的邻域内，电势一阶连续可微，该处没有导致法向场跳跃的奇异电荷层。证明或说明电势在对称面两侧满足怎样的对称性，据此确定
\(x=0\) 上的法向导数 \(\partial\varphi/\partial x\)
及相应的法向电场分量。

\par\noindent\begin{minipage}{\linewidth}
\textbf{解答}

\begin{center}
\includegraphics[width=\linewidth,height=.40\textheight,keepaspectratio]{\detokenize{assets/module-02/examples/M2-006.pdf}}
\par\smallskip{\small 图 M2-006：镜面对称使法向斜率为零}
\end{center}
\end{minipage}\par

图 M2-006 的势曲线关于 \(x=0\) 镜像，中央绿色切线水平。由偶函数关系对
\(x\)
求导，左右斜率相反，在对称面相接处只能为零，故该面可使用齐次法向导数条件。

若电荷与边界关于平面 \(x=0\) 镜面对称，则

\[
\varphi(-x,y,z)=\varphi(x,y,z),
\]

从而在对称面上

\[
\left.\frac{\partial\varphi}{\partial x}\right|_{x=0}=0.
\]

因此对称面常可作为齐次诺伊曼边界使用。
若对称面本身带奇异面电荷，则只能由对称性推出两侧法向场互为反号；例如
\(\varphi\propto-|x|\) 虽为偶函数，面上法向导数却不连续，不能直接置零。

\subsection{八、静电题中方程和方法的选择}\label{m2:ux516bux9759ux7535ux9898ux4e2dux65b9ux7a0bux548cux65b9ux6cd5ux7684ux9009ux62e9}

\par\noindent\begin{minipage}{\linewidth}
\begin{center}
\includegraphics[width=\linewidth,height=.40\textheight,keepaspectratio]{\detokenize{assets/module-02/original/tikz-m2-topic8-v1-07-solution-method-flow-pdf-v11.jpg}}
\par\smallskip{\small 图示：静电题应先分区和检查对称性，再在高斯定律、库仑积分、泊松方程与拉普拉斯边值问题之间选择}
\end{center}
\end{minipage}\par

\subsubsection{1. 先分区}\label{m2:ux5148ux5206ux533a}

标出体电荷区、无源区、点线面奇点、导体表面、有限边界和无穷远边界。每个区域内分别列方程，再用连接条件把解接起来。

\subsubsection{2. 再选未知量}\label{m2:ux518dux9009ux672aux77e5ux91cf}

\begin{itemize}
\tightlist
\item
  高度对称、目标是场强或通量：优先用高斯定律直接求 \(\mathbf E\)。
\item
  已知有限电荷分布、边界不复杂：可用库仑积分或电势积分。
\item
  无源区且边界决定性强：求解 \(\nabla^2\varphi=0\)，再由
  \(\mathbf E=-\nabla\varphi\) 还原电场。
\item
  有源区且边界也重要：求解泊松方程，并与相邻无源区的拉普拉斯解连接。
\end{itemize}

\subsubsection{3. 写全条件}\label{m2:ux5199ux5168ux6761ux4ef6}

除了显式给出的边界条件，还要检查：轴线上是否有限、原点是否允许奇异、无穷远如何衰减、对称面上的法向导数是否为零、界面上是否存在面电荷。

\subsubsection{4. 多重检验}\label{m2:ux591aux91cdux68c0ux9a8c}

\begin{enumerate}
\def\labelenumi{\arabic{enumi}.}
\tightlist
\item
  \textbf{方程检验}：分区代回泊松或拉普拉斯方程。
\item
  \textbf{界面检验}：检查 \(\varphi\)、切向场和法向跳跃。
\item
  \textbf{积分检验}：用高斯定律核对包围总电荷，用环量核对静电无旋性。
\item
  \textbf{极限检验}：考察
  \(r\to0\)、\(r\to\infty\)、几何参数趋于特殊值时是否合理。
\item
  \textbf{量纲检验}：\(\varphi\) 应为伏特，\(\mathbf E\) 应为伏特每米。
\end{enumerate}

\subsection{九、静电基本方程的使用条件}\label{m2:ux4e5dux9759ux7535ux57faux672cux65b9ux7a0bux7684ux4f7fux7528ux6761ux4ef6}

\begin{enumerate}
\def\labelenumi{\arabic{enumi}.}
\tightlist
\item
  \textbf{把整个空间都写成拉普拉斯方程}：必须先挖去奇点并分区；有电荷处应使用泊松方程或分布源。
\item
  \textbf{认为高斯定律总能直接求场}：高斯定律总成立，但缺乏足够对称性时，闭合面上的场强未知，积分式不能直接解出局域场。
\item
  \textbf{默认无穷远电势为零}：无限长线电荷和无限大平面电荷通常不允许这样选取零点。
\item
  \textbf{由局部无旋直接断言全局路径无关}：非单连通域还要检查不可缩回路的环量。
\item
  \textbf{只解微分方程，不写边界条件}：泊松方程或拉普拉斯方程本身只给出一族解。
\item
  \textbf{混淆法向方向}：写跳跃条件前必须先规定 \(\hat{\mathbf n}\)
  从哪一侧指向哪一侧。
\item
  \textbf{把电势连续误解为法向导数连续}：有面电荷时电势仍连续，但法向导数发生跳跃。
\end{enumerate}

\subsection{十、练习与解答}\label{m2:ux5341ux7ec3ux4e60ux4e0eux89e3ux7b54}

\subsubsection{作业一：由拉普拉斯方程和高斯定律确定点电荷电势}\label{m2:ux4f5cux4e1aux4e00ux7531ux62c9ux666eux62c9ux65afux65b9ux7a0bux548cux9ad8ux65afux5b9aux5f8bux786eux5b9aux70b9ux7535ux8377ux7535ux52bf}

\textbf{题号：M2-007}

\textbf{层级：} 基础　　\textbf{方法：} 高斯定理、边值问题

真空中只有一个电荷量为 \(q\) 的点电荷，固定在原点。采用球坐标
\(r=|\mathbf r|\)，取
\(\varphi(\infty)=0\)。不直接套用或积分库仑势：先在无源区域 \(r>0\)
利用拉普拉斯方程和球对称性求一般势函数，再用包围原点的高斯面确定系数。求
\(\varphi(r)\) 与 \(\mathbf E(r)\)，并说明原点奇点怎样进入系数条件。

\par\noindent\begin{minipage}{\linewidth}
\textbf{解答}

\begin{center}
\includegraphics[width=\linewidth,height=.40\textheight,keepaspectratio]{\detokenize{assets/module-02/examples/M2-007.pdf}}
\par\smallskip{\small 图 M2-007：由拉普拉斯方程和高斯定律确定点电荷电势}
\end{center}
\end{minipage}\par

图 M2-007 以原点为中心划出 \(r>0\)
的无源区域和同心高斯面。径向拉普拉斯方程先给 \(A+B/r\)，无穷远条件消去
\(A\)，包围电荷的通量再固定 \(B\)，两项条件分别进入不同的积分常数。

原点以外没有电荷，因此在 \(r>0\) 区域中

\[
\nabla^2\varphi=0.
\]

球对称性给出 \(\varphi=\varphi(r)\)。球坐标下

\[
\frac{1}{r^2}\frac{d}{dr}
\left(r^2\frac{d\varphi}{dr}\right)=0.
\]

逐次积分：

\[
r^2\frac{d\varphi}{dr}=C_1,
\qquad
\varphi(r)=C_2-\frac{C_1}{r}.
\]

令 \(A=C_2\)、\(B=-C_1\)，则

\[
\varphi(r)=A+\frac{B}{r},
\qquad
\mathbf E=-\nabla\varphi
=\frac{B}{r^2}\mathbf e_r.
\]

由 \(\varphi(\infty)=0\) 得 \(A=0\)。再对任意半径为 \(r\)
的球面应用高斯定律：

\[
4\pi r^2\frac{B}{r^2}
=
\frac{q}{\varepsilon_0},
\]

故

\[
B=\frac{q}{4\pi\varepsilon_0}.
\]

最终

\[
\boxed{
\varphi(r)=\frac{q}{4\pi\varepsilon_0r},
\qquad
\mathbf E(r)=\frac{q}{4\pi\varepsilon_0r^2}\mathbf e_r
}.
\]

\par\noindent\begin{minipage}{\linewidth}
\begin{center}
\includegraphics[width=\linewidth,height=.40\textheight,keepaspectratio]{\detokenize{assets/module-02/original/tikz-m2-static-02-point-charge-laplace-flux-pdf-v11.jpg}}
\par\smallskip{\small 图示：点电荷外部区域满足拉普拉斯方程，奇点强度由包围小球面的高斯通量固定}
\end{center}
\end{minipage}\par

这里 \(r>0\) 内的拉普拉斯方程不能确定 \(B\)；\(B\)
保存的是被挖去原点处的源强信息，必须由小球面的高斯通量补回。

\subsubsection{作业二：均匀带电实心球的场与势}\label{m2:ux4f5cux4e1aux4e8cux5747ux5300ux5e26ux7535ux5b9eux5fc3ux7403ux7684ux573aux4e0eux52bf}

\textbf{题号：M2-008}

\textbf{层级：} 基础　　\textbf{方法：} 静电场方程

真空中半径为 \(a>0\)、球心位于原点的实心球均匀带电，体密度为有符号常量
\(\rho_0\)；球外没有其他电荷，\(r=a\)
处不另附加面电荷。取无穷远电势为零。分别求 \(0\le r<a\) 与 \(r>a\)
的电场和电势，说明球心有限性如何使用，并检验球面两侧电势和法向电场的连续性。

\par\noindent\begin{minipage}{\linewidth}
\textbf{解答}

\begin{center}
\includegraphics[width=\linewidth,height=.40\textheight,keepaspectratio]{\detokenize{assets/module-02/examples/M2-008.pdf}}
\par\smallskip{\small 图 M2-008：均匀带电实心球的场与势}
\end{center}
\end{minipage}\par

图 M2-008 的球面 \(r=a\) 把体电荷区与外部无源区分开。球内包围电荷随
\(r^3\)
增加，外部包围电荷保持不变，所以场分段为线性与平方反比；积分电势后还要在同一球面连接数值和法向场。

由球对称性，\(\mathbf E=E(r)\mathbf e_r\)。

当 \(r<a\) 时，高斯球内电荷为

\[
Q(r)=\frac{4\pi}{3}\rho_0r^3.
\]

因此

\[
E(r)4\pi r^2
=
\frac{Q(r)}{\varepsilon_0},
\]

得到

\[
\mathbf E_{\mathrm{in}}(r)
=
\frac{\rho_0r}{3\varepsilon_0}\mathbf e_r.
\]

当 \(r>a\) 时，总电荷为

\[
Q=\frac{4\pi}{3}\rho_0a^3,
\]

故

\[
\mathbf E_{\mathrm{out}}(r)
=
\frac{\rho_0a^3}{3\varepsilon_0r^2}\mathbf e_r.
\]

先由外区和无穷远条件得到

\[
\varphi_{\mathrm{out}}(r)
=
\int_r^{\infty} E_{\mathrm{out}}(s)\,ds
=
\frac{\rho_0a^3}{3\varepsilon_0r}.
\]

球面电势为

\[
\varphi(a)=\frac{\rho_0a^2}{3\varepsilon_0}.
\]

内区电势满足

\[
\varphi_{\mathrm{in}}(r)-\varphi(a)
=
\int_r^a E_{\mathrm{in}}(s)\,ds
=
\frac{\rho_0(a^2-r^2)}{6\varepsilon_0},
\]

所以

\[
\boxed{
\varphi(r)=
\begin{cases}
\displaystyle
\frac{\rho_0}{6\varepsilon_0}(3a^2-r^2), & 0\le r\le a,\\[8pt]
\displaystyle
\frac{\rho_0a^3}{3\varepsilon_0r}, & r\ge a.
\end{cases}
}
\]

并且

\[
\boxed{
\mathbf E(r)=
\begin{cases}
\displaystyle
\frac{\rho_0r}{3\varepsilon_0}\mathbf e_r, & 0\le r\le a,\\[8pt]
\displaystyle
\frac{\rho_0a^3}{3\varepsilon_0r^2}\mathbf e_r, & r\ge a.
\end{cases}
}
\]

\par\noindent\begin{minipage}{\linewidth}
\begin{center}
\includegraphics[width=\linewidth,height=.40\textheight,keepaspectratio]{\detokenize{assets/module-02/original/tikz-m2-topic8-v1-09-charged-sphere-profiles.png}}
\par\smallskip{\small 图示：均匀带电实心球内场强正比于半径、球外按平方反比衰减，电势及其斜率在球面连续}
\end{center}
\end{minipage}\par

在 \(r=a\) 处，内外电势均为
\(\rho_0a^2/(3\varepsilon_0)\)，内外径向场均为
\(\rho_0a/(3\varepsilon_0)\)。球面没有额外的 \(\delta\)
型面电荷，所以法向场不跳跃。

\subsubsection{作业三：无限大均匀带电平面}\label{m2:ux4f5cux4e1aux4e09ux65e0ux9650ux5927ux5747ux5300ux5e26ux7535ux5e73ux9762}

\textbf{题号：M2-009}

\textbf{层级：} 复赛　　\textbf{方法：} 静电场方程

整个空间为真空，只有无限大平面 \(z=0\) 带均匀面电荷密度
\(\sigma>0\)，不叠加其他外场。利用平面对称性求上下两半空间的静电场。取平面自身电势
\(\varphi(0)=0\)，求两侧
\(\varphi(z)\)，说明电势在界面是否连续，以及本无限延伸源模型为何不能把无穷远设为有限的零电势。

\par\noindent\begin{minipage}{\linewidth}
\textbf{解答}

\begin{center}
\includegraphics[width=\linewidth,height=.40\textheight,keepaspectratio]{\detokenize{assets/module-02/examples/M2-009.pdf}}
\par\smallskip{\small 图 M2-009：无限大均匀带电平面}
\end{center}
\end{minipage}\par

图 M2-009
的薄柱两端分别位于平面两侧，外法向相反而场均背离正电荷面。两盖通量相加为
\(2E\Delta S\)，决定每侧场为
\(\sigma/(2\varepsilon_0)\)；再从平面零势向两侧积分得到含 \(|z|\)
的电势。

由平移对称性与镜面对称性，电场只能垂直于平面，且两侧大小相等、方向相反：

\[
\mathbf E(z)
=
E_0\operatorname{sgn}(z)\mathbf e_z.
\]

取跨越平面的薄柱面，两个端面通量之和为 \(2E_0\Delta S\)。由高斯定律

\[
2E_0\Delta S
=
\frac{\sigma\Delta S}{\varepsilon_0},
\]

故

\[
\boxed{
\mathbf E(z)
=
\frac{\sigma}{2\varepsilon_0}
\operatorname{sgn}(z)\mathbf e_z
}.
\]

这也满足跳跃条件

\[
E_z(0^+)-E_z(0^-)
=
\frac{\sigma}{\varepsilon_0}.
\]

由 \(E_z=-d\varphi/dz\) 且 \(\varphi(0)=0\)，得到

\[
\boxed{
\varphi(z)
=
-\frac{\sigma}{2\varepsilon_0}|z|
}.
\]

\par\noindent\begin{minipage}{\linewidth}
\begin{center}
\includegraphics[width=\linewidth,height=.40\textheight,keepaspectratio]{\detokenize{assets/module-02/original/tikz-m2-topic8-v1-10-infinite-sheet-jump.png}}
\par\smallskip{\small 图示：无限带电平面的法向场在界面跳跃，电势在平面处连续并向两侧线性降低}
\end{center}
\end{minipage}\par

当 \(|z|\to\infty\) 时，\(\varphi(z)\to-\infty\)，因此不存在有限的
\(\varphi(\infty)\)。这是无限电荷分布不能机械选取无穷远零点的典型例子。

\subsubsection{作业四：同心球面间的拉普拉斯边值问题}\label{m2:ux4f5cux4e1aux56dbux540cux5fc3ux7403ux9762ux95f4ux7684ux62c9ux666eux62c9ux65afux8fb9ux503cux95eeux9898}

\textbf{题号：M2-010}

\textbf{层级：} 复赛　　\textbf{方法：} 边值问题

两个同心导体球面的半径为 \(0<a<b\)，内球面保持电势
\(V_0\)，外球面保持电势
\(0\)。两球面间为无电荷真空区域，不另加破坏球对称性的源。只求 \(a<r<b\)
内的电势和径向电场，要求同时满足两端给定电势；再讨论保持 \(a,V_0\)
不变而令 \(b\to\infty\) 时结果如何变化。

\par\noindent\begin{minipage}{\linewidth}
\textbf{解答}

\begin{center}
\includegraphics[width=\linewidth,height=.40\textheight,keepaspectratio]{\detokenize{assets/module-02/examples/M2-010.pdf}}
\par\smallskip{\small 图 M2-010：同心球面间的拉普拉斯边值问题}
\end{center}
\end{minipage}\par

图 M2-010 标出内外同心球面 \(a,b\) 和中间的无源球壳。径向通解只有常量和
\(1/r\)
两项，两球面的电势分别给出两个方程；求导后的径向场必须与同一对边界的电势高低一致。

区域内 \(\rho=0\)，且具有球对称性，因此

\[
\frac{1}{r^2}\frac{d}{dr}
\left(r^2\frac{d\varphi}{dr}\right)=0.
\]

通解为

\[
\varphi(r)=A+\frac{B}{r}.
\]

代入两个边界条件：

\[
A+\frac{B}{a}=V_0,
\qquad
A+\frac{B}{b}=0.
\]

相减得到

\[
B\left(\frac1a-\frac1b\right)=V_0,
\]

所以

\[
B=\frac{V_0ab}{b-a},
\qquad
A=-\frac{V_0a}{b-a}.
\]

于是

\[
\boxed{
\varphi(r)
=
\frac{V_0a}{b-a}
\left(\frac br-1\right)
},
\]

\[
\boxed{
\mathbf E(r)
=
-\nabla\varphi
=
\frac{V_0ab}{(b-a)r^2}\mathbf e_r
}.
\]

内球所带电荷可由高斯定律求出：

\[
Q
=
4\pi a^2\varepsilon_0E(a^+)
=
4\pi\varepsilon_0\frac{ab}{b-a}V_0.
\]

因此两同心球面的电容为

\[
C=\frac{Q}{V_0}
=
4\pi\varepsilon_0\frac{ab}{b-a}.
\]

当 \(b\to\infty\) 时，

\[
\varphi(r)\to\frac{V_0a}{r},
\qquad
C\to4\pi\varepsilon_0a,
\]

恢复为半径 \(a\)
的孤立导体球结果。这个极限同时检验了边界条件、量纲与远场行为。

\section{专题九：电场线与等势线（基础必讲）}\label{m2:ux4e13ux9898ux4e5dux7535ux573aux7ebfux4e0eux7b49ux52bfux7ebfux57faux7840ux5fc5ux8bb2}

电场线表示电场方向和通量分布，等势面表示电势的空间分布。即使没有完整解析解，也可据此判断场的方向、强弱、边界行为和特殊点。

\par\noindent\begin{minipage}{\linewidth}
\begin{center}
\includegraphics[width=\linewidth,height=.40\textheight,keepaspectratio]{\detokenize{assets/module-02/original/tikz-m2-topic9-v1-01-field-equipotential-definition-pdf-v11.jpg}}
\par\smallskip{\small 图示：电场线是电场的积分曲线，等势面是电势的水平集，在正则点处二者正交且电势沿场线下降}
\end{center}
\end{minipage}\par

\subsection{一、电场线、等势面与正则点}\label{m2:ux4e00ux7535ux573aux7ebfux7b49ux52bfux9762ux4e0eux6b63ux5219ux70b9}

\subsubsection{1.
电场线的参数定义}\label{m2:ux7535ux573aux7ebfux7684ux53c2ux6570ux5b9aux4e49}

设一条曲线写成

\[
\mathbf r=\mathbf r(s).
\]

若曲线上每一点的切向量都与当地电场同向，则它是按电场方向定向的电场线：

\[
\boxed{
\frac{d\mathbf r}{ds}
=
\alpha(s)\mathbf E[\mathbf r(s)],
\qquad
\alpha(s)>0
}.
\]

\(\alpha(s)\)
只改变曲线参数化的快慢，不改变几何轨迹。若取弧长为参数，并且
\(\mathbf E\neq\mathbf0\)，可令

\[
\alpha(s)=\frac{1}{|\mathbf E|},
\qquad
\frac{d\mathbf r}{ds}
=
\frac{\mathbf E}{|\mathbf E|}.
\]

不关心方向时，也可写成

\[
d\mathbf r\times\mathbf E=\mathbf0.
\]

电场线的定义要求电场在曲线附近有意义。点电荷处 \(\mathbf E\)
奇异，零场点处
\(\mathbf E=\mathbf0\)，都不能直接用``唯一切线方向''描述。

\subsubsection{2.
等势面是电势的水平集}\label{m2:ux7b49ux52bfux9762ux662fux7535ux52bfux7684ux6c34ux5e73ux96c6}

等势面由

\[
\boxed{\varphi(\mathbf r)=C}
\]

定义。在二维截面中，它表现为等势线。

若某点满足

\[
\nabla\varphi\neq\mathbf0,
\]

则该点称为电势的正则点。由隐函数定理，正则点附近的等势集合是一张光滑曲面，其法向方向为
\(\nabla\varphi\)。

若 \(\nabla\varphi=\mathbf0\)，等价于
\(\mathbf E=\mathbf0\)，该点是电势的临界点。此处等势线可能改变拓扑、相交成分界结构，不能套用普通正交轨线图像。

\subsubsection{3. 正交性推导}\label{m2:ux6b63ux4ea4ux6027ux63a8ux5bfc}

沿等势面内任意切向位移 \(d\mathbf r_t\)，有

\[
d\varphi
=
\nabla\varphi\cdot d\mathbf r_t
=0.
\]

因此 \(\nabla\varphi\) 垂直于等势面。又因为

\[
\mathbf E=-\nabla\varphi,
\]

所以在 \(\mathbf E\neq\mathbf0\) 的正则点，

\[
\boxed{\mathbf E\perp\text{等势面}}.
\]

负号还给出了方向：电场由高电势指向低电势。

\textbf{例 M2-011：电场线与等势面正交}

在直角坐标系中，某区域的静电场为 \(\mathbf E=E_0\mathbf e_x\)，\(E_0>0\)
为常量，且场不随 \(y,z\)
改变。求一个相容的电势，写出电场线和等势面的几何方程，明确场线的有向切线方向，并检验电场线与等势面在交点处的夹角。

\par\noindent\begin{minipage}{\linewidth}
\textbf{解答}

\begin{center}
\includegraphics[width=\linewidth,height=.40\textheight,keepaspectratio]{\detokenize{assets/module-02/examples/M2-011.pdf}}
\par\smallskip{\small 图 M2-011：电场线与等势面正交}
\end{center}
\end{minipage}\par

图 M2-011 中竖直等势面与水平场线相互垂直。由 \(-\partial_x\varphi=E_0\)
得电势，再由场线切向条件保持 \(y,z\)
不变，两类曲线的不同常量分别标记不同场线和等势面。

若

\[
\mathbf E=E_0\mathbf e_x,
\]

则

\[
\varphi=-E_0x+C.
\]

电场线是

\[
y=\text{常量},\qquad z=\text{常量},
\]

等势面是

\[
x=\text{常量}.
\]

二者处处正交，且沿 \(+x\) 方向电势线性降低。

\subsection{二、电场线微分方程}\label{m2:ux4e8cux7535ux573aux7ebfux5faeux5206ux65b9ux7a0b}

\subsubsection{1.
直角坐标形式}\label{m2:ux76f4ux89d2ux5750ux6807ux5f62ux5f0f}

在直角坐标中

\[
d\mathbf r
=
dx\,\mathbf e_x
+dy\,\mathbf e_y
+dz\,\mathbf e_z,
\]

而

\[
\mathbf E
=
E_x\mathbf e_x
+E_y\mathbf e_y
+E_z\mathbf e_z.
\]

由 \(d\mathbf r\parallel\mathbf E\) 得

\[
\boxed{
\frac{dx}{E_x}
=
\frac{dy}{E_y}
=
\frac{dz}{E_z}
}.
\]

在二维平面中，若 \(E_x\neq0\)，可写为

\[
\boxed{
\frac{dy}{dx}
=
\frac{E_y(x,y)}{E_x(x,y)}
}.
\]

若某个场分量为零，不应机械地除以零。例如 \(E_y=0\) 可能意味着场线满足
\(y=\text{常量}\)；这时回到参数方程判断最稳妥。

\subsubsection{2.
正交曲线坐标中的尺度因子}\label{m2:ux6b63ux4ea4ux66f2ux7ebfux5750ux6807ux4e2dux7684ux5c3aux5ea6ux56e0ux5b50}

在一般正交曲线坐标 \((q_1,q_2,q_3)\) 中，

\[
d\mathbf r
=
h_1dq_1\,\mathbf e_1
+h_2dq_2\,\mathbf e_2
+h_3dq_3\,\mathbf e_3.
\]

所以场线方程是

\[
\boxed{
\frac{h_1dq_1}{E_1}
=
\frac{h_2dq_2}{E_2}
=
\frac{h_3dq_3}{E_3}
}.
\]

\par\noindent\begin{minipage}{\linewidth}
\begin{center}
\includegraphics[width=\linewidth,height=.40\textheight,keepaspectratio]{\detokenize{assets/module-02/theory/M2-H-coordinate.pdf}}
\par\smallskip{\small 图示：直角、圆柱和球坐标中的电场线方程必须包含相应线元的尺度因子}
\end{center}
\end{minipage}\par

圆柱坐标 \((\rho,\phi,z)\) 中

\[
\boxed{
\frac{d\rho}{E_{\rho}}
=
\frac{\rho\,d\phi}{E_{\phi}}
=
\frac{dz}{E_z}
};
\]

球坐标 \((r,\theta,\phi)\) 中

\[
\boxed{
\frac{dr}{E_r}
=
\frac{r\,d\theta}{E_{\theta}}
=
\frac{r\sin\theta\,d\phi}{E_{\phi}}
}.
\]

漏掉 \(\rho\)、\(r\) 或 \(r\sin\theta\)，得到的曲线方向会发生实质错误。

\subsubsection{3.
已知电势时的正交轨线法}\label{m2:ux5df2ux77e5ux7535ux52bfux65f6ux7684ux6b63ux4ea4ux8f68ux7ebfux6cd5}

若已经求得 \(\varphi(x,y)\)，等势线直接由

\[
\varphi(x,y)=C
\]

给出。沿等势线全微分：

\[
\frac{\partial\varphi}{\partial x}dx
+\frac{\partial\varphi}{\partial y}dy
=0.
\]

若 \(\partial\varphi/\partial y\neq0\)，等势线斜率为

\[
\left.\frac{dy}{dx}\right|_{\mathrm{\text{等势}}}
=
-\frac{\varphi_x}{\varphi_y}.
\]

而电场线斜率为

\[
\left.\frac{dy}{dx}\right|_{\mathrm{\text{场线}}}
=
\frac{E_y}{E_x}
=
\frac{\varphi_y}{\varphi_x}.
\]

因此在两个斜率都有限且非零时，

\[
\left.\frac{dy}{dx}\right|_{\mathrm{\text{等势}}}
\left.\frac{dy}{dx}\right|_{\mathrm{\text{场线}}}
=-1.
\]

所谓``取负倒数''只是正交性的坐标表达。遇到水平或竖直切线时，应直接使用梯度和切向量的点积为零，而不要强行使用斜率。

\subsection{三、电通量管与场强图像}\label{m2:ux4e09ux7535ux901aux91cfux7ba1ux4e0eux573aux5f3aux56feux50cf}

\subsubsection{1.
无源电场管的通量守恒}\label{m2:ux65e0ux6e90ux7535ux573aux7ba1ux7684ux901aux91cfux5b88ux6052}

取一束电场线围成管状侧面，再用两个横截面 \(\Sigma_1,\Sigma_2\)
封闭。由于侧面处 \(\mathbf E\) 与曲面相切，

\[
\Phi_{\mathrm{\text{侧}}}=0.
\]

若两个截面的法向都取沿电场管推进的同一方向，则高斯定律给出

\[
\int_{\Sigma_2}\mathbf E\cdot d\mathbf S
-
\int_{\Sigma_1}\mathbf E\cdot d\mathbf S
=
\frac{Q_{\mathrm{\text{管段}}}}{\varepsilon_0}.
\]

管段内无电荷时，

\[
\boxed{
\int_{\Sigma_1}\mathbf E\cdot d\mathbf S
=
\int_{\Sigma_2}\mathbf E\cdot d\mathbf S
}.
\]

若截面推进时跨过点电荷、线电荷或面电荷，通量会按包围电荷发生跃变，不能继续使用同一个常数。

\subsubsection{2.
等势面疏密与场强}\label{m2:ux7b49ux52bfux9762ux758fux5bc6ux4e0eux573aux5f3a}

沿等势面法向取两个相邻等势面，电势差为 \(\Delta\varphi\)，法向距离为
\(\Delta n\)。在足够小的局域范围内，

\[
|\mathbf E|
\approx
\frac{|\Delta\varphi|}{\Delta n}.
\]

因此，当相邻等势面取相同的电势差时，间距越小，场越强。

若电场管截面上场近似均匀且垂直于截面，则通量守恒简化为

\[
E_1A_1=E_2A_2.
\]

所以同一等通量管越细，场通常越强。

\par\noindent\begin{minipage}{\linewidth}
\begin{center}
\includegraphics[width=\linewidth,height=.40\textheight,keepaspectratio]{\detokenize{assets/module-02/original/tikz-m2-topic9-v1-04-spacing-flux-tube-pdf-v11.jpg}}
\par\smallskip{\small 图示：相邻等势面间距反映局域场强，同一无源电场管的截面通量保持不变}
\end{center}
\end{minipage}\par

图中``场线疏密''只有在每两条相邻绘制场线代表相同通量时才有定量意义。任意多画或少画几条线，不会改变真实电场。

\subsection{四、电场线的全局规则与特殊点}\label{m2:ux56dbux7535ux573aux7ebfux7684ux5168ux5c40ux89c4ux5219ux4e0eux7279ux6b8aux70b9}

\par\noindent\begin{minipage}{\linewidth}
\begin{center}
\includegraphics[width=\linewidth,height=.40\textheight,keepaspectratio]{\detokenize{assets/module-02/theory/M2-H-field-rules.pdf}}
\par\smallskip{\small 图示：光滑非零静电场的积分曲线切向唯一且不能闭合，零场点附近会出现分界线结构}
\end{center}
\end{minipage}\par

\subsubsection{1.
正则点处电场线不相交}\label{m2:ux6b63ux5219ux70b9ux5904ux7535ux573aux7ebfux4e0dux76f8ux4ea4}

若 \(\mathbf E(\mathbf r)\) 在某邻域连续可微且
\(\mathbf E\neq\mathbf0\)，参数方程

\[
\frac{d\mathbf r}{ds}=\mathbf E(\mathbf r)
\]

满足局部解的唯一性条件。通过同一初始点只能有一条积分曲线，因此两条不同电场线不能在正则点相交。

点电荷处场奇异，许多场线可以从同一源点发出；零场点处方向未定义，若干分界线可以以该点为共同极限。这些都不违反正则点的唯一性结论。

\subsubsection{2.
非零静电场线不能闭合}\label{m2:ux975eux96f6ux9759ux7535ux573aux7ebfux4e0dux80fdux95edux5408}

沿按电场方向参数化的场线，

\[
\frac{d\varphi}{ds}
=
\nabla\varphi\cdot\frac{d\mathbf r}{ds}
=
-\alpha(s)|\mathbf E|^2.
\]

在 \(\mathbf E\neq\mathbf0\) 处，

\[
\frac{d\varphi}{ds}<0.
\]

若场线闭合，走一周后回到同一点，电势应恢复原值；但上式要求电势严格下降，产生矛盾。因此静电场线不能形成非零闭合环路。

\subsubsection{3.
场线的起点、终点与无穷远}\label{m2:ux573aux7ebfux7684ux8d77ux70b9ux7ec8ux70b9ux4e0eux65e0ux7a77ux8fdc}

电场线可从正电荷出发，终止于负电荷；也可从无穷远进入或延伸到无穷远。对于局域电荷系统，若总电荷为零，则穿过足够大闭合面的净通量为零；无穷远处向外与向内的通量必须相互抵消。

零场点附近常有分界线，把趋向不同终点的场线族分开。判断分界线时，要同时利用对称性、局域线性化和全局通量。

\subsubsection{4.
电场线不是粒子轨迹}\label{m2:ux7535ux573aux7ebfux4e0dux662fux7c92ux5b50ux8f68ux8ff9}

电场线满足一阶几何方程

\[
\frac{d\mathbf r}{ds}\parallel\mathbf E,
\]

而带电粒子的实际运动满足

\[
m\frac{d^2\mathbf r}{dt^2}=q\mathbf E.
\]

粒子速度由初始速度和过去的受力共同决定，通常不与瞬时电场平行。只有在强阻尼、特殊初始条件或只讨论初始加速度方向时，粒子轨迹才可能近似沿场线。

\textbf{例 M2-012：零场点附近的双曲线场线}

给定与 \(z\) 无关的二维场
\(\mathbf E=a(x\mathbf e_x-y\mathbf e_y)\)，其中 \(a>0\) 为常量，研究
\(xy\)
子午截面。建立场线的有向切线条件并求一般场线方程，单独说明坐标轴上的特殊场线；判断原点是否为零场点，分析邻近不同象限中的场向及鞍形结构。

\par\noindent\begin{minipage}{\linewidth}
\textbf{解答}

\begin{center}
\includegraphics[width=\linewidth,height=.40\textheight,keepaspectratio]{\detokenize{assets/module-02/examples/M2-012.pdf}}
\par\smallskip{\small 图 M2-012：零场点附近的双曲线场线}
\end{center}
\end{minipage}\par

图 M2-012 的绿色双曲线在各象限沿红色场箭头延伸。将切向斜率写成
\(dy/dx=E_y/E_x\) 并积分得
\(xy=C\)，同时要单独保留两个坐标轴；原点是零场点，不能在此套用正则场线的唯一切向结论。

设

\[
\mathbf E=a(x\mathbf e_x-y\mathbf e_y),
\qquad a>0.
\]

原点是零场点。场线方程

\[
\frac{dy}{dx}=-\frac{y}{x}
\]

给出

\[
xy=C.
\]

\(x\) 轴与 \(y\)
轴是分界线；一般场线是双曲线。不同场线族可共同趋近零场点，但不能把原点当作具有唯一切向的正则点。

\subsection{五、三维对象与二维截面的对应}\label{m2:ux4e94ux4e09ux7ef4ux5bf9ux8c61ux4e0eux4e8cux7ef4ux622aux9762ux7684ux5bf9ux5e94}

\subsubsection{1.
点电荷：等势球面}\label{m2:ux70b9ux7535ux8377ux7b49ux52bfux7403ux9762}

点电荷 \(q\) 的电势与电场为

\[
\varphi(r)=\frac{q}{4\pi\varepsilon_0r},
\qquad
\mathbf E(r)=\frac{q}{4\pi\varepsilon_0r^2}\mathbf e_r.
\]

因此

\[
\varphi=C
\quad\Longleftrightarrow\quad
r=\text{常量}.
\]

真正的等势面是同心球面；过点电荷的平面只截出同心圆。

\par\noindent\begin{minipage}{\linewidth}
\begin{center}
\includegraphics[width=\linewidth,height=.40\textheight,keepaspectratio]{\detokenize{assets/module-02/original/tikz-m2-topic9-v1-05-point-charge-sphere-section-pdf-v11.jpg}}
\par\smallskip{\small 图示：点电荷的三维等势面是同心球面，二维同心圆只是球面的平面截线}
\end{center}
\end{minipage}\par

对 \(q>0\)，电场沿球面外法向穿出，电势随 \(r\) 增大而降低；对
\(q<0\)，电场方向反向。

\subsubsection{2.
无限长线电荷：等势圆柱面}\label{m2:ux65e0ux9650ux957fux7ebfux7535ux8377ux7b49ux52bfux5706ux67f1ux9762}

设线电荷沿 \(z\) 轴，线密度为 \(\lambda\)。圆柱对称性要求

\[
\mathbf E=E(\rho)\mathbf e_{\rho}.
\]

高斯定律给出

\[
\mathbf E(\rho)
=
\frac{\lambda}{2\pi\varepsilon_0\rho}\mathbf e_{\rho}.
\]

等势面是

\[
\rho=\text{常量}
\]

的同轴圆柱面。在垂直导线的横截面中，它们才表现为同心圆。

\par\noindent\begin{minipage}{\linewidth}
\begin{center}
\includegraphics[width=\linewidth,height=.40\textheight,keepaspectratio]{\detokenize{assets/module-02/original/tikz-m2-topic9-v1-08-line-charge-cylinder-section.png}}
\par\smallskip{\small 图示：无限长线电荷在三维中具有同轴圆柱等势面，横截面上的同心圆不能误认为二维点电荷模型}
\end{center}
\end{minipage}\par

点电荷的场按 \(1/r^2\) 衰减，线电荷的场按 \(1/\rho\)
衰减。二维图像相似并不意味着三维源和衰减律相同。

\subsection{六、电场线与等势线的画法和检查}\label{m2:ux516dux7535ux573aux7ebfux4e0eux7b49ux52bfux7ebfux7684ux753bux6cd5ux548cux68c0ux67e5}

\par\noindent\begin{minipage}{\linewidth}
\begin{center}
\includegraphics[width=\linewidth,height=.40\textheight,keepaspectratio]{\detokenize{assets/module-02/original/tikz-m2-topic9-v1-10-solution-workflow-pdf-v11.jpg}}
\par\smallskip{\small 图示：求场线和等势线应依次确定维数与奇点、求场或势、写曲线方程、定常数并检查特殊点}
\end{center}
\end{minipage}\par

\begin{enumerate}
\def\labelenumi{\arabic{enumi}.}
\tightlist
\item
  \textbf{确定空间维数}：区分真正的二维问题、轴对称子午面和三维截面。
\item
  \textbf{标出奇点与临界点}：点线面电荷处场可能奇异，\(\mathbf E=\mathbf0\)
  处方向未定义。
\item
  \textbf{选择入口}：已知 \(\mathbf E\) 时写积分曲线方程；已知
  \(\varphi\) 时先画水平集，再取正交轨线。
\item
  \textbf{利用对称性}：对称轴、对称面和反对称面常直接给出特殊场线或等势线。
\item
  \textbf{选择积分工具}：普通微分方程难积时，可寻找等通量管、轴对称通量函数或平面辅助函数。
\item
  \textbf{确定积分常数}：用给定初始点、所含通量、边界电势或场线端点确定。
\item
  \textbf{检查方向}：场线必须由高电势指向低电势，并满足正源、负源与无穷远之间的通量平衡。
\item
  \textbf{检查正交与极限}：正交性只在正则点检验；还要考察源点、零场点和无穷远行为。
\end{enumerate}

常见错误包括：

\begin{itemize}
\tightlist
\item
  把二维截面中的圆直接称为三维等势面。
\item
  在圆柱或球坐标场线方程中漏掉尺度因子。
\item
  在 \(\mathbf E=\mathbf0\) 处仍要求唯一场线方向。
\item
  用任意绘制的线条疏密代替定量的 \(|\mathbf E|\)。
\item
  把电场线当作带电粒子的实际运动轨迹。
\item
  让同一无源电场管的通量常数跨过电荷源而不发生变化。
\end{itemize}

\subsection{七、练习与解答}\label{m2:ux4e03ux7ec3ux4e60ux4e0eux89e3ux7b54}

\subsubsection{作业一：鞍形静电场的场线与等势线}\label{m2:ux4f5cux4e1aux4e00ux978dux5f62ux9759ux7535ux573aux7684ux573aux7ebfux4e0eux7b49ux52bfux7ebf}

\textbf{题号：M2-013}

\textbf{层级：} 基础　　\textbf{方法：} 静电场方程

在真空区域内给定
\(\mathbf E(x,y)=a(x\mathbf e_x-y\mathbf e_y)\)，\(a>0\)，场与 \(z\)
无关。先检验静电无旋条件和普通无源条件，再求含任意加常数的电势。求
\(xy\)
截面中的场线和等势线，明确其取向，并分析原点零场附近两族曲线的鞍形结构。

\par\noindent\begin{minipage}{\linewidth}
\textbf{解答}

\begin{center}
\includegraphics[width=\linewidth,height=.40\textheight,keepaspectratio]{\detokenize{assets/module-02/examples/M2-013.pdf}}
\par\smallskip{\small 图 M2-013：鞍形静电场的场线与等势线}
\end{center}
\end{minipage}\par

图 M2-013 的场箭头沿 \(x\) 远离原点、沿 \(y\)
指向原点，反映两个分量符号相反。先核对散度与旋度，再分别由切向方程和负梯度积分求场线与等势线；两个曲线族在普通点正交，在原点共同退化。

先计算散度与旋度：

\[
\nabla\cdot\mathbf E
=
\frac{\partial(ax)}{\partial x}
+\frac{\partial(-ay)}{\partial y}
=a-a=0,
\]

\[
(\nabla\times\mathbf E)_z
=
\frac{\partial E_y}{\partial x}
-\frac{\partial E_x}{\partial y}
=0.
\]

因此它可作为某无源区域内的静电场。

由

\[
E_x=-\frac{\partial\varphi}{\partial x}=ax
\]

得

\[
\varphi=-\frac a2x^2+f(y).
\]

再由

\[
E_y=-\frac{\partial\varphi}{\partial y}=-ay
\]

得到

\[
f'(y)=ay,
\qquad
f(y)=\frac a2y^2+C_0.
\]

所以

\[
\boxed{
\varphi(x,y)
=
\frac a2(y^2-x^2)+C_0
}.
\]

场线满足

\[
\frac{dy}{dx}
=
\frac{-ay}{ax}
=-\frac yx.
\]

分离变量：

\[
\frac{dy}{y}=-\frac{dx}{x},
\]

故

\[
\ln|xy|=\text{常量},
\]

即

\[
\boxed{xy=C_1}.
\]

等势线满足

\[
\boxed{y^2-x^2=C_2}.
\]

两族双曲线互相正交。\(x\) 轴和 \(y\) 轴分别是 \(xy=0\)
的两组分界线。原点处

\[
\mathbf E(0,0)=\mathbf0,
\]

电势既不是普通极大值也不是普通极小值，而是鞍点：沿 \(x\) 轴电势下降，沿
\(y\) 轴电势上升。因此不能给原点指定唯一场线方向。

\subsubsection{作业二：等量异号点电荷的电场线}\label{m2:ux4f5cux4e1aux4e8cux7b49ux91cfux5f02ux53f7ux70b9ux7535ux8377ux7684ux7535ux573aux7ebf}

\textbf{题号：M2-014}

\textbf{层级：} 基础　　\textbf{方法：} 静电场方程

真空中两个固定点电荷 \(+Q\) 与 \(-Q\) 分别位于
\((-l,0,0)\)、\((l,0,0)\)，其中 \(Q>0,l>0\)，没有其他外场。系统绕 \(x\)
轴旋转对称。取包含 \(x\)
轴的任意子午面，用场点到两电荷的距离或相应极角表示电场线方程；写明沿场线的方向，并区分一般场线与轴线上的特殊分支。

\par\noindent\begin{minipage}{\linewidth}
\textbf{解答}

\begin{center}
\includegraphics[width=\linewidth,height=.40\textheight,keepaspectratio]{\detokenize{assets/module-02/examples/M2-014.pdf}}
\par\smallskip{\small 图 M2-014：等量异号点电荷的电场线}
\end{center}
\end{minipage}\par

图 M2-014
的正电荷在左、负电荷在右，场线沿从正到负的方向绕轴对称展开。对子午面采用两源的有向通量函数相加，令其保持常量；积分时保留两电荷相反的符号，不能把两个局部极角直接当作同一角。

在子午面中用 \((x,\rho)\) 表示点 \(P\)，其中 \(\rho\) 是到 \(x\)
轴的距离。记

\[
r_+
=
\sqrt{(x+l)^2+\rho^2},
\qquad
r_-
=
\sqrt{(x-l)^2+\rho^2}.
\]

令 \(\theta_+\) 为从 \(+Q\) 指向 \(P\) 的连线与 \(+x\)
轴的夹角，\(\theta_-\) 为从 \(-Q\) 指向 \(P\) 的连线与 \(-x\)
轴的夹角，则

\[
\cos\theta_+
=
\frac{x+l}{r_+},
\qquad
\cos\theta_-
=
\frac{l-x}{r_-}.
\]

\par\noindent\begin{minipage}{\linewidth}
\begin{center}
\includegraphics[width=\linewidth,height=.40\textheight,keepaspectratio]{\detokenize{assets/module-02/original/tikz-m2-topic9-v1-06-point-dipole-flux-function-pdf-v11.jpg}}
\par\smallskip{\small 图示：等量异号点电荷的子午面场线绕对称轴生成电场管，垂直轴线的圆盘通量可作为场线标签}
\end{center}
\end{minipage}\par

取过 \(P\) 且垂直于 \(x\) 轴的圆盘，半径为 \(\rho\)，法向取 \(+x\)
方向。对位于轴上 \(x=a\) 的点电荷 \(q\)，该圆盘的通量为

\[
\Phi_q(x,\rho)
=
\frac{q}{2\varepsilon_0}
\left[
\operatorname{sgn}(x-a)
-
\frac{x-a}{\sqrt{(x-a)^2+\rho^2}}
\right].
\]

在两电荷之间 \(-l<x<l\)，\(+Q\) 的贡献为

\[
\Phi_+
=
\frac{Q}{2\varepsilon_0}
(1-\cos\theta_+),
\]

\(-Q\) 的贡献为

\[
\Phi_-
=
\frac{Q}{2\varepsilon_0}
(1-\cos\theta_-).
\]

因此总通量

\[
\Phi
=
\frac{Q}{2\varepsilon_0}
(2-\cos\theta_+-\cos\theta_-).
\]

当点 \(P\) 沿同一条子午面场线移动时，这条场线绕 \(x\)
轴生成电场管侧面，侧面通量为零；只要不跨过电荷，圆盘通量保持不变。所以

\[
\cos\theta_++\cos\theta_-=C.
\]

写成坐标形式：

\[
\boxed{
\frac{x+l}{\sqrt{(x+l)^2+\rho^2}}
+
\frac{l-x}{\sqrt{(x-l)^2+\rho^2}}
=C
}.
\]

在任意过 \(x\) 轴的平面内，可把 \(\rho\) 换成 \(|y|\)。两电荷之间的轴线
\(\rho=0\) 对应 \(C=2\)；离轴越远，\(C\) 越小。

若截面跨过点电荷，\(\operatorname{sgn}(x-a)\) 项发生变化，通量按
\(q/\varepsilon_0\)
跃变。因此``通量常数''只在每一段无源电场管内保持不变。

\subsubsection{作业三：等量异号平行线电荷的场线与等势线}\label{m2:ux4f5cux4e1aux4e09ux7b49ux91cfux5f02ux53f7ux5e73ux884cux7ebfux7535ux8377ux7684ux573aux7ebfux4e0eux7b49ux52bfux7ebf}

\textbf{题号：M2-015}

\textbf{层级：} 基础　　\textbf{方法：} 静电场方程

真空中两条无限长直线电荷平行于 \(z\) 轴，在横截面 \(xy\)
平面的位置分别为 \((-l,0)\)、\((l,0)\)，线密度依次为
\(+\lambda,-\lambda\)，其中 \(l>0,\lambda>0\)。令 \(r_+,r_-\)
分别为场点到两条线的垂直距离，并允许选择电势加常数。求横截面中的电势、等势线及电场线，说明两类曲线的几何形状和场向。

\par\noindent\begin{minipage}{\linewidth}
\textbf{解答}

\begin{center}
\includegraphics[width=\linewidth,height=.40\textheight,keepaspectratio]{\detokenize{assets/module-02/examples/M2-015.pdf}}
\par\smallskip{\small 图 M2-015：等量异号平行线电荷的场线与等势线}
\end{center}
\end{minipage}\par

图 M2-015 的两条线电荷在截面中表现为两个点，源强分别为
\(+\lambda,-\lambda\)。电势由两段距离的对数差给出，等势条件固定距离比，场线则由调和共轭或正交性确定；二维线源不能直接套三维点源的通量函数。

记到两条导线的距离为

\[
r_+
=
\sqrt{(x+l)^2+y^2},
\qquad
r_-
=
\sqrt{(x-l)^2+y^2}.
\]

单根无限线电荷的电势只确定到常数，但异号线电荷叠加后，公共参考常数相消。可写成

\[
\boxed{
\varphi(x,y)
=
\frac{\lambda}{2\pi\varepsilon_0}
\ln\frac{r_-}{r_+}
}.
\]

令 \(\varphi\) 为常量，得到

\[
\frac{r_-}{r_+}=K,
\qquad K>0.
\]

平方并整理：

\[
(x-l)^2+y^2
=
K^2[(x+l)^2+y^2].
\]

当 \(K\neq1\) 时，

\[
\left[
x-\frac{l(1+K^2)}{1-K^2}
\right]^2+y^2
=
\left[
\frac{2lK}{|1-K^2|}
\right]^2.
\]

所以等势线是阿波罗尼斯圆，其圆心和半径分别为

\[
x_K=\frac{l(1+K^2)}{1-K^2},
\qquad
R_K=\frac{2lK}{|1-K^2|}.
\]

当 \(K=1\) 时，

\[
x=0,
\]

即两导线的中垂线，也是零电势线。

下面求电场线。在不跨过两条线电荷的单连通区域内连续选取角度分支，并定义两条连线相对于
\(+x\) 轴的方向角

\[
\theta_+
=
\arg[(x+l)+iy],
\qquad
\theta_-
=
\arg[(x-l)+iy].
\]

利用

\[
\nabla\theta
=
\mathbf e_z\times\nabla\ln r,
\]

定义辅助函数

\[
\psi
=
\frac{\lambda}{2\pi\varepsilon_0}
(\theta_--\theta_+).
\]

则

\[
\nabla\psi
=
\mathbf e_z\times\nabla\varphi.
\]

因为 \(\mathbf E=-\nabla\varphi\)，

\[
\mathbf E\cdot\nabla\psi=0.
\]

所以 \(\psi=\text{常量}\) 的曲线与电场相切，即

\[
\theta_--\theta_+=\beta.
\]

这表示固定线段的两个端点对动点所张的角保持不变，其轨迹是通过两个源点的圆弧。该圆族可写为

\[
\boxed{
x^2+y^2-2by=l^2
},
\]

其中 \(b\) 是场线参数；\(x\) 轴是退化的极限场线。

\par\noindent\begin{minipage}{\linewidth}
\begin{center}
\includegraphics[width=\linewidth,height=.40\textheight,keepaspectratio]{\detokenize{assets/module-02/original/tikz-m2-topic9-v1-07-opposite-line-charge-circles-pdf-v11.jpg}}
\par\smallskip{\small 图示：异号平行线电荷的场线是通过两个源点的圆弧，等势线是距离比固定的阿波罗尼斯圆}
\end{center}
\end{minipage}\par

场线从 \(+\lambda\) 指向
\(-\lambda\)。两族圆在每个正则交点互相正交，这也由
\(\mathbf E=-\nabla\varphi\) 自动保证。

\subsubsection{作业四：无限长线电荷的场与参考电势}\label{m2:ux4f5cux4e1aux56dbux65e0ux9650ux957fux7ebfux7535ux8377ux7684ux573aux4e0eux53c2ux8003ux7535ux52bf}

\textbf{题号：M2-016}

\textbf{层级：} 复赛　　\textbf{方法：} 静电场方程

真空中只有一根沿 \(z\) 轴延伸的无限长均匀直线电荷，线密度为有符号常量
\(\lambda\)。采用柱坐标，\(\rho>0\) 表示距轴距离，\(\mathbf e_{\rho}\)
指向离轴方向。由柱对称性求电场；给定有限参考距离 \(\rho_0>0\) 及
\(\varphi(\rho_0)=0\)，求 \(\varphi(\rho)\)，并解释为何不能令
\(\varphi(\infty)=0\)。

\par\noindent\begin{minipage}{\linewidth}
\textbf{解答}

\begin{center}
\includegraphics[width=\linewidth,height=.40\textheight,keepaspectratio]{\detokenize{assets/module-02/examples/M2-016.pdf}}
\par\smallskip{\small 图 M2-016：无限长线电荷的场与参考电势}
\end{center}
\end{minipage}\par

图 M2-016 的同轴高斯柱侧面提供通量，端面法向与径向场垂直。先求
\(E_{\rho}\propto1/\rho\)，再从有限参考圆柱 \(\rho_0\)
积分；对数随距离无界，因而不能再要求无穷远处电势为有限零值。

取与导线同轴、半径为 \(\rho\)、长度为 \(L\) 的圆柱高斯面。由圆柱对称性，

\[
\mathbf E=E(\rho)\mathbf e_{\rho}.
\]

电场与上下底面平行，底面通量为零；侧面上电场大小处处相同且沿法向。于是

\[
E(\rho)(2\pi\rho L)
=
\frac{\lambda L}{\varepsilon_0}.
\]

因此

\[
\boxed{
\mathbf E(\rho)
=
\frac{\lambda}{2\pi\varepsilon_0\rho}\mathbf e_{\rho}
}.
\]

若规定 \(\varphi(\rho_0)=0\)，则

\[
\varphi(\rho)-\varphi(\rho_0)
=
-\int_{\rho_0}^{\rho}
\frac{\lambda}{2\pi\varepsilon_0\rho'}
\,d\rho'.
\]

积分得到

\[
\boxed{
\varphi(\rho)
=
-\frac{\lambda}{2\pi\varepsilon_0}
\ln\frac{\rho}{\rho_0}
}.
\]

\par\noindent\begin{minipage}{\linewidth}
\begin{center}
\includegraphics[width=\linewidth,height=.40\textheight,keepaspectratio]{\detokenize{assets/module-02/original/tikz-m2-topic9-v1-09-line-charge-potential-reference-safe-v7.png}}
\par\smallskip{\small 图示：无限长线电荷的电势只能相对有限半径参考圆柱面定义，正线电荷的电势随径向距离对数下降}
\end{center}
\end{minipage}\par

若试图令 \(\rho_0\to\infty\)，则

\[
\ln\frac{\rho}{\rho_0}\to-\infty,
\]

任意有限 \(\rho\)
处的电势都会发散。根本原因不是无穷远``靠近导线''，而是圆柱通量只在面积
\(2\pi\rho L\) 上扩散，使

\[
E(\rho)\propto\frac1\rho,
\]

从有限半径积分到无穷远时产生对数发散。因而只有电势差或相对于有限参考圆柱面的电势有意义。

\section{专题十：电偶极子（基础必讲）}\label{m2:ux4e13ux9898ux5341ux7535ux5076ux6781ux5b50ux57faux7840ux5fc5ux8bb2}

\begin{quote}
阅读准备：远场展开先找无量纲小参数 当源区尺度为 \(\ell\)、观察距离为
\(r\) 且 \(\ell/r\ll1\) 时，才可按 \(\ell/r\)
作泰勒展开。保留到一阶意味着误差通常从 \((\ell/r)^2\)
开始；对称性可能使某些低阶项自动抵消。

\textbf{例 M2-017：远场展开中两项单极贡献抵消}
\end{quote}

设 \(a>0\)，\(r>a\)，考察远区条件 \(a/r\ll1\) 下的差
\(1/(r-a)-1/(r+a)\)。将两项分别按 \(a/r\)
展开，写出相减后的最低非零阶，说明各自的 \(r^{-1}\)
首项为什么抵消，并判断差式相对于单个单极势的衰减阶次。

\par\noindent\begin{minipage}{\linewidth}
\textbf{解答}

\begin{center}
\includegraphics[width=\linewidth,height=.40\textheight,keepaspectratio]{\detokenize{assets/module-02/examples/M2-017.pdf}}
\par\smallskip{\small 图 M2-017：远场展开中两项单极贡献抵消}
\end{center}
\end{minipage}\par

图 M2-017 的有限电荷对沿轴线分布，远处场点满足
\(r\gg a\)。两段距离分别为 \(r-a\) 与 \(r+a\)，一阶展开中的 \(1/r\)
抵消，而 \(a/r^2\) 相加，留下偶极阶；近场仍应使用原来的精确距离。

\(1/(r\pm a)=r^{-1}(1\pm a/r)^{-1}\approx r^{-1}(1\mp a/r)\)。两项相减时单极项抵消，留下与
\(a/r^2\) 同阶的偶极项。

电偶极子描述总电荷为零而偶极矩非零的电荷分布。有限电荷对是实际源；在
\(r\gg\ell\) 时可作点偶极远场近似；令 \(\ell\to0\)、\(q\to\infty\)
而保持 \(q\ell\)
不变，才得到理想点偶极。三种模型的适用范围、误差阶数和奇点结构不同。

\par\noindent\begin{minipage}{\linewidth}
\begin{center}
\includegraphics[width=\linewidth,height=.40\textheight,keepaspectratio]{\detokenize{assets/module-02/original/tikz-m2-topic10-v2-01-model-levels.png}}
\par\smallskip{\small 图示：有限电荷对、远场点偶极近似与理想点偶极的三个模型层次}
\end{center}
\end{minipage}\par

\subsection{一、有限电荷对与偶极矩}\label{m2:ux4e00ux6709ux9650ux7535ux8377ux5bf9ux4e0eux5076ux6781ux77e9}

\subsubsection{1.
有限电偶极子}\label{m2:ux6709ux9650ux7535ux5076ux6781ux5b50}

两个等量异号点电荷 \(+q\) 与 \(-q\) 相距
\(\ell\)，便构成一个\textbf{有限电偶极子}。取两电荷中点为原点，并记

\[
\mathbf r_+=\frac{\boldsymbol{\ell}}2,
\qquad
\mathbf r_-=-\frac{\boldsymbol{\ell}}2,
\]

其中 \(\boldsymbol{\ell}\) 从 \(-q\) 指向 \(+q\)。偶极矩定义为

\[
\boxed{\mathbf p=q\boldsymbol{\ell}}.
\]

\(\mathbf p\) 的方向始终从负电荷指向正电荷，量纲为
\(\mathrm{C\cdot m}\)。

\begin{quote}
\(r\gg\ell\)
不是``成为偶极子''的条件，而是把有限电荷对近似为点偶极的条件。
\end{quote}

\subsubsection{2.
中性系统的偶极矩与原点无关}\label{m2:ux4e2dux6027ux7cfbux7edfux7684ux5076ux6781ux77e9ux4e0eux539fux70b9ux65e0ux5173}

对离散电荷系，一般定义

\[
\mathbf p=\sum_i q_i\mathbf r_i.
\]

若把坐标原点平移 \(\mathbf a\)，则
\(\mathbf r_i'=\mathbf r_i-\mathbf a\)，从而

\[
\mathbf p'
=\sum_iq_i(\mathbf r_i-\mathbf a)
=\mathbf p-\mathbf a\sum_iq_i.
\]

因此，当系统总电荷 \(\sum_iq_i=0\)
时，\(\mathbf p'=\mathbf p\)；若总电荷不为零，偶极矩一般依赖原点。

\textbf{例 M2-018：中性电荷对的偶极矩不随原点平移}

在真空中，点电荷 \(+q\) 位于 \(z=a\)，\(-q\) 位于 \(z=-a\)，其中
\(q>0,a>0\)。先以两电荷中点为原点求系统的总电荷与电偶极矩。再把原点平移任意常矢量
\(\mathbf b\)，令新位置坐标为
\(\mathbf r_i'=\mathbf r_i-\mathbf b\)，重新计算偶极矩，并说明中性条件怎样影响平移结果。

\par\noindent\begin{minipage}{\linewidth}
\textbf{解答}

\begin{center}
\includegraphics[width=\linewidth,height=.40\textheight,keepaspectratio]{\detokenize{assets/module-02/examples/M2-018.pdf}}
\par\smallskip{\small 图 M2-018：中性电荷对的偶极矩不随原点平移}
\end{center}
\end{minipage}\par

图 M2-018 中偶极矩由负电荷指向正电荷。两项位置加权后均沿 \(+z\)，得到
\(2qa\mathbf e_z\)；平移原点带来的改变量为平移矢量乘总电荷，而此处总电荷为零，所以偶极矩不变。

\(+q\) 位于 \(z=a\)，\(-q\) 位于 \(z=-a\)，则

\[
\mathbf p=q(a\mathbf e_z)-q(-a\mathbf e_z)=2qa\mathbf e_z.
\]

无论把原点沿任意方向平移多少，只要同时平移两电荷的坐标，结果仍为
\(2qa\mathbf e_z\)。

\subsubsection{3. 理想点偶极}\label{m2:ux7406ux60f3ux70b9ux5076ux6781}

令

\[
\ell\to0,
\qquad
q\to\infty,
\qquad
\mathbf p=q\boldsymbol{\ell}=\text{有限常矢量},
\]

所得集中在一点、只保留偶极矩信息的理想化源称为\textbf{理想点偶极}。理想点偶极在原点是奇异源；其普通函数形式的电势和电场只在
\(r>0\) 处使用。

\subsection{二、精确电势与远场展开}\label{m2:ux4e8cux7cbeux786eux7535ux52bfux4e0eux8fdcux573aux5c55ux5f00}

记

\[
k=\frac1{4\pi\varepsilon_0},
\qquad
r_+=\left|\mathbf r-\frac{\boldsymbol{\ell}}2\right|,
\qquad
r_-=\left|\mathbf r+\frac{\boldsymbol{\ell}}2\right|.
\]

有限电荷对在场点 \(\mathbf r\) 处的\textbf{精确电势}为

\[
\boxed{
\varphi_{\mathrm{\text{精确}}}(\mathbf r)
=kq\left(\frac1{r_+}-\frac1{r_-}\right)
}.
\]

相应的精确电场为

\[
\boxed{
\mathbf E_{\mathrm{\text{精确}}}(\mathbf r)
=kq\left[
\frac{\mathbf r-\boldsymbol{\ell}/2}{|\mathbf r-\boldsymbol{\ell}/2|^3}
-\frac{\mathbf r+\boldsymbol{\ell}/2}{|\mathbf r+\boldsymbol{\ell}/2|^3}
\right]
}.
\]

这两式除电荷所在点外均成立，并不要求 \(r\gg\ell\)。

\par\noindent\begin{minipage}{\linewidth}
\begin{center}
\includegraphics[width=\linewidth,height=.40\textheight,keepaspectratio]{\detokenize{assets/module-02/original/tikz-m2-topic10-v2-02-far-expansion-pdf-v11.jpg}}
\par\smallskip{\small 图示：在电荷对中点作对称泰勒展开得到远场偶极电势}
\end{center}
\end{minipage}\par

\subsubsection{1.
对称泰勒展开}\label{m2:ux5bf9ux79f0ux6cf0ux52d2ux5c55ux5f00}

当 \(r\gg\ell\) 时，以中点为展开中心：

\[
\frac1{|\mathbf r\mp\boldsymbol{\ell}/2|}
=\frac1r
\pm\frac{\boldsymbol{\ell}\cdot\hat{\mathbf r}}{2r^2}
+\frac{3(\boldsymbol{\ell}\cdot\hat{\mathbf r})^2-\ell^2}{8r^3}
+O\!\left(\frac{\ell^3}{r^4}\right).
\]

正、负电荷的零阶项相消；由于电荷分布关于中点反对称，二阶项也相消。于是

\[
\frac1{r_+}-\frac1{r_-}
=\frac{\boldsymbol{\ell}\cdot\hat{\mathbf r}}{r^2}
+O\!\left(\frac{\ell^3}{r^4}\right),
\]

从而

\[
\boxed{
\varphi(\mathbf r)
=k\frac{\mathbf p\cdot\hat{\mathbf r}}{r^2}
=k\frac{\mathbf p\cdot\mathbf r}{r^3}
}.
\]

若首项不因方向而恰好为零，则相对误差为

\[
O\!\left(\frac{\ell^2}{r^2}\right).
\]

在赤道面
\(\mathbf p\cdot\hat{\mathbf r}=0\)，偶极首项本身为零，不能再用``相对误差''衡量；此时应回到更高阶项或精确式判断。

\subsubsection{2.
电势的角分布}\label{m2:ux7535ux52bfux7684ux89d2ux5206ux5e03}

取极轴沿 \(\mathbf p\)，令 \(\theta\) 为 \(\mathbf p\) 与 \(\mathbf r\)
的夹角，则

\[
\boxed{
\varphi(r,\theta)=k\frac{p\cos\theta}{r^2}
}.
\]

\par\noindent\begin{minipage}{\linewidth}
\begin{center}
\includegraphics[width=\linewidth,height=.40\textheight,keepaspectratio]{\detokenize{assets/module-02/original/tikz-m2-topic10-v3-03-potential-sign.png}}
\par\smallskip{\small 图示：固定半径上偶极电势随角度的正负和余弦变化}
\end{center}
\end{minipage}\par

因此：

\begin{itemize}
\tightlist
\item
  正轴方向 \(\theta=0\)，电势为正且绝对值最大；
\item
  赤道面 \(\theta=\pi/2\)，电势为零；
\item
  负轴方向 \(\theta=\pi\)，电势为负且绝对值最大。
\end{itemize}

\textbf{例 M2-019：偶极势在赤道两侧等大反号}

理想电偶极子位于真空原点，其固定偶极矩为
\(\mathbf p=p\mathbf e_z\)，\(p>0\)，取无穷远电势为零。采用以 \(+z\)
为极轴的球坐标，在同一半径 \(r>0\) 上取两个场点，极角分别为
\(60^{\circ}\) 与
\(120^{\circ}\)。比较两点的电势大小及符号，并解释它们关于赤道面的对应关系。

\par\noindent\begin{minipage}{\linewidth}
\textbf{解答}

\begin{center}
\includegraphics[width=\linewidth,height=.40\textheight,keepaspectratio]{\detokenize{assets/module-02/examples/M2-019.pdf}}
\par\smallskip{\small 图 M2-019：偶极势在赤道两侧等大反号}
\end{center}
\end{minipage}\par

图 M2-019 的两点在同一圆周上，分别位于赤道面上下。极角从 \(+z\)
轴量起，两个余弦分别为
\(1/2,-1/2\)；径向系数相同，所以电势等大反号，而不是电场矢量也简单逐项反号。

在同一半径 \(r\) 上，\(\theta=60^{\circ}\) 与 \(120^{\circ}\)
两点的电势分别为

\[
\varphi_{60^{\circ}}=\frac{kp}{2r^2},
\qquad
\varphi_{120^{\circ}}=-\frac{kp}{2r^2}.
\]

两点电势等大反号，反映了偶极势关于赤道面的反对称性。

\subsection{三、偶极电场的分量推导}\label{m2:ux4e09ux5076ux6781ux7535ux573aux7684ux5206ux91cfux63a8ux5bfc}

\subsubsection{1.
坐标无关的矢量式}\label{m2:ux5750ux6807ux65e0ux5173ux7684ux77e2ux91cfux5f0f}

由

\[
\varphi=k(\mathbf p\cdot\mathbf r)r^{-3}
\]

以及 \(\mathbf p\) 为常矢量，有

\[
\nabla(\mathbf p\cdot\mathbf r)=\mathbf p,
\qquad
\nabla r^{-3}=-3r^{-5}\mathbf r.
\]

使用乘积法则：

\[
\begin{aligned}
\nabla\varphi
&=k\left[
\mathbf p\,r^{-3}
-3(\mathbf p\cdot\mathbf r)\mathbf r\,r^{-5}
\right],\\
\mathbf E
&=-\nabla\varphi.
\end{aligned}
\]

所以，在 \(r>0\) 处

\[
\boxed{
\mathbf E(\mathbf r)
=\frac1{4\pi\varepsilon_0r^3}
\left[3(\mathbf p\cdot\hat{\mathbf r})\hat{\mathbf r}-\mathbf p\right]
}.
\]

\subsubsection{2. 球坐标分量}\label{m2:ux7403ux5750ux6807ux5206ux91cf}

当极轴沿 \(\mathbf p\) 时，电势与方位角无关。球坐标梯度给出

\[
\mathbf E
=-\frac{\partial\varphi}{\partial r}\mathbf e_r
-\frac1r\frac{\partial\varphi}{\partial\theta}\mathbf e_{\theta}.
\]

逐项求导：

\[
\begin{aligned}
E_r
&=-\frac{\partial}{\partial r}
\left(k\frac{p\cos\theta}{r^2}\right)
=\frac{2kp\cos\theta}{r^3},\\
E_{\theta}
&=-\frac1r\frac{\partial}{\partial\theta}
\left(k\frac{p\cos\theta}{r^2}\right)
=\frac{kp\sin\theta}{r^3},\\
E_{\phi}&=0.
\end{aligned}
\]

故

\[
\boxed{
E_r=\frac{2p\cos\theta}{4\pi\varepsilon_0r^3},
\qquad
E_{\theta}=\frac{p\sin\theta}{4\pi\varepsilon_0r^3}
}.
\]

场强大小为

\[
\boxed{
|\mathbf E|
=\frac{p}{4\pi\varepsilon_0r^3}
\sqrt{1+3\cos^2\theta}
}.
\]

\par\noindent\begin{minipage}{\linewidth}
\begin{center}
\includegraphics[width=\linewidth,height=.40\textheight,keepaspectratio]{\detokenize{assets/module-02/original/tikz-m2-topic10-v2-04-field-components-margin12-v11.png}}
\par\smallskip{\small 图示：偶极电场的径向分量、角向分量及坐标无关矢量式}
\end{center}
\end{minipage}\par

电势按 \(r^{-2}\) 衰减，电场按 \(r^{-3}\) 衰减。这里的 \(r^{-3}\)
是点偶极的远场规律，不适用于紧邻有限电荷对中某个点电荷的区域。

\subsubsection{3.
轴线与赤道面检验}\label{m2:ux8f74ux7ebfux4e0eux8d64ux9053ux9762ux68c0ux9a8c}

\par\noindent\begin{minipage}{\linewidth}
\begin{center}
\includegraphics[width=\linewidth,height=.40\textheight,keepaspectratio]{\detokenize{assets/module-02/original/tikz-m2-topic10-v2-05-axis-equator-pdf-v11.jpg}}
\par\smallskip{\small 图示：偶极轴线和赤道面上的电场方向与大小}
\end{center}
\end{minipage}\par

在两侧轴线上，电场都沿 \(\mathbf p\)：

\[
\boxed{
\mathbf E_{\mathrm{\text{轴}}}
=\frac{2p}{4\pi\varepsilon_0r^3}\hat{\mathbf p}
}.
\]

在赤道面上，电场与 \(\mathbf p\) 反向：

\[
\boxed{
\mathbf E_{\mathrm{\text{赤道}}}
=-\frac{p}{4\pi\varepsilon_0r^3}\hat{\mathbf p}
}.
\]

同一半径处，轴线场强与赤道面场强之比为 \(2:1\)。

\begin{quote}
赤道面上 \(\varphi=0\)，但
\(\mathbf E\neq\mathbf0\)。电势为零只说明该点的电势数值为零，并不说明电势的空间梯度为零。
\end{quote}

\subsection{四、场线、通量函数与等势面}\label{m2:ux56dbux573aux7ebfux901aux91cfux51fdux6570ux4e0eux7b49ux52bfux9762}

\subsubsection{1.
场线微分方程}\label{m2:ux573aux7ebfux5faeux5206ux65b9ux7a0b}

在子午面内，场线切向微元为

\[
d\mathbf s=\mathbf e_r\,dr+\mathbf e_{\theta}\,r\,d\theta.
\]

场线上 \(d\mathbf s\parallel\mathbf E\)，所以

\[
\frac{dr}{r\,d\theta}
=\frac{E_r}{E_{\theta}}
=2\cot\theta.
\]

分离变量并积分：

\[
\frac{dr}{r}=2\cot\theta\,d\theta,
\]

\[
\ln r=2\ln(\sin\theta)+\text{常数}.
\]

因此场线族为

\[
\boxed{r=C\sin^2\theta},
\qquad C>0.
\]

\subsubsection{2.
极冠通量函数}\label{m2:ux6781ux51a0ux901aux91cfux51fdux6570}

取半径为 \(r\) 的球面，从正轴到极角 \(\theta\) 的球面极冠通量为

\[
\begin{aligned}
\Psi(r,\theta)
&=\int_0^{2\pi}\int_0^{\theta}
E_r(r,\theta')r^2\sin\theta'\,d\theta'\,d\phi\\
&=\frac{p}{2\varepsilon_0r}\sin^2\theta.
\end{aligned}
\]

轴对称场中，同一条场线旋转所得通量管的 \(\Psi\) 保持不变，因此

\[
\frac{\sin^2\theta}{r}=\text{常数},
\]

再次得到 \(r=C\sin^2\theta\)。

\par\noindent\begin{minipage}{\linewidth}
\begin{center}
\includegraphics[width=\linewidth,height=.40\textheight,keepaspectratio]{\detokenize{assets/module-02/original/tikz-m2-topic10-v3-06-field-lines-flux.png}}
\par\smallskip{\small 图示：偶极场线方程与球面极冠通量函数的对应关系}
\end{center}
\end{minipage}\par

\textbf{例 M2-020：同一偶极场线上半径随极角变化}

真空原点有理想偶极子 \(\mathbf p=p\mathbf e_z\)，\(p>0\)。在以 \(+z\)
为极轴的子午面内，一条非轴线场线与赤道面 \(\theta=\pi/2\) 相交于
\(r=r_0>0\)。利用偶极场线关系确定同一分支在 \(\theta=30^{\circ}\)
处的半径，说明积分常数如何由赤道交点确定。

\par\noindent\begin{minipage}{\linewidth}
\textbf{解答}

\begin{center}
\includegraphics[width=\linewidth,height=.40\textheight,keepaspectratio]{\detokenize{assets/module-02/examples/M2-020.pdf}}
\par\smallskip{\small 图 M2-020：同一偶极场线上半径随极角变化}
\end{center}
\end{minipage}\par

图 M2-020 的赤道交点固定积分常数 \(C=r_0\)。从 \(+z\)
轴量到另一个标记点的极角为 \(30^{\circ}\)，代入同一条场线的
\(\sin^2\theta\)，得到径向长度缩为 \(r_0/4\)，不是横坐标缩为四分之一。

某条场线在赤道面与 \(r=r_0\) 相交。因为 \(\theta=\pi/2\) 时
\(C=r_0\)，故在 \(\theta=30^{\circ}\) 处

\[
r=r_0\sin^230^{\circ}=\frac{r_0}{4}.
\]

\subsubsection{3.
等势面与场线正交}\label{m2:ux7b49ux52bfux9762ux4e0eux573aux7ebfux6b63ux4ea4}

非零等势面满足

\[
\varphi_0=k\frac{p\cos\theta}{r^2},
\]

即

\[
r^2=\frac{kp}{\varphi_0}\cos\theta,
\]

其中右端必须为正。由于
\(\mathbf E=-\nabla\varphi\)，电场线在每一点都与等势面正交。

理想点偶极场线在图中汇聚到原点，是因为原点被压缩成奇点；有限电荷对的真实场线从
\(+q\) 出发并终止于
\(-q\)。场线只表示场的方向，不应无条件理解为带电粒子的运动轨迹。

\subsection{五、外电场中的永久偶极}\label{m2:ux4e94ux5916ux7535ux573aux4e2dux7684ux6c38ux4e45ux5076ux6781}

以下各式中的 \(\mathbf E\)
都表示\textbf{外电场}，不包括偶极自身产生的电场。

\subsubsection{1.
匀强外场中的合力与力矩}\label{m2:ux5300ux5f3aux5916ux573aux4e2dux7684ux5408ux529bux4e0eux529bux77e9}

在匀强外场中，两端受力分别为

\[
\mathbf F_+=q\mathbf E,
\qquad
\mathbf F_-=-q\mathbf E,
\]

故合力为零。关于偶极中心的合力矩为

\[
\begin{aligned}
\boldsymbol\tau
&=\frac{\boldsymbol{\ell}}2\times q\mathbf E
+\left(-\frac{\boldsymbol{\ell}}2\right)\times(-q\mathbf E)\\
&=q\boldsymbol{\ell}\times\mathbf E.
\end{aligned}
\]

因此

\[
\boxed{
\mathbf F_{\mathrm{\text{合}}}=\mathbf0,
\qquad
\boldsymbol\tau=\mathbf p\times\mathbf E
}.
\]

若 \(\beta\) 为 \(\mathbf p\) 与 \(\mathbf E\) 的夹角，则力矩大小为

\[
\tau=pE\sin\beta.
\]

相应的取向势能为

\[
\boxed{U=-\mathbf p\cdot\mathbf E=-pE\cos\beta}.
\]

\par\noindent\begin{minipage}{\linewidth}
\begin{center}
\includegraphics[width=\linewidth,height=.40\textheight,keepaspectratio]{\detokenize{assets/module-02/original/tikz-m2-topic10-v2-07-uniform-field.png}}
\par\smallskip{\small 图示：永久偶极在匀强外场中的力偶、力矩、势能和稳定性}
\end{center}
\end{minipage}\par

\(\beta=0\) 是稳定平衡，\(\beta=\pi\) 是不稳定平衡。

\textbf{例 M2-021：偶极转矩形成小角度恢复}

一个刚性永久电偶极子在匀强静电场 \(\mathbf E\)
中绕其固定中心转动，偶极矩大小为 \(p>0\)，外场大小为
\(E>0\)，绕垂直于转动平面的轴的转动惯量为
\(I>0\)。忽略摩擦及其他力矩。以偶极矩与外场夹角描述运动，确定稳定平衡位置，并求该位置附近小角度振动的角频率。

\par\noindent\begin{minipage}{\linewidth}
\textbf{解答}

\begin{center}
\includegraphics[width=\linewidth,height=.40\textheight,keepaspectratio]{\detokenize{assets/module-02/examples/M2-021.pdf}}
\par\smallskip{\small 图 M2-021：偶极转矩形成小角度恢复}
\end{center}
\end{minipage}\par

图 M2-021 的偶极矩相对水平电场偏转
\(\beta\)，两端电力形成使夹角减小的转矩。角运动方程中的负号由这一恢复方向决定，小角度取
\(\sin\beta\simeq\beta\) 后，与简谐运动标准式比较读出角频率。

若刚性偶极绕中心转动的转动惯量为 \(I\)，在 \(\beta\ll1\) 时

\[
I\ddot\beta=-pE\sin\beta\approx-pE\beta,
\]

所以小振动角频率为

\[
\omega=\sqrt{\frac{pE}{I}}.
\]

\subsubsection{2.
缓慢变化外场中的净力}\label{m2:ux7f13ux6162ux53d8ux5316ux5916ux573aux4e2dux7684ux51c0ux529b}

设外场在偶极尺度 \(\ell\) 上变化很小。两端合力为

\[
\mathbf F
=q\left[
\mathbf E\left(\mathbf r+\frac{\boldsymbol{\ell}}2\right)
-\mathbf E\left(\mathbf r-\frac{\boldsymbol{\ell}}2\right)
\right].
\]

在偶极中心作对称展开：

\[
\begin{aligned}
\mathbf E\left(\mathbf r\pm\frac{\boldsymbol{\ell}}2\right)
={}&\mathbf E(\mathbf r)
\pm\frac12(\boldsymbol{\ell}\cdot\nabla)\mathbf E
+\frac18(\boldsymbol{\ell}\cdot\nabla)^2\mathbf E\\
&\pm\frac1{48}(\boldsymbol{\ell}\cdot\nabla)^3\mathbf E
+\frac1{384}(\boldsymbol{\ell}\cdot\nabla)^4\mathbf E
+O\!\left(\ell^5\nabla^5\mathbf E\right).
\end{aligned}
\]

相减后偶次项抵消，因此

\[
\mathbf F
=(\mathbf p\cdot\nabla)\mathbf E
+\frac q{24}(\boldsymbol{\ell}\cdot\nabla)^3\mathbf E
+O\!\left(q\ell^5\nabla^5\mathbf E\right).
\]

当外场在尺度 \(\ell\)
上缓慢变化时，三阶方向导数项相对主项小两个尺度因子，最低阶结果才是

\[
\boxed{\mathbf F\simeq(\mathbf p\cdot\nabla)\mathbf E}.
\]

对固定的 \(\mathbf p\)，矢量恒等式给出

\[
\nabla(\mathbf p\cdot\mathbf E)
=(\mathbf p\cdot\nabla)\mathbf E
+\mathbf p\times(\nabla\times\mathbf E).
\]

若外场为静电场，\(\nabla\times\mathbf E=\mathbf0\)，并只保留上述最低阶项，则

\[
\boxed{
\mathbf F\simeq\nabla(\mathbf p\cdot\mathbf E)
=-\nabla U
}.
\]

\par\noindent\begin{minipage}{\linewidth}
\begin{center}
\includegraphics[width=\linewidth,height=.40\textheight,keepaspectratio]{\detokenize{assets/module-02/original/tikz-m2-topic10-v2-08-nonuniform-field.png}}
\par\smallskip{\small 图示：非匀强外场中偶极两端受力不完全抵消并产生净力}
\end{center}
\end{minipage}\par

\textbf{例 M2-022：非均匀场使两端电力不能抵消}

一个尺寸很小的刚性永久偶极子，其偶极矩固定为
\(\mathbf p=p\mathbf e_x\)，\(p>0\)。中心附近给定外场
\(\mathbf E(x)=(E_0+\gamma x)\mathbf e_x\)，\(E_0,\gamma\)
为常量；偶极尺寸远小于外场显著变化的长度尺度，忽略它对外场的反作用。将偶极子视为两端异号电荷，保留尺寸的一阶效应求净力，并说明非均匀场为何使两端电力不能抵消。

\par\noindent\begin{minipage}{\linewidth}
\textbf{解答}

\begin{center}
\includegraphics[width=\linewidth,height=.40\textheight,keepaspectratio]{\detokenize{assets/module-02/examples/M2-022.pdf}}
\par\smallskip{\small 图 M2-022：非均匀场使两端电力不能抵消}
\end{center}
\end{minipage}\par

图 M2-022 的右侧场箭头更长，示意
\(\gamma>0\)。有限偶极两端的反向电力不再等大，差值在小间距极限成为
\(p\partial_xE_x\)，故净力为 \(p\gamma\mathbf e_x\)；均匀部分 \(E_0\)
只产生彼此抵消的平动力。

若 \(\mathbf p=p\mathbf e_x\)，且局部外场为

\[
\mathbf E(x)=(E_0+\gamma x)\mathbf e_x,
\]

则

\[
\mathbf F=p\gamma\mathbf e_x.
\]

当 \(\gamma>0\) 时，永久偶极被拉向场较强的一侧。

\subsubsection{3.
永久偶极与感应偶极的能量差别}\label{m2:ux6c38ux4e45ux5076ux6781ux4e0eux611fux5e94ux5076ux6781ux7684ux80fdux91cfux5deeux522b}

对于与外场无关、大小固定的永久偶极，

\[
U_{\mathrm{\text{永久}}}=-\mathbf p\cdot\mathbf E.
\]

对于各向同性线性可极化粒子，若

\[
\mathbf p_{\mathrm{\text{感应}}}=\alpha\mathbf E,
\]

偶极矩是随外场从零逐步建立的，故

\[
U_{\mathrm{\text{感应}}}
=-\int_{\mathbf0}^{\mathbf E}\mathbf p(\mathbf E')\cdot d\mathbf E'
=-\frac12\alpha E^2.
\]

因此

\[
\boxed{
\mathbf F_{\mathrm{\text{感应}}}
=\frac12\alpha\nabla E^2
}.
\]

感应偶极问题中不能把 \(U=-\mathbf p\cdot\mathbf E\) 直接代入
\(\mathbf p=\alpha\mathbf E\)，否则会多出一倍。

\subsection{六、理想点偶极的分布源表示（拔高衔接）}\label{m2:ux516dux7406ux60f3ux70b9ux5076ux6781ux7684ux5206ux5e03ux6e90ux8868ux793aux62d4ux9ad8ux8854ux63a5}

有限电荷对的电荷密度可写为

\[
\rho(\mathbf r)
=q\delta^{(3)}\!\left(\mathbf r-\frac{\boldsymbol{\ell}}2\right)
-q\delta^{(3)}\!\left(\mathbf r+\frac{\boldsymbol{\ell}}2\right).
\]

在 \(\ell\to0\) 时展开：

\[
\delta^{(3)}\!\left(\mathbf r\mp\frac{\boldsymbol{\ell}}2\right)
=\delta^{(3)}(\mathbf r)
\mp\frac12\boldsymbol{\ell}\cdot\nabla\delta^{(3)}(\mathbf r)
+O(\ell^2).
\]

零阶项相消，得到理想点偶极的分布源

\[
\boxed{
\rho_{\mathbf p}(\mathbf r)
=-\mathbf p\cdot\nabla\delta^{(3)}(\mathbf r)
}.
\]

\par\noindent\begin{minipage}{\linewidth}
\begin{center}
\includegraphics[width=\linewidth,height=.40\textheight,keepaspectratio]{\detokenize{assets/module-02/original/tikz-m2-topic10-v2-09-distribution-source-pdf-v11.jpg}}
\par\smallskip{\small 图示：有限电荷对收缩为理想点偶极及其德尔塔分布导数表示}
\end{center}
\end{minipage}\par

把它代入电势积分，并按分布导数作分部积分：

\[
\begin{aligned}
\varphi(\mathbf r)
&=k\int
\frac{\rho_{\mathbf p}(\mathbf r')}{|\mathbf r-\mathbf r'|}\,dV'\\
&=k\frac{\mathbf p\cdot\mathbf r}{r^3},
\qquad r>0.
\end{aligned}
\]

若只讨论
\(r>0\)，前述偶极电场公式已经完整。若把电场看作定义在全空间的分布，则还需写成

\[
\mathbf E
=\operatorname{PV}\!\left\{
\frac{3(\mathbf p\cdot\hat{\mathbf r})\hat{\mathbf r}-\mathbf p}
{4\pi\varepsilon_0r^3}
\right\}
-\frac{\mathbf p}{3\varepsilon_0}\delta^{(3)}(\mathbf r),
\]

其中 \(\operatorname{PV}\)
表示主值。最后一项只在涉及原点的分布积分时出现；普通远场计算不要把它误加到
\(r>0\) 的场中。

\subsection{七、电偶极子问题的近似条件}\label{m2:ux4e03ux7535ux5076ux6781ux5b50ux95eeux9898ux7684ux8fd1ux4f3cux6761ux4ef6}

\par\noindent\begin{minipage}{\linewidth}
\begin{center}
\includegraphics[width=\linewidth,height=.40\textheight,keepaspectratio]{\detokenize{assets/module-02/original/tikz-m2-topic10-v2-10-workflow.png}}
\par\smallskip{\small 图示：电偶极问题中的模型、电势、场线和外场力学}
\end{center}
\end{minipage}\par

\begin{enumerate}
\def\labelenumi{\arabic{enumi}.}
\tightlist
\item
  \textbf{先判模型。} 有限电荷对本身就是偶极子；只有在 \(r\gg\ell\)
  时才能用点偶极远场式。
\item
  \textbf{源场优先求电势。} 正负电荷的标量电势先相减，再取
  \(\mathbf E=-\nabla\varphi\)，通常比直接叠加三维矢量场稳定。
\item
  \textbf{牢记偶极矩方向。} \(\mathbf p\) 从 \(-q\) 指向
  \(+q\)；轴线两侧的电场都沿 \(\mathbf p\)，赤道面电场沿
  \(-\mathbf p\)。
\item
  \textbf{不要把零电势当成零电场。} 赤道面上
  \(\varphi=0\)，但法向梯度不为零。
\item
  \textbf{检查衰减阶数。} 偶极势按 \(r^{-2}\) 衰减，偶极场按 \(r^{-3}\)
  衰减。
\item
  \textbf{区分自身场与外场。}
  \(\boldsymbol\tau=\mathbf p\times\mathbf E\)、\(U=-\mathbf p\cdot\mathbf E\)
  中的 \(\mathbf E\) 是外场。
\item
  \textbf{区分永久偶极与感应偶极。} 永久偶极能量没有
  \(1/2\)；线性感应偶极能量有 \(1/2\)。
\item
  \textbf{奇点处另行处理。} \(r>0\) 的普通场公式不能直接代入
  \(r=0\)；涉及全空间分布时才需要主值与接触项。
\end{enumerate}

\subsection{八、练习与解答}\label{m2:ux516bux7ec3ux4e60ux4e0eux89e3ux7b54}

\subsubsection{作业一：由有限电荷对推出偶极远场}\label{m2:ux4f5cux4e1aux4e00ux7531ux6709ux9650ux7535ux8377ux5bf9ux63a8ux51faux5076ux6781ux8fdcux573a}

\textbf{题号：M2-023}

\textbf{层级：} 基础　　\textbf{方法：} 多极展开

真空中两固定电荷 \(+q,-q\) 分别位于
\(+\boldsymbol\ell/2,-\boldsymbol\ell/2\)，其中
\(q>0\)，\(\boldsymbol\ell\)
从负电荷指向正电荷，\(\ell=|\boldsymbol\ell|\)。场点位置为
\(\mathbf r\)，取 \(\varphi(\infty)=0\)。写出精确电势，再在 \(r\gg\ell\)
下保留最低非零阶，推导偶极势、以 \(\boldsymbol\ell\)
为极轴的球坐标电场分量，以及不依赖坐标选择的电场矢量式。

\par\noindent\begin{minipage}{\linewidth}
\textbf{解答}

\begin{center}
\includegraphics[width=\linewidth,height=.40\textheight,keepaspectratio]{\detokenize{assets/module-02/examples/M2-023.pdf}}
\par\smallskip{\small 图 M2-023：由有限电荷对推出偶极远场}
\end{center}
\end{minipage}\par

图 M2-023 用两段源点到场点的距离构成精确电势。对这两段距离同时按
\(\ell/r\) 展开，零阶项相消、一次项组成
\(\mathbf p\cdot\mathbf r/r^3\)；再求负梯度时必须同时微分径向幂和角向因子。

精确电势为

\[
\varphi_{\mathrm{\text{精确}}}
=kq\left(
\frac1{|\mathbf r-\boldsymbol{\ell}/2|}
-\frac1{|\mathbf r+\boldsymbol{\ell}/2|}
\right).
\]

在 \(r\gg\ell\) 时，

\[
\frac1{|\mathbf r\mp\boldsymbol{\ell}/2|}
=\frac1r
\pm\frac{\boldsymbol{\ell}\cdot\hat{\mathbf r}}{2r^2}
+O\!\left(\frac{\ell^2}{r^3}\right).
\]

两式相减并利用 \(\mathbf p=q\boldsymbol{\ell}\)：

\[
\varphi
=k\frac{\mathbf p\cdot\hat{\mathbf r}}{r^2}
=k\frac{p\cos\theta}{r^2}.
\]

取极轴沿 \(\mathbf p\)，则

\[
E_r=-\frac{\partial\varphi}{\partial r}
=\frac{2kp\cos\theta}{r^3},
\]

\[
E_{\theta}=-\frac1r\frac{\partial\varphi}{\partial\theta}
=\frac{kp\sin\theta}{r^3},
\qquad E_{\phi}=0.
\]

另一方面，写成
\(\varphi=k(\mathbf p\cdot\mathbf r)r^{-3}\)，直接求梯度得

\[
\boxed{
\mathbf E
=\frac{k}{r^3}
\left[3(\mathbf p\cdot\hat{\mathbf r})\hat{\mathbf r}-\mathbf p\right]
}.
\]

对关于中点反对称的电荷对，下一非零修正来自三阶项；一般方向上偶极首项的相对误差为
\(O((\ell/r)^2)\)。

\subsubsection{作业二：偶极场线的两种推导}\label{m2:ux4f5cux4e1aux4e8cux5076ux6781ux573aux7ebfux7684ux4e24ux79cdux63a8ux5bfc}

\textbf{题号：M2-024}

\textbf{层级：} 基础　　\textbf{方法：} 多极展开

真空原点有 \(\mathbf p=p\mathbf e_z\)
的理想偶极子，\(p>0\)，\(k=1/(4\pi\varepsilon_0)\)。在球坐标中已知
\(E_r=2kp\cos\theta/r^3\)、\(E_{\theta}=kp\sin\theta/r^3\)、\(E_{\phi}=0\)，其中
\(r>0\)，\(\theta\) 自 \(+z\)
轴量起。分别利用子午面内场线的微分方程和以极轴为中心的极冠通量函数，推导一般非轴线场线方程，并比较两种积分常数的含义。

\par\noindent\begin{minipage}{\linewidth}
\textbf{解答}

\begin{center}
\includegraphics[width=\linewidth,height=.40\textheight,keepaspectratio]{\detokenize{assets/module-02/examples/M2-024.pdf}}
\par\smallskip{\small 图 M2-024：偶极场线的两种推导}
\end{center}
\end{minipage}\par

图 M2-024 的子午场线与球面极冠给出两条推导路线。局部切向比值决定
\(dr/d\theta\)，而极冠通量沿同一根场管保持不变；两者都得到
\(r=C\sin^2\theta\)，坐标轴上的特殊场线应另列。

由场线切向与电场平行，

\[
\frac{dr}{r\,d\theta}
=\frac{E_r}{E_{\theta}}
=2\cot\theta.
\]

积分得

\[
\ln r=2\ln\sin\theta+\text{常数},
\]

所以

\[
\boxed{r=C\sin^2\theta}.
\]

再用通量法。半径 \(r\) 球面上，从正轴到极角 \(\theta\) 的通量为

\[
\begin{aligned}
\Psi(r,\theta)
&=2\pi r^2\int_0^{\theta}
\frac{2kp\cos\theta'}{r^3}\sin\theta'\,d\theta'\\
&=\frac{p}{2\varepsilon_0r}\sin^2\theta.
\end{aligned}
\]

同一轴对称通量管上 \(\Psi\) 不变，故 \(\sin^2\theta/r\) 为常数，仍得到

\[
r=C\sin^2\theta.
\]

\subsubsection{作业三：轴线与赤道面的方向检验}\label{m2:ux4f5cux4e1aux4e09ux8f74ux7ebfux4e0eux8d64ux9053ux9762ux7684ux65b9ux5411ux68c0ux9a8c}

\textbf{题号：M2-025}

\textbf{层级：} 基础　　\textbf{方法：} 多极展开

理想电偶极子位于真空原点，固定偶极矩为
\(\mathbf p=p\mathbf e_z\)，\(p>0\)。取距原点同为 \(r>0\)
的轴线上场点和赤道面场点。分别求电场的矢量形式，比较其大小，明确轴线两端与赤道面上的场是沿
\(\mathbf p\)、反向于 \(\mathbf p\)，还是沿各自的局部径向。

\par\noindent\begin{minipage}{\linewidth}
\textbf{解答}

\begin{center}
\includegraphics[width=\linewidth,height=.40\textheight,keepaspectratio]{\detokenize{assets/module-02/examples/M2-025.pdf}}
\par\smallskip{\small 图 M2-025：轴线与赤道面的方向检验}
\end{center}
\end{minipage}\par

图 M2-025 的轴线和赤道面分别选取同一半径上的场点。轴线上两项矢量合成沿
\(\mathbf p\)，赤道面上投影项消失，只剩
\(-\mathbf p\)；公用径向系数后，大小之比为2:1，方向相反。

偶极场为

\[
\mathbf E
=\frac{k}{r^3}
\left[3(\mathbf p\cdot\hat{\mathbf r})\hat{\mathbf r}-\mathbf p\right].
\]

在正轴上，\(\hat{\mathbf r}=\mathbf e_z\)，故

\[
\mathbf E=\frac{2kp}{r^3}\mathbf e_z.
\]

在负轴上，\(\hat{\mathbf r}=-\mathbf e_z\)，且
\(\mathbf p\cdot\hat{\mathbf r}=-p\)，所以

\[
\mathbf E
=\frac{k}{r^3}
\left[3(-p)(-\mathbf e_z)-p\mathbf e_z\right]
=\frac{2kp}{r^3}\mathbf e_z.
\]

可见两侧轴线场都沿
\(\mathbf p\)。在赤道面，\(\mathbf p\cdot\hat{\mathbf r}=0\)，所以

\[
\mathbf E=-\frac{kp}{r^3}\mathbf e_z.
\]

因此

\[
|\mathbf E_{\mathrm{\text{轴}}}|:|\mathbf E_{\mathrm{\text{赤道}}}|=2:1,
\]

轴线场沿 \(\mathbf p\)，赤道面场沿 \(-\mathbf p\)。

\subsubsection{作业四：永久偶极的转动与点电荷外场中的平动}\label{m2:ux4f5cux4e1aux56dbux6c38ux4e45ux5076ux6781ux7684ux8f6cux52a8ux4e0eux70b9ux7535ux8377ux5916ux573aux4e2dux7684ux5e73ux52a8}

\textbf{题号：M2-026}

\textbf{层级：} 基础　　\textbf{方法：} 多极展开、机电耦合

一个刚性永久偶极子大小为 \(p>0\)，绕其固定中心的转动惯量为
\(I>0\)，无摩擦。

\begin{enumerate}
\def\labelenumi{(\arabic{enumi})}
\item
  将其置于匀强电场 \(\mathbf E\) 中，令夹角为
  \(\beta\)，求转动方程及稳定平衡附近的小振动周期。
\item
  另考虑真空原点固定点电荷 \(Q\) 所产生的非均匀外场。偶极中心位于
  \(\mathbf R=R\hat{\mathbf R}\)，\(R\)
  远大于偶极尺寸，计算力时偶极矩保持给定取向。求一阶偶极近似的净力，再判断
  \(Q>0\)、\(\mathbf p=p\hat{\mathbf R}\) 时的力向。
\end{enumerate}

\par\noindent\begin{minipage}{\linewidth}
\textbf{解答}

\begin{center}
\includegraphics[width=\linewidth,height=.40\textheight,keepaspectratio]{\detokenize{assets/module-02/examples/M2-026.pdf}}
\par\smallskip{\small 图 M2-026：永久偶极的转动与点电荷外场中的平动}
\end{center}
\end{minipage}\par

图 M2-026 区分偶极自身的取向角与其中心到外部点电荷的位矢。转动部分由
\(\mathbf p\times\mathbf E\)
给出，平动部分则对局部外场求空间导数；代入径向偶极时还须保留场方向随位置变化的贡献。

匀强场中力矩为 \(\boldsymbol\tau=\mathbf p\times\mathbf E\)。取使
\(\beta\) 增大的方向为正，回复力矩为

\[
\tau_{\beta}=-pE\sin\beta.
\]

故运动方程为

\[
I\ddot\beta=-pE\sin\beta.
\]

稳定平衡在 \(\beta=0\)。小角度时 \(\sin\beta\approx\beta\)：

\[
\ddot\beta+\frac{pE}{I}\beta=0.
\]

所以

\[
\omega=\sqrt{\frac{pE}{I}},
\qquad
\boxed{T=2\pi\sqrt{\frac{I}{pE}}}.
\]

点电荷产生的外场为

\[
\mathbf E_Q(\mathbf R)=kQ\frac{\mathbf R}{R^3}.
\]

永久偶极势能为

\[
U=-\mathbf p\cdot\mathbf E_Q
=-kQ\frac{\mathbf p\cdot\mathbf R}{R^3}.
\]

固定 \(\mathbf p\) 并求 \(\mathbf F=-\nabla U\)：

\[
\boxed{
\mathbf F
=\frac{kQ}{R^3}
\left[
\mathbf p-3(\mathbf p\cdot\hat{\mathbf R})\hat{\mathbf R}
\right]
}.
\]

当 \(Q>0\) 且 \(\mathbf p=p\hat{\mathbf R}\) 时，

\[
\boxed{
\mathbf F=-\frac{2kQp}{R^3}\hat{\mathbf R}
}.
\]

净力指向点电荷。此时 \(\mathbf p\parallel\mathbf E_Q\)，力矩为零。

\subsubsection{作业五：线性感应偶极的能量与受力}\label{m2:ux4f5cux4e1aux4e94ux7ebfux6027ux611fux5e94ux5076ux6781ux7684ux80fdux91cfux4e0eux53d7ux529b}

\textbf{题号：M2-027}

\textbf{层级：} 基础　　\textbf{方法：} 能量法、多极展开

一个足够小的各向同性可极化粒子处在给定静电外场中，线性、可逆响应满足
\(\mathbf p=\alpha\mathbf E\)，\(\alpha>0\)
为常量。粒子尺度远小于外场变化尺度，忽略其对外场的扰动，并以无外场时相互作用能为零。证明诱导偶极势能与力分别为
\(U=-\alpha E^2/2\)、\(\mathbf F=\alpha\nabla E^2/2\)。若中心附近
\(\mathbf E(x)=E_0(1+x/L)\mathbf e_x\)，\(L>0\)，求位于 \(x=0\)
的粒子所受力。

\par\noindent\begin{minipage}{\linewidth}
\textbf{解答}

\begin{center}
\includegraphics[width=\linewidth,height=.40\textheight,keepaspectratio]{\detokenize{assets/module-02/examples/M2-027.pdf}}
\par\smallskip{\small 图 M2-027：线性感应偶极的能量与受力}
\end{center}
\end{minipage}\par

图 M2-027
的诱导偶极矩与局部电场同向，其大小也随场变。能量需要从零极化逐步积分，因而出现
\(1/2\)；随后求梯度时对 \(E^2\)
求导，在原点读取局部斜率，得到指向强场侧的力。

把外场从零缓慢增加到 \(\mathbf E\)。偶极矩也从零增加到
\(\alpha\mathbf E\)，因此势能变化为

\[
U
=-\int_{\mathbf0}^{\mathbf E}
\mathbf p(\mathbf E')\cdot d\mathbf E'
=-\int_0^E\alpha E'\,dE'
=-\frac12\alpha E^2.
\]

于是

\[
\mathbf F=-\nabla U
=\frac12\alpha\nabla E^2.
\]

在给定的一维场中，

\[
E^2(x)=E_0^2\left(1+\frac{x}{L}\right)^2,
\]

所以

\[
\left.\frac{dE^2}{dx}\right|_{x=0}
=\frac{2E_0^2}{L}.
\]

最终

\[
\boxed{
\mathbf F(0)=\frac{\alpha E_0^2}{L}\mathbf e_x
}.
\]

粒子被拉向场强较大的一侧。

\section{专题十一：多极展开（基础必讲与拔高选讲）}\label{m2:ux4e13ux9898ux5341ux4e00ux591aux6781ux5c55ux5f00ux57faux7840ux5fc5ux8bb2ux4e0eux62d4ux9ad8ux9009ux8bb2}

多极展开按无量纲小量 \(a/r\)
将局域电荷分布的远场分成单极、偶极、四极等阶次。低阶项记录整体电荷信息，高阶项逐步反映源区形状；真实电荷的位置并未改变。

本专题统一用 \(L=0,1,2,\ldots\) 表示多极展开的阶数，用 \(\ell\)
表示具体电荷构型中的长度，避免两个符号混淆。

第一轮学习应掌握低阶多极项、对称性筛选和远场阶数判断；严格余项界、原点平移、球谐矩以及非轴对称四极张量用于决赛层面的系统化与误差控制，可在第二轮学习。

\subsection{一、源区、收敛条件与截断误差}\label{m2:ux4e00ux6e90ux533aux6536ux655bux6761ux4ef6ux4e0eux622aux65adux8befux5dee}

设电荷分布 \(\rho(\mathbf r')\) 完全包含在以所选原点 \(O\)
为球心、半径为 \(a\) 的球内，即所有源点都满足 \(r'\le a\)。场点为
\(\mathbf r\)，精确电势为

\[
\varphi(\mathbf r)
=\frac1{4\pi\varepsilon_0}
\int\frac{\rho(\mathbf r')}{|\mathbf r-\mathbf r'|}\,dV'.
\]

\par\noindent\begin{minipage}{\linewidth}
\begin{center}
\includegraphics[width=\linewidth,height=.40\textheight,keepaspectratio]{\detokenize{assets/module-02/original/tikz-m2-topic11-v3-01-source-region-pdf-v11.jpg}}
\par\smallskip{\small 图示：局域源区、源点、场点以及多极展开的两个距离条件}
\end{center}
\end{minipage}\par

\subsubsection{1.
库仑核的精确勒让德级数}\label{m2:ux5e93ux4ed1ux6838ux7684ux7cbeux786eux52d2ux8ba9ux5fb7ux7ea7ux6570}

令 \(\gamma\) 为 \(\mathbf r\) 与 \(\mathbf r'\) 的夹角，则

\[
|\mathbf r-\mathbf r'|
=r\sqrt{1-2\frac{r'}r\cos\gamma+\frac{r'^2}{r^2}}.
\]

利用勒让德多项式的母函数

\[
\frac1{\sqrt{1-2tx+t^2}}
=\sum_{L=0}^{\infty}t^LP_L(x),
\qquad |t|<1,
\]

取 \(t=r'/r\)、\(x=\cos\gamma\)，得到

\[
\boxed{
\frac1{|\mathbf r-\mathbf r'|}
=\sum_{L=0}^{\infty}
\frac{r'^L}{r^{L+1}}P_L(\cos\gamma)
},
\qquad r>r'.
\]

若 \(r>a\)，则对源区内所有点都有 \(r'/r<1\)，上式在整个源区成立。于是

\[
\boxed{
\varphi(\mathbf r)
=\frac1{4\pi\varepsilon_0}
\sum_{L=0}^{\infty}
\frac1{r^{L+1}}
\int\rho(\mathbf r')r'^LP_L(\cos\gamma)\,dV'
}.
\]

这里要区分两个条件：

\begin{itemize}
\tightlist
\item
  \(r>a\)：级数可以收敛地表示精确电势；
\item
  \(r\gg a\)：\(a/r\) 很小，只保留少数低阶项才具有良好精度。
\end{itemize}

\subsubsection{2.
截断余项的上界}\label{m2:ux622aux65adux4f59ux9879ux7684ux4e0aux754c}

记

\[
\eta=\frac ar<1.
\]

若库仑核保留到 \(L\) 阶，由于 \(|P_n(x)|\le1\)，余项满足

\[
\begin{aligned}
|R_L|
&\le\frac1r\sum_{n=L+1}^{\infty}\left(\frac{r'}r\right)^n\\
&\le\frac1r\frac{\eta^{L+1}}{1-\eta}.
\end{aligned}
\]

定义源区的总绝对电荷量

\[
Q_{\text{绝}}=\int|\rho(\mathbf r')|\,dV',
\]

则电势截断误差满足

\[
\boxed{
|\Delta\varphi_L|
\le
\frac1{4\pi\varepsilon_0}
\frac{Q_{\text{绝}}}{r}
\frac{\eta^{L+1}}{1-\eta}
}.
\]

\par\noindent\begin{minipage}{\linewidth}
\begin{center}
\includegraphics[width=\linewidth,height=.40\textheight,keepaspectratio]{\detokenize{assets/module-02/theory/unified7-kernel-error.pdf}}
\par\smallskip{\small 图示：库仑核的低阶勒让德项、衰减阶次和截断误差界}
\end{center}
\end{minipage}\par

\textbf{例 M2-028：截断界控制绝对误差尺度}

一个局域电荷分布完全包含在以展开原点为中心、半径为 \(a\)
的球内。定义绝对电荷总量
\(Q_{\rm abs}=\int|\rho(\mathbf r')|\,dV'\)（点电荷情形为
\(\sum_i|q_i|\)）。观察点到原点距离为 \(r>a\)，满足
\(a/r=0.1\)。多极展开保留到 \(L=2\)
的四极项，利用几何级数型截断界估计被遗漏电势的绝对上界，以
\(Q_{\rm abs}/(4\pi\varepsilon_0r)\) 为尺度给出结果。

\par\noindent\begin{minipage}{\linewidth}
\textbf{解答}

\begin{center}
\includegraphics[width=\linewidth,height=.40\textheight,keepaspectratio]{\detokenize{assets/module-02/examples/M2-028.pdf}}
\par\smallskip{\small 图 M2-028：截断界控制绝对误差尺度}
\end{center}
\end{minipage}\par

图 M2-028 把全部源限制在半径 \(a\) 内，场点位于外部。遗漏级数从
\(L+1=3\) 阶开始，用 \(\eta=a/r\)
控制等比尾和；其尺度使用各电荷绝对值之和，因此不能把该界直接解释成实际总电势的相对误差。

若 \(a/r=0.1\)，保留到四极项 \(L=2\)，则

\[
\frac{\eta^3}{1-\eta}
=\frac{10^{-3}}{0.9}
=\frac1{900}.
\]

这说明绝对误差不超过

\[
\frac1{900}
\frac{Q_{\text{绝}}}{4\pi\varepsilon_0r}.
\]

这不是总电势的相对误差：若正负电荷强烈抵消，或观测方向恰好位于某阶项的角向节点，总电势可能远小于上述电势尺度。

\subsection{二、单极项、偶极项与四极项（基础必讲）}\label{m2:ux4e8cux5355ux6781ux9879ux5076ux6781ux9879ux4e0eux56dbux6781ux9879ux57faux7840ux5fc5ux8bb2}

展开库仑核至二阶：

\[
\frac1{|\mathbf r-\mathbf r'|}
=\frac1r
+\frac{r'\cos\gamma}{r^2}
+\frac{r'^2(3\cos^2\gamma-1)}{2r^3}
+O\!\left(\frac{a^3}{r^4}\right).
\]

记 \(\mathbf n=\hat{\mathbf r}\)，则

\[
r'\cos\gamma=\mathbf n\cdot\mathbf r'.
\]

\subsubsection{1. 单极矩}\label{m2:ux5355ux6781ux77e9}

总电荷为

\[
\boxed{Q=\int\rho(\mathbf r')\,dV'}.
\]

单极电势为

\[
\boxed{
\varphi^{(0)}=\frac1{4\pi\varepsilon_0}\frac Qr
}.
\]

\textbf{例 M2-029：净电荷决定最远区的首项}

真空中有两个固定点电荷 \(+3q\) 与 \(-q\)，其中
\(q>0\)，它们都位于某个有限源区中，源区之外没有其他电荷或外场。取无穷远电势为零，在观察距离远大于源区尺度时，确定电势的最低非零多极阶和对应首项，说明为什么只凭总电荷就能确定最远区的主要衰减形式。

\par\noindent\begin{minipage}{\linewidth}
\textbf{解答}

\begin{center}
\includegraphics[width=\linewidth,height=.40\textheight,keepaspectratio]{\detokenize{assets/module-02/examples/M2-029.pdf}}
\par\smallskip{\small 图 M2-029：净电荷决定最远区的首项}
\end{center}
\end{minipage}\par

图 M2-029 中两电荷代数和为
\(2q\)，不因摆放位置而改变。远区展开先保留这一总量除以
\(r\)，电荷间距和方向只影响下一阶及以上项；最远区不能因异号电荷同时存在而默认单极项消失。

若局域体系由 \(+3q\) 与 \(-q\) 组成，则
\(Q=2q\)。无论两电荷如何摆放，只要观测距离远大于源区尺度，最低阶电势必为

\[
\varphi^{(0)}=\frac{1}{4\pi\varepsilon_0}\frac{2q}{r}.
\]

电荷间距和方向只从偶极项及更高阶项开始进入。

\subsubsection{2. 偶极矩}\label{m2:ux5076ux6781ux77e9}

电偶极矩为

\[
\boxed{
\mathbf p=\int\rho(\mathbf r')\mathbf r'\,dV'
}.
\]

偶极电势为

\[
\boxed{
\varphi^{(1)}
=\frac1{4\pi\varepsilon_0}
\frac{\mathbf p\cdot\mathbf n}{r^2}
}.
\]

\textbf{例 M2-030：中性体系从偶极阶开始}

真空中一个局域电荷对由 \(-q\) 与 \(+q\)
构成，\(q>0\)；从负电荷到正电荷的位移为非零矢量
\(\mathbf d\)，令中点为展开原点。在 \(r\gg|\mathbf d|\)
的远区，求总电荷及最低非零多极矩，写出首项电势的角向依赖与随 \(r\)
的衰减阶数。

\par\noindent\begin{minipage}{\linewidth}
\textbf{解答}

\begin{center}
\includegraphics[width=\linewidth,height=.40\textheight,keepaspectratio]{\detokenize{assets/module-02/examples/M2-030.pdf}}
\par\smallskip{\small 图 M2-030：中性体系从偶极阶开始}
\end{center}
\end{minipage}\par

图 M2-030 的绿色位移从负端指向正端，故
\(\mathbf p=q\mathbf d\)。两电荷的 \(1/r\)
项相消，第一项由该有向位移在观测方向上的投影决定；反转位移会反转偶极首项的符号。

\(-q\) 指向 \(+q\) 的位移为 \(\mathbf d\) 时，

\[
Q=0,
\qquad
\mathbf p=q\mathbf d.
\]

因此单极项消失，远场从 \(r^{-2}\) 的偶极电势开始；把 \(\mathbf d\)
方向反转会使整个偶极首项变号。

\subsubsection{3. 四极矩张量}\label{m2:ux56dbux6781ux77e9ux5f20ux91cf}

本讲义固定采用对称无迹四极矩张量

\[
\boxed{
\mathcal Q_{ij}
=\int\rho(\mathbf r')
\left(3x_i'x_j'-r'^2\delta_{ij}\right)dV'
}.
\]

它满足

\[
\mathcal Q_{ij}=\mathcal Q_{ji},
\qquad
\mathcal Q_{xx}+\mathcal Q_{yy}+\mathcal Q_{zz}=0,
\]

因此一般只有五个独立分量。使用重复指标求和约定，有

\[
\begin{aligned}
n_i\mathcal Q_{ij}n_j
&=\int\rho(\mathbf r')
\left[3(\mathbf n\cdot\mathbf r')^2-r'^2\right]dV'.
\end{aligned}
\]

所以四极电势为

\[
\boxed{
\varphi^{(2)}
=\frac1{4\pi\varepsilon_0}
\frac{n_i\mathcal Q_{ij}n_j}{2r^3}
}.
\]

\textbf{例 M2-031：三个电荷也能构成线型四极子}

在真空 \(z\) 轴上，\(z=\ell\) 与 \(z=-\ell\) 处各固定一个
\(+q\)，原点处固定 \(-2q\)，其中
\(q>0,\ell>0\)，取无穷远电势为零。以原点为展开中心、\(+z\)
为极轴，求总电荷和电偶极矩，判断首个非零多极阶；给出轴对称四极系数
\(A_2=\sum_iq_i z_i^2\) 或等价四极矩，说明远场首项怎样衰减。

\par\noindent\begin{minipage}{\linewidth}
\textbf{解答}

\begin{center}
\includegraphics[width=\linewidth,height=.40\textheight,keepaspectratio]{\detokenize{assets/module-02/examples/M2-031.pdf}}
\par\smallskip{\small 图 M2-031：三个电荷也能构成线型四极子}
\end{center}
\end{minipage}\par

图 M2-031
的上下正电荷互相抵消偶极矩，中央负电荷再消去总电荷。把三个位置代入同一个无迹张量定义，得到纵向与横向分量之比为
\(-2\)；因此远区首项为四极阶，而非由电荷个数命名。

在 \(z=\pm\ell\) 各放置 \(+q\)，在原点放置 \(-2q\)。则

\[
Q=0,
\qquad
\mathbf p=\mathbf0,
\]

而

\[
\mathcal Q_{zz}=4q\ell^2,
\qquad
\mathcal Q_{xx}=\mathcal Q_{yy}=-2q\ell^2.
\]

故最低非零远场是四极项，电势按 \(r^{-3}\)
衰减。这也说明``四极子''是最低非零矩量的名称，并不要求恰有四个电荷。

有些资料把

\[
\Theta_{ij}
=\int\rho\left(x_i'x_j'-\frac13r'^2\delta_{ij}\right)dV'
\]

称为四极矩，此时
\(\mathcal Q_{ij}=3\Theta_{ij}\)。比较不同资料时，必须同时核对定义和电势公式中的系数。

\par\noindent\begin{minipage}{\linewidth}
\begin{center}
\includegraphics[width=\linewidth,height=.40\textheight,keepaspectratio]{\detokenize{assets/module-02/original/tikz-m2-topic11-v3-04-quadrupole-tensor-margin12-v11.png}}
\par\smallskip{\small 图示：对称无迹四极矩张量及其沿观测方向的投影}
\end{center}
\end{minipage}\par

综合前三阶：

\[
\boxed{
\varphi(\mathbf r)
=\frac1{4\pi\varepsilon_0}
\left[
\frac Qr
+\frac{\mathbf p\cdot\mathbf n}{r^2}
+\frac{n_i\mathcal Q_{ij}n_j}{2r^3}
+\cdots
\right]
}.
\]

\par\noindent\begin{minipage}{\linewidth}
\begin{center}
\includegraphics[width=\linewidth,height=.40\textheight,keepaspectratio]{\detokenize{assets/module-02/original/tikz-m2-topic11-v3-03-leading-order-pdf-v11.jpg}}
\par\smallskip{\small 图示：总电荷、偶极矩与四极矩中第一个非零量决定远场主导阶}
\end{center}
\end{minipage}\par

{\def\LTcaptype{none} % do not increment counter
\begin{longtable}[]{@{}
  >{\raggedright\arraybackslash}p{(\linewidth - 8\tabcolsep) * \real{0.2000}}
  >{\raggedright\arraybackslash}p{(\linewidth - 8\tabcolsep) * \real{0.2000}}
  >{\raggedright\arraybackslash}p{(\linewidth - 8\tabcolsep) * \real{0.2000}}
  >{\raggedright\arraybackslash}p{(\linewidth - 8\tabcolsep) * \real{0.2000}}
  >{\raggedright\arraybackslash}p{(\linewidth - 8\tabcolsep) * \real{0.2000}}@{}}
\toprule\noalign{}
\begin{minipage}[b]{\linewidth}\raggedright
展开阶数
\end{minipage} & \begin{minipage}[b]{\linewidth}\raggedright
矩量
\end{minipage} & \begin{minipage}[b]{\linewidth}\raggedright
电势衰减
\end{minipage} & \begin{minipage}[b]{\linewidth}\raggedright
电场衰减
\end{minipage} & \begin{minipage}[b]{\linewidth}\raggedright
成为主导项的条件
\end{minipage} \\
\midrule\noalign{}
\endhead
\bottomrule\noalign{}
\endlastfoot
\(L=0\) & \(Q\) & \(r^{-1}\) & \(r^{-2}\) & \(Q\ne0\) \\
\(L=1\) & \(\mathbf p\) & \(r^{-2}\) & \(r^{-3}\) &
\(Q=0,\ \mathbf p\ne0\) \\
\(L=2\) & \(\mathcal Q_{ij}\) & \(r^{-3}\) & \(r^{-4}\) &
\(Q=0,\ \mathbf p=0,\ \mathcal Q\ne0\) \\
\end{longtable}
}

``四极子''不是``必须有四个点电荷''的意思。严格判据是：总电荷和偶极矩消失，而四极矩不消失。三个点电荷同样可以构成线型四极子。

\subsection{三、坐标原点平移与多极矩重分配（拔高选讲）}\label{m2:ux4e09ux5750ux6807ux539fux70b9ux5e73ux79fbux4e0eux591aux6781ux77e9ux91cdux5206ux914dux62d4ux9ad8ux9009ux8bb2}

设新原点相对旧原点的位移为 \(\mathbf a\)，则源点坐标变为

\[
\mathbf r_i'=\mathbf r_i-\mathbf a.
\]

\par\noindent\begin{minipage}{\linewidth}
\begin{center}
\includegraphics[width=\linewidth,height=.40\textheight,keepaspectratio]{\detokenize{assets/module-02/original/tikz-m2-topic11-v3-05-origin-shift-pdf-v11.jpg}}
\par\smallskip{\small 图示：同一电荷分布在原点平移前后的源点坐标与多极矩变换}
\end{center}
\end{minipage}\par

总电荷显然不变：

\[
\boxed{Q'=Q}.
\]

偶极矩变为

\[
\begin{aligned}
\mathbf p'
&=\int\rho(\mathbf r')\bigl(\mathbf r'-\mathbf a\bigr)dV'\\
&=\mathbf p-Q\mathbf a,
\end{aligned}
\]

即

\[
\boxed{\mathbf p'=\mathbf p-Q\mathbf a}.
\]

对四极矩逐项展开可得

\[
\boxed{
\begin{aligned}
\mathcal Q_{ij}'
= {}&\mathcal Q_{ij}
-3(a_ip_j+a_jp_i)
+2(\mathbf a\cdot\mathbf p)\delta_{ij}\\
&+Q(3a_ia_j-a^2\delta_{ij}).
\end{aligned}
}
\]

由此得到：

\begin{enumerate}
\def\labelenumi{\arabic{enumi}.}
\tightlist
\item
  \(Q\) 与原点无关；
\item
  当 \(Q=0\) 时，\(\mathbf p\) 与原点无关；
\item
  当 \(Q=0\) 且 \(\mathbf p=0\) 时，\(\mathcal Q_{ij}\) 与原点无关。
\end{enumerate}

若 \(Q\ne0\)，可以选取

\[
\mathbf a=\frac{\mathbf p}{Q}
\]

使新原点下
\(\mathbf p'=0\)。这只是把偶极修正吸收到原点选择中；它不意味着真实电荷集中在该点，也不保证四极矩和更高阶矩消失。

\textbf{例 M2-032：等效中心可以位于实际源区之外}

真空 \(x\) 轴上有两个固定电荷：\(x=\ell\) 处为 \(+3q\)，\(x=-\ell\) 处为
\(-q\)，其中 \(q>0,\ell>0\)。先以 \(x=0\)
为原点求总电荷和偶极矩，再把展开原点平移到 \(x=x_c\)，用新坐标
\(x_i'=x_i-x_c\) 求使偶极矩为零的
\(x_c\)。比较该中心与两个实际电荷位置，说明``等效中心''是否必须落在源点之间。

\par\noindent\begin{minipage}{\linewidth}
\textbf{解答}

\begin{center}
\includegraphics[width=\linewidth,height=.40\textheight,keepaspectratio]{\detokenize{assets/module-02/examples/M2-032.pdf}}
\par\smallskip{\small 图 M2-032：等效中心可以位于实际源区之外}
\end{center}
\end{minipage}\par

图 M2-032 的新原点 \(O_c\)
位于右侧真实电荷之外。先计算代数总电荷和一次矩，再用 \(p_x-Qa_x=0\)
选取平移量；这里消去的是展开中的偶极项，没有把真实电荷移到 \(O_c\)。

\(+3q\) 位于 \(x=\ell\)，\(-q\) 位于 \(x=-\ell\)。以原点为参考，

\[
Q=2q,
\qquad
p_x=3q\ell+q\ell=4q\ell.
\]

使偶极矩消失的新原点位于

\[
a_x=\frac{p_x}{Q}=2\ell.
\]

该点甚至可以位于原电荷分布之外，进一步说明它是展开意义下的等效中心，而不是物质中心。

\subsection{四、对称性筛选与球谐展开（基础筛选与拔高球谐）}\label{m2:ux56dbux5bf9ux79f0ux6027ux7b5bux9009ux4e0eux7403ux8c10ux5c55ux5f00ux57faux7840ux7b5bux9009ux4e0eux62d4ux9ad8ux7403ux8c10}

在计算多极矩之前，先用对称性判断哪些阶必为零，通常能显著缩短计算。

\par\noindent\begin{minipage}{\linewidth}
\begin{center}
\includegraphics[width=\linewidth,height=.40\textheight,keepaspectratio]{\detokenize{assets/module-02/original/tikz-m2-topic11-v3-06-symmetry-filter-pdf-v11.jpg}}
\par\smallskip{\small 图示：反演对称性和轴对称性对多极阶数与球谐分量的筛选}
\end{center}
\end{minipage}\par

\subsubsection{1. 反演对称性}\label{m2:ux53cdux6f14ux5bf9ux79f0ux6027}

球谐函数在空间反演下满足

\[
Y_{Lm}(-\Omega)=(-1)^LY_{Lm}(\Omega).
\]

因此：

\begin{itemize}
\tightlist
\item
  若 \(\rho(-\mathbf r)=\rho(\mathbf r)\)，所有奇数阶 \(L\) 消失；
\item
  若 \(\rho(-\mathbf r)=-\rho(\mathbf r)\)，所有偶数阶 \(L\) 消失。
\end{itemize}

注意，``某些阶因对称性消失''不等于``余下最低阶一定非零''，还要继续计算或利用其他对称性判断。

\subsubsection{2.
一般球谐展开}\label{m2:ux4e00ux822cux7403ux8c10ux5c55ux5f00}

库仑核的球谐加法定理为

\[
\frac1{|\mathbf r-\mathbf r'|}
=4\pi\sum_{L=0}^{\infty}\sum_{m=-L}^{L}
\frac1{2L+1}
\frac{r'^L}{r^{L+1}}
Y_{Lm}(\Omega)Y_{Lm}^*(\Omega'),
\qquad r>r'.
\]

定义球多极矩

\[
\boxed{
q_{Lm}
=\int\rho(\mathbf r')r'^LY_{Lm}^*(\Omega')\,dV'
},
\]

便有

\[
\boxed{
\varphi(\mathbf r)
=\frac1{\varepsilon_0}
\sum_{L=0}^{\infty}\sum_{m=-L}^{L}
\frac{q_{Lm}}{(2L+1)r^{L+1}}Y_{Lm}(\Omega)
}.
\]

这一定义采用归一化球谐函数；若球谐归一化约定改变，\(q_{Lm}\)
的系数也会随之改变。

\subsubsection{3. 轴对称情形}\label{m2:ux8f74ux5bf9ux79f0ux60c5ux5f62}

若源绕 \(z\) 轴旋转不变，则只有 \(m=0\) 分量。定义

\[
A_L
=\int\rho(\mathbf r')r'^LP_L(\cos\theta')\,dV',
\]

电势化为

\[
\boxed{
\varphi(r,\theta)
=\frac1{4\pi\varepsilon_0}
\sum_{L=0}^{\infty}
\frac{A_L}{r^{L+1}}P_L(\cos\theta)
}.
\]

其中

\[
A_0=Q,
\qquad
A_1=p_z,
\qquad
A_2=\frac12\mathcal Q_{zz}.
\]

\subsection{五、线型电四极子（基础必讲与拔高选讲）}\label{m2:ux4e94ux7ebfux578bux7535ux56dbux6781ux5b50ux57faux7840ux5fc5ux8bb2ux4e0eux62d4ux9ad8ux9009ux8bb2}

考虑三个点电荷：\(+q\) 位于 \(z=\ell\)，\(-2q\) 位于原点，\(+q\) 位于
\(z=-\ell\)。

\par\noindent\begin{minipage}{\linewidth}
\begin{center}
\includegraphics[width=\linewidth,height=.40\textheight,keepaspectratio]{\detokenize{assets/module-02/original/tikz-m2-topic11-v3-07-linear-quadrupole.png}}
\par\smallskip{\small 图示：线型三电荷的总电荷、偶极矩、四极张量与远场电势}
\end{center}
\end{minipage}\par

\subsubsection{1.
精确电势与阶数筛选}\label{m2:ux7cbeux786eux7535ux52bfux4e0eux9636ux6570ux7b5bux9009}

精确电势为

\[
\begin{aligned}
\varphi_{\mathrm{\text{精确}}}(r,\theta)
=\frac{q}{4\pi\varepsilon_0}
\Bigg[&
\frac1{\sqrt{r^2+\ell^2-2r\ell\cos\theta}}\\
&+\frac1{\sqrt{r^2+\ell^2+2r\ell\cos\theta}}
-\frac2r
\Bigg].
\end{aligned}
\]

系统总电荷与偶极矩均为零：

\[
Q=q-2q+q=0,
\]

\[
\mathbf p=q\ell\mathbf e_z+q(-\ell\mathbf e_z)=\mathbf0.
\]

电荷分布满足反演偶对称，故所有奇数阶 \(L\) 消失。对
\(r>\ell\)，精确级数可写为

\[
\boxed{
\varphi(r,\theta)
=\frac{2q}{4\pi\varepsilon_0}
\sum_{n=1}^{\infty}
\frac{\ell^{2n}}{r^{2n+1}}P_{2n}(\cos\theta)
}.
\]

首项对应 \(n=1\)、\(L=2\)：

\[
\boxed{
\varphi^{(2)}(r,\theta)
=\frac{q\ell^2}{4\pi\varepsilon_0r^3}
(3\cos^2\theta-1)
}.
\]

下一修正为 \(L=4\)，所以在四极首项不为零的一般方向上，相对修正为
\(O((\ell/r)^2)\)。

\subsubsection{2.
四极张量与电场}\label{m2:ux56dbux6781ux5f20ux91cfux4e0eux7535ux573a}

逐个点电荷求和：

\[
\boxed{
\mathcal Q
=q\ell^2
\begin{pmatrix}
-2&0&0\\
0&-2&0\\
0&0&4
\end{pmatrix}
}.
\]

令

\[
C=\frac{q\ell^2}{4\pi\varepsilon_0},
\]

则 \(\varphi^{(2)}=C(3\cos^2\theta-1)r^{-3}\)。求梯度得

\[
\boxed{
E_r
=\frac{3q\ell^2}{4\pi\varepsilon_0r^4}
(3\cos^2\theta-1)
},
\]

\[
\boxed{
E_{\theta}
=\frac{6q\ell^2}{4\pi\varepsilon_0r^4}
\sin\theta\cos\theta
},
\qquad E_{\phi}=0.
\]

电势按 \(r^{-3}\) 衰减，电场按 \(r^{-4}\) 衰减。

\subsubsection{3.
四极首项的零势锥}\label{m2:ux56dbux6781ux9996ux9879ux7684ux96f6ux52bfux9525}

四极首项在

\[
3\cos^2\theta_0-1=0
\]

处为零，即

\[
\boxed{
\cos^2\theta_0=\frac13,
\qquad
\theta_0\approx54.7^{\circ}, 125.3^{\circ}
}.
\]

子午截面中的四条射线绕 \(z\) 轴旋转，形成两张圆锥面。

\par\noindent\begin{minipage}{\linewidth}
\begin{center}
\includegraphics[width=\linewidth,height=.40\textheight,keepaspectratio]{\detokenize{assets/module-02/original/tikz-m2-topic11-v3-08-zero-cone-field-pdf-v11.jpg}}
\par\smallskip{\small 图示：线型四极首项的正负区域、零势锥和非零角向电场}
\end{center}
\end{minipage}\par

在该角度，四极近似给出

\[
\varphi^{(2)}=0,
\qquad
E_r^{(2)}=0,
\]

但

\[
E_{\theta}^{(2)}
=\frac{6q\ell^2}{4\pi\varepsilon_0r^4}
\sin\theta_0\cos\theta_0
\ne0.
\]

所以零电势锥不是零电场面。对有限三电荷系统，\(L=4,6,\ldots\)
项仍然存在，精确零势面会相对四极锥发生 \(O((\ell/r)^2)\) 的偏移；在
\(\theta_0\) 方向不能用四极首项计算总电势的相对误差。

\subsection{六、非轴对称四极子（拔高选讲）}\label{m2:ux516dux975eux8f74ux5bf9ux79f0ux56dbux6781ux5b50ux62d4ux9ad8ux9009ux8bb2}

在 \(xy\) 平面内，令 \(+q\) 位于 \((\pm a,0,0)\)，\(-q\) 位于
\((0,\pm a,0)\)。该系统有

\[
Q=0,
\qquad
\mathbf p=\mathbf0.
\]

四极张量为

\[
\boxed{
\mathcal Q
=6qa^2
\begin{pmatrix}
1&0&0\\
0&-1&0\\
0&0&0
\end{pmatrix}
}.
\]

因为

\[
n_x=\sin\theta\cos\phi,
\qquad
n_y=\sin\theta\sin\phi,
\]

所以

\[
\begin{aligned}
n_i\mathcal Q_{ij}n_j
&=6qa^2(n_x^2-n_y^2)\\
&=6qa^2\sin^2\theta\cos2\phi.
\end{aligned}
\]

远场四极电势为

\[
\boxed{
\varphi^{(2)}(r,\theta,\phi)
=\frac{3qa^2}{4\pi\varepsilon_0r^3}
\sin^2\theta\cos2\phi
}.
\]

\par\noindent\begin{minipage}{\linewidth}
\begin{center}
\includegraphics[width=\linewidth,height=.40\textheight,keepaspectratio]{\detokenize{assets/module-02/original/tikz-m2-topic11-v3-09-planar-quadrupole-safe-v7.png}}
\par\smallskip{\small 图示：平面四电荷的四极张量及赤道面四瓣角分布}
\end{center}
\end{minipage}\par

在赤道面 \(\theta=\pi/2\)，电势随 \(\cos2\phi\) 变化：沿 \(\pm x\)
方向为正，沿 \(\pm y\) 方向为负；\(x=\pm y\)
对应的竖直平面是角向节点面。这个例子说明，一般四极矩必须用张量或
\(q_{2m}\) 描述，不能只用一个标量``\(Q_2\)''。

\subsection{七、多极展开的使用条件}\label{m2:ux4e03ux591aux6781ux5c55ux5f00ux7684ux4f7fux7528ux6761ux4ef6}

\par\noindent\begin{minipage}{\linewidth}
\begin{center}
\includegraphics[width=\linewidth,height=.40\textheight,keepaspectratio]{\detokenize{assets/module-02/original/tikz-m2-topic11-v3-10-workflow-pdf-v11.jpg}}
\par\smallskip{\small 图示：多极展开中的源区、对称性、矩量和截断误差}
\end{center}
\end{minipage}\par

\begin{enumerate}
\def\labelenumi{\arabic{enumi}.}
\tightlist
\item
  \textbf{先圈定源区。} 选择原点并给出包围半径 \(a\)；若场点不满足
  \(r>a\)，外部多极级数不能覆盖整个源区。
\item
  \textbf{区分收敛与近似。} \(r>a\) 保证可展开，\(a/r\ll1\)
  才保证低阶截断有效。
\item
  \textbf{先用对称性。} 反演、镜像和轴对称常能直接消去整组多极项。
\item
  \textbf{再找第一个非零矩。}
  点电荷个数不能决定``几极子''，低阶矩是否消失才是判据。
\item
  \textbf{固定四极矩约定。} 若定义中是否含因子 \(3\)
  不同，电势公式前的系数也不同。
\item
  \textbf{检查原点依赖。} \(Q\) 总是不变；\(\mathbf p\) 只在 \(Q=0\)
  时不变；\(\mathcal Q\) 只在 \(Q=0\) 且 \(\mathbf p=0\) 时不变。
\item
  \textbf{不要忽略角向节点。}
  全局主导项可能在某些方向恰好为零，下一阶会在这些方向局部主导。
\item
  \textbf{检查衰减幂次。} \(L\) 阶电势按 \(r^{-L-1}\) 衰减，相应电场按
  \(r^{-L-2}\) 衰减。
\item
  \textbf{近场回到精确式。}
  多极展开不能描述紧邻某个真实点电荷的库仑奇点结构。
\end{enumerate}

\subsection{八、练习与解答}\label{m2:ux516bux7ec3ux4e60ux4e0eux89e3ux7b54-1}

\subsubsection{作业一：线型电四极子的远场}\label{m2:ux4f5cux4e1aux4e00ux7ebfux578bux7535ux56dbux6781ux5b50ux7684ux8fdcux573a}

\textbf{题号：M2-033}

\textbf{层级：} 复赛　　\textbf{方法：} 多极展开

所有电荷位于真空，\(q>0,\ell>0\)；以原点为球坐标中心，\(+z\)
为极轴，取无穷远零势。远场比较使用 \(r\gg\ell\)。

三个点电荷 \(+q,-2q,+q\) 分别位于 \(z=\ell,0,-\ell\)。求：

\begin{enumerate}
\def\labelenumi{\arabic{enumi}.}
\tightlist
\item
  精确电势；
\item
  第一个非零多极项及其相对修正阶数；
\item
  四极近似下的 \(E_r\)、\(E_{\theta}\)；
\item
  四极首项的零势锥，并说明该锥面上电场是否为零。
\end{enumerate}

\par\noindent\begin{minipage}{\linewidth}
\textbf{解答}

\begin{center}
\includegraphics[width=\linewidth,height=.40\textheight,keepaspectratio]{\detokenize{assets/module-02/examples/M2-033.pdf}}
\par\smallskip{\small 图 M2-033：线型电四极子的远场}
\end{center}
\end{minipage}\par

图 M2-033
的三个源沿同一轴线且关于中心对称，奇数阶全部抵消。精确距离展开后取第一个非零偶数阶，再对径向和极角分别求导；在四极电势为零的锥面上，角向导数仍可非零，所以电场不随电势首项一起消失。

精确电势为

\[
\begin{aligned}
\varphi_{\mathrm{\text{精确}}}
=\frac{q}{4\pi\varepsilon_0}
\Bigg[&
\frac1{\sqrt{r^2+\ell^2-2r\ell\cos\theta}}
+\frac1{\sqrt{r^2+\ell^2+2r\ell\cos\theta}}
-\frac2r
\Bigg].
\end{aligned}
\]

总电荷为零，偶极矩也为零；反演偶对称又使所有奇数阶消失。因此展开从
\(L=2\) 开始：

\[
\varphi
=\frac{2q}{4\pi\varepsilon_0}
\sum_{n=1}^{\infty}
\frac{\ell^{2n}}{r^{2n+1}}P_{2n}(\cos\theta).
\]

首项为

\[
\boxed{
\varphi^{(2)}
=\frac{q\ell^2}{4\pi\varepsilon_0r^3}
(3\cos^2\theta-1)
}.
\]

下一项是 \(L=4\)，在首项不为零的一般方向上相对修正为 \(O((\ell/r)^2)\)。

求梯度：

\[
\boxed{
E_r
=\frac{3q\ell^2}{4\pi\varepsilon_0r^4}
(3\cos^2\theta-1)
},
\]

\[
\boxed{
E_{\theta}
=\frac{6q\ell^2}{4\pi\varepsilon_0r^4}
\sin\theta\cos\theta
}.
\]

零势角满足

\[
\cos^2\theta_0=\frac13.
\]

在该角度 \(\varphi^{(2)}=E_r^{(2)}=0\)，但
\(E_{\theta}^{(2)}\ne0\)，所以零势锥不是零场面。有限体系的高阶项还会使精确零势面略微偏离该锥面。

\subsubsection{作业二：两个等量正电荷的单极项与四极修正}\label{m2:ux4f5cux4e1aux4e8cux4e24ux4e2aux7b49ux91cfux6b63ux7535ux8377ux7684ux5355ux6781ux9879ux4e0eux56dbux6781ux4feeux6b63}

\textbf{题号：M2-034}

\textbf{层级：} 复赛　　\textbf{方法：} 多极展开

电荷位于真空，\(q>0,a>0\)；以两电荷中点为展开原点，\(+z\)
为极轴，\(\theta\) 自该轴量起，零势取在无穷远。

两个点电荷 \(+q\) 位于 \(z=\pm a\)。求远场电势至 \(r^{-3}\)
阶，并比较轴线和赤道面上四极修正相对单极项的大小。

\par\noindent\begin{minipage}{\linewidth}
\textbf{解答}

\begin{center}
\includegraphics[width=\linewidth,height=.40\textheight,keepaspectratio]{\detokenize{assets/module-02/examples/M2-034.pdf}}
\par\smallskip{\small 图 M2-034：两个等量正电荷的单极项与四极修正}
\end{center}
\end{minipage}\par

图 M2-034
的两个同号电荷使总电荷加倍，而位置的一次矩抵消。以同一个半径比较轴向和赤道方向时，四极项的
\(P_2\) 因子分别为1与 \(-1/2\)，由此决定相对单极项的修正符号和大小。

精确电势为

\[
\varphi
=\frac{q}{4\pi\varepsilon_0}
\left[
\frac1{|\mathbf r-a\mathbf e_z|}
+\frac1{|\mathbf r+a\mathbf e_z|}
\right].
\]

两项相加时奇数阶消失：

\[
\varphi
=\frac{2q}{4\pi\varepsilon_0}
\left[
\frac1r
+\frac{a^2}{r^3}P_2(\cos\theta)
+O\!\left(\frac{a^4}{r^5}\right)
\right].
\]

因此

\[
\boxed{
\varphi
=\frac1{4\pi\varepsilon_0}
\left[
\frac{2q}{r}
+\frac{qa^2}{r^3}(3\cos^2\theta-1)
\right]
+O\!\left(\frac{qa^4}{4\pi\varepsilon_0r^5}\right)
}.
\]

轴线上 \(\cos\theta=\pm1\)，故

\[
\frac{\varphi^{(2)}}{\varphi^{(0)}}
=\frac{2qa^2/r^3}{2q/r}
=\frac{a^2}{r^2}.
\]

赤道面上 \(\cos\theta=0\)，故

\[
\frac{\varphi^{(2)}}{\varphi^{(0)}}
=-\frac{qa^2/r^3}{2q/r}
=-\frac12\frac{a^2}{r^2}.
\]

该系统虽然偶极矩为零，但总电荷不为零，因此远场仍由单极项主导，四极项只是修正。

\subsubsection{作业三：平面四电荷的四极张量}\label{m2:ux4f5cux4e1aux4e09ux5e73ux9762ux56dbux7535ux8377ux7684ux56dbux6781ux5f20ux91cf}

\textbf{题号：M2-035}

\textbf{层级：} 复赛　　\textbf{方法：} 多极展开

四个电荷位于真空，\(q>0,a>0\)，以给定坐标原点展开。采用对称无迹约定
\(\mathcal Q_{ij}=\sum_{\nu} q_{\nu}(3x_{\nu i}x_{\nu j}-r_{\nu}^2\delta_{ij})\)，场点方向为
\(\mathbf n=\mathbf r/r\)。

\(+q\) 位于 \((\pm a,0,0)\)，\(-q\) 位于 \((0,\pm a,0)\)。求
\(Q\)、\(\mathbf p\)、\(\mathcal Q_{ij}\) 和远场四极电势。

\par\noindent\begin{minipage}{\linewidth}
\textbf{解答}

\begin{center}
\includegraphics[width=\linewidth,height=.40\textheight,keepaspectratio]{\detokenize{assets/module-02/examples/M2-035.pdf}}
\par\smallskip{\small 图 M2-035：平面四电荷的四极张量}
\end{center}
\end{minipage}\par

图 M2-035 中 \(x\) 轴正电荷与 \(y\)
轴负电荷具有相同半径。总电荷和偶极矩均为零，四极张量的 \(xx,yy\)
分量等大反号，\(zz\) 分量为零；与观测方向二次投影后得到四瓣角分布。

总电荷为

\[
Q=2q-2q=0.
\]

每对同号电荷关于原点对称，所以

\[
\mathbf p=\mathbf0.
\]

按

\[
\mathcal Q_{ij}
=\sum_s q_s(3x_{s,i}x_{s,j}-r_s^2\delta_{ij})
\]

逐项求和。所有非对角元均为零，并有

\[
\mathcal Q_{xx}=6qa^2,
\qquad
\mathcal Q_{yy}=-6qa^2,
\qquad
\mathcal Q_{zz}=0.
\]

因此

\[
\boxed{
\mathcal Q
=6qa^2
\begin{pmatrix}
1&0&0\\
0&-1&0\\
0&0&0
\end{pmatrix}
}.
\]

代入四极电势：

\[
\begin{aligned}
\varphi^{(2)}
&=\frac1{4\pi\varepsilon_0}
\frac{3qa^2}{r^3}(n_x^2-n_y^2)\\
&=\boxed{
\frac{3qa^2}{4\pi\varepsilon_0r^3}
\sin^2\theta\cos2\phi
}.
\end{aligned}
\]

\subsubsection{作业四：证明四极矩的原点不变条件}\label{m2:ux4f5cux4e1aux56dbux8bc1ux660eux56dbux6781ux77e9ux7684ux539fux70b9ux4e0dux53d8ux6761ux4ef6}

\textbf{题号：M2-036}

\textbf{层级：} 复赛　　\textbf{方法：} 多极展开

用源坐标 \(\mathbf r_{\nu}\) 定义
\(\mathbf p=\sum q_{\nu}\mathbf r_{\nu}\) 及
\(\mathcal Q_{ij}=\sum q_{\nu}(3x_{\nu i}x_{\nu j}-r_{\nu}^2\delta_{ij})\)；平移后
\(\mathbf r_{\nu}'=\mathbf r_{\nu}-\mathbf a\)。连续分布可将求和改为积分。

已知一个局域电荷系统满足
\(Q=0\)、\(\mathbf p=0\)。证明平移坐标原点后，其对称无迹四极矩
\(\mathcal Q_{ij}\) 不变。

\par\noindent\begin{minipage}{\linewidth}
\textbf{解答}

\begin{center}
\includegraphics[width=\linewidth,height=.40\textheight,keepaspectratio]{\detokenize{assets/module-02/examples/M2-036.pdf}}
\par\smallskip{\small 图 M2-036：证明四极矩的原点不变条件}
\end{center}
\end{minipage}\par

图 M2-036
的坐标平移只改变同一组源点的位置坐标。展开新四极矩时，附加项分别正比于总电荷和偶极矩；由于题设两者为零，所有平移项同时消失，不能把这一不变性推广到任意带净电荷体系。

原点平移 \(\mathbf a\) 后，\(x_i'=x_i-a_i\)。于是

\[
\begin{aligned}
\mathcal Q_{ij}'
=\int\rho\Bigl[&3(x_i-a_i)(x_j-a_j)\\
&-|\mathbf r-\mathbf a|^2\delta_{ij}\Bigr]dV.
\end{aligned}
\]

展开并整理：

\[
\begin{aligned}
\mathcal Q_{ij}'
= {}&\mathcal Q_{ij}
-3(a_ip_j+a_jp_i)
+2(\mathbf a\cdot\mathbf p)\delta_{ij}\\
&+Q(3a_ia_j-a^2\delta_{ij}).
\end{aligned}
\]

代入 \(Q=0\)、\(\mathbf p=0\)，立即得到

\[
\boxed{\mathcal Q_{ij}'=\mathcal Q_{ij}}.
\]

若只满足 \(Q=0\) 而 \(\mathbf p\ne0\)，四极矩一般仍会随原点改变。

\subsubsection{作业五：截断误差界的使用}\label{m2:ux4f5cux4e1aux4e94ux622aux65adux8befux5deeux754cux7684ux4f7fux7528}

\textbf{题号：M2-037}

\textbf{层级：} 复赛　　\textbf{方法：} 多极展开

源区球心即多极展开原点，\(r>a\)；定义
\(Q_{\rm \text{绝}}=\int|\rho|\,dV\)，点电荷情形改为
\(\sum|q_i|\)。这里要求电势绝对误差相对于指定正尺度的界。

某局域电荷分布包含在半径 \(a\) 内，场点满足 \(a/r=1/5\)。若保留到四极项
\(L=2\)，估计遗漏电势相对于尺度

\[
\frac{Q_{\text{绝}}}{4\pi\varepsilon_0r}
\]

的上界，并说明它是否等于总电势的相对误差。

\par\noindent\begin{minipage}{\linewidth}
\textbf{解答}

\begin{center}
\includegraphics[width=\linewidth,height=.40\textheight,keepaspectratio]{\detokenize{assets/module-02/examples/M2-037.pdf}}
\par\smallskip{\small 图 M2-037：截断误差界的使用}
\end{center}
\end{minipage}\par

图 M2-037 将源区与外部场点分开，控制参数为
\(\eta=1/5\)。保留四极项后，从三阶开始的尾和上界为
\(\eta^3/(1-\eta)=0.01\)；该数乘以绝对电荷尺度，不能除以可能因相消而很小的总势当作同一相对误差。

取 \(\eta=1/5\)、\(L=2\)：

\[
\frac{\eta^{L+1}}{1-\eta}
=\frac{(1/5)^3}{1-1/5}
=\frac{1/125}{4/5}
=\frac1{100}.
\]

所以

\[
\boxed{
|\Delta\varphi_2|
\le
\frac1{100}
\frac{Q_{\text{绝}}}{4\pi\varepsilon_0r}
}.
\]

这只是以总绝对电荷构造的绝对误差界，不一定是总电势的相对误差。若低阶项因正负电荷抵消或角向节点而很小，总电势的相对误差可能远大于
\(1\%\)，此时应直接比较下一非零项或使用精确式。

\section{专题十二：电偶极子与多极展开练习（拔高选讲）}\label{m2:ux4e13ux9898ux5341ux4e8cux7535ux5076ux6781ux5b50ux4e0eux591aux6781ux5c55ux5f00ux7ec3ux4e60ux62d4ux9ad8ux9009ux8bb2}

以下练习检验多极展开中的几个基本判断：远区条件是否成立，最低非零多极矩是哪一阶，以及角向节点是否同时消去电场主阶。带净电荷时，还要区分坐标原点与等效电荷中心。

\par\noindent\begin{minipage}{\linewidth}
\begin{center}
\includegraphics[width=\linewidth,height=.40\textheight,keepaspectratio]{\detokenize{assets/module-02/original/tikz-m2-topic12-v8-01-margin12-v11.png}}
\par\smallskip{\small 图示：多极综合题的适用范围与判阶}
\end{center}
\end{minipage}\par

\subsection{一、远区判阶：适用范围先于矩量计算}\label{m2:ux4e00ux8fdcux533aux5224ux9636ux9002ux7528ux8303ux56f4ux5148ux4e8eux77e9ux91cfux8ba1ux7b97}

设局域电荷分布完全包含在以所选原点 \(O\) 为球心、半径为 \(a\)
的球内。对场点 \(\mathbf r=r\mathbf n\)，有

\[
\varphi(\mathbf r)
=\frac{1}{4\pi\varepsilon_0}
\left[
\frac{Q}{r}
+\frac{\mathbf p\cdot\mathbf n}{r^2}
+\frac{n_i\mathcal Q_{ij}n_j}{2r^3}
+O\!\left(\frac{Q_{\mathrm{\text{绝}}}a^3}{r^4}\right)
\right],
\]

其中

\[
Q=\int\rho\,dV,
\qquad
\mathbf p=\int\rho\mathbf r'\,dV,
\]

\[
\mathcal Q_{ij}
=\int\rho(3x_i'x_j'-r'^2\delta_{ij})\,dV,
\qquad
Q_{\mathrm{\text{绝}}}=\int|\rho|\,dV.
\]

这里必须区分两个条件：

\begin{enumerate}
\def\labelenumi{\arabic{enumi}.}
\tightlist
\item
  \(r>a\) 保证外部多极级数对整个源区成立；
\item
  \(a/r\ll1\) 才保证只保留少数低阶项通常具有较高精度。
\end{enumerate}

若第一个非零多极矩的阶数为 \(L\)，则一般方向上

\[
\varphi^{(L)}\propto r^{-L-1},
\qquad
\mathbf E^{(L)}\propto r^{-L-2}.
\]

因此判阶顺序是

\[
Q\longrightarrow\mathbf p\longrightarrow\mathcal Q_{ij}
\longrightarrow\text{更高阶矩}.
\]

\par\noindent\begin{minipage}{\linewidth}
\begin{center}
\includegraphics[width=\linewidth,height=.40\textheight,keepaspectratio]{\detokenize{assets/module-02/original/tikz-m2-topic12-v8-02.png}}
\par\smallskip{\small 图示：单极、偶极与四极主导条件及其衰减阶次}
\end{center}
\end{minipage}\par

{\def\LTcaptype{none} % do not increment counter
\begin{longtable}[]{@{}
  >{\raggedright\arraybackslash}p{(\linewidth - 6\tabcolsep) * \real{0.2500}}
  >{\raggedright\arraybackslash}p{(\linewidth - 6\tabcolsep) * \real{0.2500}}
  >{\raggedright\arraybackslash}p{(\linewidth - 6\tabcolsep) * \real{0.2500}}
  >{\raggedright\arraybackslash}p{(\linewidth - 6\tabcolsep) * \real{0.2500}}@{}}
\toprule\noalign{}
\begin{minipage}[b]{\linewidth}\raggedright
第一个非零矩
\end{minipage} & \begin{minipage}[b]{\linewidth}\raggedright
主导电势
\end{minipage} & \begin{minipage}[b]{\linewidth}\raggedright
主导电场
\end{minipage} & \begin{minipage}[b]{\linewidth}\raggedright
必要条件
\end{minipage} \\
\midrule\noalign{}
\endhead
\bottomrule\noalign{}
\endlastfoot
单极矩 \(Q\) & \(r^{-1}\) & \(r^{-2}\) & \(Q\ne0\) \\
偶极矩 \(\mathbf p\) & \(r^{-2}\) & \(r^{-3}\) &
\(Q=0,\ \mathbf p\ne0\) \\
四极矩 \(\mathcal Q\) & \(r^{-3}\) & \(r^{-4}\) &
\(Q=0,\ \mathbf p=0,\ \mathcal Q\ne0\) \\
\end{longtable}
}

\textbf{例 M2-038：偶极项为零不等于纯四极子}

真空 \(z\) 轴的 \(z=\pm a\) 处各固定一个电荷 \(+q\)，其中
\(a>0,q>0\)。以中点为展开原点，取 \(\varphi(\infty)=0\)，只考察
\(r\gg a\)
的远区。利用关于原点的反演对称性筛选多极阶数，求总电荷与偶极矩，并判断偶极项消失是否足以将该体系称为纯四极子。

\par\noindent\begin{minipage}{\linewidth}
\textbf{解答}

\begin{center}
\includegraphics[width=\linewidth,height=.40\textheight,keepaspectratio]{\detokenize{assets/module-02/examples/M2-038.pdf}}
\par\smallskip{\small 图 M2-038：偶极项为零不等于纯四极子}
\end{center}
\end{minipage}\par

图 M2-038
的两电荷同号且关于原点对称，间距标记不是非零偶极矩。反演对称性使奇数阶消失，但总电荷仍为
\(2q\)；所以场从单极阶开始，四极只是下一修正。

两个 \(+q\) 位于 \(z=\pm a\)，虽然关于原点的偶极矩为零，但总电荷为
\(2q\)，所以远区仍由单极项主导。偶极项消失只意味着下一修正从四极项开始，并不意味着该体系是``纯四极子''。

若体系具有反演对称性，还可在计算前筛去整组阶数：

\[
\rho(-\mathbf r)=\rho(\mathbf r)
\quad\Longrightarrow\quad
L=1,3,5,\ldots\text{ 全部消失},
\]

\[
\rho(-\mathbf r)=-\rho(\mathbf r)
\quad\Longrightarrow\quad
L=0,2,4,\ldots\text{ 全部消失}.
\]

对称性只能说明某些阶必为零；余下的最低阶是否非零，仍需继续计算。

\subsection{二、带净电荷体系：等效电荷中心}\label{m2:ux4e8cux5e26ux51c0ux7535ux8377ux4f53ux7cfbux7b49ux6548ux7535ux8377ux4e2dux5fc3}

设新原点 \(O'\) 相对旧原点 \(O\) 的位移为
\(\mathbf a\)。同一源点在新坐标中的位置为

\[
\mathbf r_i'=\mathbf r_i-\mathbf a.
\]

总电荷不变，而偶极矩变为

\[
\begin{aligned}
\mathbf p'
&=\sum_iq_i\mathbf r_i'\\
&=\sum_iq_i(\mathbf r_i-\mathbf a)\\
&=\mathbf p-Q\mathbf a.
\end{aligned}
\]

即

\[
\boxed{\mathbf p'=\mathbf p-Q\mathbf a}.
\]

\par\noindent\begin{minipage}{\linewidth}
\begin{center}
\includegraphics[width=\linewidth,height=.40\textheight,keepaspectratio]{\detokenize{assets/module-02/original/tikz-m2-topic12-v8-03.png}}
\par\smallskip{\small 图示：原点平移时源点坐标与偶极矩的变化}
\end{center}
\end{minipage}\par

当 \(Q\ne0\) 时，可以选择

\[
\boxed{\mathbf a_c=\frac{\mathbf p}{Q}}
\]

使新原点下

\[
\mathbf p^{(c)}=\mathbf0.
\]

\(O_c\)
称为带净电荷体系的等效电荷中心。它的含义是：在单极项和偶极项的精度内，原体系与位于
\(O_c\) 的点电荷 \(Q\) 具有相同远场。

要严格验证这一点，考虑位于 \(\mathbf a_c\) 的点电荷 \(Q\)：

\[
\varphi_c(\mathbf r)
=\frac{1}{4\pi\varepsilon_0}
\frac{Q}{|\mathbf r-\mathbf a_c|}.
\]

当 \(r\gg a_c\) 时，

\[
\frac{1}{|\mathbf r-\mathbf a_c|}
=\frac1r
+\frac{\mathbf a_c\cdot\mathbf n}{r^2}
+O\!\left(\frac{a_c^2}{r^3}\right).
\]

于是

\[
\varphi_c(\mathbf r)
=\frac{1}{4\pi\varepsilon_0}
\left[
\frac Qr
+\frac{Q\mathbf a_c\cdot\mathbf n}{r^2}
+O\!\left(\frac{Qa_c^2}{r^3}\right)
\right].
\]

代入 \(Q\mathbf a_c=\mathbf p\)，单极项与偶极项恰好同原体系一致。

\textbf{例 M2-039：等效中心可以位于实际源区之外}

在真空中，\(+3q\) 固定于 \(x=\ell\)，\(-q\) 固定于 \(x=-\ell\)，其中
\(q>0,\ell>0\)。以 \(x=0\)
为初始展开原点，计算总电荷和电偶极矩。若新原点沿 \(x\) 轴平移到
\(x_c\)，求使新偶极矩为零的唯一
\(x_c\)，并说明所得中心在实际两电荷区间内还是区间外。

\par\noindent\begin{minipage}{\linewidth}
\textbf{解答}

\begin{center}
\includegraphics[width=\linewidth,height=.40\textheight,keepaspectratio]{\detokenize{assets/module-02/examples/M2-039.pdf}}
\par\smallskip{\small 图 M2-039：等效中心可以位于实际源区之外}
\end{center}
\end{minipage}\par

图 M2-039 的 \(O_c\) 由一次矩与总电荷之比确定，位置为
\(2\ell\)。这一点落在原源区外不违反任何物理条件，因为它是远区展开的参考中心，并非电荷的真实支撑点。

\(+3q\) 位于 \(x=\ell\)，\(-q\) 位于 \(x=-\ell\)。以中点为原点，

\[
Q=2q,
\qquad
p_x=3q\ell+q\ell=4q\ell,
\]

所以

\[
x_c=\frac{p_x}{Q}=2\ell.
\]

等效中心位于原电荷区之外并不矛盾，因为它是远区展开的一阶等效点，不是实际电荷所在的位置。

\subsection{三、消去偶极项以后：残余四极矩与收敛半径}\label{m2:ux4e09ux6d88ux53bbux5076ux6781ux9879ux4ee5ux540eux6b8bux4f59ux56dbux6781ux77e9ux4e0eux6536ux655bux534aux5f84}

等效电荷中心只消去偶极项，不能消去体系的全部形状信息。四极矩在原点平移下满足

\[
\begin{aligned}
\mathcal Q_{ij}'={}&\mathcal Q_{ij}
-3(a_ip_j+a_jp_i)
+2(\mathbf a\cdot\mathbf p)\delta_{ij}\\
&+Q(3a_ia_j-a^2\delta_{ij}).
\end{aligned}
\]

令 \(\mathbf a=\mathbf a_c=\mathbf p/Q\)，得到等效中心处的残余四极矩

\[
\boxed{
\mathcal Q_{ij}^{(c)}
=\mathcal Q_{ij}
-\frac{3p_ip_j-p^2\delta_{ij}}{Q}
}.
\]

\par\noindent\begin{minipage}{\linewidth}
\begin{center}
\includegraphics[width=\linewidth,height=.40\textheight,keepaspectratio]{\detokenize{assets/module-02/original/tikz-m2-topic12-v8-04.png}}
\par\smallskip{\small 图示：等效电荷中心消去偶极项但保留残余四极矩}
\end{center}
\end{minipage}\par

若用

\[
\mathbf R=\mathbf r-\mathbf a_c,
\qquad
R=|\mathbf R|,
\qquad
\mathbf N=\frac{\mathbf R}{R}
\]

表示场点相对等效中心的位置，则

\[
\varphi(\mathbf r)
=\frac{1}{4\pi\varepsilon_0}
\left[
\frac QR
+\frac{N_i\mathcal Q_{ij}^{(c)}N_j}{2R^3}
+O\!\left(\frac{Q_{\mathrm{\text{绝}}}b^3}{R^4}\right)
\right],
\]

其中 \(b\) 是以 \(O_c\) 为球心包围全部源电荷所需的半径。

这里有一个容易忽略的事实：改变展开中心会同时改变源区半径。消去偶极项不一定提高截断精度；若
\(O_c\) 距离真实电荷区较远，新的 \(b\) 可能比旧的 \(a\) 大得多。

\textbf{例 M2-040：消去偶极项后仍有残余四极矩}

真空中 \(+3q\) 位于 \(x=\ell\)、\(-q\) 位于 \(x=-\ell\)，其中
\(q>0,\ell>0\)。现将多极展开中心移到 \(x_c=2\ell\)，以 \(x_i'=x_i-x_c\)
表示源点的新坐标。求相对于新中心的总电荷、偶极矩及最低残余高阶矩，判断消去偶极项后，该体系能否在整个外部空间严格等同于位于新中心的单个点电荷。

\par\noindent\begin{minipage}{\linewidth}
\textbf{解答}

\begin{center}
\includegraphics[width=\linewidth,height=.40\textheight,keepaspectratio]{\detokenize{assets/module-02/examples/M2-040.pdf}}
\par\smallskip{\small 图 M2-040：消去偶极项后仍有残余四极矩}
\end{center}
\end{minipage}\par

图 M2-040 保留两枚真实电荷，同时标出移到 \(2\ell\)
的展开中心。新中心处一次矩消失，但对相对新中心的坐标重新求二次矩仍得到非零张量；所以远区局部切线近似指向该中心，不代表近场完全成为点电荷场。

旧原点下

\[
\mathcal Q
=q\ell^2
\begin{pmatrix}
4&0&0\\
0&-2&0\\
0&0&-2
\end{pmatrix},
\qquad
\mathbf p=4q\ell\mathbf e_x,
\qquad
Q=2q.
\]

代入残余四极矩公式，得到

\[
\mathcal Q^{(c)}
=q\ell^2
\begin{pmatrix}
-12&0&0\\
0&6&0\\
0&0&6
\end{pmatrix}
\ne0.
\]

因此该体系并不严格等同于位于 \(x=2\ell\) 的点电荷
\(2q\)；两者只在单极和偶极精度内等效。

对于沿 \(x\) 轴的带净电荷体系，若极轴也取为 \(x\) 轴，则保留至偶极项时

\[
E_r
=\frac{1}{4\pi\varepsilon_0}
\left(
\frac{Q}{r^2}
+\frac{2p\cos\theta}{r^3}
\right),
\]

\[
E_{\theta}
=\frac{1}{4\pi\varepsilon_0}
\frac{p\sin\theta}{r^3}.
\]

场线局部切线相对旧原点径向方向的偏角 \(\beta\) 满足

\[
\tan\beta=\frac{E_{\theta}}{E_r}.
\]

在 \(r\) 足够大时，

\[
\beta\simeq\frac{p\sin\theta}{Qr}
=\frac{a_c\sin\theta}{r}.
\]

这正是从 \(O_c\)
指向远处场点的直线相对旧原点径向方向的几何偏角。因此远区场线的局部切线反向延长后，近似交于等效电荷中心。

\par\noindent\begin{minipage}{\linewidth}
\begin{center}
\includegraphics[width=\linewidth,height=.40\textheight,keepaspectratio]{\detokenize{assets/module-02/original/tikz-m2-topic12-v8-05.png}}
\par\smallskip{\small 图示：远区局部切线反向延长至等效电荷中心}
\end{center}
\end{minipage}\par

虚线只是远区切线的反向延长，不代表真实近场场线。真实场线仍由实际电荷分布决定；残余四极矩和更高阶矩会使它偏离严格直线。

\subsection{四、中性体系、对称性与角向节点}\label{m2:ux56dbux4e2dux6027ux4f53ux7cfbux5bf9ux79f0ux6027ux4e0eux89d2ux5411ux8282ux70b9}

当 \(Q=0\) 时，

\[
\mathbf p'=\mathbf p.
\]

因此中性体系的偶极矩与原点无关。只要
\(\mathbf p\ne0\)，就不可能通过移动原点消去偶极主项。

若进一步有
\(\mathbf p=0\)，则四极矩也与原点无关。由此可见，第一个非零多极矩具有内禀意义：所有更低阶矩都为零时，它不再依赖原点选择。

\par\noindent\begin{minipage}{\linewidth}
\begin{center}
\includegraphics[width=\linewidth,height=.40\textheight,keepaspectratio]{\detokenize{assets/module-02/original/tikz-m2-topic12-v8-06.png}}
\par\smallskip{\small 图示：同号电荷、异号电荷与线型三电荷的远区主导项对比}
\end{center}
\end{minipage}\par

\subsubsection{1.
全局主导项与方向主导项}\label{m2:ux5168ux5c40ux4e3bux5bfcux9879ux4e0eux65b9ux5411ux4e3bux5bfcux9879}

``四极项是全局首项''不等于``四极电势在每个方向都非零''。以线型三电荷

\[
+q\text{ 位于 }z=\ell,
\qquad
-2q\text{ 位于 }z=0,
\qquad
+q\text{ 位于 }z=-\ell
\]

为例。对 \(r>\ell\)，精确外部级数为

\[
\varphi(r,\theta)
=\frac{2q}{4\pi\varepsilon_0}
\sum_{n=1}^{\infty}
\frac{\ell^{2n}}{r^{2n+1}}P_{2n}(\cos\theta).
\]

一般方向上的首项是

\[
\varphi^{(2)}
=\frac{2q\ell^2}{4\pi\varepsilon_0r^3}
P_2(\cos\theta)
=\frac{q\ell^2}{4\pi\varepsilon_0r^3}
(3\cos^2\theta-1).
\]

当

\[
\cos^2\theta_0=\frac13
\]

时，\(P_2(\cos\theta_0)=0\)。此时不能说``电势就是零''，而要保留下一个允许的偶数阶：

\[
P_4(\cos\theta_0)=-\frac7{18}.
\]

所以在这条方向上

\[
\boxed{
\varphi(r,\theta_0)
=-\frac{7q\ell^4}{9(4\pi\varepsilon_0)r^5}
+O\!\left(\frac{q\ell^6}{\varepsilon_0r^7}\right)
}.
\]

\par\noindent\begin{minipage}{\linewidth}
\begin{center}
\includegraphics[width=\linewidth,height=.40\textheight,keepaspectratio]{\detokenize{assets/module-02/original/tikz-m2-topic12-v8-07.png}}
\par\smallskip{\small 图示：线型四极子的角向节点与节点方向的下一阶电势}
\end{center}
\end{minipage}\par

但是，四极近似的角向电场为

\[
E_{\theta}^{(2)}
=\frac{6q\ell^2}{4\pi\varepsilon_0r^4}
\sin\theta\cos\theta.
\]

在 \(\theta=\theta_0\) 处它仍不为零。因此：

\[
\boxed{
\text{某阶电势在一个方向为零，不能推出该阶电场也为零。}
}
\]

\subsubsection{2.
节点附近的误差判断}\label{m2:ux8282ux70b9ux9644ux8fd1ux7684ux8befux5deeux5224ux65ad}

若把库仑核截断到 \(L\) 阶，记 \(\eta=a/r<1\)，则绝对误差满足

\[
\boxed{
|\Delta\varphi_L|
\le
\frac{1}{4\pi\varepsilon_0}
\frac{Q_{\mathrm{\text{绝}}}}{r}
\frac{\eta^{L+1}}{1-\eta}
}.
\]

在一般方向上，可以把下一非零阶与首项比较，估计相对误差；在首项的角向节点附近，分母趋近于零，这种相对误差估计会失去意义，应改用绝对误差或直接多保留几阶。

\subsection{五、从多极电势到场线几何}\label{m2:ux4e94ux4eceux591aux6781ux7535ux52bfux5230ux573aux7ebfux51e0ux4f55}

对轴对称的单一 \(L\) 阶理想多极，写成

\[
\varphi_L(r,\theta)
=\frac{C_L}{r^{L+1}}P_L(\cos\theta).
\]

求梯度得到

\[
E_r
=\frac{(L+1)C_L}{r^{L+2}}P_L(\cos\theta),
\]

\[
E_{\theta}
=\frac{C_L}{r^{L+2}}
\sin\theta\,P_L'(\cos\theta).
\]

在子午面内，场线切向与电场平行，因此

\[
\boxed{
\frac{dr}{r\,d\theta}
=\frac{E_r}{E_{\theta}}
=\frac{(L+1)P_L(\cos\theta)}
{\sin\theta\,P_L'(\cos\theta)}
}.
\]

\subsubsection{1. 单极场}\label{m2:ux5355ux6781ux573a}

\(L=0\) 时 \(E_{\theta}=0\)，场线是从中心向外或向内的径向射线。

\subsubsection{2. 偶极场}\label{m2:ux5076ux6781ux573a}

\(L=1\) 时 \(P_1(\cos\theta)=\cos\theta\)，于是

\[
\frac{dr}{r}=2\cot\theta\,d\theta.
\]

积分得到

\[
\boxed{r=C\sin^2\theta}.
\]

\subsubsection{3. 四极场}\label{m2:ux56dbux6781ux573a}

\(L=2\) 时

\[
P_2(\cos\theta)=\frac12(3\cos^2\theta-1),
\]

因而在 \(\sin\theta\cos\theta\ne0\) 的区域

\[
\frac{dr}{r}
=\left(\cot\theta-\frac12\tan\theta\right)d\theta.
\]

积分得到

\[
\boxed{
r=C\sin\theta\sqrt{|\cos\theta|}
}.
\]

\par\noindent\begin{minipage}{\linewidth}
\begin{center}
\includegraphics[width=\linewidth,height=.40\textheight,keepaspectratio]{\detokenize{assets/module-02/original/tikz-m2-topic12-v8-08.png}}
\par\smallskip{\small 图示：单极、偶极和四极理想场线的子午截面}
\end{center}
\end{minipage}\par

推导中除去了 \(\sin\theta=0\) 或 \(\cos\theta=0\)
的情形，所以对四极场还应单独补上对称轴和赤道面内的径向场线。

这些公式描述的是单一理想多极项。有限电荷体系只在远区近似服从相应形状；接近源区后必须回到精确电场。理想多极场线在原点汇合，是因为原点被压缩为奇异源，并不表示静电场线可以在光滑无源区闭合。

场线给出电场方向，也不等同于带电粒子的运动轨迹。粒子轨迹还取决于初速度、质量和速度随时间的变化。

\subsection{六、均匀带电圆环}\label{m2:ux516dux5747ux5300ux5e26ux7535ux5706ux73af}

半径为 \(a\) 的均匀细圆环位于 \(xy\) 平面，总电荷为
\(Q\)，圆心取为原点。该分布关于原点反演对称，又关于 \(z\)
轴旋转对称，因此：

\[
\mathbf p=0,
\]

所有奇数阶消失，且只有 \(m=0\) 的轴对称分量保留。

轴对称多极系数为

\[
A_L=\int\rho(\mathbf r')r'^L
P_L(\cos\theta')\,dV'.
\]

圆环上 \(r'=a\)、\(\theta'=\pi/2\)，故

\[
A_0=Q,
\]

\[
A_2=Qa^2P_2(0)=-\frac12Qa^2.
\]

远区电势至四极修正为

\[
\boxed{
\varphi(r,\theta)
=\frac{1}{4\pi\varepsilon_0}
\left[
\frac Qr
-\frac{Qa^2}{2r^3}P_2(\cos\theta)
+O\!\left(\frac{Qa^4}{r^5}\right)
\right]
}.
\]

\par\noindent\begin{minipage}{\linewidth}
\begin{center}
\includegraphics[width=\linewidth,height=.40\textheight,keepaspectratio]{\detokenize{assets/module-02/original/tikz-m2-topic12-v8-09.png}}
\par\smallskip{\small 图示：均匀带电圆环的多极矩与轴线、赤道方向检验}
\end{center}
\end{minipage}\par

轴线上 \(r=z\)、\(\theta=0\)，精确电势为

\[
\varphi_{\mathrm{\text{轴}}}(z)
=\frac{1}{4\pi\varepsilon_0}
\frac{Q}{\sqrt{z^2+a^2}}.
\]

当 \(z\gg a\) 时，

\[
\varphi_{\mathrm{\text{轴}}}(z)
=\frac{Q}{4\pi\varepsilon_0z}
\left[
1-\frac12\frac{a^2}{z^2}
+O\!\left(\frac{a^4}{z^4}\right)
\right],
\]

与 \(P_2(1)=1\) 的多极结果完全一致。

在赤道远区，\(P_2(0)=-1/2\)，四极修正变为正值：

\[
\varphi_{\mathrm{\text{赤道}}}(r)
=\frac{1}{4\pi\varepsilon_0}
\left[
\frac Qr
+\frac{Qa^2}{4r^3}
+O\!\left(\frac{Qa^4}{r^5}\right)
\right].
\]

轴线和赤道方向的修正符号相反，反映了圆环相对点电荷的扁平形状。

\subsection{七、多极展开的选阶与误差检查}\label{m2:ux4e03ux591aux6781ux5c55ux5f00ux7684ux9009ux9636ux4e0eux8befux5deeux68c0ux67e5}

\par\noindent\begin{minipage}{\linewidth}
\begin{center}
\includegraphics[width=\linewidth,height=.40\textheight,keepaspectratio]{\detokenize{assets/module-02/original/tikz-m2-topic12-v8-10-margin12-v11.png}}
\par\smallskip{\small 图示：多极展开的选原点、判阶与误差检查}
\end{center}
\end{minipage}\par

\begin{enumerate}
\def\labelenumi{\arabic{enumi}.}
\tightlist
\item
  \textbf{先圈定源区。}
  没有给出源区半径和场点位置，就不能判断外部展开是否收敛。
\item
  \textbf{区分 \(r>a\) 与 \(r\gg a\)。}
  前者是可展开条件，后者是低阶近似条件。
\item
  \textbf{先用对称性，再计算矩量。}
  对称性可以筛阶，但不能自动保证余下最低阶非零。
\item
  \textbf{只由低阶矩命名。}
  三个点电荷可以是单极主导、偶极主导，也可以是四极主导。
\item
  \textbf{带净电荷时检查原点依赖。} \(Q\ne0\) 时 \(\mathbf p\)
  依赖原点；\(\mathbf a_c=\mathbf p/Q\) 只消去偶极项。
\item
  \textbf{等效中心不是真实电荷中心。}
  它可以位于源区之外，也不保证新的源区半径更小。
\item
  \textbf{中性体系不能靠平移消去偶极矩。} \(Q=0\) 时 \(\mathbf p\)
  与原点无关。
\item
  \textbf{固定四极矩约定。} 本讲义使用对称无迹张量
  \(\mathcal Q_{ij}=\int\rho(3x_i'x_j'-r'^2\delta_{ij})dV\)。
\item
  \textbf{区分全局首项与方向首项。}
  首项在角向节点为零时，要继续保留下一阶。
\item
  \textbf{电势为零不等于电场为零。} 还必须检查电势的角向和径向导数。
\item
  \textbf{节点附近不要滥用相对误差。} 应改用绝对误差或增加截断阶数。
\item
  \textbf{理想场线只用于远区。}
  接近真实电荷时必须返回精确场，场线也不是粒子轨迹。
\end{enumerate}

\subsection{八、练习与解答}\label{m2:ux516bux7ec3ux4e60ux4e0eux89e3ux7b54-2}

\subsubsection{作业一：三电荷体系的等效中心与残余四极矩}\label{m2:ux4f5cux4e1aux4e00ux4e09ux7535ux8377ux4f53ux7cfbux7684ux7b49ux6548ux4e2dux5fc3ux4e0eux6b8bux4f59ux56dbux6781ux77e9}

\textbf{题号：M2-041}

\textbf{层级：} 复赛　　\textbf{方法：} 多极展开、等效变换

所有电荷位于真空，\(\ell>0\)；球坐标以旧原点为中心，\(\theta\) 从 \(+x\)
轴量起，零势取在无穷远，远区须位于相应源区球之外。轴对称系数约定
\(A_L=\sum_iq_i x_i^L\)。

在 \(x\) 轴上，\(x=-\ell\) 处放置 \(-q\)，\(x=0\) 处放置
\(-q\)，\(x=\ell\) 处放置 \(+3q\)，其中
\(q>0\)。以原点为展开中心，极轴取为 \(x\) 轴。

\begin{enumerate}
\def\labelenumi{\arabic{enumi}.}
\tightlist
\item
  求 \(Q\)、\(\mathbf p\) 和远区电势至 \(r^{-3}\) 阶；
\item
  只保留单极与偶极项，求
  \(E_r\)、\(E_{\theta}\)，并说明远区场线切线反向延长后的交点；
\item
  求等效电荷中心及该中心处的残余四极矩；
\item
  比较以旧原点和等效中心展开时的源区半径，说明消去偶极项是否一定提高精度。
\end{enumerate}

\par\noindent\begin{minipage}{\linewidth}
\textbf{解答}

\begin{center}
\includegraphics[width=\linewidth,height=.40\textheight,keepaspectratio]{\detokenize{assets/module-02/examples/M2-041.pdf}}
\par\smallskip{\small 图 M2-041：三电荷体系的等效中心与残余四极矩}
\end{center}
\end{minipage}\par

图 M2-041 的三枚有符号电荷都要参加 \(Q,p_x,A_2\)
的求和。由一次矩确定等效中心后，应重新量取每个源点距新中心的最大距离；即使消去了偶极项，新的展开小参数也未必更小。

总电荷为

\[
Q=-q-q+3q=q.
\]

偶极矩沿 \(x\) 轴：

\[
\mathbf p
=\left[(-q)(-\ell)+(-q)(0)+(3q)\ell\right]\mathbf e_x
=4q\ell\mathbf e_x.
\]

对轴上点电荷，轴对称系数

\[
A_2=\sum_iq_ix_i^2
=(-q)\ell^2+0+(3q)\ell^2
=2q\ell^2.
\]

因此

\[
\boxed{
\varphi(r,\theta)
=\frac{1}{4\pi\varepsilon_0}
\left[
\frac qr
+\frac{4q\ell\cos\theta}{r^2}
+\frac{2q\ell^2}{r^3}P_2(\cos\theta)
+O\!\left(\frac{q\ell^3}{r^4}\right)
\right]
}.
\]

也可写成

\[
\varphi
=\frac{1}{4\pi\varepsilon_0}
\left[
\frac qr
+\frac{4q\ell\cos\theta}{r^2}
+\frac{q\ell^2(3\cos^2\theta-1)}{r^3}
+\cdots
\right].
\]

只保留单极与偶极项，求梯度得

\[
\boxed{
E_r
=\frac{1}{4\pi\varepsilon_0}
\left(
\frac q{r^2}
+\frac{8q\ell\cos\theta}{r^3}
\right)
},
\]

\[
\boxed{
E_{\theta}
=\frac{1}{4\pi\varepsilon_0}
\frac{4q\ell\sin\theta}{r^3}
}.
\]

远区偏角为

\[
\beta\simeq\frac{E_{\theta}}{E_r}
\simeq\frac{4\ell\sin\theta}{r}.
\]

若切线反向延长交于 \(x=h\)，几何上有

\[
\beta\simeq\frac{h\sin\theta}{r},
\]

所以

\[
\boxed{h=4\ell}.
\]

这也由等效中心公式直接得到：

\[
\boxed{
\mathbf a_c=\frac{\mathbf p}{Q}=4\ell\mathbf e_x
}.
\]

旧原点下的四极张量为

\[
\mathcal Q
=q\ell^2
\begin{pmatrix}
4&0&0\\
0&-2&0\\
0&0&-2
\end{pmatrix}.
\]

利用

\[
\mathcal Q_{ij}^{(c)}
=\mathcal Q_{ij}
-\frac{3p_ip_j-p^2\delta_{ij}}{Q},
\]

得到

\[
\boxed{
\mathcal Q^{(c)}
=q\ell^2
\begin{pmatrix}
-28&0&0\\
0&14&0\\
0&0&14
\end{pmatrix}
}.
\]

也可以直接用新坐标检验。相对 \(x_c=4\ell\)，三个电荷的位置依次是

\[
-5\ell,
\qquad
-4\ell,
\qquad
-3\ell.
\]

因此新中心下的源区半径为

\[
b=5\ell,
\]

而旧原点下只需

\[
a=\ell.
\]

新原点虽使偶极项消失，却把包围源区所需的半径从 \(\ell\) 增大到
\(5\ell\)。所以消去偶极项并不保证低阶截断更精确；还必须比较新的 \(b/R\)
与旧的 \(a/r\)。

\subsubsection{作业二：中性体系的偶极矩不能由平移消去}\label{m2:ux4f5cux4e1aux4e8cux4e2dux6027ux4f53ux7cfbux7684ux5076ux6781ux77e9ux4e0dux80fdux7531ux5e73ux79fbux6d88ux53bb}

\textbf{题号：M2-042}

\textbf{层级：} 复赛　　\textbf{方法：} 多极展开

所有电荷位于真空，\(q>0,\ell>0\)；以 \(+x\) 为极轴，令
\(A_L=\sum_iq_i x_i^L\)，新坐标为 \(x_i'=x_i-a\)。电势零点在无穷远。

在 \(x\) 轴上，\(x=2\ell\) 处放置 \(+q\)，\(x=-\ell\) 处放置
\(+q\)，原点处放置 \(-2q\)。

\begin{enumerate}
\def\labelenumi{\arabic{enumi}.}
\tightlist
\item
  求 \(Q\)、\(p_x\) 与原点下的 \(A_2\)；
\item
  写出远区电势至四极项；
\item
  把原点沿 \(x\) 轴平移 \(a\)，求 \(p_x'\) 与 \(A_2'\)；
\item
  能否选择原点同时消去偶极项和四极项？
\end{enumerate}

\par\noindent\begin{minipage}{\linewidth}
\textbf{解答}

\begin{center}
\includegraphics[width=\linewidth,height=.40\textheight,keepaspectratio]{\detokenize{assets/module-02/examples/M2-042.pdf}}
\par\smallskip{\small 图 M2-042：中性体系的偶极矩不能由平移消去}
\end{center}
\end{minipage}\par

图 M2-042 的三源总电荷为零，但不具有消去一次矩的对称性。平移时偶极矩因
\(Q=0\)
保持不变，而四极系数仍会变化；因此可以调某些高阶项，却不能同时把非零偶极项消去。

总电荷为

\[
Q=q+q-2q=0.
\]

偶极矩为

\[
p_x=q(2\ell)+q(-\ell)+(-2q)(0)=q\ell.
\]

所以该体系是偶极主导，而不是四极主导。

轴对称四极系数为

\[
A_2
=\sum_iq_ix_i^2
=q(2\ell)^2+q(-\ell)^2
=5q\ell^2.
\]

故

\[
\boxed{
\varphi(r,\theta)
=\frac{1}{4\pi\varepsilon_0}
\left[
\frac{q\ell\cos\theta}{r^2}
+\frac{5q\ell^2}{r^3}P_2(\cos\theta)
+O\!\left(\frac{q\ell^3}{r^4}\right)
\right]
}.
\]

原点平移后

\[
p_x'=p_x-Qa=q\ell.
\]

这说明中性体系的非零偶极矩不能由平移消去。

对 \(A_2\)，直接展开：

\[
\begin{aligned}
A_2'
&=\sum_iq_i(x_i-a)^2\\
&=\sum_iq_ix_i^2-2a\sum_iq_ix_i+a^2\sum_iq_i\\
&=A_2-2ap_x+a^2Q.
\end{aligned}
\]

代入本题结果：

\[
\boxed{A_2'=5q\ell^2-2aq\ell}.
\]

选择

\[
a=\frac52\ell
\]

可以使 \(A_2'=0\)，但仍有

\[
p_x'=q\ell\ne0.
\]

因此不能同时消去偶极项和四极项；无论选择什么原点，远区首项始终是同一个偶极项。

\subsubsection{作业三：线型四极子的节点方向}\label{m2:ux4f5cux4e1aux4e09ux7ebfux578bux56dbux6781ux5b50ux7684ux8282ux70b9ux65b9ux5411}

\textbf{题号：M2-043}

\textbf{层级：} 复赛　　\textbf{方法：} 多极展开、渐近展开

电荷位于真空，\(q>0,\ell>0\)；采用以原点为中心、\(+z\)
为极轴的球坐标，取 \(\varphi(\infty)=0\)。多极级数用于
\(r>\ell\)，渐近比较用于 \(r\gg\ell\)。

三个点电荷 \(+q,-2q,+q\) 分别位于 \(z=\ell,0,-\ell\)。

\begin{enumerate}
\def\labelenumi{\arabic{enumi}.}
\tightlist
\item
  写出精确外部多极级数；
\item
  求一般方向上的首项；
\item
  在 \(\cos^2\theta_0=1/3\) 的方向上，求电势的第一个非零项；
\item
  判断该方向上的电场是否也从相同高阶开始。
\end{enumerate}

\par\noindent\begin{minipage}{\linewidth}
\textbf{解答}

\begin{center}
\includegraphics[width=\linewidth,height=.40\textheight,keepaspectratio]{\detokenize{assets/module-02/examples/M2-043.pdf}}
\par\smallskip{\small 图 M2-043：线型四极子的节点方向}
\end{center}
\end{minipage}\par

图 M2-043 的节点锥使
\(P_2(\cos\theta_0)=0\)，只消去该方向上的四极电势值。随后应从精确偶数阶级数继续读取四阶项；电场还包含四极项的角导数，其起始阶数不能只根据电势在锥面上为零判断。

系统关于原点反演偶对称，只有偶数阶。又有 \(Q=0\)，所以

\[
\boxed{
\varphi(r,\theta)
=\frac{2q}{4\pi\varepsilon_0}
\sum_{n=1}^{\infty}
\frac{\ell^{2n}}{r^{2n+1}}P_{2n}(\cos\theta)
},
\qquad r>\ell.
\]

一般方向的首项为

\[
\boxed{
\varphi^{(2)}
=\frac{2q\ell^2}{4\pi\varepsilon_0r^3}
P_2(\cos\theta)
}.
\]

在 \(\cos^2\theta_0=1/3\) 处，

\[
P_2(\cos\theta_0)=0,
\]

而

\[
P_4(\cos\theta_0)
=\frac18(35\cos^4\theta_0-30\cos^2\theta_0+3)
=-\frac7{18}.
\]

因此节点方向上的电势从 \(L=4\) 开始：

\[
\boxed{
\varphi(r,\theta_0)
=-\frac{7q\ell^4}{9(4\pi\varepsilon_0)r^5}
+O(r^{-7})
}.
\]

四极近似的电场分量为

\[
E_r^{(2)}
=\frac{6q\ell^2}{4\pi\varepsilon_0r^4}
P_2(\cos\theta),
\]

\[
E_{\theta}^{(2)}
=\frac{6q\ell^2}{4\pi\varepsilon_0r^4}
\sin\theta\cos\theta.
\]

在 \(\theta_0\) 处，

\[
E_r^{(2)}=0,
\qquad
E_{\theta}^{(2)}\ne0.
\]

所以该方向上的电势虽从 \(r^{-5}\) 项开始，电场仍保留 \(r^{-4}\)
的四极主阶。

\subsubsection{作业四：均匀带电圆环的四极修正}\label{m2:ux4f5cux4e1aux56dbux5747ux5300ux5e26ux7535ux5706ux73afux7684ux56dbux6781ux4feeux6b63}

\textbf{题号：M2-044}

\textbf{层级：} 复赛　　\textbf{方法：} 多极展开

圆环圆心为原点，\(a>0\)，电荷分布为真空中的理想均匀线源；球坐标极轴取
\(+z\)，\(\theta\) 从该轴量起，取无穷远零势。轴对称系数采用
\(A_L=\int r'^LP_L(\cos\theta')\,dq\)。

半径为 \(a\) 的均匀细圆环位于 \(xy\) 平面，总电荷为 \(Q\)。

\begin{enumerate}
\def\labelenumi{\arabic{enumi}.}
\tightlist
\item
  用对称性判断哪些低阶矩消失；
\item
  求 \(A_2\) 并写出远区电势至 \(r^{-3}\) 阶；
\item
  用轴线精确电势检验四极项；
\item
  判断轴线与赤道远区的四极修正符号。
\end{enumerate}

\par\noindent\begin{minipage}{\linewidth}
\textbf{解答}

\begin{center}
\includegraphics[width=\linewidth,height=.40\textheight,keepaspectratio]{\detokenize{assets/module-02/examples/M2-044.pdf}}
\par\smallskip{\small 图 M2-044：均匀带电圆环的四极修正}
\end{center}
\end{minipage}\par

图 M2-044
的圆环在赤道面内均匀分布，轴对称性使横向偶极矩抵消，反演对称性消去所有奇数阶。二阶角因子在源环上取
\(P_2(0)=-1/2\)，由此得到四极修正；轴上精确势的展开提供独立检验。

圆环关于原点反演对称，所以所有奇数阶消失；关于 \(z\)
轴旋转对称，所以只有 \(m=0\) 分量。总电荷 \(Q\ne0\)，因此单极项主导，而

\[
\mathbf p=0.
\]

圆环上 \(r'=a\)、\(\cos\theta'=0\)，故

\[
A_2
=\int dq\,a^2P_2(0)
=-\frac12Qa^2.
\]

于是

\[
\boxed{
\varphi(r,\theta)
=\frac{1}{4\pi\varepsilon_0}
\left[
\frac Qr
-\frac{Qa^2}{2r^3}P_2(\cos\theta)
+O\!\left(\frac{Qa^4}{r^5}\right)
\right]
}.
\]

轴线上精确电势为

\[
\varphi_{\mathrm{\text{轴}}}(z)
=\frac{Q}{4\pi\varepsilon_0\sqrt{z^2+a^2}}.
\]

展开得

\[
\varphi_{\mathrm{\text{轴}}}(z)
=\frac{Q}{4\pi\varepsilon_0z}
\left[
1-\frac12\frac{a^2}{z^2}
+\frac38\frac{a^4}{z^4}
+\cdots
\right].
\]

因为轴线上 \(P_2(1)=1\)，其中 \(-Qa^2/(2z^3)\)
正是多极展开给出的四极修正。

赤道面上 \(P_2(0)=-1/2\)，故

\[
\Delta\varphi_{\mathrm{\text{赤道}}}^{(2)}
=\frac{Qa^2}{4(4\pi\varepsilon_0)r^3}>0
\]

当 \(Q>0\)；轴线上相应修正为负。两者符号相反。

\subsubsection{作业五：轴对称偶极与四极场线}\label{m2:ux4f5cux4e1aux4e94ux8f74ux5bf9ux79f0ux5076ux6781ux4e0eux56dbux6781ux573aux7ebf}

\textbf{题号：M2-045}

\textbf{层级：} 复赛　　\textbf{方法：} 多极展开

取球坐标的极轴为该轴对称源的对称轴，\(\theta\) 从正极轴量起，\(P_L\)
为勒让德多项式，\(C_L\ne0\) 为给定常量；只在 \(r>0\)
的理想多极场区域求场线。

轴对称单一 \(L\) 阶理想多极电势为

\[
\varphi_L=\frac{C_L}{r^{L+1}}P_L(\cos\theta).
\]

\begin{enumerate}
\def\labelenumi{\arabic{enumi}.}
\tightlist
\item
  推导子午面内的场线微分方程；
\item
  分别求 \(L=1\) 与 \(L=2\) 的场线；
\item
  说明推导中遗漏了哪些特殊场线，并解释这些公式为什么不能直接用于有限电荷体系的近场。
\end{enumerate}

\par\noindent\begin{minipage}{\linewidth}
\textbf{解答}

\begin{center}
\includegraphics[width=\linewidth,height=.40\textheight,keepaspectratio]{\detokenize{assets/module-02/examples/M2-045.pdf}}
\par\smallskip{\small 图 M2-045：轴对称偶极与四极场线}
\end{center}
\end{minipage}\par

图 M2-045 的偶极与四极流线具有不同角向结构，推导时均令局部切向与
\(E_r\mathbf e_r+E_{\theta}\mathbf e_{\theta}\)
平行。除积分所得一般曲线外，还需核对推导中除以零而丢失的轴线、赤道线或角节点。

由 \(\mathbf E=-\nabla\varphi\)，

\[
E_r
=-\frac{\partial\varphi_L}{\partial r}
=\frac{(L+1)C_L}{r^{L+2}}P_L(\cos\theta),
\]

\[
E_{\theta}
=-\frac1r\frac{\partial\varphi_L}{\partial\theta}
=\frac{C_L}{r^{L+2}}
\sin\theta\,P_L'(\cos\theta).
\]

场线切向满足

\[
\frac{dr}{r\,d\theta}
=\frac{E_r}{E_{\theta}},
\]

所以

\[
\boxed{
\frac{dr}{r\,d\theta}
=\frac{(L+1)P_L(\cos\theta)}
{\sin\theta\,P_L'(\cos\theta)}
}.
\]

对 \(L=1\)，

\[
P_1(\cos\theta)=\cos\theta,
\]

从而

\[
\frac{dr}{r}=2\cot\theta\,d\theta.
\]

积分得

\[
\boxed{r=C\sin^2\theta}.
\]

对 \(L=2\)，

\[
P_2(\cos\theta)=\frac12(3\cos^2\theta-1),
\qquad
P_2'(\cos\theta)=3\cos\theta.
\]

所以

\[
\begin{aligned}
\frac{dr}{r}
&=\frac{P_2(\cos\theta)}{\sin\theta\cos\theta}
\,d\theta\\
&=\left(\cot\theta-\frac12\tan\theta\right)d\theta.
\end{aligned}
\]

积分：

\[
\ln r
=\ln\sin\theta
+\frac12\ln|\cos\theta|
+\text{常数},
\]

故

\[
\boxed{
r=C\sin\theta\sqrt{|\cos\theta|}
}.
\]

推导中除以了 \(\sin\theta\)；四极情形还除以了
\(\cos\theta\)。因此对称轴是应单独补上的特殊场线，四极场的赤道面内径向线也应单独补上。

这些曲线来自单一理想多极项，而有限电荷体系还含其他阶，并在真实电荷处具有各自的库仑奇点。因此它们只能描述远区几何，不能替代有限体系的近场精确场线。

\section{专题十三：导体静电平衡与边界条件（基础必讲）}\label{m2:ux4e13ux9898ux5341ux4e09ux5bfcux4f53ux9759ux7535ux5e73ux8861ux4e0eux8fb9ux754cux6761ux4ef6ux57faux7840ux5fc5ux8bb2}

导体静电问题的关键不在于记住若干孤立公式，而在于同时处理三类约束：导体材料内的平衡条件、导体表面的局部边界条件，以及接地、孤立或导线连接所给出的整体约束。三者缺一不可。

\par\noindent\begin{minipage}{\linewidth}
\begin{center}
\includegraphics[width=\linewidth,height=.40\textheight,keepaspectratio]{\detokenize{assets/module-02/original/tikz-m2-topic13-v8-01.png}}
\par\smallskip{\small 图示：自由电荷重排并建立静电平衡}
\end{center}
\end{minipage}\par

\subsection{一、静电平衡的定义与基本结论}\label{m2:ux4e00ux9759ux7535ux5e73ux8861ux7684ux5b9aux4e49ux4e0eux57faux672cux7ed3ux8bba}

导体达到静电平衡，是指导体中的自由电荷不再发生宏观定向迁移，电荷分布不再随时间变化。对电导率为有限非零值的欧姆导体，局部关系为

\[
\mathbf J=\gamma\mathbf E.
\]

静电平衡时 \(\mathbf J=\mathbf 0\)，因而导体材料内部必须满足

\[
\boxed{\mathbf E_{\text{材料内}}=\mathbf 0.}
\]

由 \(\mathbf E=-\nabla\varphi\) 立即得到

\[
\nabla\varphi=\mathbf 0,
\qquad
\boxed{\varphi=C.}
\]

因此，同一块\textbf{连通导体}的材料内部及其全部表面具有同一电势。若几块导体由理想导线连接，它们合起来构成一个连通导体，平衡后也具有同一电势。彼此绝缘的不同导体则可以具有不同的电势常数
\(C_1,C_2,\ldots\)。

在导体材料内部的普通点，高斯定律给出

\[
\rho=\varepsilon_0\nabla\cdot\mathbf E=0.
\]

所以静电平衡时，导体所带的净电荷只能分布在表面。这里的``内部''是指导体材料的正则区域；表面、尖端和棱边处允许出现面电荷，理想尖端处还可能出现场的奇异增强。

\textbf{例 M2-046：连通导体最终相同的是电势}

两个彼此分离的导体最初处于静电平衡，电势分别为
\(V_1,V_2\)，两者采用同一参考零点。现用理想细导线连接两导体，整个系统不接地，待电荷重新分布达到静电平衡。写出连接后的共同约束，说明最终相等的是电势还是各自的电荷量，并指出导线内电场与该结论的关系。

\par\noindent\begin{minipage}{\linewidth}
\textbf{解答}

\begin{center}
\includegraphics[width=\linewidth,height=.40\textheight,keepaspectratio]{\detokenize{assets/module-02/examples/M2-046.pdf}}
\par\smallskip{\small 图 M2-046：连通导体最终相同的是电势}
\end{center}
\end{minipage}\par

图 M2-046
的导线连接两导体表面，静电平衡时导线内部不能维持切向电场，所以两端电势最终相同。电荷如何重新分配还取决于球面大小、相互位置和外部边界，不能由等势条件推出电荷相等。

两个相互绝缘的导体最初电势分别为 \(V_1\)、\(V_2\)。接通导线后，只要
\(V_1\ne V_2\)，导线中就会存在电场并发生电荷迁移；最终条件不是``两个导体电荷相等''，而是

\[
V_1'=V_2'.
\]

电荷如何分配取决于两导体的几何形状、相互位置和周围边界。

\subsection{二、导体表面的边界条件}\label{m2:ux4e8cux5bfcux4f53ux8868ux9762ux7684ux8fb9ux754cux6761ux4ef6}

设单位法向量 \(\mathbf n\) 从导体材料指向外部真空，\(\sigma\)
为该处面电荷密度。必须区分表面内、外两侧的极限：

\[
\mathbf E_{\text{内}}
=\lim_{h\to0^+}\mathbf E(\mathbf r-h\mathbf n),
\qquad
\mathbf E_{\text{外}}
=\lim_{h\to0^+}\mathbf E(\mathbf r+h\mathbf n).
\]

理想面电荷处的电场一般不连续，因此``表面上的电场''若未说明取哪一侧极限，就不是完整表述。

\subsubsection{1.
法向分量：高斯薄柱}\label{m2:ux6cd5ux5411ux5206ux91cfux9ad8ux65afux8584ux67f1}

取跨过表面的极薄高斯柱，端面面积为
\(\Delta S\)。当柱高趋于零时，侧面通量趋于零，高斯定律给出

\[
(\mathbf E_{\text{外}}-\mathbf E_{\text{内}})\cdot\mathbf n\,\Delta S
=\frac{\sigma\Delta S}{\varepsilon_0}.
\]

因此一般的法向跃变条件是

\[
\boxed{(\mathbf E_{\text{外}}-\mathbf E_{\text{内}})\cdot\mathbf n
=\frac{\sigma}{\varepsilon_0}.}
\]

导体材料内 \(\mathbf E_{\text{内}}=0\)，所以

\[
\boxed{E_{n,\text{外}}=\frac{\sigma}{\varepsilon_0}.}
\]

\subsubsection{2.
切向分量：极薄回路}\label{m2:ux5207ux5411ux5206ux91cfux6781ux8584ux56deux8def}

取跨过表面的极小矩形回路，两条长边平行于表面。静电场满足

\[
\oint\mathbf E\cdot d\mathbf l=0.
\]

令矩形高度趋于零，短边贡献消失，得到

\[
(\mathbf E_{t,\text{外}}-\mathbf E_{t,\text{内}})\cdot\Delta\mathbf l=0.
\]

于是

\[
\boxed{\mathbf E_{t,\text{外}}=\mathbf E_{t,\text{内}}.}
\]

导体内切向场为零，故导体外侧紧贴表面的电场只能沿法线：

\[
\boxed{\mathbf E_{t,\text{外}}=\mathbf 0.}
\]

\par\noindent\begin{minipage}{\linewidth}
\begin{center}
\includegraphics[width=\linewidth,height=.40\textheight,keepaspectratio]{\detokenize{assets/module-02/original/tikz-m2-topic13-v8-02.png}}
\par\smallskip{\small 图示：用高斯薄柱和极薄回路分别导出法向与切向边界条件}
\end{center}
\end{minipage}\par

\textbf{例 M2-047：带符号的表面电荷决定场向}

一个处于静电平衡的导体，其表面邻接真空。取表面小面元，令 \(\mathbf n\)
从导体实体指向外部真空，该处实际面电荷密度为有符号量
\(\sigma\)。由导体内部零场及界面条件判断紧邻表面外侧电场的方向，分别讨论
\(\sigma>0\)、\(\sigma<0\)，同时说明是否存在切向电场。

\par\noindent\begin{minipage}{\linewidth}
\textbf{解答}

\begin{center}
\includegraphics[width=\linewidth,height=.40\textheight,keepaspectratio]{\detokenize{assets/module-02/examples/M2-047.pdf}}
\par\smallskip{\small 图 M2-047：带符号的表面电荷决定场向}
\end{center}
\end{minipage}\par

图 M2-047
的橙色箭头规定从导体指向真空的外法向。边界式中的法向分量带符号，正的面电荷对应蓝色向外场；若面电荷变负，场箭头反转，但法向定义保持不变。

若某处 \(\sigma>0\)，则 \(E_{n,\text{外}}>0\)，外侧电场沿外法线；若
\(\sigma<0\)，则外侧电场与外法线反向。结论来自带符号的边界条件，而不是单纯记忆``电场垂直表面''。

\subsubsection{3.
电势连续，法向导数跃变}\label{m2:ux7535ux52bfux8fdeux7eedux6cd5ux5411ux5bfcux6570ux8dc3ux53d8}

对有限面电荷密度且不存在理想偶极层的普通导体表面，跨越厚度 \(2h\)
的短线段有

\[
\varphi(\mathbf r+h\mathbf n)-\varphi(\mathbf r-h\mathbf n)
=-\int_{-h}^{h}\mathbf E\cdot\mathbf n\,ds.
\]

虽然电场的法向分量有有限跃变，但令 \(h\to0\) 时上式积分趋于零，因此

\[
\boxed{\varphi_{\text{外}}=\varphi_{\text{内}}=C.}
\]

另一方面，由 \(\mathbf E=-\nabla\varphi\) 可得

\[
\boxed{\sigma
=-\varepsilon_0\frac{\partial\varphi_{\text{外}}}{\partial n}.}
\]

这里的偏导沿从导体指向真空的外法线计算。电势连续并不意味着法向导数连续；面电荷正是由法向导数的跃变所承担。

\par\noindent\begin{minipage}{\linewidth}
\begin{center}
\includegraphics[width=\linewidth,height=.40\textheight,keepaspectratio]{\detokenize{assets/module-02/original/tikz-m2-topic13-v8-03.png}}
\par\smallskip{\small 图示：电势在导体表面连续而法向电场发生跃变}
\end{center}
\end{minipage}\par

\textbf{例 M2-048：导体表面外侧的电势斜率}

一个静电平衡导体在所考察处表面可视为无限大平面，真空侧紧邻表面的均匀面电荷密度为
\(\sigma>0\)。令法向 \(\mathbf n\) 从导体指向真空，并以
\(\partial/\partial n\)
表示沿该方向的导数。利用界面条件求导体外侧的电场及电势法向导数，解释二者符号之间的关系。

\par\noindent\begin{minipage}{\linewidth}
\textbf{解答}

\begin{center}
\includegraphics[width=\linewidth,height=.40\textheight,keepaspectratio]{\detokenize{assets/module-02/examples/M2-048.pdf}}
\par\smallskip{\small 图 M2-048：导体表面外侧的电势斜率}
\end{center}
\end{minipage}\par

图 M2-048 的下侧是零场导体材料，真空侧的法向场由面电荷跳跃直接确定。再用
\(E_n=-\partial_n\varphi\)
读取斜率符号，正面电荷处沿外法向电势降低，不能把场强和电势导数写成同号。

若真空侧表面电荷密度为均匀的 \(\sigma>0\)，则

\[
\mathbf E_{\text{外}}=\frac{\sigma}{\varepsilon_0}\mathbf n,
\qquad
\frac{\partial\varphi_{\text{外}}}{\partial n}=-\frac{\sigma}{\varepsilon_0}.
\]

沿外法线方向电势降低，与 \(\mathbf E=-\nabla\varphi\) 一致。

\subsection{三、接地、孤立与导线连接}\label{m2:ux4e09ux63a5ux5730ux5b64ux7acbux4e0eux5bfcux7ebfux8fdeux63a5}

``导体等势''只给出电势在每个连通导体上为常数，并不自动给出这个常数的数值。完整边值问题还要判断导体的连接状态。

\par\noindent\begin{minipage}{\linewidth}
\begin{center}
\includegraphics[width=\linewidth,height=.40\textheight,keepaspectratio]{\detokenize{assets/module-02/original/tikz-m2-topic13-v8-04.png}}
\par\smallskip{\small 图示：接地、孤立和导线连接对应不同的整体约束}
\end{center}
\end{minipage}\par

\subsubsection{1. 接地导体}\label{m2:ux63a5ux5730ux5bfcux4f53}

以大地电势为零时，接地导体满足

\[
\boxed{\varphi=0.}
\]

它的总电荷通常不是预先给定的，因为电荷可以从大地流入或流向大地。接地是\textbf{电势约束}，不是``导体不带电''；处于外电场中的接地导体完全可以带有非零感应电荷。

\subsubsection{2. 孤立导体}\label{m2:ux5b64ux7acbux5bfcux4f53}

孤立导体满足

\[
\varphi=C,
\qquad
\int_S\sigma\,dS=Q.
\]

其中总电荷 \(Q\) 给定，而电势常数 \(C\) 由求解结果确定。孤立且中性只表示
\(Q=0\)，并不表示表面处处有 \(\sigma=0\)；外场仍可使正、负感应电荷分离。

\subsubsection{3.
未接地的连接导体}\label{m2:ux672aux63a5ux5730ux7684ux8fdeux63a5ux5bfcux4f53}

若干导体由理想导线连接但整体未接地时，

\[
\varphi_1=\varphi_2=\cdots=C,
\]

同时整个连接系统的总电荷守恒。各部分电荷一般不守恒，因为电荷可以沿导线重新分配。

\subsection{四、封闭空腔、感应电荷与静电屏蔽}\label{m2:ux56dbux5c01ux95edux7a7aux8154ux611fux5e94ux7535ux8377ux4e0eux9759ux7535ux5c4fux853d}

\subsubsection{1.
无电荷封闭空腔}\label{m2:ux65e0ux7535ux8377ux5c01ux95edux7a7aux8154}

设导体内部有一个封闭空腔，腔内没有电荷。腔壁属于同一连通导体，故其电势为同一常数
\(C\)。令

\[
u=\varphi-C.
\]

空腔内满足

\[
\nabla^2u=0,
\qquad
u|_{\partial V}=0.
\]

利用恒等式
\(\nabla\cdot(u\nabla u)=|\nabla u|^2+u\nabla^2u\)，在空腔体积 \(V\)
内积分：

\[
\int_V|\nabla u|^2\,dV
=\oint_{\partial V}u\frac{\partial u}{\partial n}\,dS
-\int_Vu\nabla^2u\,dV=0.
\]

被积函数非负，故 \(\nabla u=0\)，从而

\[
\boxed{\varphi=C,\qquad \mathbf E_{\text{腔}}=0.}
\]

因此，外部静电场不能进入封闭的无电荷空腔。这个结论与导体是否接地无关。

\subsubsection{2.
空腔内有电荷}\label{m2:ux7a7aux8154ux5185ux6709ux7535ux8377}

若一个空腔内的总电荷为
\(q\)，在包围该空腔且完全位于导体材料中的闭合面上有
\(\mathbf E=0\)，高斯定律给出

\[
q+Q_{\text{内壁}}=0.
\]

所以

\[
\boxed{Q_{\text{内壁}}=-q.}
\]

若导体本身的总电荷为 \(Q_{\text{导体}}\)，由电荷守恒

\[
Q_{\text{导体}}=Q_{\text{内壁}}+Q_{\text{外壁}},
\]

得到

\[
\boxed{Q_{\text{外壁}}=Q_{\text{导体}}+q.}
\]

这些式子只确定内、外壁的\textbf{总电荷}。若电荷偏离空腔对称中心，内壁面电荷密度通常不均匀；不能仅凭高斯定律求出其局部分布。

\par\noindent\begin{minipage}{\linewidth}
\begin{center}
\includegraphics[width=\linewidth,height=.40\textheight,keepaspectratio]{\detokenize{assets/module-02/original/tikz-m2-topic13-v8-05.png}}
\par\smallskip{\small 图示：无电荷空腔与带电空腔的电场和壁面总电荷}
\end{center}
\end{minipage}\par

\textbf{例 M2-049：中性导体空腔的电荷守恒}

一个最初电中性的孤立导体含有一处完全封闭的空腔，腔内放置不接触内壁的点电荷
\(q\)。系统达到静电平衡，导体未接地，也没有电荷穿过导体外边界。求空腔内壁和导体外表面的总电荷，说明结果是否依赖点电荷在腔内的具体位置或空腔形状；不要求面密度的逐点分布。

\par\noindent\begin{minipage}{\linewidth}
\textbf{解答}

\begin{center}
\includegraphics[width=\linewidth,height=.40\textheight,keepaspectratio]{\detokenize{assets/module-02/examples/M2-049.pdf}}
\par\smallskip{\small 图 M2-049：中性导体空腔的电荷守恒}
\end{center}
\end{minipage}\par

图 M2-049
中在导体材料内作包围空腔的高斯面，通量为零，因此腔内电荷与内壁电荷相消。导体本身的电荷总量保持为零，故外壁必须留有
\(+q\)；材料内无场不等于导体外部也无场。

此时 \(Q_{\text{导体}}=0\)，故

\[
Q_{\text{内壁}}=-q,
\qquad
Q_{\text{外壁}}=+q.
\]

导体材料内仍有 \(\mathbf E=0\)，但外表面的 \(+q\) 一般会在外部产生电场。

\subsubsection{3.
必须区分屏蔽的两个方向}\label{m2:ux5fc5ux987bux533aux5206ux5c4fux853dux7684ux4e24ux4e2aux65b9ux5411}

\par\noindent\begin{minipage}{\linewidth}
\begin{center}
\includegraphics[width=\linewidth,height=.40\textheight,keepaspectratio]{\detokenize{assets/module-02/original/tikz-m2-topic13-v8-06.png}}
\par\smallskip{\small 图示：外场进入空腔与腔内电荷影响外部是两个不同问题}
\end{center}
\end{minipage}\par

\begin{itemize}
\tightlist
\item
  \textbf{外场能否进入无电荷空腔：}不能。封闭空腔的边界是常电势面，腔内无源，故腔内场为零；孤立和接地都成立。
\item
  \textbf{腔内电荷能否影响外部：}取决于外表面的整体约束。孤立中性导体的外壁总电荷为
  \(+q\)，外部一般存在电场。
\item
  \textbf{外壳接地且外部没有其他电荷时：}外表面电势为零，无穷远电势也为零。由外部区域的唯一性，外部电势处处为零，故外表面
  \(\sigma=0\)、外部电场为零。此时导体从大地取得总电荷
  \(-q\)，它全部表现为内壁总电荷。
\item
  \textbf{外部另有电荷时：}即使外壳接地，外表面仍会因外部电荷产生感应电荷；``外表面
  \(\sigma=0\)''不再成立。
\end{itemize}

\subsection{五、导体表面的电场力与电压力}\label{m2:ux4e94ux5bfcux4f53ux8868ux9762ux7684ux7535ux573aux529bux4e0eux7535ux538bux529b}

不能直接把 \(\sigma\mathbf E_{\text{外}}\)
当作面电荷所受的单位面积电场力，因为 \(\mathbf E_{\text{外}}\)
中包含了该面元自身产生的场。把局部表面近似为平面面电荷层，其自场在两侧大小相等、方向相反，因此去掉自场后，作用于该面元的外场为两侧极限的平均：

\[
\boxed{\mathbf E_{\text{去自场}}
=\frac{\mathbf E_{\text{外}}+\mathbf E_{\text{内}}}{2}.}
\]

导体内 \(\mathbf E_{\text{内}}=0\)，于是

\[
\mathbf E_{\text{去自场}}
=\frac{\mathbf E_{\text{外}}}{2}
=\frac{\sigma}{2\varepsilon_0}\mathbf n.
\]

表面单位面积所受电场力为

\[
\boxed{\mathbf f_s
=\sigma\mathbf E_{\text{去自场}}
=\frac{\sigma^2}{2\varepsilon_0}\mathbf n.}
\]

相应的向外电压力为

\[
\boxed{p_{\text{电}}
=\frac{\sigma^2}{2\varepsilon_0}
=\frac{\varepsilon_0E_{\text{外}}^2}{2}.}
\]

无论 \(\sigma\) 正负，\(\sigma^2\)
都为正，故静电压力总是把导体表面向外拉伸。该结果也可由真空中的麦克斯韦应力张量得到。

\par\noindent\begin{minipage}{\linewidth}
\begin{center}
\includegraphics[width=\linewidth,height=.40\textheight,keepaspectratio]{\detokenize{assets/module-02/original/tikz-m2-topic13-v8-07.png}}
\par\smallskip{\small 图示：表面电荷受力应使用两侧电场的平均值}
\end{center}
\end{minipage}\par

\textbf{例 M2-050：局部电压力沿外法向}

静电平衡导体的一小块表面邻接真空，该处电荷面密度为
\(\sigma=2.0\times10^{-6}\,\mathrm{C/m^2}\)。面元面积为
\(\Delta S\)，足够小以便把 \(\sigma\) 和外法向 \(\mathbf n\)
视为常量。求该面元受到的静电力大小及方向，说明计算电压力时为什么不能直接把包含面元自场的全部外侧场作为作用场。

\par\noindent\begin{minipage}{\linewidth}
\textbf{解答}

\begin{center}
\includegraphics[width=\linewidth,height=.40\textheight,keepaspectratio]{\detokenize{assets/module-02/examples/M2-050.pdf}}
\par\smallskip{\small 图 M2-050：局部电压力沿外法向}
\end{center}
\end{minipage}\par

图 M2-050 的橙色小片选定面积
\(\Delta S\)，红色力沿外法向。先由面密度求压力
\(\sigma^2/(2\varepsilon_0)\)，再乘局部面积；由于压力含平方，面电荷变号不反转该局部压力方向。

若某处 \(\sigma=2.0\times10^{-6}\ \mathrm{C/m^2}\)，面积为 \(\Delta S\)
的小面元所受法向力近似为

\[
\Delta\mathbf F
=\frac{\sigma^2}{2\varepsilon_0}\Delta S\,\mathbf n.
\]

面积足够小时才可把 \(\sigma\) 视为常量；对整个曲面求合力必须积分
\(\int_S\mathbf f_s\,dS\)。

\subsection{六、曲率、尖端效应和近似条件}\label{m2:ux516dux66f2ux7387ux5c16ux7aefux6548ux5e94ux548cux8fd1ux4f3cux6761ux4ef6}

``越尖处电荷越多''只能作为常见趋势，不能当作脱离整体边界条件的局部定理。一个可严格计算的近似模型是：两个相距很远的导体球，半径分别为
\(R_1\)、\(R_2\)，用很细的长导线连接。忽略两球间的相互影响和导线电容，则等势条件给出

\[
\frac{Q_1}{4\pi\varepsilon_0R_1}
=\frac{Q_2}{4\pi\varepsilon_0R_2},
\]

所以

\[
\frac{Q_1}{Q_2}=\frac{R_1}{R_2}.
\]

两球表面电荷密度和场强之比为

\[
\boxed{
\frac{\sigma_1}{\sigma_2}
=\frac{E_1}{E_2}
=\frac{R_2}{R_1}.}
\]

在这个模型中，较小半径对应较大的表面电荷密度和场强。但对一般不规则导体，\(\sigma\)
是整个边值问题的解，还取决于其他导体、外加电场、接地状态和远处边界，不能只看某一点的局部曲率。

\par\noindent\begin{minipage}{\linewidth}
\begin{center}
\includegraphics[width=\linewidth,height=.40\textheight,keepaspectratio]{\detokenize{assets/module-02/original/tikz-m2-topic13-v8-08.png}}
\par\smallskip{\small 图示：连接双球模型与一般几何中的曲率边界}
\end{center}
\end{minipage}\par

\textbf{例 M2-051：小球电荷较少却有更强表面场}

真空中两个相距很远的导体球半径分别为 \(R\) 和
\(3R\)，\(R>0\)。两球由足够细的导线相连，系统不接地，带总电荷
\(Q\ne0\)。忽略导线电容、导线带电及两球相互感应，使每球可按孤立球近似。求平衡时两球电荷之比和紧邻表面的电场大小之比，并判断小球的电荷较少是否必然意味着表面场较弱。

\par\noindent\begin{minipage}{\linewidth}
\textbf{解答}

\begin{center}
\includegraphics[width=\linewidth,height=.40\textheight,keepaspectratio]{\detokenize{assets/module-02/examples/M2-051.pdf}}
\par\smallskip{\small 图 M2-051：小球电荷较少却有更强表面场}
\end{center}
\end{minipage}\par

图 M2-051 的两球半径按1:3画出，细线使表面电势相同。用 \(Q_i/R_i\)
相等联立总电荷守恒，先确定分配比例；表面场另除一次半径，所以较小球的总电荷较少，表面场反而更大。

若系统总电荷为 \(Q\)，则

\[
Q_1=\frac{Q}{4},
\qquad
Q_2=\frac{3Q}{4},
\]

但小球表面场强是大球的三倍：

\[
\frac{E_1}{E_2}=3.
\]

这说明``小球总电荷较少''与``小球表面场更强''并不矛盾。

\subsection{七、孤立带电导体球：分段场与电势}\label{m2:ux4e03ux5b64ux7acbux5e26ux7535ux5bfcux4f53ux7403ux5206ux6bb5ux573aux4e0eux7535ux52bf}

半径为 \(R\) 的孤立导体球带总电荷
\(Q\)，取无穷远电势为零。球对称性和高斯定律给出

\[
\mathbf E(r)=
\begin{cases}
\mathbf 0, & 0\le r<R,\\[4pt]
\dfrac{Q}{4\pi\varepsilon_0r^2}\mathbf e_r, & r>R.
\end{cases}
\]

电势在 \(r=R\) 连续：

\[
\varphi(r)=
\begin{cases}
\dfrac{Q}{4\pi\varepsilon_0R}, & 0\le r\le R,\\[6pt]
\dfrac{Q}{4\pi\varepsilon_0r}, & r\ge R.
\end{cases}
\]

表面电荷密度、电容和静电能分别为

\[
\sigma=\frac{Q}{4\pi R^2},
\qquad
C=\frac{Q}{\varphi(R)}=4\pi\varepsilon_0R,
\qquad
U=\frac{Q^2}{2C}=\frac{Q^2}{8\pi\varepsilon_0R}.
\]

在表面处，电势连续；场强却从内侧的 \(0\) 跳到外侧的
\(Q/(4\pi\varepsilon_0R^2)\)。作图时不能用一条连续曲线把这两个单侧极限连接起来。

\par\noindent\begin{minipage}{\linewidth}
\begin{center}
\includegraphics[width=\linewidth,height=.40\textheight,keepaspectratio]{\detokenize{assets/module-02/original/tikz-m2-topic13-v8-09.png}}
\par\smallskip{\small 图示：孤立带电导体球的空间场、电场强度和电势分段图}
\end{center}
\end{minipage}\par

\subsection{八、导体静电题的求解顺序}\label{m2:ux516bux5bfcux4f53ux9759ux7535ux9898ux7684ux6c42ux89e3ux987aux5e8f}

\par\noindent\begin{minipage}{\linewidth}
\begin{center}
\includegraphics[width=\linewidth,height=.40\textheight,keepaspectratio]{\detokenize{assets/module-02/original/tikz-m2-topic13-v8-10.png}}
\par\smallskip{\small 图示：导体静电题的区域划分、整体约束与表面量求解流程}
\end{center}
\end{minipage}\par

处理导体静电题时，可按以下次序推进：

\begin{enumerate}
\def\labelenumi{\arabic{enumi}.}
\tightlist
\item
  划分导体材料、各空腔和外部真空区域。
\item
  在导体材料内写 \(\mathbf E=0\)、\(\varphi=C_i\)。
\item
  判断每个导体是接地、孤立还是通过导线相连。
\item
  在各真空区域写泊松方程或拉普拉斯方程，并列出边界条件。
\item
  由对称性、分离变量、电像法或唯一性定理求电势。
\item
  用 \(\sigma=-\varepsilon_0\partial_n\varphi_{\text{外}}\)
  求表面电荷，再核对总电荷。
\item
  若求表面力，使用平均场或麦克斯韦应力张量，不能直接乘外侧总场。
\end{enumerate}

常见错误集中在以下边界：

\begin{itemize}
\tightlist
\item
  \(\mathbf E=0\) 只保证在导体材料内成立；带电空腔内通常有场。
\item
  ``等势''不等于``零电势''；只有选定参考并接地后才常写 \(\varphi=0\)。
\item
  接地固定电势，孤立固定总电荷；二者不能交换。
\item
  \(Q_{\text{内壁}}=-q\) 是总量关系，不直接给出局部面电荷密度。
\item
  理想面电荷处要区分
  \(\mathbf E_{\text{内}}\)、\(\mathbf E_{\text{外}}\)
  与去自场后的平均场。
\item
  ``尖端电荷密度大''有适用条件，不能替代完整的边值求解。
\end{itemize}

\subsection{九、练习与解答}\label{m2:ux4e5dux7ec3ux4e60ux4e0eux89e3ux7b54}

\subsubsection{作业一：孤立带电导体球}\label{m2:ux4f5cux4e1aux4e00ux5b64ux7acbux5e26ux7535ux5bfcux4f53ux7403}

\textbf{题号：M2-052}

\textbf{层级：} 基础　　\textbf{方法：} 静电场方程

真空中一个半径为 \(R>0\) 的孤立球形导体带有符号总电荷
\(Q\)，处于静电平衡，外部没有其他电荷或外场。以球心为原点，取
\(\varphi(\infty)=0\)。求导体内外各区域的电场和电势、球面电荷密度、相对于无穷远的自电容，以及该带电导体的静电能。

\par\noindent\begin{minipage}{\linewidth}
\textbf{解答}

\begin{center}
\includegraphics[width=\linewidth,height=.40\textheight,keepaspectratio]{\detokenize{assets/module-02/examples/M2-052.pdf}}
\par\smallskip{\small 图 M2-052：孤立带电导体球}
\end{center}
\end{minipage}\par

图 M2-052
的球面划分导体内部和真空外部。内部电势是边界常量、场为零；外部利用球对称性和无穷远零势求解，再从同一表面场依次求
\(\sigma,C,U\)，各结果应共享同一个 \(Q,R\)。

静电平衡时球内电场为零。对 \(r>R\) 的同心球面使用高斯定律：

\[
E(r)4\pi r^2=\frac{Q}{\varepsilon_0},
\]

故

\[
\mathbf E(r)=
\begin{cases}
\mathbf 0, & r<R,\\[4pt]
\dfrac{Q}{4\pi\varepsilon_0r^2}\mathbf e_r, & r>R.
\end{cases}
\]

当 \(Q<0\) 时，上式系数为负，场方向自动反向。由 \(\varphi(\infty)=0\)，

\[
\varphi(r)=\int_r^{\infty} E(r')\,dr'
=\frac{Q}{4\pi\varepsilon_0r},
\qquad r\ge R.
\]

球内与表面等势，所以

\[
\varphi(r)=\frac{Q}{4\pi\varepsilon_0R},
\qquad r\le R.
\]

再由表面外侧法向场

\[
E_n(R^+)=\frac{Q}{4\pi\varepsilon_0R^2}
\]

得到

\[
\sigma=\varepsilon_0E_n(R^+)=\frac{Q}{4\pi R^2}.
\]

球的电势为 \(V=Q/(4\pi\varepsilon_0R)\)，故

\[
C=\frac{Q}{V}=4\pi\varepsilon_0R,
\qquad
U=\frac{Q^2}{2C}=\frac{Q^2}{8\pi\varepsilon_0R}.
\]

\subsubsection{作业二：两个远离的连接导体球}\label{m2:ux4f5cux4e1aux4e8cux4e24ux4e2aux8fdcux79bbux7684ux8fdeux63a5ux5bfcux4f53ux7403}

\textbf{题号：M2-053}

\textbf{层级：} 基础　　\textbf{方法：} 静电场方程

真空中两个球形导体半径分别为 \(0<R_1<R_2\)，球心距远大于
\(R_2\)，由一根足够细的长导线连接。系统不接地，总自由电荷为
\(Q\ne0\)。忽略导线带电、电容和两球相互感应，以共同无穷远零势比较两球电势。求达到平衡时的
\(Q_1,Q_2\)、各球均匀表面电荷密度之比及外侧场强大小之比。

\par\noindent\begin{minipage}{\linewidth}
\textbf{解答}

\begin{center}
\includegraphics[width=\linewidth,height=.40\textheight,keepaspectratio]{\detokenize{assets/module-02/examples/M2-053.pdf}}
\par\smallskip{\small 图 M2-053：两个远离的连接导体球}
\end{center}
\end{minipage}\par

图 M2-053
的细导线只提供等势约束，不计自身电容。在相距很远的近似下，每球电势可用各自的
\(Q_i/(4\pi\varepsilon_0R_i)\)，而面密度还要除以
\(4\pi R_i^2\)；因此电荷比、面密度比与场强比不能混成同一比例。

连接后两球等势：

\[
\frac{Q_1}{4\pi\varepsilon_0R_1}
=\frac{Q_2}{4\pi\varepsilon_0R_2},
\qquad
Q_1+Q_2=Q.
\]

解得

\[
Q_1=\frac{R_1}{R_1+R_2}Q,
\qquad
Q_2=\frac{R_2}{R_1+R_2}Q.
\]

由于 \(\sigma_i=Q_i/(4\pi R_i^2)\)，

\[
\frac{\sigma_1}{\sigma_2}
=\frac{Q_1R_2^2}{Q_2R_1^2}
=\frac{R_2}{R_1}.
\]

又因 \(E_i=\sigma_i/\varepsilon_0\)，所以

\[
\boxed{
\frac{\sigma_1}{\sigma_2}
=\frac{E_1}{E_2}
=\frac{R_2}{R_1}>1.}
\]

小球分得的总电荷较少，但单位面积电荷更多，表面场更强。

\subsubsection{作业三：中心点电荷与球形导体壳}\label{m2:ux4f5cux4e1aux4e09ux4e2dux5fc3ux70b9ux7535ux8377ux4e0eux7403ux5f62ux5bfcux4f53ux58f3}

\textbf{题号：M2-054}

\textbf{层级：} 基础　　\textbf{方法：} 静电场方程

真空中有内半径 \(a\)、外半径 \(b\)
的同心球形导体壳，\(0<a<b\)，导体壳本身最初带总电荷
\(Q_c\)。在球心固定点电荷
\(q\)，点电荷不与壳接触，外界没有其他源。取无穷远零势，求
\(0<r<a\)、\(a<r<b\)、\(r>b\)
的电场与电势，以及内外球面的电荷密度。再将壳接地，保持中心电荷和外部无源条件不变，讨论电荷分布与外部场如何改变。

\par\noindent\begin{minipage}{\linewidth}
\textbf{解答}

\begin{center}
\includegraphics[width=\linewidth,height=.40\textheight,keepaspectratio]{\detokenize{assets/module-02/examples/M2-054.pdf}}
\par\smallskip{\small 图 M2-054：中心点电荷与球形导体壳}
\end{center}
\end{minipage}\par

图 M2-054 的内半径 \(a\) 和外半径 \(b\)
分出空腔、导体材料与外部真空。材料内的高斯面先固定内壁电荷
\(-q\)，导体总量再决定外壁；分区积分场时，在两球面连续连接电势，接地情形改由边界电势固定外部解。

球心放置电荷保持球对称。导体材料 \(a<r<b\) 内有
\(\mathbf E=0\)。在材料内取同心高斯面，包围电荷 \(q\) 与内壁电荷，故

\[
q+Q_{\text{内壁}}=0,
\qquad
Q_{\text{内壁}}=-q.
\]

导体总电荷为 \(Q_c\)，于是

\[
Q_{\text{外壁}}=Q_c+q.
\]

由球对称性，

\[
\mathbf E(r)=
\begin{cases}
\dfrac{q}{4\pi\varepsilon_0r^2}\mathbf e_r, & 0<r<a,\\[6pt]
\mathbf 0, & a<r<b,\\[4pt]
\dfrac{Q_c+q}{4\pi\varepsilon_0r^2}\mathbf e_r, & r>b.
\end{cases}
\]

两表面的电荷均匀分布：

\[
\sigma_{\text{内}}
=-\frac{q}{4\pi a^2},
\qquad
\sigma_{\text{外}}
=\frac{Q_c+q}{4\pi b^2}.
\]

外部电势为

\[
\varphi(r)=\frac{Q_c+q}{4\pi\varepsilon_0r},
\qquad r\ge b.
\]

导体材料内为常量

\[
\varphi(r)=\frac{Q_c+q}{4\pi\varepsilon_0b},
\qquad a\le r\le b.
\]

在空腔内由电势连续性得到

\[
\varphi(r)
=\frac{q}{4\pi\varepsilon_0}
\left(\frac1r-\frac1a\right)
+\frac{Q_c+q}{4\pi\varepsilon_0b},
\qquad 0<r<a.
\]

若导体壳接地且外部无其他电荷，则外表面电势与无穷远电势都为零。唯一性定理给出外部
\(\varphi=0\)，因此

\[
Q_{\text{外壁}}=0,
\qquad
Q_{\text{内壁}}=-q,
\qquad
Q_{\text{导体}}=-q.
\]

此时

\[
\varphi(r)=
\frac{q}{4\pi\varepsilon_0}
\left(\frac1r-\frac1a\right),
\qquad 0<r<a,
\]

而 \(r\ge a\) 的导体材料和外部空间中电势均为零。

\subsubsection{作业四：多个空腔的总电荷核算}\label{m2:ux4f5cux4e1aux56dbux591aux4e2aux7a7aux8154ux7684ux603bux7535ux8377ux6838ux7b97}

\textbf{题号：M2-055}

\textbf{层级：} 基础　　\textbf{方法：} 静电场方程

一个孤立、未接地的连通导体含有三个互不连通且完全封闭的空腔。三腔内分别放置不接触内壁的电荷系统，其总电荷为
\(q_1,q_2,q_3\)；导体材料本身所带总电荷为
\(Q_c\)。达到静电平衡后，求三个内壁各自的总感应电荷及外表面总电荷。说明仅有这些总量数据，是否足以确定任一壁面的逐点面电荷密度。

\par\noindent\begin{minipage}{\linewidth}
\textbf{解答}

\begin{center}
\includegraphics[width=\linewidth,height=.40\textheight,keepaspectratio]{\detokenize{assets/module-02/examples/M2-055.pdf}}
\par\smallskip{\small 图 M2-055：多个空腔的总电荷核算}
\end{center}
\end{minipage}\par

图 M2-055
的三个空腔彼此封闭，对每个空腔都可在导体材料中单独作零通量高斯面。先分别得到各内壁电荷，再用整个导体电荷守恒求外壁总量；这些积分条件不决定每个面元上的密度。

分别在导体材料内取包围每个空腔的高斯面。由于各高斯面上
\(\mathbf E=0\)，有

\[
Q_{\text{内壁},i}=-q_i,
\qquad i=1,2,3.
\]

导体总电荷等于全部内壁电荷与外壁电荷之和：

\[
Q_c=-(q_1+q_2+q_3)+Q_{\text{外壁}}.
\]

因此

\[
\boxed{Q_{\text{外壁}}=Q_c+q_1+q_2+q_3.}
\]

这些结果只确定总量。局部面电荷密度还取决于空腔形状、电荷位置、外表面形状和外部边界条件，不能由高斯定律单独确定。

\subsubsection{作业五：带电导体球的电压力}\label{m2:ux4f5cux4e1aux4e94ux5e26ux7535ux5bfcux4f53ux7403ux7684ux7535ux538bux529b}

\textbf{题号：M2-056}

\textbf{层级：} 基础　　\textbf{方法：} 静电场方程

真空中半径为 \(R>0\) 的孤立球形导体带固定总电荷
\(Q\)，外部无其他电荷或场。求静电平衡时球面的电压力及作用方向。再设球面作保持总电荷不变的均匀微小膨胀
\(dR\)，由静电能变化计算电场所做虚功，检验压力与面积、位移之间的关系。

\par\noindent\begin{minipage}{\linewidth}
\textbf{解答}

\begin{center}
\includegraphics[width=\linewidth,height=.40\textheight,keepaspectratio]{\detokenize{assets/module-02/examples/M2-056.pdf}}
\par\smallskip{\small 图 M2-056：带电导体球的电压力}
\end{center}
\end{minipage}\par

图 M2-056 中球面法向随位置径向向外，局部压力由外侧场平方确定。固定 \(Q\)
使球半径增加时能量下降，\(-dU\)
等于压力乘整个球面面积再乘径向位移，由此核对压力的系数。

球面外侧场强为

\[
E_{\text{外}}=\frac{Q}{4\pi\varepsilon_0R^2}.
\]

故电压力

\[
p=\frac{\varepsilon_0E_{\text{外}}^2}{2}
=\boxed{\frac{Q^2}{32\pi^2\varepsilon_0R^4}}.
\]

球的静电能为

\[
U(R)=\frac{Q^2}{8\pi\varepsilon_0R}.
\]

在总电荷 \(Q\) 不变时令半径增加 \(dR\)，场能变化为

\[
dU=-\frac{Q^2}{8\pi\varepsilon_0R^2}\,dR.
\]

电压力对球面向外做功

\[
dW=p\,(4\pi R^2)dR.
\]

能量守恒要求 \(dW=-dU\)，于是

\[
p\,(4\pi R^2)
=\frac{Q^2}{8\pi\varepsilon_0R^2},
\]

再次得到

\[
p=\frac{Q^2}{32\pi^2\varepsilon_0R^4}.
\]

这一检验同时说明，压力公式中的因子 \(1/2\) 不能遗漏。

\section{专题十四：静电能、电容与虚功法（基础必讲与拔高选讲）}\label{m2:ux4e13ux9898ux5341ux56dbux9759ux7535ux80fdux7535ux5bb9ux4e0eux865aux529fux6cd5ux57faux7840ux5fc5ux8bb2ux4e0eux62d4ux9ad8ux9009ux8bb2}

\begin{quote}
阅读准备：矩阵关系与``固定什么'' 多导体关系
\(\mathbf Q=\mathbf C\mathbf V\) 只是把一组线性方程集中书写；第 \(i\)
行表示 \(Q_i\)
怎样由所有导体电势决定。由能量对几何坐标求力时，偏导符号下必须注明保持电荷还是保持电势，因为外电源是否交换能量会改变所用势能函数。

\textbf{例 M2-057：固定电荷与固定电压使用不同能量约束}
\end{quote}

真空平行板电容器的有效极板面积为 \(A\)，板间距为可变参数
\(d>0\)，忽略边缘场。令机械力沿增大板间距方向的分量为
\(F_d\)。分别考虑极板与外界绝缘、自由电荷 \(Q\)
固定，以及连接理想电源、电压 \(V\)
固定两种过程。利用虚功求两种约束下的力，写明应使用的能量函数，并在固定电压过程中计入电源的能量交换。

\par\noindent\begin{minipage}{\linewidth}
\textbf{解答}

\begin{center}
\includegraphics[width=\linewidth,height=.40\textheight,keepaspectratio]{\detokenize{assets/module-02/examples/M2-057.pdf}}
\par\smallskip{\small 图 M2-057：固定电荷与固定电压使用不同能量约束}
\end{center}
\end{minipage}\par

图 M2-057
的两板受力趋于减小间距。若固定电荷，直接对场能求距离导数；若接理想恒压电源，则须把电源交换能量并入
\(-CV^2/2\)，两种约束最后都给出吸引方向。

图 M2-057 中 \(d\) 是两板间距，\(C=\varepsilon_0A/d\)。固定 \(Q\)
时，场能与力分别为

\[
U_Q=\frac{Q^2d}{2\varepsilon_0A},\qquad
F_{d,Q}=-\left.\frac{\partial U_Q}{\partial d}\right|_Q
=-\frac{Q^2}{2\varepsilon_0A}.
\]

固定 \(V\) 时，把理想电源的能量交换计入有效势能

\[
\mathcal F_V=-\frac12CV^2=-\frac{\varepsilon_0AV^2}{2d},
\qquad F_{d,V}=-\left.\frac{\partial\mathcal F_V}{\partial d}\right|_V
=-\frac{\varepsilon_0AV^2}{2d^2}.
\]

两种 \(F_d\) 都沿减小间距方向。恒压过程中，\(dQ=V\,dC\)，电源向系统做功

\[
dW_{\rm src}=V\,dQ=V^2dC,\qquad
dU_{\rm field}=\frac12V^2dC.
\]

其余一半成为场对极板所做的机械功：
\(dW_{\rm mech}=dW_{\rm src}-dU_{\rm field}=\tfrac12V^2dC=F_{d,V}\,dd\)。若漏掉电源交换，只对
\(CV^2/2\) 求负导数，便会得到错误的力方向。

静电能由电荷、场和导体几何共同决定。对称性足够强时，可以先求场再积分能量；几何参数变化时，能量导数给出力。计算前先判断外部约束：固定电荷、固定电势，或多导体的混合约束。

\par\noindent\begin{minipage}{\linewidth}
\begin{center}
\includegraphics[width=\linewidth,height=.40\textheight,keepaspectratio]{\detokenize{assets/module-02/original/tikz-m2-topic14-energy-v1-01-three-views.png}}
\par\smallskip{\small 图示：静电能可以从组装电荷、场能密度和导体电势三个角度计算}
\end{center}
\end{minipage}\par

\subsection{一、从电荷组装功得到静电能}\label{m2:ux4e00ux4eceux7535ux8377ux7ec4ux88c5ux529fux5f97ux5230ux9759ux7535ux80fd}

\subsubsection{1.
离散点电荷的相互作用能}\label{m2:ux79bbux6563ux70b9ux7535ux8377ux7684ux76f8ux4e92ux4f5cux7528ux80fd}

把 \(N\) 个点电荷从无穷远依次缓慢搬到指定位置。放置第 \(i\)
个电荷时，只计算此前已经放好的电荷在该点产生的电势，不计它自己的电势。总组装功为

\[
U
=\sum_{i<j}\frac{1}{4\pi\varepsilon_0}
\frac{q_iq_j}{r_{ij}}.
\]

若记 \(\varphi_{\ne i}(\mathbf r_i)\) 为除第 \(i\)
个电荷外其余电荷产生的电势，则也可写成

\[
\boxed{
U=\frac12\sum_{i=1}^N q_i\varphi_{\ne i}(\mathbf r_i)
}.
\]

因子 \(1/2\)
用来消除每一对电荷在求和中的重复计数。理想点电荷自身的库仑场能发散，不能把
\(q_i\varphi(\mathbf r_i)\) 中的自势直接算入上述有限相互作用能。

\textbf{例 M2-058：两点电荷只有一对相互作用}

真空中两个理想点电荷 \(q_1,q_2\) 相距
\(d>0\)，无其他电荷。将两电荷在无穷远分离时的相互作用能定义为零，只求两者之间的有限相互作用部分，不计各点电荷发散的自能。求相互作用能，说明其正负与两电荷符号的关系，并解释用逐对求和时为何不能将同一电荷对重复计数。

\par\noindent\begin{minipage}{\linewidth}
\textbf{解答}

\begin{center}
\includegraphics[width=\linewidth,height=.40\textheight,keepaspectratio]{\detokenize{assets/module-02/examples/M2-058.pdf}}
\par\smallskip{\small 图 M2-058：两点电荷只有一对相互作用}
\end{center}
\end{minipage}\par

图 M2-058 只包含一个由两枚不同点电荷组成的配对，间距为
\(d\)。把其中一枚从无穷远搬到另一枚产生的势中，直接得到
\(kq_1q_2/d\)；若使用双重求和，才需另用 \(1/2\) 消去重复计数。

相距 \(d\) 的 \(q_1,q_2\) 只有一对相互作用，因此

\[
\boxed{U=\frac{q_1q_2}{4\pi\varepsilon_0d}}.
\]

同号时 \(U>0\)，异号时
\(U<0\)；这里以两电荷相距无穷远时的相互作用能为零。

\subsubsection{2.
连续电荷分布}\label{m2:ux8fdeux7eedux7535ux8377ux5206ux5e03}

对有限尺寸、足够光滑的连续电荷分布，把电荷从零逐渐增加到最终分布。在线性真空中，若中间状态取
\(\rho_s=s\rho\)，则 \(\varphi_s=s\varphi\)。于是

\[
\begin{aligned}
U
&=\int_0^1ds\int \varphi_s(\mathbf r)\rho(\mathbf r)\,dV\\
&=\int_0^1s\,ds\int\rho\varphi\,dV.
\end{aligned}
\]

因此

\[
\boxed{U=\frac12\int\rho(\mathbf r)\varphi(\mathbf r)\,dV}.
\]

对体电荷、面电荷和线电荷分别有

\[
U=\frac12\int_V\rho\varphi\,dV,
\qquad
U=\frac12\int_S\sigma\varphi\,dS,
\qquad
U=\frac12\int_L\lambda\varphi\,dl,
\]

前提是相应积分有限，且 \(\varphi\)
包含全部其他电荷及连续分布中正确的自相互作用。理想点电荷、零半径线电荷等模型需要把发散自能与有限相互作用能分开。

\textbf{例 M2-059：等势导体的能量积分}

真空中一个孤立导体带总电荷 \(Q\)，静电平衡时整个导体电势为
\(V\)，取无穷远电势为零，外部没有其他源。设表面电荷密度为
\(\sigma(\mathbf r)\)，可不均匀。由表面电荷建立过程对应的能量积分求导体静电能，说明等势条件怎样把面积积分化为
\(Q,V\) 的表达式。

\par\noindent\begin{minipage}{\linewidth}
\textbf{解答}

\begin{center}
\includegraphics[width=\linewidth,height=.40\textheight,keepaspectratio]{\detokenize{assets/module-02/examples/M2-059.pdf}}
\par\smallskip{\small 图 M2-059：等势导体的能量积分}
\end{center}
\end{minipage}\par

图 M2-059 的整个导体表面保持同一电势 \(V\)，所以可把 \(V\)
提出表面积分。剩余的面密度积分就是总电荷 \(Q\)，再保留组装电荷能量中的
\(1/2\)，得到 \(QV/2\)。

若一块孤立导体处于电势 \(V\)，其表面总电荷为 \(Q\)，则表面各点电势相同：

\[
U=\frac12\int_S\sigma V\,dS
=\boxed{\frac12QV}.
\]

\subsection{二、从电荷能量变换到场能量}\label{m2:ux4e8cux4eceux7535ux8377ux80fdux91cfux53d8ux6362ux5230ux573aux80fdux91cf}

在区域 \(V\) 内，利用

\[
\rho=\varepsilon_0\nabla\cdot\mathbf E,
\qquad
\mathbf E=-\nabla\varphi,
\]

以及

\[
\nabla\cdot(\varphi\mathbf E)
=\mathbf E\cdot\nabla\varphi+\varphi\nabla\cdot\mathbf E
=-E^2+\varphi\nabla\cdot\mathbf E,
\]

可得

\[
\begin{aligned}
\frac12\int_V\rho\varphi\,dV
&=\frac{\varepsilon_0}{2}\int_V
\varphi\nabla\cdot\mathbf E\,dV\\
&=\frac{\varepsilon_0}{2}\int_VE^2\,dV
+\frac{\varepsilon_0}{2}\oint_{\partial V}
\varphi\mathbf E\cdot\mathbf n\,dS.
\end{aligned}
\]

等价地，因 \(\mathbf E\cdot\mathbf n=-\partial_n\varphi\)，

\[
\boxed{
\frac12\int_V\rho\varphi\,dV
=\frac{\varepsilon_0}{2}\int_V|\nabla\varphi|^2\,dV
-\frac{\varepsilon_0}{2}\oint_{\partial V}
\varphi\frac{\partial\varphi}{\partial n}\,dS
}.
\]

\par\noindent\begin{minipage}{\linewidth}
\begin{center}
\includegraphics[width=\linewidth,height=.40\textheight,keepaspectratio]{\detokenize{assets/module-02/original/tikz-m2-topic14-energy-v1-02-boundary-term.png}}
\par\smallskip{\small 图示：分部积分把电荷能量分成区域内场能和边界项}
\end{center}
\end{minipage}\par

若取全空间，电荷局域分布且无穷远处
\(\varphi=O(r^{-1})\)、\(E=O(r^{-2})\)，则半径为 \(R\) 的外球面边界项为
\(O(R^{-1})\)，在 \(R\to\infty\) 时消失，于是

\[
\boxed{
U=\frac{\varepsilon_0}{2}\int_{\mathbb R^3}E^2\,dV
},
\qquad
\boxed{u_E=\frac12\varepsilon_0E^2}.
\]

场能密度 \(u_E\)
非负，但这不与异号点电荷的有限相互作用能为负矛盾：理想点电荷各自具有发散的正自能；从总场能中减去固定的发散自能后，剩余的交叉项可以为负。

\textbf{例 M2-060：平行板场能等于能量密度乘体积}

真空平行板电容器两极板有效面积均为 \(A\)，板间距为 \(d>0\)，电势差为
\(V\)。忽略边缘效应，将板间电场视为均匀，能量主要存于体积 \(Ad\)
内。利用真空场能密度求总场能，并与电容器的整体能量表达式比较，明确所用电场大小与几何参数的关系。

\par\noindent\begin{minipage}{\linewidth}
\textbf{解答}

\begin{center}
\includegraphics[width=\linewidth,height=.40\textheight,keepaspectratio]{\detokenize{assets/module-02/examples/M2-060.pdf}}
\par\smallskip{\small 图 M2-060：平行板场能等于能量密度乘体积}
\end{center}
\end{minipage}\par

图 M2-060 的近似匀强场仅占两板之间体积 \(Ad\)。用 \(E=V/d\)
代入场能密度，再乘这一体积，即可与由电容定义得到的 \(CV^2/2\)
对照；边缘场已在题设中忽略。

忽略边缘效应，两板面积为 \(A\)、间距为 \(d\)，板间场强 \(E=V/d\)。场能为

\[
U=\frac12\varepsilon_0E^2Ad
=\frac12\frac{\varepsilon_0A}{d}V^2.
\]

由 \(C=\varepsilon_0A/d\)，得到 \(U=CV^2/2\)。

\subsection{三、单导体电容、双导体电容与电容系数矩阵}\label{m2:ux4e09ux5355ux5bfcux4f53ux7535ux5bb9ux53ccux5bfcux4f53ux7535ux5bb9ux4e0eux7535ux5bb9ux7cfbux6570ux77e9ux9635}

\subsubsection{1.
单导体和双导体电容}\label{m2:ux5355ux5bfcux4f53ux548cux53ccux5bfcux4f53ux7535ux5bb9}

孤立导体以无穷远为电势零点时，电容定义为

\[
\boxed{C=\frac{Q}{V}}.
\]

两导体电容器带等量异号电荷 \(+Q,-Q\)，两导体电势分别为 \(V_1,V_2\)，则

\[
\boxed{C=\frac{Q}{V_1-V_2}}.
\]

在线性真空中，电容只由导体几何和介质决定，与 \(Q,V\)
的大小无关。其能量有三种等价写法：

\[
\boxed{
U=\frac{Q^2}{2C}
=\frac12Q(V_1-V_2)
=\frac12C(V_1-V_2)^2
}.
\]

\textbf{例 M2-061：孤立球的电容与自能}

真空中一个孤立导体球半径为 \(R>0\)，带总电荷
\(Q\)，外界没有其他电荷或场，取无穷远电势为零。由球面电势求相对于无穷远的自电容，并计算静电能；说明能量是储存在导体内部还是外部场区。

\par\noindent\begin{minipage}{\linewidth}
\textbf{解答}

\begin{center}
\includegraphics[width=\linewidth,height=.40\textheight,keepaspectratio]{\detokenize{assets/module-02/examples/M2-061.pdf}}
\par\smallskip{\small 图 M2-061：孤立球的电容与自能}
\end{center}
\end{minipage}\par

图 M2-061
的场箭头只位于球面外侧，内部是等势导体。由表面半径处的库仑势读取
\(V\)，再按 \(Q/V\) 定义电容，最后代入
\(QV/2\)；这样得到的能量含全空间外部场的自能。

半径为 \(R\) 时，球面电势为 \(V=Q/(4\pi\varepsilon_0R)\)，故

\[
\boxed{C=4\pi\varepsilon_0R},
\qquad
\boxed{U=\frac{Q^2}{8\pi\varepsilon_0R}}.
\]

\subsubsection{2.
多导体的电容系数矩阵（拔高选讲）}\label{m2:ux591aux5bfcux4f53ux7684ux7535ux5bb9ux7cfbux6570ux77e9ux9635ux62d4ux9ad8ux9009ux8bb2}

取一个参考导体接地，或规定无穷远电势为零。其余 \(N\) 个导体的电势为
\(V_i\)、电荷为 \(Q_i\)。泊松方程的线性性给出

\[
\boxed{Q_i=\sum_{j=1}^NC_{ij}V_j}.
\]

格林互易关系保证

\[
\boxed{C_{ij}=C_{ji}}.
\]

在通常约定下，\(C_{ii}>0\)，而 \(i\ne j\) 时
\(C_{ij}\le0\)。若把参考导体也纳入完整矩阵，则每一行元素之和为零；若只列出相对于接地参考体的
\(N\times N\) 子矩阵，行和一般不为零，不能机械套用``行和为零''。

多导体静电能为

\[
\boxed{
U=\frac12\sum_iQ_iV_i
=\frac12\sum_{i,j}C_{ij}V_iV_j
=\frac12\mathbf V^{\mathsf T}\mathbf C\mathbf V
}.
\]

在参考电势固定后，\(\mathbf C\) 对独立导体电势是正定的，因此
\(U>0\)，除非所有 \(V_i=0\)。

\par\noindent\begin{minipage}{\linewidth}
\begin{center}
\includegraphics[width=\linewidth,height=.40\textheight,keepaspectratio]{\detokenize{assets/module-02/original/tikz-m2-topic14-energy-v1-03-capacitance-matrix.png}}
\par\smallskip{\small 图示：单导体、双导体和多导体电容的电荷电势关系}
\end{center}
\end{minipage}\par

\textbf{例 M2-062：完整二导体矩阵保留规范零模}

考虑只保留两导体之间互电容的理想二导体电容器，电容为
\(C>0\)，不另列接地参考导体。两导体电势为
\(V_1,V_2\)，电荷只以等量异号形式存储。把两导体同时列入电荷---电势关系，写出
\(2\times2\)
电容矩阵；讨论将两个电势同时加上任意常数后电荷是否改变，以及矩阵对应的零模含义。

\par\noindent\begin{minipage}{\linewidth}
\textbf{解答}

\begin{center}
\includegraphics[width=\linewidth,height=.40\textheight,keepaspectratio]{\detokenize{assets/module-02/examples/M2-062.pdf}}
\par\smallskip{\small 图 M2-062：完整二导体矩阵保留规范零模}
\end{center}
\end{minipage}\par

图 M2-062 的两块板都保留在未知电势向量中，电荷只随差值 \(V_1-V_2\)
变化。给两个电势一起加常数不改变任何板间电压，所以矩阵乘 \((1,1)\)
为零，行和为零正是这一自由度的代数表示。

若把两块导体都列入且不另取接地参考体，则

\[
\begin{pmatrix}Q_1\\Q_2\end{pmatrix}
=C
\begin{pmatrix}1&-1\\-1&1\end{pmatrix}
\begin{pmatrix}V_1\\V_2\end{pmatrix}.
\]

因此 \(Q_1=C(V_1-V_2)=-Q_2\)，矩阵各行之和为零；整体平移 \(V_1,V_2\)
不改变电荷和电场。

\subsection{四、虚功法：固定电荷与固定电势必须分开}\label{m2:ux56dbux865aux529fux6cd5ux56faux5b9aux7535ux8377ux4e0eux56faux5b9aux7535ux52bfux5fc5ux987bux5206ux5f00}

设导体几何由广义坐标 \(x\) 描述，电容为 \(C(x)\)。规定广义力 \(F_x\)
的正方向为 \(x\) 增大的方向。

\subsubsection{1. 固定电荷}\label{m2:ux56faux5b9aux7535ux8377}

系统与外界绝缘，\(Q\) 固定时，场能

\[
U_Q(x)=\frac{Q^2}{2C(x)}.
\]

准静态位移 \(dx\) 中，电场对机械部分做功 \(F_xdx=-dU_Q\)，所以

\[
\boxed{
F_x=-\left(\frac{\partial U}{\partial x}\right)_Q
=\frac{Q^2}{2C^2}\frac{dC}{dx}
}.
\]

\subsubsection{2. 固定电势}\label{m2:ux56faux5b9aux7535ux52bf}

电容器接在理想恒压电源上时，\(V\) 固定而 \(Q=C(x)V\)
会改变。此时不能只对场能 \(U_{\rm \text{场}}=CV^2/2\)
取负导数，因为电源还要转移电荷并交换能量。

电源维持电势 \(V\) 时，有效机械势能为

\[
\mathcal G
=U_{\rm \text{场}}-VQ
=\frac12CV^2-CV^2
=-\frac12CV^2.
\]

因此

\[
\boxed{
F_x=-\left(\frac{\partial\mathcal G}{\partial x}\right)_V
=\frac12V^2\frac{dC}{dx}
}.
\]

\par\noindent\begin{minipage}{\linewidth}
\begin{center}
\includegraphics[width=\linewidth,height=.40\textheight,keepaspectratio]{\detokenize{assets/module-02/original/tikz-m2-topic14-energy-v3-04-fixed-q-v-safe-v7.png}}
\par\smallskip{\small 图示：固定电荷时只有场能变化，固定电势时还必须计入电源交换的能量}
\end{center}
\end{minipage}\par

两式都表明系统趋向于使电容增大的方向，但它们对应不同的控制条件，数值一般不同。

\textbf{例 M2-063：板间吸引使距离减小}

真空平行板电容器面积为 \(A\)，板间距为 \(d>0\)，忽略边缘效应，电容
\(C(d)=\varepsilon_0A/d\)。以增大 \(d\) 的方向为广义力 \(F_d\)
的正方向。分别在自由电荷 \(Q\) 固定、电压 \(V\) 由理想电源固定时用虚功求
\(F_d\)，说明电源功应计入哪个过程，并据符号判断两板相吸还是相斥。

\par\noindent\begin{minipage}{\linewidth}
\textbf{解答}

\begin{center}
\includegraphics[width=\linewidth,height=.40\textheight,keepaspectratio]{\detokenize{assets/module-02/examples/M2-063.pdf}}
\par\smallskip{\small 图 M2-063：板间吸引使距离减小}
\end{center}
\end{minipage}\par

图 M2-063 的距离 \(d\) 取增大板间距为正，而红色吸引力指向减小
\(d\)。对两种约束各自使用适当能量，均得到负广义力；固定电压情形不能只对正的场能
\(CV^2/2\) 取负导数。

取 \(x=d\)，\(C=\varepsilon_0A/d\)，故
\(dC/dd=-\varepsilon_0A/d^2<0\)。于是

\[
\boxed{F_d=-\frac{Q^2}{2\varepsilon_0A}}
\quad(Q\text{ 固定}),
\]

\[
\boxed{F_d=-\frac{\varepsilon_0AV^2}{2d^2}}
\quad(V\text{ 固定}).
\]

负号表示力使两板间距减小，即两板相互吸引。

\subsubsection{3.
多导体推广（拔高选讲）}\label{m2:ux591aux5bfcux4f53ux63a8ux5e7fux62d4ux9ad8ux9009ux8bb2}

若各导体电荷固定，可用电势系数矩阵 \(\mathbf P=\mathbf C^{-1}\)：

\[
U_Q=\frac12\mathbf Q^{\mathsf T}\mathbf P\mathbf Q,
\qquad
F_x=-\frac12\mathbf Q^{\mathsf T}
\frac{\partial\mathbf P}{\partial x}\mathbf Q.
\]

若各导体电势由独立理想电源固定，则

\[
\mathcal G_V=-\frac12\mathbf V^{\mathsf T}\mathbf C\mathbf V,
\qquad
\boxed{
F_x=\frac12\mathbf V^{\mathsf T}
\frac{\partial\mathbf C}{\partial x}\mathbf V
}.
\]

若只有部分电荷、部分电势固定，应先把相应自由度消去，再对与实际约束匹配的能量函数求导。

\subsection{五、麦克斯韦应力张量与导体表面电压力（拔高选讲）}\label{m2:ux4e94ux9ea6ux514bux65afux97e6ux5e94ux529bux5f20ux91cfux4e0eux5bfcux4f53ux8868ux9762ux7535ux538bux529bux62d4ux9ad8ux9009ux8bb2}

静电真空中的麦克斯韦应力张量为

\[
\boxed{
\mathbf T
=\varepsilon_0
\left(
\mathbf E\mathbf E-\frac12E^2\mathbf I
\right)
}.
\]

对包围物体并完全位于真空中的闭曲面 \(S\)，电磁力为

\[
\boxed{
\mathbf F=\oint_S\mathbf T\cdot\mathbf n\,dS
}.
\]

导体表面外侧电场纯法向，\(\mathbf E=E_n\mathbf n\)。应力张量给出的法向拉力密度为

\[
\boxed{
p=\frac12\varepsilon_0E_n^2
=\frac{\sigma^2}{2\varepsilon_0}
}.
\]

这里的 \(E_n\)
是导体外侧极限。若误把带电面自身在两侧产生的场相加，就会漏掉或多出因子
\(1/2\)。

\par\noindent\begin{minipage}{\linewidth}
\begin{center}
\includegraphics[width=\linewidth,height=.40\textheight,keepaspectratio]{\detokenize{assets/module-02/original/tikz-m2-topic14-energy-v2-05-stress-pressure.png}}
\par\smallskip{\small 图示：导体表面外侧法向电场产生电压力，局部应力与整体虚功可以互相检验}
\end{center}
\end{minipage}\par

\textbf{例 M2-064：球面虚位移检验电压力}

真空中一个半径为 \(R\) 的孤立导体球带固定电荷
\(Q\)，无外加场，取无穷远零势。先由表面电荷密度求球面的电压力；再令球面作保持
\(Q\) 不变的均匀虚位移
\(dR\)，用静电能变化与压力在球面上所做总功核对压力大小及方向。

\par\noindent\begin{minipage}{\linewidth}
\textbf{解答}

\begin{center}
\includegraphics[width=\linewidth,height=.40\textheight,keepaspectratio]{\detokenize{assets/module-02/examples/M2-064.pdf}}
\par\smallskip{\small 图 M2-064：球面虚位移检验电压力}
\end{center}
\end{minipage}\par

图 M2-064
的球面各点具有相同场强和压力，但法向方向随位置改变。这里做的是各面元沿本地方向同时外移的虚功，不是求整体平动力；把压力乘以
\(4\pi R^2dR\) 才能与固定电荷的能量变化比较。

半径为 \(R\)、电荷为 \(Q\) 时，

\[
E_n=\frac{Q}{4\pi\varepsilon_0R^2},
\]

所以

\[
\boxed{
p=\frac{Q^2}{32\pi^2\varepsilon_0R^4}
}.
\]

另一方面，固定 \(Q\) 时 \(U=Q^2/(8\pi\varepsilon_0R)\)。令球面向外移动
\(dR\)，有

\[
p(4\pi R^2)dR=-dU,
\]

得到同一结果。

\subsection{六、固定电荷与固定电势条件的选择}\label{m2:ux516dux56faux5b9aux7535ux8377ux4e0eux56faux5b9aux7535ux52bfux6761ux4ef6ux7684ux9009ux62e9}

\par\noindent\begin{minipage}{\linewidth}
\begin{center}
\includegraphics[width=\linewidth,height=.40\textheight,keepaspectratio]{\detokenize{assets/module-02/original/tikz-m2-topic14-energy-v1-06-method-map.png}}
\par\smallskip{\small 图示：由控制量、几何和目标物理量选择静电能计算方法}
\end{center}
\end{minipage}\par

\begin{enumerate}
\def\labelenumi{\arabic{enumi}.}
\tightlist
\item
  \textbf{先写控制条件。} ``电容器''三个字不足以决定能量导数；必须写明
  \(Q\) 固定还是 \(V\) 固定。
\item
  \textbf{点电荷先排除自能。} 计算相互作用能时用
  \(\varphi_{\ne i}\)；镜像法中真实电荷与感应电荷系统的能量也不能直接等同于真实电荷和像电荷的两体能量。
\item
  \textbf{场能积分要覆盖全部有场区域。}
  只积电容器板间区域依赖于忽略边缘场的近似。
\item
  \textbf{有限区域保留边界项。} \(\frac12\int\rho\varphi\) 与
  \(\frac12\varepsilon_0\int E^2\)
  只有在边界项消失或已被正确处理时才直接相等。
\item
  \textbf{压力是局部标量，合力是矢量积分。}
  球面各处电压力均向外，但整个孤立球所受外部合力仍为零。
\item
  \textbf{量纲检查。} 电容单位为
  \(\mathrm F=\mathrm C/\mathrm V\)；\(Q^2/C\)、\(CV^2\) 与
  \(\varepsilon_0E^2V_{\text{体积}}\) 都应为焦耳。
\end{enumerate}

\subsection{七、练习与解答}\label{m2:ux4e03ux7ec3ux4e60ux4e0eux89e3ux7b54-1}

\subsubsection{作业一：平行板电容器的能量与吸力}\label{m2:ux4f5cux4e1aux4e00ux5e73ux884cux677fux7535ux5bb9ux5668ux7684ux80fdux91cfux4e0eux5438ux529b}

\textbf{题号：M2-065}

\textbf{层级：} 基础　　\textbf{方法：} 能量法

真空中两块平行导体板的有效面积均为 \(A>0\)，间距为
\(d>0\)；忽略边缘效应和板外场能，板间带等量异号自由电荷。

\begin{enumerate}
\def\labelenumi{(\arabic{enumi})}
\item
  求电容及用 \(Q\) 或 \(V\) 表示的场能。
\item
  以增大 \(d\) 为广义坐标正向，分别求孤立固定电荷
  \(Q\)、连接理想恒压源固定电压 \(V\) 时的广义力。
\item
  用导体表面的电压力复核吸力大小，并说明固定电压时电源的能量交换。
\end{enumerate}

\par\noindent\begin{minipage}{\linewidth}
\textbf{解答}

\begin{center}
\includegraphics[width=\linewidth,height=.40\textheight,keepaspectratio]{\detokenize{assets/module-02/examples/M2-065.pdf}}
\par\smallskip{\small 图 M2-065：平行板电容器的能量与吸力}
\end{center}
\end{minipage}\par

图 M2-065
的板距是机械坐标，正方向使两板分离。场能积分得到电容后，分别固定 \(Q\)
和 \(V\)
求相应势能导数；电压力的大小乘板面积应与两种力公式在同一状态下相符。

板间场均匀，

\[
C=\frac{\varepsilon_0A}{d}.
\]

固定 \(Q\) 时

\[
U_Q=\frac{Q^2d}{2\varepsilon_0A},
\qquad
\boxed{F_d=-\frac{Q^2}{2\varepsilon_0A}}.
\]

固定 \(V\) 时使用含电源交换的有效势能，

\[
\mathcal G_V=-\frac{\varepsilon_0AV^2}{2d},
\qquad
\boxed{F_d=-\frac{\varepsilon_0AV^2}{2d^2}}.
\]

板外侧近似无场，板间场为 \(E=Q/(\varepsilon_0A)=V/d\)。电压力

\[
p=\frac12\varepsilon_0E^2.
\]

一块板所受吸力大小为 \(pA\)，分别给出上述两式的绝对值。

\subsubsection{作业二：同轴圆柱电容器}\label{m2:ux4f5cux4e1aux4e8cux540cux8f74ux5706ux67f1ux7535ux5bb9ux5668}

\textbf{题号：M2-066}

\textbf{层级：} 基础　　\textbf{方法：} 能量法

真空中两根很长的同轴导体圆柱构成电容器，内柱外半径为 \(a\)，外柱内半径为
\(b\)，\(0<a<b\)；有效长度 \(L\)
远大于横向尺寸，忽略端部效应。内柱每单位长度带自由电荷
\(\lambda\)，外柱面对内柱的表面带相应负电荷，板间无其他源。

\begin{enumerate}
\def\labelenumi{(\arabic{enumi})}
\item
  求内柱相对外柱的电势差、有效长度 \(L\) 的电容和单位长度场能。
\item
  求内柱外表面电压力。
\item
  在 \(a<\rho<b\) 的场区积分能量密度，复核单位长度场能。
\end{enumerate}

\par\noindent\begin{minipage}{\linewidth}
\textbf{解答}

\begin{center}
\includegraphics[width=\linewidth,height=.40\textheight,keepaspectratio]{\detokenize{assets/module-02/examples/M2-066.pdf}}
\par\smallskip{\small 图 M2-066：同轴圆柱电容器}
\end{center}
\end{minipage}\par

图 M2-066 的同轴圆柱把径向积分限制在
\(a<\rho<b\)，端部效应忽略。高斯柱给
\(E_{\rho}\propto1/\rho\)，积分电势或能量都产生同一个
\(\ln(b/a)\)；电压力则读取内柱外表面 \(\rho=a\) 的局部场。

在 \(a<r<b\) 中，高斯定律给出

\[
E(r)=\frac{\lambda}{2\pi\varepsilon_0r}.
\]

故

\[
V_a-V_b
=\int_a^bE(r)\,dr
=\frac{\lambda}{2\pi\varepsilon_0}\ln\frac ba.
\]

总电荷 \(Q=\lambda L\)，所以

\[
\boxed{C=\frac{2\pi\varepsilon_0L}{\ln(b/a)}}.
\]

单位长度能量为

\[
\boxed{
\frac UL
=\frac12\lambda(V_a-V_b)
=\frac{\lambda^2}{4\pi\varepsilon_0}\ln\frac ba
}.
\]

内柱表面外侧场强为 \(E(a)=\lambda/(2\pi\varepsilon_0a)\)，故电压力

\[
\boxed{p(a)=\frac{\lambda^2}{8\pi^2\varepsilon_0a^2}}.
\]

它沿内柱材料的外法线，表现为局部向外的张应力；绕圆周积分后的合力矢量为零。

场能积分为

\[
\begin{aligned}
\frac UL
&=\int_a^b\frac12\varepsilon_0E^2(2\pi r)\,dr\\
&=\frac{\lambda^2}{4\pi\varepsilon_0}
\int_a^b\frac{dr}{r},
\end{aligned}
\]

与前式相同。

\subsubsection{作业三：孤立带电导体球的三重检验}\label{m2:ux4f5cux4e1aux4e09ux5b64ux7acbux5e26ux7535ux5bfcux4f53ux7403ux7684ux4e09ux91cdux68c0ux9a8c}

\textbf{题号：M2-067}

\textbf{层级：} 基础　　\textbf{方法：} 能量法

真空中半径为 \(R>0\) 的孤立导体球带总电荷 \(Q\)，外部无其他源，取
\(\varphi(\infty)=0\)。分别利用
\(U=QV/2\)、真空场能密度的全空间积分求静电能；再保持 \(Q\)
固定，使球面发生均匀微小半径变化，用虚功求表面电压力，并与局部表面场计算比较。

\par\noindent\begin{minipage}{\linewidth}
\textbf{解答}

\begin{center}
\includegraphics[width=\linewidth,height=.40\textheight,keepaspectratio]{\detokenize{assets/module-02/examples/M2-067.pdf}}
\par\smallskip{\small 图 M2-067：孤立带电导体球的三重检验}
\end{center}
\end{minipage}\par

图 M2-067 的球面外侧是场能积分区域。先由表面电势作 \(QV/2\)，再对外域的
\(E^2\)
作体积分，所得自能应一致；固定电荷下对半径求导，并用球面积换成局部压力，完成第三项检验。

球面电势与电容为

\[
V=\frac{Q}{4\pi\varepsilon_0R},
\qquad
C=4\pi\varepsilon_0R.
\]

所以

\[
U=\frac12QV
=\boxed{\frac{Q^2}{8\pi\varepsilon_0R}}.
\]

球外 \(E=Q/(4\pi\varepsilon_0r^2)\)，球内无场，因此

\[
\begin{aligned}
U
&=\int_R^{\infty}\frac12\varepsilon_0E^2(4\pi r^2)\,dr\\
&=\frac{Q^2}{8\pi\varepsilon_0}
\int_R^{\infty}\frac{dr}{r^2}
=\frac{Q^2}{8\pi\varepsilon_0R}.
\end{aligned}
\]

固定 \(Q\) 时

\[
p(4\pi R^2)=-\frac{dU}{dR},
\]

故

\[
\boxed{p=\frac{Q^2}{32\pi^2\varepsilon_0R^4}}
=\frac12\varepsilon_0E^2(R^+).
\]

\subsubsection{作业四：两极板电荷、储能与电势零点}\label{m2:ux4f5cux4e1aux56dbux4e24ux6781ux677fux7535ux8377ux50a8ux80fdux4e0eux7535ux52bfux96f6ux70b9}

\textbf{题号：M2-068}

\textbf{层级：} 基础　　\textbf{方法：} 平行板电场、能量法

真空中两块平行金属板的正对面积均为 \(S\)，板间距为
\(d\)，忽略边缘效应，只考虑两板之间的电场。两板所带总电荷为零；电容为
\(C_0=\varepsilon_0S/d\)。相对于同一电势零点，两板电势分别为
\(V_1,V_2\)，记板间电压 \(U_{12}=V_1-V_2\)。

\begin{enumerate}
\def\labelenumi{(\arabic{enumi})}
\item
  先以 \(V_1>V_2\) 为例，在图中指出电场方向和两板电荷的正负；求极板电荷
  \(Q_1,Q_2\)。
\item
  仅更换电势零点，使两个电势数值同时增加同一常数
  \(V_s\)，实际极板和外部连接均不改变。板间电压、电场和电荷会改变吗？以
  \(5\,V,2\,V\) 同时变为 \(15\,V,12\,V\) 说明。
\item
  求板间静电能 \(W\)，说明计算需要的是两个电势各自的数值，还是它们的差。
\end{enumerate}

\par\noindent\begin{minipage}{\linewidth}
\textbf{解答}

\begin{center}
\includegraphics[width=\linewidth,height=.40\textheight,keepaspectratio]{\detokenize{assets/module-02/examples/M2-068.pdf}}
\par\smallskip{\small 图 M2-068：两极板电荷、储能与电势零点}
\end{center}
\end{minipage}\par

图 M2-068
左侧的1号板电势较高，蓝色电场箭头从1号板指向2号板。忽略边缘效应时，

\[
E=\frac{V_1-V_2}{d},\qquad Q_1=\varepsilon_0ES
=C_0(V_1-V_2),\qquad Q_2=-Q_1.
\]

因此左板带正电、右板带等量负电。若
\(V_1<V_2\)，电场箭头和电荷正负同时反转；上面的带符号电荷公式仍成立。

图右侧把两板电势同时加上 \(10\,V\)，但电势差仍是 \(3\,V\)。一般地，

\[
U_{12}^{\prime}=(V_1+V_s)-(V_2+V_s)=U_{12}.
\]

这只是更换电势的数值零点，板间电场和两板电荷均不变，并不是把电容器另接到某个电源。

左图两板之间的场区体积为 \(Sd\)，故

\[
W=\frac12\varepsilon_0E^2Sd
=\boxed{\frac12C_0(V_1-V_2)^2}.
\]

电荷、板间电场和储能都由电势差决定；给两个电势同时加上常数不会改变这些物理量。本题只需电容的定义和均匀场能量公式。

\subsubsection{作业五：变重叠面积电容器}\label{m2:ux4f5cux4e1aux4e94ux53d8ux91cdux53e0ux9762ux79efux7535ux5bb9ux5668}

\textbf{题号：M2-069}

\textbf{层级：} 基础　　\textbf{方法：} 能量法

真空平行板电容器的两板宽度为 \(b>0\)，间距为固定值 \(d>0\)。沿 \(x\)
方向相对平移时，两板的重叠长度为 \(x>0\)，有效面积为
\(bx\)；只在重叠区足够大、可忽略边缘效应的范围讨论。以增大重叠长度为正向，分别求固定极板电荷
\(Q\) 和连接理想电源固定电压 \(V\) 时的水平力，并写明能量法所用约束。

\par\noindent\begin{minipage}{\linewidth}
\textbf{解答}

\begin{center}
\includegraphics[width=\linewidth,height=.40\textheight,keepaspectratio]{\detokenize{assets/module-02/examples/M2-069.pdf}}
\par\smallskip{\small 图 M2-069：变重叠面积电容器}
\end{center}
\end{minipage}\par

图 M2-069 的机械位移增加两板重叠长度，场区面积为 \(bx\)，故电容与 \(x\)
成正比。用与外电路约束相应的能量求导，力均倾向增加重叠；固定电荷与固定电压的大小随
\(x\) 的变化规律不同。

重叠面积 \(A=bx\)，故

\[
C(x)=\frac{\varepsilon_0bx}{d},
\qquad
\frac{dC}{dx}=\frac{\varepsilon_0b}{d}.
\]

固定 \(Q\) 时

\[
\boxed{
F_x=\frac{Q^2}{2C^2}\frac{\varepsilon_0b}{d}
=\frac{Q^2d}{2\varepsilon_0bx^2}
}.
\]

固定 \(V\) 时

\[
\boxed{
F_x=\frac12V^2\frac{\varepsilon_0b}{d}
}.
\]

两种情况下力都使重叠面积增大。固定 \(V\) 的理想模型给出与 \(x\)
无关的力；接近板端时边缘场不可忽略，这一结论会出现修正。

\section{专题十五：静电场唯一性定理（基础必讲）}\label{m2:ux4e13ux9898ux5341ux4e94ux9759ux7535ux573aux552fux4e00ux6027ux5b9aux7406ux57faux7840ux5fc5ux8bb2}

静电边值问题还要确认候选解是否就是物理解。若候选电势在研究区域内具有正确的源，满足微分方程、全部边界条件和必要的无穷远条件，唯一性定理便排除另一个不同的物理解。

\par\noindent\begin{minipage}{\linewidth}
\begin{center}
\includegraphics[width=\linewidth,height=.40\textheight,keepaspectratio]{\detokenize{assets/module-02/original/tikz-m2-topic15-v8-01.png}}
\par\smallskip{\small 图示：泊松方程、边界数据、无穷远条件和奇点共同确定静电解}
\end{center}
\end{minipage}\par

\subsection{一、边值问题由哪些数据组成}\label{m2:ux4e00ux8fb9ux503cux95eeux9898ux7531ux54eaux4e9bux6570ux636eux7ec4ux6210}

静电场满足

\[
\nabla\times\mathbf E=\mathbf 0,
\qquad
\nabla\cdot\mathbf E=\frac{\rho}{\varepsilon_0}.
\]

在单连通静电区域内可引入电势

\[
\mathbf E=-\nabla\varphi,
\]

从而得到泊松方程

\[
\boxed{\nabla^2\varphi=-\frac{\rho}{\varepsilon_0}.}
\]

无电荷区域满足拉普拉斯方程

\[
\boxed{\nabla^2\varphi=0.}
\]

微分方程只规定电势在每一点附近怎样弯曲，并不能单独确定整个区域中的电势。一个完整的静电边值问题至少要说明：

\begin{enumerate}
\def\labelenumi{\arabic{enumi}.}
\tightlist
\item
  求解区域 \(V\) 及其边界 \(S\)；
\item
  区域内的体电荷密度以及点电荷、线电荷等奇点；
\item
  边界上给出的电势或法向导数；
\item
  无界区域中的无穷远行为和电势参考零点。
\end{enumerate}

\textbf{例 M2-070：同一微分方程仍有一族直线解}

在一维无源真空区间 \(0<x<L\) 内，电势只依赖 \(x\)，满足
\(d^2\varphi/dx^2=0\)，尚未给定端点电势。求一般解，指出其中有多少个待定常数；解释只知道微分方程为何不能确定唯一场，并说明给定两端电势后如何固定这两个常数。

\par\noindent\begin{minipage}{\linewidth}
\textbf{解答}

\begin{center}
\includegraphics[width=\linewidth,height=.40\textheight,keepaspectratio]{\detokenize{assets/module-02/examples/M2-070.pdf}}
\par\smallskip{\small 图 M2-070：同一微分方程仍有一族直线解}
\end{center}
\end{minipage}\par

图 M2-070
的几条直线都具有零曲率，因而都满足同一无源方程。方程本身只给出两个积分常数，边界数据才决定选哪一条；不能只验证域内二阶导数为零就宣布得到物理解。

一维无源方程

\[
\frac{d^2\varphi}{dx^2}=0
\]

的通解为

\[
\varphi(x)=Ax+B.
\]

只有给出足够的边界数据后，常数 \(A,B\)
才能确定。因此``满足拉普拉斯方程''远远不足以证明候选解正确。

\subsection{二、三类常用边界条件与定理表述}\label{m2:ux4e8cux4e09ux7c7bux5e38ux7528ux8fb9ux754cux6761ux4ef6ux4e0eux5b9aux7406ux8868ux8ff0}

设 \(\mathbf n\) 为求解区域 \(V\) 边界的外法线。

\subsubsection{1.
第一类边界条件}\label{m2:ux7b2cux4e00ux7c7bux8fb9ux754cux6761ux4ef6}

在全部边界上给定电势：

\[
\varphi|_S=f(\mathbf r).
\]

这类条件也称狄利克雷边界条件。它固定了电势的参考零点，因此电势本身唯一。

\subsubsection{2.
第二类边界条件}\label{m2:ux7b2cux4e8cux7c7bux8fb9ux754cux6761ux4ef6}

在全部边界上给定法向导数：

\[
\left.\frac{\partial\varphi}{\partial n}\right|_S=g(\mathbf r).
\]

由于
\(E_n=-\partial_n\varphi\)，它等价于给定边界法向电场。这类条件也称诺伊曼边界条件。

把泊松方程在 \(V\) 内积分并使用高斯定理，得到

\[
\oint_S\frac{\partial\varphi}{\partial n}\,dS
=\int_V\nabla^2\varphi\,dV
=-\frac1{\varepsilon_0}\int_V\rho\,dV.
\]

所以纯第二类边界数据必须满足相容条件

\[
\boxed{
\oint_S g\,dS
=-\frac{Q_V}{\varepsilon_0}.}
\]

不满足该条件时，不是``解不唯一''，而是\textbf{根本不存在解}。满足相容条件时，电势可整体加上任意常数，但电场唯一。

\subsubsection{3.
混合边界条件}\label{m2:ux6df7ux5408ux8fb9ux754cux6761ux4ef6}

把边界分成互不重叠的两部分 \(S_D\) 与 \(S_N\)：

\[
S=S_D\cup S_N,
\qquad
S_D\cap S_N=\varnothing,
\]

并分别给定

\[
\varphi|_{S_D}=f,
\qquad
\partial_n\varphi|_{S_N}=g.
\]

若 \(S_D\) 含有非零面积的边界部分，通常能够固定加性常数，电势唯一。

\par\noindent\begin{minipage}{\linewidth}
\begin{center}
\includegraphics[width=\linewidth,height=.40\textheight,keepaspectratio]{\detokenize{assets/module-02/original/tikz-m2-topic15-v8-02.png}}
\par\smallskip{\small 图示：第一类、第二类和混合边界条件给出的信息不同}
\end{center}
\end{minipage}\par

\subsubsection{4. 唯一性定理}\label{m2:ux552fux4e00ux6027ux5b9aux7406}

设 \(V\)
为连通区域，电势在除已明确处理的奇点外足够光滑，并给定区域内相同的电荷分布。

\begin{itemize}
\tightlist
\item
  给定第一类边界条件时，泊松方程的解唯一。
\item
  给定满足相容条件的纯第二类边界条件时，任意两个解至多相差一个常数，电场唯一。
\item
  给定第一类与第二类的分段混合边界，且第一类部分非空时，解唯一。
\end{itemize}

``连通''不能省略。若区域有多个互不连通的分支，纯第二类问题可在每个分支上分别加一个常数。

\textbf{例 M2-071：常值球面边界直接给出常值解}

半径为 \(R>0\)
的球形区域内为真空、没有体电荷或点奇点，球面上处处给定相同电势
\(V_0\)。要求电势在整个球内有限。利用边界条件与唯一性定理求球内电势和电场，说明常值边界如何限制内部的非均匀调和分量。

\par\noindent\begin{minipage}{\linewidth}
\textbf{解答}

\begin{center}
\includegraphics[width=\linewidth,height=.40\textheight,keepaspectratio]{\detokenize{assets/module-02/examples/M2-071.pdf}}
\par\smallskip{\small 图 M2-071：常值球面边界直接给出常值解}
\end{center}
\end{minipage}\par

图 M2-071 的整个球面保持同一个
\(V_0\)。常值函数在球内满足拉普拉斯方程，也逐点满足全部边界；由唯一性排除再叠加其他调和函数的可能，球内负梯度因此处处为零。

球内无电荷，整个球面给定 \(\varphi=V_0\)。常数函数 \(\varphi=V_0\)
满足拉普拉斯方程和全部第一类边界条件；由唯一性定理，它就是唯一解，因此球内
\(\mathbf E=\mathbf0\)。

\subsection{三、唯一性定理的证明}\label{m2:ux4e09ux552fux4e00ux6027ux5b9aux7406ux7684ux8bc1ux660e}

\subsubsection{1.
构造两个解的差}\label{m2:ux6784ux9020ux4e24ux4e2aux89e3ux7684ux5dee}

设 \(\varphi_1\) 与 \(\varphi_2\) 是同一边值问题的两个候选解：

\[
\nabla^2\varphi_1=-\frac{\rho}{\varepsilon_0},
\qquad
\nabla^2\varphi_2=-\frac{\rho}{\varepsilon_0}.
\]

令

\[
u=\varphi_1-\varphi_2.
\]

相减后电荷源抵消：

\[
\boxed{\nabla^2u=0.}
\]

第一类边界上有 \(u=0\)；第二类边界上有 \(\partial_n u=0\)。

\subsubsection{2.
建立第一格林恒等式}\label{m2:ux5efaux7acbux7b2cux4e00ux683cux6797ux6052ux7b49ux5f0f}

散度的乘积法则给出

\[
\nabla\cdot(u\nabla u)
=(\nabla u)\cdot(\nabla u)+u\nabla^2u
=|\nabla u|^2+u\nabla^2u.
\]

在 \(V\) 内积分并使用高斯定理：

\[
\oint_Su\frac{\partial u}{\partial n}\,dS
=\int_V|\nabla u|^2\,dV
+\int_Vu\nabla^2u\,dV.
\]

因为 \(\nabla^2u=0\)，得到证明的核心等式

\[
\boxed{
\int_V|\nabla u|^2\,dV
=\oint_Su\frac{\partial u}{\partial n}\,dS.}
\]

\par\noindent\begin{minipage}{\linewidth}
\begin{center}
\includegraphics[width=\linewidth,height=.40\textheight,keepaspectratio]{\detokenize{assets/module-02/original/tikz-m2-topic15-v8-03.png}}
\par\smallskip{\small 图示：唯一性定理的能量积分证明}
\end{center}
\end{minipage}\par

\subsubsection{3.
检查边界积分}\label{m2:ux68c0ux67e5ux8fb9ux754cux79efux5206}

\par\noindent\begin{minipage}{\linewidth}
\begin{center}
\includegraphics[width=\linewidth,height=.40\textheight,keepaspectratio]{\detokenize{assets/module-02/original/tikz-m2-topic15-v8-04.png}}
\par\smallskip{\small 图示：三类边界条件使边界积分分别消失}
\end{center}
\end{minipage}\par

\begin{itemize}
\tightlist
\item
  第一类边界：\(u=0\)，故 \(u\,\partial_n u=0\)。
\item
  第二类边界：\(\partial_n u=0\)，故 \(u\,\partial_n u=0\)。
\item
  混合边界：在 \(S_D\) 上第一因子为零，在 \(S_N\)
  上第二因子为零，故两段积分分别为零。
\end{itemize}

因此

\[
\int_V|\nabla u|^2\,dV=0.
\]

由于 \(|\nabla u|^2\ge0\)，并且 \(u\) 足够光滑，只能有

\[
\nabla u=\mathbf 0.
\]

在连通区域内

\[
u=C.
\]

第一类或含非空第一类部分的混合边界给出 \(C=0\)，所以
\(\varphi_1=\varphi_2\)。纯第二类边界允许任意 \(C\)，但

\[
-\nabla\varphi_1=-\nabla\varphi_2,
\]

故电场仍唯一。

这不是反证法，而是把``两个解可能不同''直接压缩为``它们只能相差常数''的能量积分证明。

\textbf{例 M2-072：零拉普拉斯不等于零梯度}

在三维右手直角坐标系中给定光滑标量场 \(u=x^2+y^2-2z^2\)。分别直接计算
\(\nabla u\)、\(\nabla^2u\) 和 \(\nabla\cdot(u\nabla u)\)，核对恒等式
\(\nabla\cdot(u\nabla u)=|\nabla u|^2+u\nabla^2u\)。据此判断
\(\nabla^2u=0\) 是否意味着梯度在各点均为零。

\par\noindent\begin{minipage}{\linewidth}
\textbf{解答}

\begin{center}
\includegraphics[width=\linewidth,height=.40\textheight,keepaspectratio]{\detokenize{assets/module-02/examples/M2-072.pdf}}
\par\smallskip{\small 图 M2-072：零拉普拉斯不等于零梯度}
\end{center}
\end{minipage}\par

图 M2-072
的蓝色与橙色坐标剖面具有相反曲率，三个方向的二阶导数相加为零。乘积求导时仍有来自梯度平方的项，所以
\(u\) 调和并不使 \(\nabla\cdot(u\nabla u)\)
为零；应逐项核对同一组二次系数。

取调和函数

\[
u=x^2+y^2-2z^2,
\qquad
\nabla^2u=2+2-4=0.
\]

一方面

\[
|\nabla u|^2
=(2x)^2+(2y)^2+(-4z)^2
=4x^2+4y^2+16z^2.
\]

另一方面直接计算也有

\[
\nabla\cdot(u\nabla u)
=4x^2+4y^2+16z^2,
\]

与 \(\nabla\cdot(u\nabla u)=|\nabla u|^2+u\nabla^2u\) 一致。

\subsection{四、纯第二类问题：相容性与加性常数}\label{m2:ux56dbux7eafux7b2cux4e8cux7c7bux95eeux9898ux76f8ux5bb9ux6027ux4e0eux52a0ux6027ux5e38ux6570}

\par\noindent\begin{minipage}{\linewidth}
\begin{center}
\includegraphics[width=\linewidth,height=.40\textheight,keepaspectratio]{\detokenize{assets/module-02/original/tikz-m2-topic15-v8-05.png}}
\par\smallskip{\small 图示：纯第二类问题的总通量相容条件和加性常数}
\end{center}
\end{minipage}\par

第二类边界的两项结论必须分开：

\begin{enumerate}
\def\labelenumi{\arabic{enumi}.}
\tightlist
\item
  \textbf{相容条件控制存在性。} 边界总法向导数必须与区域总电荷匹配。
\item
  \textbf{加性常数控制唯一性。} 若 \(\varphi\) 是解，则 \(\varphi+C\)
  具有相同的法向导数和相同的电场。
\end{enumerate}

\textbf{例 M2-073：相同端点斜率仍不能固定零势}

一维无源区域为 \(0<x<L\)，\(L>0\)，电势满足
\(\varphi''(x)=0\)。先规定两端沿同一 \(+x\) 方向的导数均为常数 \(a\)，即
\(\varphi'(0)=\varphi'(L)=a\)，求解并说明剩余的不确定性；再设两端导数分别为
\(a,b\) 且
\(a\ne b\)，判断是否存在满足所有条件的解。这里的导数不分别改取两端外法向。

\par\noindent\begin{minipage}{\linewidth}
\textbf{解答}

\begin{center}
\includegraphics[width=\linewidth,height=.40\textheight,keepaspectratio]{\detokenize{assets/module-02/examples/M2-073.pdf}}
\par\smallskip{\small 图 M2-073：相同端点斜率仍不能固定零势}
\end{center}
\end{minipage}\par

图 M2-073
的两条平行势线具有相同斜率而不同高度，对应同一电场和不同零点。无源方程迫使整段斜率恒定，所以两端给不同斜率时不是常数未定，而是数据不相容。

对 \(0<x<L\) 的一维无源问题，若给定

\[
\varphi'(0)=a,
\qquad
\varphi'(L)=a,
\]

则 \(\varphi=ax+C\)，常数 \(C\) 无法由斜率确定。若却给出
\(\varphi'(0)=a\)、\(\varphi'(L)=b\ne a\)，则与 \(\varphi''=0\)
不相容，不存在解。

\subsection{五、给定孤立导体总电荷时的唯一性}\label{m2:ux4e94ux7ed9ux5b9aux5b64ux7acbux5bfcux4f53ux603bux7535ux8377ux65f6ux7684ux552fux4e00ux6027}

导体边界上常常不知道电势数值，只知道每一块孤立导体是等势体且总电荷给定。这不是普通的逐点第二类条件，但仍可证明电场唯一。

设场区域中有若干互不接触的导体。两候选解之差为 \(u\)。在第 \(i\)
个导体表面 \(S_i\) 上，因两候选解各自等势，故

\[
u=C_i.
\]

取 \(\mathbf n\) 为场区域的外法线。在内边界 \(S_i\) 上，\(\mathbf n\)
指向导体内部。两候选解在该导体上的总电荷相同，因此

\[
\Delta Q_i
=\varepsilon_0\oint_{S_i}\frac{\partial u}{\partial n}\,dS
=0.
\]

第一格林恒等式的导体边界贡献为

\[
\oint_{S_i}u\frac{\partial u}{\partial n}\,dS
=C_i\oint_{S_i}\frac{\partial u}{\partial n}\,dS
=\frac{C_i\Delta Q_i}{\varepsilon_0}=0.
\]

若外边界电势固定，或无穷远项按规定消失，则

\[
\int_V|\nabla u|^2\,dV=0.
\]

于是

\[
\boxed{\mathbf E_1=\mathbf E_2.}
\]

这就是导体问题常用的第二唯一性结论：给定体电荷、导体几何和每个孤立导体的总电荷，并给出外部参考条件后，电场唯一；若没有固定电势零点，电势仍可整体加常数。

\par\noindent\begin{minipage}{\linewidth}
\begin{center}
\includegraphics[width=\linewidth,height=.40\textheight,keepaspectratio]{\detokenize{assets/module-02/original/tikz-m2-topic15-v8-06.png}}
\par\smallskip{\small 图示：给定各孤立导体总电荷时的唯一性证明}
\end{center}
\end{minipage}\par

若几块导体由理想导线连接，应把它们视为同一个连通导体，只规定整个连接系统的总电荷，而不是分别固定每一部分的电荷。

\subsection{六、无界区域、奇点与界面}\label{m2:ux516dux65e0ux754cux533aux57dfux5947ux70b9ux4e0eux754cux9762}

\subsubsection{1. 无界区域}\label{m2:ux65e0ux754cux533aux57df}

无界区域不能直接把``无穷远边界项为零''当作显然。严格做法是先用半径为
\(R\) 的大球面截断区域，在有限区域中使用格林恒等式，再令
\(R\to\infty\)。候选解必须满足足够强的无穷远衰减，使

\[
\oint_{S_R}u\frac{\partial u}{\partial n}\,dS\to0.
\]

对于局域电荷分布，通常规定 \(\varphi\to0\) 或规定一致的远场渐近行为。

\subsubsection{2.
点电荷等奇点}\label{m2:ux70b9ux7535ux8377ux7b49ux5947ux70b9}

若物理区域内有点电荷，电势在电荷位置不光滑。应先挖去包围每个奇点的小球，在剩余光滑区域中积分，再令小球半径趋于零。

两个候选解必须包含\textbf{相同位置、相同电量}的点电荷。这样它们的库仑奇异项在差值
\(u\) 中抵消，\(u\)
在原奇点附近保持正则，小球边界项才可消失。若候选解删去、增加或移动了物理区域内的真实电荷，它们根本不是同一边值问题的两个解。

\par\noindent\begin{minipage}{\linewidth}
\begin{center}
\includegraphics[width=\linewidth,height=.40\textheight,keepaspectratio]{\detokenize{assets/module-02/original/tikz-m2-topic15-v8-07.png}}
\par\smallskip{\small 图示：无界区域与点电荷奇点应先截断和挖除再取极限}
\end{center}
\end{minipage}\par

\subsubsection{3.
最大值原理的直观理解}\label{m2:ux6700ux5927ux503cux539fux7406ux7684ux76f4ux89c2ux7406ux89e3}

在有界连通区域内，并且函数能够连续延拓到边界时，非恒定调和函数不能在内部取得最大值或最小值，极值只能出现在边界。若差值
\(u\) 满足拉普拉斯方程且边界上 \(u=0\)，那么内部既不能有 \(u>0\)
的最大值，也不能有 \(u<0\) 的最小值，因此只能有 \(u=0\)。

最大值原理为第一类边界唯一性提供了直观解释；能量积分法则更容易推广到混合边界和介质问题。

\subsubsection{4.
分段线性介质中的推广}\label{m2:ux5206ux6bb5ux7ebfux6027ux4ecbux8d28ux4e2dux7684ux63a8ux5e7f}

若介电常数为空间中分段取正值的 \(\varepsilon(\mathbf r)\)，方程写成

\[
\nabla\cdot(\varepsilon\nabla\varphi)=-\rho_f.
\]

两个解之差满足

\[
\nabla\cdot(\varepsilon\nabla u)=0.
\]

在各介质区域分别积分；若两候选解满足相同的电势连续条件和相同的界面自由电荷条件，内部界面的贡献成对抵消，得到

\[
\int_V\varepsilon|\nabla u|^2\,dV=0.
\]

只要 \(\varepsilon>0\)，仍有 \(\nabla u=0\)。

\par\noindent\begin{minipage}{\linewidth}
\begin{center}
\includegraphics[width=\linewidth,height=.40\textheight,keepaspectratio]{\detokenize{assets/module-02/original/tikz-m2-topic15-v8-08.png}}
\par\smallskip{\small 图示：最大值原理和正介电常数下的加权能量证明}
\end{center}
\end{minipage}\par

\subsection{七、唯一性定理为什么允许电像法}\label{m2:ux4e03ux552fux4e00ux6027ux5b9aux7406ux4e3aux4ec0ux4e48ux5141ux8bb8ux7535ux50cfux6cd5}

电像法不要求虚构区域中的电荷是真实存在的。它只要求在真正研究的\textbf{物理区域}内，构造出的电势同时满足：

\begin{enumerate}
\def\labelenumi{\arabic{enumi}.}
\tightlist
\item
  正确的真实电荷与奇点；
\item
  正确的泊松方程或拉普拉斯方程；
\item
  导体表面上的指定电势；
\item
  正确的无穷远条件。
\end{enumerate}

以 \(z=0\) 的无限大接地平面和位于 \((0,0,d)\) 的点电荷 \(+q\)
为例。只研究 \(z>0\)。在构造区域 \(z<0\) 的镜像位置 \((0,0,-d)\) 放置
\(-q\)，得到

\[
\varphi(x,y,z)
=\frac{q}{4\pi\varepsilon_0}
\left[
\frac{1}{\sqrt{x^2+y^2+(z-d)^2}}
-\frac{1}{\sqrt{x^2+y^2+(z+d)^2}}
\right],
\qquad z>0.
\]

在 \(z=0\) 上，两距离相等，故
\(\varphi=0\)；无穷远处电势趋零；物理区域内只保留了真实电荷
\(+q\)。因此唯一性定理保证该表达式就是 \(z>0\) 中的物理解。

\par\noindent\begin{minipage}{\linewidth}
\begin{center}
\includegraphics[width=\linewidth,height=.40\textheight,keepaspectratio]{\detokenize{assets/module-02/original/tikz-m2-topic15-v8-09.png}}
\par\smallskip{\small 图示：接地平面问题的物理区域、构造区域和镜像电荷}
\end{center}
\end{minipage}\par

镜像电荷只是一种计算工具。唯一性定理只保证两个问题在 \(z>0\)
物理区域中的电势相同，并不声称导体内部真的存在 \(-q\)。

\subsection{八、候选电势的检查顺序}\label{m2:ux516bux5019ux9009ux7535ux52bfux7684ux68c0ux67e5ux987aux5e8f}

\par\noindent\begin{minipage}{\linewidth}
\begin{center}
\includegraphics[width=\linewidth,height=.40\textheight,keepaspectratio]{\detokenize{assets/module-02/original/tikz-m2-topic15-v8-10.png}}
\par\smallskip{\small 图示：候选解必须同时通过区域、源、方程、边界和无穷远检验}
\end{center}
\end{minipage}\par

验证试探解、分离变量解或电像法构造时，依次检查：

\begin{enumerate}
\def\labelenumi{\arabic{enumi}.}
\tightlist
\item
  划定真正需要求解的物理区域。
\item
  核对区域内真实电荷的位置、总量和奇点强度。
\item
  在每个光滑子区域核对泊松方程或拉普拉斯方程。
\item
  逐段核对全部边界条件和介质界面条件。
\item
  核对无穷远行为、电势参考零点和周期性等整体条件。
\end{enumerate}

常见错误包括：

\begin{itemize}
\tightlist
\item
  只验证导体表面是等势面，却没有核对这个常数是否为题目指定的电势。
\item
  只验证拉普拉斯方程，却在物理区域内增加、删去或移动了真实电荷。
\item
  忘记纯第二类边界的总通量相容条件。
\item
  把``电势至多相差常数''误说成``电势严格相同''。
\item
  在无界区域中漏掉无穷远条件。
\item
  直接跨过点电荷或面电荷奇点使用要求光滑性的积分恒等式。
\item
  给定孤立导体总电荷时，误把未知的导体电势当成已知数。
\end{itemize}

\subsection{九、练习与解答}\label{m2:ux4e5dux7ec3ux4e60ux4e0eux89e3ux7b54-1}

\subsubsection{作业一：一维泊松方程的第一类边值问题}\label{m2:ux4f5cux4e1aux4e00ux4e00ux7ef4ux6ccaux677eux65b9ux7a0bux7684ux7b2cux4e00ux7c7bux8fb9ux503cux95eeux9898}

\textbf{题号：M2-074}

\textbf{层级：} 复赛　　\textbf{方法：} 边值问题

真空中的一维平板区域为 \(0<x<L\)，\(L>0\)，物理量与 \(y,z\)
无关。区内均匀体电荷密度为 \(\rho_0\)，两端保持零电势，因而
\(\varphi''=-\rho_0/\varepsilon_0\)、\(\varphi(0)=\varphi(L)=0\)。求区内
\(\varphi(x)\) 和
\(E_x(x)\)，核对边界数据，并通过两解之差说明解为何唯一。

\par\noindent\begin{minipage}{\linewidth}
\textbf{解答}

\begin{center}
\includegraphics[width=\linewidth,height=.40\textheight,keepaspectratio]{\detokenize{assets/module-02/examples/M2-074.pdf}}
\par\smallskip{\small 图 M2-074：一维泊松方程的第一类边值问题}
\end{center}
\end{minipage}\par

图 M2-074
的两端零势把电势曲线固定在同一基线上，均匀体源使其具有固定负曲率。积分所得二次式必须同时通过两端，场零点对应中央水平切线；差值函数的零边界再保证唯一性。

积分两次：

\[
\varphi(x)=-\frac{\rho_0}{2\varepsilon_0}x^2+Ax+B.
\]

由 \(\varphi(0)=0\) 得 \(B=0\)；由 \(\varphi(L)=0\) 得

\[
A=\frac{\rho_0L}{2\varepsilon_0}.
\]

所以

\[
\boxed{
\varphi(x)=\frac{\rho_0}{2\varepsilon_0}x(L-x).}
\]

电场为

\[
\boxed{
E_x(x)=-\frac{d\varphi}{dx}
=\frac{\rho_0}{\varepsilon_0}
\left(x-\frac L2\right).}
\]

若另有一解 \(\tilde\varphi\)，则差值 \(u=\varphi-\tilde\varphi\) 满足
\(u''=0\) 且
\(u(0)=u(L)=0\)。线性函数只有在两端都为零时才恒为零，因此解唯一。这是一维情形下唯一性定理的直接表现。

\subsubsection{作业二：均匀带电球内的纯第二类条件}\label{m2:ux4f5cux4e1aux4e8cux5747ux5300ux5e26ux7535ux7403ux5185ux7684ux7eafux7b2cux4e8cux7c7bux6761ux4ef6}

\textbf{题号：M2-075}

\textbf{层级：} 基础　　\textbf{方法：} 静电场方程

以原点为球心、半径为 \(R>0\) 的真空球形区域内，体电荷密度为常量
\(\rho_0\)，球内没有其他点奇点。仅在整个球面规定外法向导数
\(\partial\varphi/\partial n=g\)，其中 \(\mathbf n=\mathbf e_r\)，\(g\)
为常量。求这一纯诺伊曼问题的相容条件，并给出球心有限的球对称电势；讨论
\(\rho_0\ne0\) 时误设 \(g=0\) 是否可能，以及电势是否仍可加常数。

\par\noindent\begin{minipage}{\linewidth}
\textbf{解答}

\begin{center}
\includegraphics[width=\linewidth,height=.40\textheight,keepaspectratio]{\detokenize{assets/module-02/examples/M2-075.pdf}}
\par\smallskip{\small 图 M2-075：均匀带电球内的纯第二类条件}
\end{center}
\end{minipage}\par

图 M2-075
的球面外法向为径向，恒定法向导数乘球面积给出边界总量。它必须等于负的包围体电荷除以
\(\varepsilon_0\)，先决定允许的 \(g\)；满足后径向势仍可整体加常数。

区域总电荷为

\[
Q_V=\frac{4\pi R^3}{3}\rho_0.
\]

纯第二类相容条件为

\[
4\pi R^2g=-\frac{Q_V}{\varepsilon_0},
\]

因此

\[
\boxed{g=-\frac{\rho_0R}{3\varepsilon_0}.}
\]

球对称泊松方程为

\[
\frac1{r^2}\frac{d}{dr}
\left(r^2\frac{d\varphi}{dr}\right)
=-\frac{\rho_0}{\varepsilon_0}.
\]

积分得

\[
r^2\frac{d\varphi}{dr}
=-\frac{\rho_0}{3\varepsilon_0}r^3+A.
\]

为使 \(r=0\) 处电势有限，必须有 \(A=0\)。再次积分：

\[
\boxed{
\varphi(r)=-\frac{\rho_0r^2}{6\varepsilon_0}+C.}
\]

其中 \(C\) 任意，而电场

\[
\mathbf E(r)=\frac{\rho_0r}{3\varepsilon_0}\mathbf e_r
\]

唯一。若 \(\rho_0\ne0\) 却给定
\(g=0\)，边界总通量为零而区域总电荷不为零，违反相容条件，因此无解。

\subsubsection{作业三：乘积恒等式的手算核对}\label{m2:ux4f5cux4e1aux4e09ux4e58ux79efux6052ux7b49ux5f0fux7684ux624bux7b97ux6838ux5bf9}

\textbf{题号：M2-076}

\textbf{层级：} 基础　　\textbf{方法：} 静电场方程

在三维直角坐标系中取 \(u(x,y,z)=x^2+y^2-2z^2\)，所有导数作用于一般点
\((x,y,z)\)。一条路径先将 \(u\nabla u\)
展开为三个分量再求散度，另一条路径分别计算 \(|\nabla u|^2\) 与
\(u\nabla^2u\)。完成两种计算并验证乘积恒等式，明确调和条件与梯度非零并不矛盾。

\par\noindent\begin{minipage}{\linewidth}
\textbf{解答}

\begin{center}
\includegraphics[width=\linewidth,height=.40\textheight,keepaspectratio]{\detokenize{assets/module-02/examples/M2-076.pdf}}
\par\smallskip{\small 图 M2-076：乘积恒等式的手算核对}
\end{center}
\end{minipage}\par

图 M2-076
展示同一调和二次函数的分量计算。先把三个二阶导数加成零，再完整展开
\(\nabla u\) 的模方；直接求 \(\nabla\cdot(u\nabla u)\)
时保留乘积的两次作用，才能得到同样的三个平方项。

先求

\[
\nabla u=(2x,2y,-4z),
\qquad
\nabla^2u=2+2-4=0.
\]

所以右侧为

\[
|\nabla u|^2+u\nabla^2u
=4x^2+4y^2+16z^2.
\]

左侧为

\[
\begin{aligned}
\nabla\cdot(u\nabla u)
&=\frac{\partial(2xu)}{\partial x}
+\frac{\partial(2yu)}{\partial y}
+\frac{\partial(-4zu)}{\partial z}\\
&=(2u+4x^2)+(2u+4y^2)+(-4u+16z^2)\\
&=4x^2+4y^2+16z^2.
\end{aligned}
\]

两侧相等。这里 \(u\nabla^2u=0\)
是因为所选函数恰好调和，不是乘积恒等式中永远可以删去的一项。

\subsubsection{作业四：无限大接地平面前的点电荷}\label{m2:ux4f5cux4e1aux56dbux65e0ux9650ux5927ux63a5ux5730ux5e73ux9762ux524dux7684ux70b9ux7535ux8377}

\textbf{题号：M2-077}

\textbf{层级：} 复赛　　\textbf{方法：} 电像法

无限大接地导体平面位于 \(z=0\)，真空半空间 \(z>0\) 中只有一个点电荷
\(+q\)，位置为 \((0,0,d)\)，其中
\(q>0,d>0\)。要求电势在无穷远衰减。用电像法求真实半空间的电势、导体表面感应面密度、感应总电荷及点电荷受到的力。再求与诱导电荷有关的有限相互作用能，以点电荷远离平面时为零，不计点电荷本身发散的自能。

\par\noindent\begin{minipage}{\linewidth}
\textbf{解答}

\begin{center}
\includegraphics[width=\linewidth,height=.40\textheight,keepaspectratio]{\detokenize{assets/module-02/examples/M2-077.pdf}}
\par\smallskip{\small 图 M2-077：无限大接地平面前的点电荷}
\end{center}
\end{minipage}\par

图 M2-077
中真实电荷与反射像分居接地平面两侧，平面各点到两源等距。确定势后从真空侧取法向导数求面密度，再对整个平面积分；作用力由像场给出，而感应能仍保留组装过程中的
\(1/2\)。

在构造区域 \(z<0\) 的 \((0,0,-d)\) 放置镜像电荷
\(-q\)。记柱坐标径向距离为 \(\varrho=\sqrt{x^2+y^2}\)，则

\[
\boxed{
\varphi(\varrho,z)
=\frac{q}{4\pi\varepsilon_0}
\left[
\frac1{\sqrt{\varrho^2+(z-d)^2}}
-\frac1{\sqrt{\varrho^2+(z+d)^2}}
\right],\quad z>0.}
\]

它在 \(z=0\) 上为零，在无穷远处趋零，并在物理区域中保留了正确的 \(+q\)
奇点，故由唯一性定理可知它是唯一物理解。

导体指向真空的外法线为 \(+\mathbf e_z\)。表面电荷密度为

\[
\sigma(\varrho)
=-\varepsilon_0
\left.\frac{\partial\varphi}{\partial z}\right|_{z=0^+}
=\boxed{
-\frac{qd}{2\pi(\varrho^2+d^2)^{3/2}}}.
\]

感应总电荷为

\[
\begin{aligned}
Q_{\text{感应}}
&=\int_0^{\infty}\sigma(\varrho)\,2\pi\varrho\,d\varrho\\
&=-qd\int_0^{\infty}
\frac{\varrho\,d\varrho}{(\varrho^2+d^2)^{3/2}}
=-q.
\end{aligned}
\]

真实电荷所受力等于镜像电荷在真实电荷位置产生的库仑力。两者距离为
\(2d\)，所以

\[
\boxed{
\mathbf F
=-\frac{q^2}{16\pi\varepsilon_0d^2}\mathbf e_z.}
\]

系统静电能不能直接取真实电荷与镜像电荷相互作用能，而应乘 \(1/2\)：

\[
U=\frac12q\varphi_{\mathrm{ind}}(0,0,d)
=\frac12q
\left(-\frac{q}{4\pi\varepsilon_0(2d)}\right)
=\boxed{-\frac{q^2}{16\pi\varepsilon_0d}}.
\]

检验

\[
-\frac{dU}{dd}
=-\frac{q^2}{16\pi\varepsilon_0d^2}
\]

与力的 \(z\) 分量一致。

\subsubsection{作业五：一维混合边界问题}\label{m2:ux4f5cux4e1aux4e94ux4e00ux7ef4ux6df7ux5408ux8fb9ux754cux95eeux9898}

\textbf{题号：M2-078}

\textbf{层级：} 基础　　\textbf{方法：} 静电场方程

在一维无源真空区间 \(0<x<L\) 内，电势与其余坐标无关。给定混合边界条件
\(\varphi(0)=V_0\)、\(E_x(L)=E_0\)，其中 \(L>0\)，\(V_0,E_0\)
为有符号常量，且
\(E_x=-d\varphi/dx\)。求区间内电势；再设存在两组解，利用它们的差满足的方程和边界条件说明唯一性。

\par\noindent\begin{minipage}{\linewidth}
\textbf{解答}

\begin{center}
\includegraphics[width=\linewidth,height=.40\textheight,keepaspectratio]{\detokenize{assets/module-02/examples/M2-078.pdf}}
\par\smallskip{\small 图 M2-078：一维混合边界问题}
\end{center}
\end{minipage}\par

图 M2-078 的一端给电势、另一端给电场。直线势的斜率先由右端 \(E_0\)
的负号固定，高度再由左端 \(V_0\)
固定；对两候选解作差时，零值与零导数分别消去两个常数。

由 \(\varphi''=0\)，设

\[
\varphi(x)=Ax+B.
\]

第一类边界给出 \(B=V_0\)。又因 \(E_x=-\varphi'\)，第二类边界给出

\[
A=-E_0.
\]

所以

\[
\boxed{\varphi(x)=V_0-E_0x.}
\]

若有两个解，其差 \(u\) 满足

\[
u''=0,
\qquad
u(0)=0,
\qquad
u'(L)=0.
\]

写成 \(u=ax+b\)，由 \(u(0)=0\) 得 \(b=0\)，由 \(u'(L)=0\) 得 \(a=0\)，故
\(u\equiv0\)。混合边界中的第一类部分固定了加性常数，因此电势唯一。

\section{专题十六：静电场的局域展开与平衡稳定性（拔高选讲）}\label{m2:ux4e13ux9898ux5341ux516dux9759ux7535ux573aux7684ux5c40ux57dfux5c55ux5f00ux4e0eux5e73ux8861ux7a33ux5b9aux6027ux62d4ux9ad8ux9009ux8bb2}

静电平衡点附近的运动不需要每次都从完整库仑力重新计算。把电势在平衡点附近展开，最低阶非零项就给出了局域势能形状、恢复力和小振动频率。真正需要谨慎区分的是：自由粒子与受约束粒子的平衡条件不同，电势与势能的极值还要结合电荷符号，而二阶项为零时不能停止计算。

\par\noindent\begin{minipage}{\linewidth}
\begin{center}
\includegraphics[width=\linewidth,height=.40\textheight,keepaspectratio]{\detokenize{assets/module-02/original/tikz-m2-topic16-v8-01.png}}
\par\smallskip{\small 图示：平衡点附近的电势展开、Hessian 与主轴}
\end{center}
\end{minipage}\par

\subsection{一、自由平衡与受约束平衡}\label{m2:ux4e00ux81eaux7531ux5e73ux8861ux4e0eux53d7ux7ea6ux675fux5e73ux8861}

对电荷量为 \(q\ne0\) 的自由点粒子，静电力为

\[
\mathbf F=q\mathbf E=-q\nabla\varphi.
\]

自由平衡要求三个方向的合力都为零：

\[
\boxed{\mathbf E(O)=\mathbf 0}
\qquad\Longleftrightarrow\qquad
\boxed{\nabla\varphi(O)=\mathbf 0.}
\]

若粒子被限制在曲线或曲面上，约束力可以抵消法向电场力。此时只要求允许位移方向上的力为零，空间总电场不一定为零。

例如，粒子沿曲线 \(\mathbf r(s)\) 运动，\(s\) 为弧长，单位切向量为

\[
\mathbf t=\frac{d\mathbf r}{ds}.
\]

势能 \(U=q\varphi\) 沿曲线的一阶导数为

\[
\frac{dU}{ds}
=q\nabla\varphi\cdot\mathbf t
=-q\mathbf E\cdot\mathbf t.
\]

所以受约束平衡只要求

\[
\boxed{\mathbf E\cdot\mathbf t=0.}
\]

把``受约束平衡''直接写成 \(\mathbf E=0\)，会漏掉一大类竞赛模型。

\subsection{二、电势的局域展开与主轴}\label{m2:ux4e8cux7535ux52bfux7684ux5c40ux57dfux5c55ux5f00ux4e0eux4e3bux8f74}

以平衡点 \(O\) 为原点，记微小位移

\[
\boldsymbol{\xi}=(x,y,z)^{\mathsf T}.
\]

若电势在原点附近至少三次可微，泰勒展开为

\[
\varphi(\boldsymbol{\xi})
=\varphi_0
+(\nabla\varphi)_0\cdot\boldsymbol{\xi}
+\frac12\boldsymbol{\xi}^{\mathsf T}H\boldsymbol{\xi}
+O(|\boldsymbol{\xi}|^3),
\]

其中二阶导数矩阵

\[
H_{ij}
=\left.\frac{\partial^2\varphi}
{\partial x_i\partial x_j}\right|_O
\]

称为 Hessian 矩阵。混合偏导可交换时

\[
H^{\mathsf T}=H.
\]

它是实对称矩阵，因此一定存在正交矩阵 \(R\)，使

\[
R^{\mathsf T}HR
=\operatorname{diag}(\lambda_1,\lambda_2,\lambda_3).
\]

\(\lambda_i\) 是三个主曲率。选取对应的正交主轴 \(\xi_1,\xi_2,\xi_3\)
后，二阶展开自动没有交叉项：

\[
\varphi
=\varphi_0
+\frac12
\left(\lambda_1\xi_1^2
+\lambda_2\xi_2^2
+\lambda_3\xi_3^2\right)
+O(|\boldsymbol{\xi}|^3).
\]

主轴的存在来自实对称矩阵的正交对角化，并不要求系统本身具有轴对称或镜像对称。几何对称性只能帮助我们更快地识别这些主轴。

\subsubsection{无电荷区的迹约束}\label{m2:ux65e0ux7535ux8377ux533aux7684ux8ff9ux7ea6ux675f}

若平衡点附近没有电荷，电势满足

\[
\nabla^2\varphi=0.
\]

在平衡点处

\[
\boxed{
\operatorname{tr}H
=\lambda_1+\lambda_2+\lambda_3=0.}
\]

因此三个非零主曲率不可能全为正，也不可能全为负。二阶项非退化时，电势必在某些方向向上弯、另一些方向向下弯。

\par\noindent\begin{minipage}{\linewidth}
\begin{center}
\includegraphics[width=\linewidth,height=.40\textheight,keepaspectratio]{\detokenize{assets/module-02/original/tikz-m2-topic16-v8-02.png}}
\par\smallskip{\small 图示：用三个主轴截面表示无源静电势的鞍形约束}
\end{center}
\end{minipage}\par

\textbf{例 M2-079：迹为零的二次势呈鞍形}

某无源真空区域的局部电势为 \(\varphi=\varphi_0+A(x^2+y^2-2z^2)/2\)，其中
\(A>0\)，\(\varphi_0\) 为常量。以原点为观察中心，比较沿 \(x,y,z\)
三个轴的电势变化，检验拉普拉斯方程，并判断原点是局部最大、局部最小还是鞍形驻点。

\par\noindent\begin{minipage}{\linewidth}
\textbf{解答}

\begin{center}
\includegraphics[width=\linewidth,height=.40\textheight,keepaspectratio]{\detokenize{assets/module-02/examples/M2-079.pdf}}
\par\smallskip{\small 图 M2-079：迹为零的二次势呈鞍形}
\end{center}
\end{minipage}\par

图 M2-079 的横向势剖面向上、轴向剖面向下，Hessian
三个本征值和为零。原点在不同方向呈不同极值性质，故不是三维势的局部极值；测试电荷的稳定方向还会随其符号翻转。

若

\[
\varphi=\varphi_0+\frac A2(x^2+y^2-2z^2),
\qquad A>0,
\]

则

\[
H=\operatorname{diag}(A,A,-2A),
\qquad
\operatorname{tr}H=0.
\]

电势沿 \(x,y\) 方向取局部极小，沿 \(z\)
方向取局部极大，所以原点是鞍形平衡点，不是三维极值点。

\subsection{三、势能、线性运动与稳定性判据}\label{m2:ux4e09ux52bfux80fdux7ebfux6027ux8fd0ux52a8ux4e0eux7a33ux5b9aux6027ux5224ux636e}

测试电荷的势能为

\[
U=q\varphi.
\]

在自由平衡点附近，主轴方向上的势能增量为

\[
\Delta U
=\frac q2\sum_i\lambda_i\xi_i^2
+O(|\boldsymbol{\xi}|^3).
\]

电场和运动方程的一阶近似分别为

\[
\mathbf E=-H\boldsymbol{\xi}+O(|\boldsymbol{\xi}|^2),
\]

\[
m\ddot\xi_i=-q\lambda_i\xi_i.
\]

于是每个主方向可分别判断：

\begin{itemize}
\item
  \(q\lambda_i>0\)：存在恢复力，线性运动为简谐振动，

  \[
  \omega_i=\sqrt{\frac{q\lambda_i}{m}}.
  \]
\item
  \(q\lambda_i<0\)：位移受到同向放大，解含指数增长项，线性不稳定。
\item
  \(q\lambda_i=0\)：二阶恢复力消失，必须继续检查更高阶项；不能把
  \(\omega_i=0\) 当作稳定结论。
\end{itemize}

\par\noindent\begin{minipage}{\linewidth}
\begin{center}
\includegraphics[width=\linewidth,height=.40\textheight,keepaspectratio]{\detokenize{assets/module-02/original/tikz-m2-topic16-v8-03.png}}
\par\smallskip{\small 图示：势能主曲率的正负与恢复、失稳和高阶判定}
\end{center}
\end{minipage}\par

自由粒子要在三维线性稳定，必须要求 \(qH\) 正定，即三个 \(q\lambda_i\)
全为正。但无源区满足

\[
q(\lambda_1+\lambda_2+\lambda_3)=0,
\]

所以三个 \(q\lambda_i\)
不可能同时为正。改变测试电荷的正负只会同时翻转三个方向，不会消除符号冲突。

\subsection{四、恩绍定理的适用范围}\label{m2:ux56dbux6069ux7ecdux5b9aux7406ux7684ux9002ux7528ux8303ux56f4}

\subsubsection{1. 定理内容}\label{m2:ux5b9aux7406ux5185ux5bb9}

在由固定静止电荷产生的纯静电场中，若测试点附近没有电荷源，一个自由点电荷不可能获得三维完全稳定的孤立平衡。

二阶非退化时，这由 \(\operatorname{tr}H=0\)
立即得到。更一般地，调和函数的强最大值原理说明：非恒定调和函数不能在区域内部取得严格局部最大值或最小值。正电荷稳定需要
\(\varphi\) 取严格局部最小，负电荷稳定需要 \(\varphi\)
取严格局部最大，两者都不可能。

\subsubsection{2.
二阶项全零并不构成反例}\label{m2:ux4e8cux9636ux9879ux5168ux96f6ux5e76ux4e0dux6784ux6210ux53cdux4f8b}

考虑

\[
\varphi=B(x^3-3xy^2)=Br^3\cos3\theta.
\]

它满足

\[
\nabla^2\varphi=0.
\]

原点处

\[
\nabla\varphi(O)=0,
\qquad
H(O)=0.
\]

二阶近似完全平坦，但三阶项沿六个扇区交替变号，因此原点既不是最大值也不是最小值。

\par\noindent\begin{minipage}{\linewidth}
\begin{center}
\includegraphics[width=\linewidth,height=.40\textheight,keepaspectratio]{\detokenize{assets/module-02/original/tikz-m2-topic16-v8-04.png}}
\par\smallskip{\small 图示：退化调和函数的二阶项为零但高阶项沿不同方向变号}
\end{center}
\end{minipage}\par

如果一个调和函数在内部真的取得局部最大值或最小值，强最大值原理要求它在相应连通邻域内为常数。此时没有恢复力，只能得到中性平坦区域，而不是孤立束缚。

\subsubsection{3.
使用恩绍定理前要检查假设}\label{m2:ux4f7fux7528ux6069ux7ecdux5b9aux7406ux524dux8981ux68c0ux67e5ux5047ux8bbe}

\par\noindent\begin{minipage}{\linewidth}
\begin{center}
\includegraphics[width=\linewidth,height=.40\textheight,keepaspectratio]{\detokenize{assets/module-02/original/tikz-m2-topic16-v8-05.png}}
\par\smallskip{\small 图示：恩绍定理的条件和不适用情形}
\end{center}
\end{minipage}\par

恩绍定理讨论的是：

\begin{itemize}
\tightlist
\item
  自由点粒子；
\item
  三维完全稳定；
\item
  静态电场；
\item
  平衡点邻域无电荷；
\item
  没有额外机械约束、主动反馈或时间变化场。
\end{itemize}

只要改变其中某项，就不能直接照搬结论。例如：

\begin{itemize}
\tightlist
\item
  细杆、导轨或曲面约束只要求允许方向稳定；
\item
  交变电场可产生动态平均束缚；
\item
  主动反馈不再是保守静电边值问题；
\item
  平衡点附近存在连续电荷时，\(\nabla^2\varphi\) 不再为零。
\end{itemize}

\textbf{例 M2-080：连续电荷区允许三个方向同时恢复}

某真空区域内给定电势
\(\varphi(\mathbf r)=\alpha r^2\)，\(r=|\mathbf r|\)，\(\alpha\ne0\)
为常量。质量为 \(m>0\)、有符号电荷量为 \(q\ne0\)
的测试粒子可在该区域自由运动，不扰动原电荷分布，也不受其他力。求产生该势的体电荷密度及粒子受力，判断原点何时稳定，并在稳定情形求小振动角频率。

\par\noindent\begin{minipage}{\linewidth}
\textbf{解答}

\begin{center}
\includegraphics[width=\linewidth,height=.40\textheight,keepaspectratio]{\detokenize{assets/module-02/examples/M2-080.pdf}}
\par\smallskip{\small 图 M2-080：连续电荷区允许三个方向同时恢复}
\end{center}
\end{minipage}\par

图 M2-080 的箭头表示 \(q\alpha>0\)
时三维都指向原点的力。负梯度给线性恢复项，拉普拉斯则揭示这里需要非零连续体电荷；因此这种各向同性稳定不违反无源区域的恩绍结论。

若局部电势为

\[
\varphi=\alpha r^2,
\]

则

\[
\nabla^2\varphi=6\alpha,
\qquad
\rho=-6\varepsilon_0\alpha.
\]

测试电荷受到

\[
\mathbf F=-2q\alpha\mathbf r.
\]

当 \(q\alpha>0\) 时三个方向都有恢复力，角频率为

\[
\omega=\sqrt{\frac{2q\alpha}{m}}.
\]

这里平衡点邻域有连续电荷，已经不属于恩绍定理的无源假设。

\subsection{五、约束运动的二阶判据}\label{m2:ux4e94ux7ea6ux675fux8fd0ux52a8ux7684ux4e8cux9636ux5224ux636e}

对曲线 \(\mathbf r(s)\) 上的粒子，势能二阶导数为

\[
\frac{d^2U}{ds^2}
=q\left[
\mathbf t^{\mathsf T}H\mathbf t
+\nabla\varphi\cdot\frac{d^2\mathbf r}{ds^2}
\right].
\]

这个式子包含两个来源：

\begin{enumerate}
\def\labelenumi{\arabic{enumi}.}
\tightlist
\item
  电势本身在切向方向的曲率 \(\mathbf t^{\mathsf T}H\mathbf t\)；
\item
  曲线转弯时，非零法向场与曲率共同产生的项。
\end{enumerate}

若平衡点本身满足 \(\mathbf E(O)=0\)，则
\(\nabla\varphi(O)=0\)，第二项消失。对通过该点的直杆也有
\(d^2\mathbf r/ds^2=0\)。此时定义

\[
\boxed{
\lambda_{\mathrm{eff}}
=\mathbf t^{\mathsf T}H\mathbf t,}
\]

并有

\[
\Delta U
=\frac12q\lambda_{\mathrm{eff}}h^2+O(h^3),
\]

\[
m\ddot h=-q\lambda_{\mathrm{eff}}h.
\]

所以直杆上的稳定条件和小振动角频率为

\[
\boxed{q\lambda_{\mathrm{eff}}>0,}
\]

\[
\boxed{
\omega=\sqrt{\frac{q\lambda_{\mathrm{eff}}}{m}}.}
\]

\par\noindent\begin{minipage}{\linewidth}
\begin{center}
\includegraphics[width=\linewidth,height=.40\textheight,keepaspectratio]{\detokenize{assets/module-02/original/tikz-m2-topic16-v8-06-margin12-v11.png}}
\par\smallskip{\small 图示：把三维 Hessian 投影到直杆切向方向}
\end{center}
\end{minipage}\par

\textbf{例 M2-081：约束只要求切向合力为零}

半径为 \(R>0\)、圆心在原点的光滑绝缘圆环固定在 \(xy\) 平面。质量为
\(m>0\)、电荷为 \(q>0\)
的小珠只能沿环运动，忽略重力、摩擦和珠自身对外场的影响。给定匀强外场
\(\mathbf E=E_0\mathbf e_x\)，\(E_0>0\)；取环上位置
\(x=R\cos\theta\)。求稳定平衡位置及小振动角频率，并说明平衡时电场力是否必须为零。

\par\noindent\begin{minipage}{\linewidth}
\textbf{解答}

\begin{center}
\includegraphics[width=\linewidth,height=.40\textheight,keepaspectratio]{\detokenize{assets/module-02/examples/M2-081.pdf}}
\par\smallskip{\small 图 M2-081：约束只要求切向合力为零}
\end{center}
\end{minipage}\par

图 M2-081
中绿色点位于沿电场方向的圆环端点。此处切向电力为零，但电场本身不为零；法向电力由约束承担。把弧长写成
\(R\theta\)，势能二次项与转动动能比较后读出频率。

设均匀场 \(\mathbf E=E_0\mathbf e_x\)，带电小珠只能在半径 \(R\)
的光滑圆环上运动，令 \(x=R\cos\theta\)。电势和势能为

\[
\varphi=-E_0R\cos\theta,
\qquad
U=-qE_0R\cos\theta.
\]

当 \(qE_0>0\) 时，\(\theta=0\)
为稳定平衡。此处空间电场并不为零，只是电场力沿圆环切向的分量为零，法向分量由约束力平衡。对小角度展开：

\[
U\approx-qE_0R+\frac12qE_0R\theta^2.
\]

由动能 \(T=mR^2\dot\theta^2/2\) 得

\[
\omega=\sqrt{\frac{qE_0}{mR}}.
\]

\subsection{六、两等量正电荷与倾斜细杆}\label{m2:ux516dux4e24ux7b49ux91cfux6b63ux7535ux8377ux4e0eux503eux659cux7ec6ux6746}

两个固定点电荷 \(+Q\) 位于 \(z=\pm l\)，原点处的电场因对称性为零。质量为
\(m\)、电荷量为 \(q\) 的小珠只能沿通过原点、与 \(z\) 轴夹角为 \(\theta\)
的直杆运动。取 \(0\le\theta\le\pi/2\)。

\par\noindent\begin{minipage}{\linewidth}
\begin{center}
\includegraphics[width=\linewidth,height=.40\textheight,keepaspectratio]{\detokenize{assets/module-02/original/tikz-m2-topic16-v8-07.png}}
\par\smallskip{\small 图示：两等量正电荷、倾斜细杆与位移分量}
\end{center}
\end{minipage}\par

\subsubsection{1.
原点附近的空间电势}\label{m2:ux539fux70b9ux9644ux8fd1ux7684ux7a7aux95f4ux7535ux52bf}

记

\[
k=\frac1{4\pi\varepsilon_0}.
\]

在 \(x\) 轴上

\[
\varphi(x,0,0)
=\frac{2kQ}{\sqrt{l^2+x^2}}
\approx\frac{2kQ}{l}-\frac{kQ}{l^3}x^2.
\]

与

\[
\varphi=\varphi_0+\frac12\lambda_x x^2+\cdots
\]

比较，得到

\[
\lambda_x=\lambda_y=-\frac{2kQ}{l^3}.
\]

由迹为零

\[
\lambda_z=-2\lambda_x=\frac{4kQ}{l^3}.
\]

令

\[
A=\frac{2kQ}{l^3}
=\frac{Q}{2\pi\varepsilon_0l^3}>0,
\]

则

\[
H=\operatorname{diag}(-A,-A,2A).
\]

这说明原点附近沿 \(z\) 轴为电势极小方向，沿任意横向为电势极大方向。对
\(q>0\)，轴向有恢复力，横向则失稳。

\subsubsection{2.
投影到细杆方向}\label{m2:ux6295ux5f71ux5230ux7ec6ux6746ux65b9ux5411}

细杆单位切向量为

\[
\mathbf t=(\sin\theta,0,\cos\theta)^{\mathsf T}.
\]

有效曲率为

\[
\begin{aligned}
\lambda_{\mathrm{eff}}
&=\mathbf t^{\mathsf T}H\mathbf t\\
&=A(2\cos^2\theta-\sin^2\theta)\\
&=A(3\cos^2\theta-1).
\end{aligned}
\]

当 \(q,Q>0\) 时，严格的二阶稳定条件为

\[
\lambda_{\mathrm{eff}}>0,
\]

即

\[
\boxed{\tan^2\theta<2.}
\]

在约定的角度范围内

\[
\boxed{
0\le\theta<\theta_c,
\qquad
\theta_c=\arctan\sqrt2
=\arccos\frac1{\sqrt3}.}
\]

对应的小振动角频率为

\[
\boxed{
\omega
=\sqrt{
\frac{qQ(2\cos^2\theta-\sin^2\theta)}
{2\pi\varepsilon_0ml^3}}.}
\]

\subsubsection{3.
临界角必须检查四阶项}\label{m2:ux4e34ux754cux89d2ux5fc5ux987bux68c0ux67e5ux56dbux9636ux9879}

杆上距原点的有向坐标为 \(h\)。精确电势为

\[
\varphi(h)
=kQ\left[
\frac1{\sqrt{l^2+h^2-2lh\cos\theta}}
+\frac1{\sqrt{l^2+h^2+2lh\cos\theta}}
\right].
\]

当 \(|h|<l\) 时，用勒让德多项式展开：

\[
\varphi(h)
=\frac{2kQ}{l}
\left[
1+P_2(\cos\theta)\frac{h^2}{l^2}
+P_4(\cos\theta)\frac{h^4}{l^4}
+O\!\left(\frac{h^6}{l^6}\right)
\right].
\]

其中

\[
P_2(c)=\frac{3c^2-1}{2},
\qquad
P_4(c)=\frac{35c^4-30c^2+3}{8}.
\]

在临界角 \(\cos^2\theta_c=1/3\) 处，

\[
P_2(\cos\theta_c)=0,
\]

但

\[
P_4(\cos\theta_c)=-\frac7{18}.
\]

因此对 \(qQ>0\)，临界角处

\[
\Delta U
=-\frac{7kqQ}{9l^5}h^4+O(h^6),
\]

原点是很平的局部极大而不是稳定点。故稳定范围必须取严格不等式
\(\tan^2\theta<2\)。

\par\noindent\begin{minipage}{\linewidth}
\begin{center}
\includegraphics[width=\linewidth,height=.40\textheight,keepaspectratio]{\detokenize{assets/module-02/original/tikz-m2-topic16-v8-08.png}}
\par\smallskip{\small 图示：斜杆稳定窗口、临界四阶失稳和二阶失稳}
\end{center}
\end{minipage}\par

\subsubsection{4.
电场投影的交叉验证}\label{m2:ux7535ux573aux6295ux5f71ux7684ux4ea4ux53c9ux9a8cux8bc1}

由

\[
\mathbf E=-H\boldsymbol{\xi}
\]

和 \(\boldsymbol{\xi}=h\mathbf t\)，杆方向电场为

\[
E_{\tau}
=\mathbf t\cdot\mathbf E
=-\lambda_{\mathrm{eff}}h.
\]

切向力

\[
F_{\tau}=qE_{\tau}
=-q\lambda_{\mathrm{eff}}h.
\]

它与势能法给出的恢复力完全一致。

\par\noindent\begin{minipage}{\linewidth}
\begin{center}
\includegraphics[width=\linewidth,height=.40\textheight,keepaspectratio]{\detokenize{assets/module-02/original/tikz-m2-topic16-v8-09.png}}
\par\smallskip{\small 图示：势能曲率法与电场投影法的符号对应}
\end{center}
\end{minipage}\par

\subsection{七、导体静电题的求解顺序}\label{m2:ux4e03ux5bfcux4f53ux9759ux7535ux9898ux7684ux6c42ux89e3ux987aux5e8f}

\par\noindent\begin{minipage}{\linewidth}
\begin{center}
\includegraphics[width=\linewidth,height=.40\textheight,keepaspectratio]{\detokenize{assets/module-02/original/tikz-m2-topic16-v8-10.png}}
\par\smallskip{\small 图示：平衡点稳定性的判断}
\end{center}
\end{minipage}\par

处理局域稳定性问题时，建议按以下顺序：

\begin{enumerate}
\def\labelenumi{\arabic{enumi}.}
\tightlist
\item
  判断粒子是自由运动，还是只能沿曲线、曲面运动。
\item
  写正确的平衡条件：自由粒子要求
  \(\mathbf E=0\)，受约束粒子只要求切向力为零。
\item
  展开电势或电场，确定最低阶非零项。
\item
  对二阶项写 Hessian；利用实对称性选取主轴。
\item
  在无源区先使用 \(\operatorname{tr}H=0\) 或最大值原理检查恩绍限制。
\item
  对受约束运动，把 Hessian 投影到允许方向；若平衡处 \(\mathbf E\ne0\)
  且约束弯曲，还要保留曲率项。
\item
  只有二次势能严格为正时，才可直接写简谐振动频率；二阶项为零必须继续高阶。
\end{enumerate}

常见错误包括：

\begin{itemize}
\tightlist
\item
  看见 \(\mathbf E=0\) 就断言稳定，没有检查位移后的力。
\item
  直接判断 \(\varphi\) 的极值，忘记势能是 \(U=q\varphi\)。
\item
  认为交叉项只能靠几何对称性消去；实际上应使用 Hessian 的正交主轴。
\item
  在含电荷区域仍使用 \(\operatorname{tr}H=0\)。
\item
  把 \(\omega=0\) 当成稳定，而没有检查三阶、四阶项。
\item
  对受约束粒子仍强行要求 \(\mathbf E=0\)。
\item
  混淆细杆与 \(z\) 轴、与横轴之间的夹角。
\end{itemize}

\subsection{八、练习与解答}\label{m2:ux516bux7ec3ux4e60ux4e0eux89e3ux7b54-3}

\subsubsection{作业一：二次鞍形势中的自由粒子}\label{m2:ux4f5cux4e1aux4e00ux4e8cux6b21ux978dux5f62ux52bfux4e2dux7684ux81eaux7531ux7c92ux5b50}

\textbf{题号：M2-082}

\textbf{层级：} 决赛　　\textbf{方法：} 稳定性分析、能量法

在一块无源真空区域内，给定局部电势
\(\varphi=\varphi_0+A(x^2+y^2-2z^2)/2\)，\(A>0\)。质量为 \(m>0\)
的非相对论测试粒子不受约束或其他力，电荷量为
\(q\ne0\)，且不扰动原场。分别讨论 \(q>0\)、\(q<0\)
时原点附近三个坐标方向的线性稳定性，求稳定方向的小振动角频率，并判断能否得到三维稳定平衡。

\par\noindent\begin{minipage}{\linewidth}
\textbf{解答}

\begin{center}
\includegraphics[width=\linewidth,height=.40\textheight,keepaspectratio]{\detokenize{assets/module-02/examples/M2-082.pdf}}
\par\smallskip{\small 图 M2-082：二次鞍形势中的自由粒子}
\end{center}
\end{minipage}\par

图 M2-082 的二次势在 \(x,y,z\) 三方向的曲率为
\(A,A,-2A\)。乘测试电荷后才是势能曲率，所以正负电荷的稳定方向互换；每个正势能曲率除以质量给相应的频率平方。

Hessian 为

\[
H=\operatorname{diag}(A,A,-2A).
\]

三个主方向运动满足

\[
m\ddot x=-qAx,
\qquad
m\ddot y=-qAy,
\qquad
m\ddot z=2qAz.
\]

若 \(q>0\)，\(x,y\) 方向有恢复力，

\[
\omega_x=\omega_y=\sqrt{\frac{qA}{m}},
\]

而 \(z\) 方向指数失稳。若 \(q<0\)，\(x,y\) 方向失稳，\(z\)
方向有恢复力，

\[
\omega_z=\sqrt{\frac{2|q|A}{m}}.
\]

无论 \(q\)
取何符号，都至少有一个不稳定方向，所以原点不是自由粒子的三维稳定平衡。

\subsubsection{作业二：两等量正电荷与倾斜细杆}\label{m2:ux4f5cux4e1aux4e8cux4e24ux7b49ux91cfux6b63ux7535ux8377ux4e0eux503eux659cux7ec6ux6746}

\textbf{题号：M2-083}

\textbf{层级：} 决赛　　\textbf{方法：} 能量法

真空中两个固定正点电荷 \(+Q\) 位于
\(z=\pm l\)，\(Q>0,l>0\)。一根光滑绝缘直杆通过原点并与 \(z\) 轴成锐角
\(\theta\)；杆上有质量 \(m>0\)、电荷 \(q>0\)
的小珠，仅能沿杆滑动。忽略重力、摩擦及小珠对固定电荷的影响。判断原点平衡的稳定角度范围，求稳定范围内的小振动角频率，并在临界角处用更高阶势能判断稳定性。

\par\noindent\begin{minipage}{\linewidth}
\textbf{解答}

\begin{center}
\includegraphics[width=\linewidth,height=.40\textheight,keepaspectratio]{\detokenize{assets/module-02/examples/M2-083.pdf}}
\par\smallskip{\small 图 M2-083：两等量正电荷与倾斜细杆}
\end{center}
\end{minipage}\par

图 M2-083 的细杆方向把三维 Hessian
投影到允许的一维位移。先沿杆代入两源的精确距离并展开，二次系数决定通常稳定范围；临界角的二次项为零时，还须继续检验四次项而不能直接判稳定。

由对称性，原点处 \(\mathbf E=0\)。在横轴上

\[
\varphi(x,0,0)
=\frac{2kQ}{\sqrt{l^2+x^2}}
\approx\frac{2kQ}{l}-\frac{kQ}{l^3}x^2.
\]

因此横向主曲率为

\[
\lambda_x=\lambda_y=-\frac{2kQ}{l^3}.
\]

由拉普拉斯方程

\[
\lambda_z=-\lambda_x-\lambda_y
=\frac{4kQ}{l^3}.
\]

沿杆单位向量 \(\mathbf t=(\sin\theta,0,\cos\theta)^{\mathsf T}\)，故

\[
\begin{aligned}
\lambda_{\mathrm{eff}}
&=\lambda_x\sin^2\theta
+\lambda_z\cos^2\theta\\
&=\frac{2kQ}{l^3}
(2\cos^2\theta-\sin^2\theta).
\end{aligned}
\]

因 \(q>0\)，稳定要求 \(\lambda_{\mathrm{eff}}>0\)：

\[
\boxed{\tan^2\theta<2.}
\]

即

\[
\boxed{0\le\theta<\arctan\sqrt2.}
\]

小振动角频率为

\[
\begin{aligned}
\omega
&=\sqrt{\frac{q\lambda_{\mathrm{eff}}}{m}}\\
&=\boxed{
\sqrt{
\frac{qQ(2\cos^2\theta-\sin^2\theta)}
{2\pi\varepsilon_0ml^3}}}.
\end{aligned}
\]

临界角满足 \(\cos^2\theta_c=1/3\)。二阶项为零，但四阶系数

\[
P_4(1/\sqrt3)=-\frac7{18}<0.
\]

因此

\[
\Delta U
=-\frac{7kqQ}{9l^5}h^4+O(h^6)<0
\]

对足够小的非零 \(h\) 成立，原点为局部极大，临界角处不稳定。

\subsubsection{作业三：退化调和平衡点}\label{m2:ux4f5cux4e1aux4e09ux9000ux5316ux8c03ux548cux5e73ux8861ux70b9}

\textbf{题号：M2-084}

\textbf{层级：} 决赛　　\textbf{方法：} 能量法

某真空区域内给定与 \(z\) 无关的电势 \(\varphi=B(x^3-3xy^2)\)，\(B\ne0\)
为常量。一个非零电荷的测试粒子可在原点附近自由运动，忽略其他力和自身对源的反作用。核对势的调和性，求原点的电场及电势
Hessian
矩阵，并判断原点是否稳定；若二阶判据退化，应利用更高阶局部电势变化判断。

\par\noindent\begin{minipage}{\linewidth}
\textbf{解答}

\begin{center}
\includegraphics[width=\linewidth,height=.40\textheight,keepaspectratio]{\detokenize{assets/module-02/examples/M2-084.pdf}}
\par\smallskip{\small 图 M2-084：退化调和平衡点}
\end{center}
\end{minipage}\par

图 M2-084
的原点处梯度和二阶项都消失，但三阶角向分布仍有正负交替。沿不同方向考察势能即可找到下降通道；因此零电场与零
Hessian 只是退化，不能单独推出稳定。

梯度为

\[
\nabla\varphi
=\bigl(3B(x^2-y^2),-6Bxy,0\bigr).
\]

所以

\[
\nabla\varphi(O)=0,
\qquad
\mathbf E(O)=0.
\]

二阶导数均正比于 \(x\) 或 \(y\)，故

\[
H(O)=0.
\]

二阶线性稳定判据完全失效。沿 \(y=0\) 有

\[
\varphi=Bx^3,
\]

它在 \(x\) 穿过零点时改变符号；更一般地

\[
\varphi=Br^3\cos3\theta
\]

沿相邻角扇区交替正负。因此原点不是电势极值，对任意非零 \(q\)
都不能形成严格稳定束缚。

\subsubsection{作业四：连续电荷区中的各向同性平衡}\label{m2:ux4f5cux4e1aux56dbux8fdeux7eedux7535ux8377ux533aux4e2dux7684ux5404ux5411ux540cux6027ux5e73ux8861}

\textbf{题号：M2-085}

\textbf{层级：} 决赛　　\textbf{方法：} 能量法

某区域内物理介质为真空，电势给定为 \(\varphi=\alpha r^2\)，其中
\(r=\sqrt{x^2+y^2+z^2}\)，\(\alpha\ne0\)。测试粒子的质量为
\(m>0\)、电荷为有符号量
\(q\ne0\)，它可自由运动且不改变源分布。求区域内的体电荷密度，分析原点平衡的三维稳定性，在稳定情形求小振动角频率，并说明本题为何不满足无源静电区域的前提。

\par\noindent\begin{minipage}{\linewidth}
\textbf{解答}

\begin{center}
\includegraphics[width=\linewidth,height=.40\textheight,keepaspectratio]{\detokenize{assets/module-02/examples/M2-085.pdf}}
\par\smallskip{\small 图 M2-085：连续电荷区中的各向同性平衡}
\end{center}
\end{minipage}\par

图 M2-085
的径向势对三个方向给出相同曲率，所需源由泊松方程确定。恢复力条件取决于
\(q\alpha\) 而非只看
\(\alpha\)，稳定时三个频率相同；此处有连续源，不能调用无源稳定性禁戒。

由泊松方程

\[
\nabla^2\varphi=6\alpha
=-\frac{\rho}{\varepsilon_0},
\]

得到

\[
\boxed{\rho=-6\varepsilon_0\alpha.}
\]

电场和力为

\[
\mathbf E=-2\alpha\mathbf r,
\qquad
\mathbf F=-2q\alpha\mathbf r.
\]

当 \(q\alpha>0\) 时，三个方向均有恢复力，原点稳定；运动方程

\[
m\ddot{\mathbf r}=-2q\alpha\mathbf r
\]

给出

\[
\boxed{\omega=\sqrt{\frac{2q\alpha}{m}}.}
\]

当 \(q\alpha<0\) 时原点不稳定。这里
\(\rho\ne0\)，所以稳定情形不违背恩绍定理。

\subsubsection{作业五：均匀电场中的圆环约束}\label{m2:ux4f5cux4e1aux4e94ux5747ux5300ux7535ux573aux4e2dux7684ux5706ux73afux7ea6ux675f}

\textbf{题号：M2-086}

\textbf{层级：} 决赛　　\textbf{方法：} 能量法

真空中半径为 \(R>0\)、圆心在原点的光滑绝缘圆环固定于 \(xy\) 平面。质量为
\(m>0\)、电荷为 \(q\ne0\)
的小珠只允许沿圆环运动；不计重力、摩擦及小珠自身对外场的影响。外场为
\(\mathbf E=E_0\mathbf e_x\)，\(E_0>0\)，用 \(x=R\cos\theta\)
参数化珠的位置。求全部平衡位置，按 \(q\)
的符号讨论稳定性，求稳定位置附近的小振动角频率。

\par\noindent\begin{minipage}{\linewidth}
\textbf{解答}

\begin{center}
\includegraphics[width=\linewidth,height=.40\textheight,keepaspectratio]{\detokenize{assets/module-02/examples/M2-086.pdf}}
\par\smallskip{\small 图 M2-086：均匀电场中的圆环约束}
\end{center}
\end{minipage}\par

图 M2-086 的允许路径是一条圆环，不是三维自由位移。先在环上把位置代为
\(x=R\cos\theta\)，再用势能的一阶和二阶角导数判平衡与稳定；弧长动力学中的质量因子为
\(mR^2\)。

取圆环参数

\[
x=R\cos\theta,
\qquad
y=R\sin\theta.
\]

均匀场电势可取

\[
\varphi=-E_0x=-E_0R\cos\theta.
\]

势能为

\[
U=-qE_0R\cos\theta.
\]

平衡条件

\[
\frac{dU}{d\theta}=qE_0R\sin\theta=0
\]

给出

\[
\theta=0,\pi.
\]

二阶导数为

\[
\frac{d^2U}{d\theta^2}
=qE_0R\cos\theta.
\]

若 \(qE_0>0\)，\(\theta=0\) 稳定、\(\theta=\pi\) 不稳定；若
\(qE_0<0\)，二者角色交换。在稳定点附近令小偏角为 \(\eta\)，则

\[
U=U_{\min}+\frac12|qE_0|R\eta^2+O(\eta^4).
\]

动能为

\[
T=\frac12mR^2\dot\eta^2.
\]

所以

\[
\boxed{
\omega=\sqrt{\frac{|qE_0|}{mR}}.}
\]

稳定点处空间电场仍为非零常量，电场力的法向分量由圆环约束力平衡；这正是受约束平衡与自由平衡的区别。

\section{第二单元：边值方程与解析方法}\label{m2:ux7b2cux4e8cux5355ux5143ux8fb9ux503cux65b9ux7a0bux4e0eux89e3ux6790ux65b9ux6cd5}

本单元从专题十七进入一般边值问题，依次使用分离变量、电像、特殊函数和保角变换。

\section{专题十七：静电边值问题的方程与边界条件（基础必讲）}\label{m2:ux4e13ux9898ux5341ux4e03ux9759ux7535ux8fb9ux503cux95eeux9898ux7684ux65b9ux7a0bux4e0eux8fb9ux754cux6761ux4ef6ux57faux7840ux5fc5ux8bb2}

静电边值问题由求解区域、源、控制方程和边界数据共同确定。偏微分方程规定区域内部的局域规律，边界条件选出具体解，唯一性定理用于确认候选解。分离变量、电像法、格林函数和数值方法放在后续专题与附录中。

\par\noindent\begin{minipage}{\linewidth}
\begin{center}
\includegraphics[width=\linewidth,height=.40\textheight,keepaspectratio]{\detokenize{assets/module-02/theory/M2-H-boundary.pdf}}
\par\smallskip{\small 图示：静电边值问题的区域、源、控制方程与边界数据}
\end{center}
\end{minipage}\par

\subsection{一、控制方程必须连同求解区域一起说明}\label{m2:ux4e00ux63a7ux5236ux65b9ux7a0bux5fc5ux987bux8fdeux540cux6c42ux89e3ux533aux57dfux4e00ux8d77ux8bf4ux660e}

在真空静电场中，

\[
\nabla\times\mathbf E=0,
\qquad
\nabla\cdot\mathbf E=\frac{\rho}{\varepsilon_0}.
\]

第一式保证在单连通静电区域内可以引入电势

\[
\mathbf E=-\nabla\varphi.
\]

代入高斯定律：

\[
\nabla\cdot(-\nabla\varphi)
=\frac{\rho}{\varepsilon_0},
\]

得到泊松方程

\[
\boxed{
\nabla^2\varphi=-\frac{\rho}{\varepsilon_0}
}.
\]

在没有体电荷且不包含奇点的子区域内，\(\rho=0\)，方程退化为

\[
\boxed{\nabla^2\varphi=0},
\]

这就是拉普拉斯方程。满足拉普拉斯方程的函数称为调和函数。

\subsubsection{1.
``无源区域''不等于``无电场区域''}\label{m2:ux65e0ux6e90ux533aux57dfux4e0dux7b49ux4e8eux65e0ux7535ux573aux533aux57df}

\(\nabla^2\varphi=0\) 只表示求解区域内部没有电荷源，不表示 \(\varphi\)
为常数或
\(\mathbf E=0\)。区域外的电荷、导体表面以及无穷远条件都可以通过边界数据在区域内部建立非零电场。

\textbf{例 M2-087：无源并不意味着场为零}

在一个没有体电荷的真空区域，考察候选电势
\(\varphi(x,y,z)=-E_0z\)，\(E_0\ne0\)
为常量，不要求此局部模型在无穷远趋零。核对它是否满足拉普拉斯方程，求对应电场，并解释为什么``源项为零''不能推出该区域内电场处处为零。

\par\noindent\begin{minipage}{\linewidth}
\textbf{解答}

\begin{center}
\includegraphics[width=\linewidth,height=.40\textheight,keepaspectratio]{\detokenize{assets/module-02/examples/M2-087.pdf}}
\par\smallskip{\small 图 M2-087：无源并不意味着场为零}
\end{center}
\end{minipage}\par

图 M2-087
的平行场箭头无散度，但每支都非零。线性电势满足拉普拉斯方程，说明无体源只约束曲率之和；没有给出常值接地壁面时，不能据此推断整个区域场为零。

函数

\[
\varphi=-E_0z
\]

满足

\[
\nabla^2\varphi=0,
\]

但

\[
\mathbf E=E_0\mathbf e_z\ne0.
\]

\subsubsection{2.
体电荷、点电荷和面电荷要分开处理}\label{m2:ux4f53ux7535ux8377ux70b9ux7535ux8377ux548cux9762ux7535ux8377ux8981ux5206ux5f00ux5904ux7406}

\par\noindent\begin{minipage}{\linewidth}
\begin{center}
\includegraphics[width=\linewidth,height=.40\textheight,keepaspectratio]{\detokenize{assets/module-02/original/tikz-m2-topic17-v8-02.png}}
\par\smallskip{\small 图示：光滑体电荷、点电荷奇点与面电荷界面的不同处理方式}
\end{center}
\end{minipage}\par

若 \(\rho\) 是光滑体电荷密度，泊松方程在区域内部逐点成立。

点电荷 \(q\) 位于 \(\mathbf r_0\) 时，可写成分布形式

\[
\nabla^2\varphi
=-\frac{q}{\varepsilon_0}
\delta^{(3)}(\mathbf r-\mathbf r_0).
\]

在普通函数计算中，也可以把点 \(\mathbf r_0\)
从区域中挖去：挖去后的区域满足拉普拉斯方程，而小球面上的通量条件保留点电荷强度。

面电荷不应简单塞入界面两侧的普通泊松方程中，而应通过法向场或法向导数的跳跃条件处理。

在均匀线性介质中，若介电常数为 \(\varepsilon\)、自由体电荷为
\(\rho_f\)，则

\[
\nabla\cdot(\varepsilon\nabla\varphi)=-\rho_f.
\]

当 \(\varepsilon\) 在该区域内为常数时，才可化为

\[
\nabla^2\varphi=-\frac{\rho_f}{\varepsilon}.
\]

\subsection{二、边界条件：局部类型、分段组合与界面跳跃}\label{m2:ux4e8cux8fb9ux754cux6761ux4ef6ux5c40ux90e8ux7c7bux578bux5206ux6bb5ux7ec4ux5408ux4e0eux754cux9762ux8df3ux8dc3}

设求解区域为 \(\Omega\)，边界为 \(\partial\Omega\)，\(\mathbf n\)
取区域的外法向。常见局部边界条件分为三类。

\subsubsection{1.
第一类边界条件}\label{m2:ux7b2cux4e00ux7c7bux8fb9ux754cux6761ux4ef6-1}

直接给定边界电势：

\[
\boxed{
\varphi|_{\partial\Omega}=f
}.
\]

接地导体表面满足 \(f=0\)；保持在固定电势 \(V_0\) 的导体表面满足
\(f=V_0\)。

\subsubsection{2.
第二类边界条件}\label{m2:ux7b2cux4e8cux7c7bux8fb9ux754cux6761ux4ef6-1}

给定沿外法向的导数：

\[
\boxed{
\left.\frac{\partial\varphi}{\partial n}
\right|_{\partial\Omega}=g
}.
\]

由于

\[
E_n=\mathbf E\cdot\mathbf n
=-\frac{\partial\varphi}{\partial n},
\]

第二类条件等价于给定法向电场。

\subsubsection{3.
第三类边界条件}\label{m2:ux7b2cux4e09ux7c7bux8fb9ux754cux6761ux4ef6}

在同一边界点给定电势与法向导数的线性组合：

\[
\boxed{
\alpha\varphi+
\beta\frac{\partial\varphi}{\partial n}=h
}.
\]

这也称为罗宾边界条件。

\subsubsection{4.
分段混合边界}\label{m2:ux5206ux6bb5ux6df7ux5408ux8fb9ux754c}

若边界分成 \(\Gamma_D\) 与 \(\Gamma_N\)，在 \(\Gamma_D\) 上给定电势，在
\(\Gamma_N\) 上给定法向导数，则称为分段混合边界：

\[
\varphi|_{\Gamma_D}=f,
\qquad
\partial_n\varphi|_{\Gamma_N}=g.
\]

它与第三类条件不同：第三类条件是在同一点给出线性组合；分段混合条件是在不同边界段给不同类型的数据。

\par\noindent\begin{minipage}{\linewidth}
\begin{center}
\includegraphics[width=\linewidth,height=.40\textheight,keepaspectratio]{\detokenize{assets/module-02/original/tikz-m2-topic17-v8-03.png}}
\par\smallskip{\small 图示：第一类、第二类、第三类与分段混合边界条件}
\end{center}
\end{minipage}\par

\subsubsection{5.
纯第二类问题的相容条件}\label{m2:ux7eafux7b2cux4e8cux7c7bux95eeux9898ux7684ux76f8ux5bb9ux6761ux4ef6}

对泊松方程在 \(\Omega\) 上积分：

\[
\int_{\Omega}\nabla^2\varphi\,dV
=-\frac1{\varepsilon_0}
\int_{\Omega}\rho\,dV.
\]

由散度定理，

\[
\oint_{\partial\Omega}
\frac{\partial\varphi}{\partial n}\,dS
=-\frac1{\varepsilon_0}
\int_{\Omega}\rho\,dV.
\]

因此纯第二类问题必须满足

\[
\boxed{
\oint_{\partial\Omega}g\,dS
=-\frac1{\varepsilon_0}
\int_{\Omega}\rho\,dV
}.
\]

若这个条件不满足，问题无解。若满足，电势仍可整体加上任意常数，但电场不变。

\textbf{例 M2-088：边界总法向导数必须与源相容}

设 \(V\) 为有限体积的连通真空区域，闭合边界 \(S\)
分片光滑且取外法向。区内无任何电荷或奇点，要求电势足够光滑，并在整个边界指定同一正常数
\(\partial\varphi/\partial n=g_0>0\)。利用拉普拉斯方程的体积分与散度定理判断此诺伊曼边值问题是否可解，指出全部边界数据应满足的总体约束。

\par\noindent\begin{minipage}{\linewidth}
\textbf{解答}

\begin{center}
\includegraphics[width=\linewidth,height=.40\textheight,keepaspectratio]{\detokenize{assets/module-02/examples/M2-088.pdf}}
\par\smallskip{\small 图 M2-088：边界总法向导数必须与源相容}
\end{center}
\end{minipage}\par

图 M2-088
的箭头表示假定在整个闭合边界上同号的法向斜率。将它沿边界积分必为正，而无体源要求同一个积分为零；因此矛盾来自边界条件，不是求解技巧或电势零点问题。

若 \(\rho=0\)，却在整个边界上指定同号的常数
\(\partial_n\varphi=g_0>0\)，则

\[
\oint_{\partial\Omega}g_0\,dS>0,
\]

与相容条件矛盾，所以这样的纯第二类问题不存在解。

\subsubsection{6.
介质界面与导体表面}\label{m2:ux4ecbux8d28ux754cux9762ux4e0eux5bfcux4f53ux8868ux9762}

令 \(\mathbf n\) 从介质 1 指向介质
2。若界面不存在理想偶极层，静电势连续：

\[
\boxed{\varphi_2=\varphi_1}.
\]

由

\[
\mathbf n\cdot(\mathbf D_2-\mathbf D_1)=\sigma_f,
\qquad
\mathbf D=-\varepsilon\nabla\varphi,
\]

得到

\[
\boxed{
\varepsilon_2\partial_n\varphi_2
-\varepsilon_1\partial_n\varphi_1
=-\sigma_f
}.
\]

\par\noindent\begin{minipage}{\linewidth}
\begin{center}
\includegraphics[width=\linewidth,height=.40\textheight,keepaspectratio]{\detokenize{assets/module-02/original/tikz-m2-topic17-v8-04.png}}
\par\smallskip{\small 图示：两介质界面和真空导体表面的法向约定与跳跃条件}
\end{center}
\end{minipage}\par

对真空中的导体，取 \(\mathbf n\) 从导体内部指向真空外侧。导体内部

\[
\mathbf E_{\mathrm{\text{内}}}=0,
\qquad
\varphi=\text{常数},
\]

紧贴外表面有

\[
\boxed{
-\left.\frac{\partial\varphi}{\partial n}
\right|_{\mathrm{\text{外}}}
=\frac{\sigma}{\varepsilon_0}
}.
\]

该式的符号依赖 \(\mathbf n\) 的方向；写跳跃条件前必须先画出法向。

\subsubsection{7. 对称面条件}\label{m2:ux5bf9ux79f0ux9762ux6761ux4ef6}

``对称面''不能一律写成 \(\partial_n\varphi=0\)。

\begin{itemize}
\tightlist
\item
  若电荷与边界关于该面镜像偶对称，则 \(\varphi\) 为偶函数，面上
  \(\partial_n\varphi=0\)；
\item
  若两侧边界数据关于该面反对称，并把零电势取在对称面上，则 \(\varphi\)
  为奇函数，面上 \(\varphi=0\)。
\end{itemize}

\subsection{三、唯一性定理在边值问题中的调用方式}\label{m2:ux4e09ux552fux4e00ux6027ux5b9aux7406ux5728ux8fb9ux503cux95eeux9898ux4e2dux7684ux8c03ux7528ux65b9ux5f0f}

专题十五已经用格林第一恒等式给出了唯一性定理的严格证明。本专题不重复证明，而把它整理成构造候选解时必须逐项核对的条件。设
\(\varphi_1\)、\(\varphi_2\) 是两个候选解，令
\(u=\varphi_1-\varphi_2\)；只有在同一求解域内源项相同，才能先得到
\(\nabla^2u=0\)。

\par\noindent\begin{minipage}{\linewidth}
\begin{center}
\includegraphics[width=\linewidth,height=.40\textheight,keepaspectratio]{\detokenize{assets/module-02/original/tikz-m2-topic17-v8-05.png}}
\par\smallskip{\small 图示：调用唯一性定理时依次检查方程、边界与无穷远条件}
\end{center}
\end{minipage}\par

\subsubsection{1.
先检查求解域与方程}\label{m2:ux5148ux68c0ux67e5ux6c42ux89e3ux57dfux4e0eux65b9ux7a0b}

\begin{itemize}
\tightlist
\item
  两个候选解必须在\textbf{同一个连通求解域}内比较；点电荷、线电荷和导体内部等被挖去的区域必须完全一致。
\item
  候选解必须具有相同的体源与奇点强度。仅在普通点满足
  \(\nabla^2\varphi=0\)，却漏掉一个 \(\delta\) 源，不能调用唯一性定理。
\end{itemize}

\subsubsection{2.
再按边界类型检查齐次差}\label{m2:ux518dux6309ux8fb9ux754cux7c7bux578bux68c0ux67e5ux9f50ux6b21ux5dee}

\begin{itemize}
\tightlist
\item
  第一类边界：\(u=0\)，电势唯一。
\item
  纯第二类边界：\(\partial_n u=0\)，只能推出
  \(u=C\)；还要先检查总通量相容条件，电势需再指定一个参考常数。
\item
  分段混合边界：在非空的 \(\Gamma_D\) 上 \(u=0\)，在 \(\Gamma_N\) 上
  \(\partial_n u=0\)，通常可得到 \(u=0\)。
\item
  第三类边界：若
  \(\alpha u+\beta\,\partial_n u=0\)，还需检查使边界能量项具有确定符号的系数条件，例如
  \(\beta\ne0\) 且 \(\alpha/\beta\ge0\)。
\end{itemize}

\subsubsection{3.
外部区域必须补上无穷远条件}\label{m2:ux5916ux90e8ux533aux57dfux5fc5ux987bux8865ux4e0aux65e0ux7a77ux8fdcux6761ux4ef6}

无界区域要先用大球截断，再令球半径
\(R\to\infty\)。候选解必须具有题目规定的衰减或远场渐近行为，使无穷远边界项消失；否则仍可叠加常数或其他不衰减调和函数，唯一性结论并不成立。

\textbf{例 M2-089：像源只能放在真实求解域之外}

考虑真空真实求解域 \(D\)
内的静电狄利克雷问题：真实电荷分布已知，各导体边界电势已知；若 \(D\)
无界，还规定正确的无穷远行为。现用真实源与位于 \(D\)
之外的虚拟像源写出候选势
\(\widetilde\varphi\)。列出必须在真实域和边界上完成的检验，说明何时可由唯一性定理确认候选势，并解释为何不能在无真实电荷的位置向
\(D\) 内加入像源奇点。

\par\noindent\begin{minipage}{\linewidth}
\textbf{解答}

\begin{center}
\includegraphics[width=\linewidth,height=.40\textheight,keepaspectratio]{\detokenize{assets/module-02/examples/M2-089.pdf}}
\par\smallskip{\small 图 M2-089：像源只能放在真实求解域之外}
\end{center}
\end{minipage}\par

图 M2-089 的灰色半空间容纳辅助像源，真实半空间只含原来的 \(q\)。边界点
\(P\)
的两段距离相等用来逐点验证零势，再核对域内奇点和无穷远行为；这些条件全部成立后，才可用唯一性确认候选势。

若候选电势在真实区域内具有正确的源和奇点，在整段导体边界上满足给定电势，并满足正确的无穷远行为，那么即使表达式使用了真实区域外的像电荷，唯一性定理仍保证它就是该区域内的物理解。

\subsection{四、最大值原理与平均值性质}\label{m2:ux56dbux6700ux5927ux503cux539fux7406ux4e0eux5e73ux5747ux503cux6027ux8d28}

设 \(\varphi\) 在有界连通区域 \(\Omega\) 内调和，并在闭区域上连续。若
\(\varphi\) 不是常数，则其最大值和最小值只能出现在边界
\(\partial\Omega\) 上。

换言之，若

\[
m\le\varphi|_{\partial\Omega}\le M,
\]

则区域内部也有

\[
\boxed{
m\le\varphi(\mathbf r)\le M
}.
\]

\par\noindent\begin{minipage}{\linewidth}
\begin{center}
\includegraphics[width=\linewidth,height=.40\textheight,keepaspectratio]{\detokenize{assets/module-02/original/tikz-m2-topic17-v8-06.png}}
\par\smallskip{\small 图示：调和电势的内部取值被边界最大值和最小值夹住}
\end{center}
\end{minipage}\par

更强的表述是：若调和函数在内部达到最大值或最小值，它在该连通区域内只能是常数。

调和函数还满足球面平均值性质。若以 \(\mathbf r_0\) 为球心、半径 \(R\)
的闭球完全位于无源区域内，则

\[
\boxed{
\varphi(\mathbf r_0)
=\frac1{4\pi R^2}
\oint_{S_R}\varphi\,dS
}.
\]

中心电势等于任意充分小球面上的平均电势，所以它不可能严格大于或严格小于周围所有点。

\subsubsection{物理后果}\label{m2:ux7269ux7406ux540eux679c}

\begin{enumerate}
\def\labelenumi{\arabic{enumi}.}
\tightlist
\item
  无源区域内部不能出现电势的严格孤立极大值或极小值；
\item
  若无源区域全部边界都接地，且区域内没有奇点，则唯一解是 \(\varphi=0\)；
\item
  若边界电势都在 \(0\) 与 \(V_0\) 之间，内部电势不可能越出该范围；
\item
  最大值原理不能跨过点电荷奇点、面电荷界面或未指定的无穷远边界直接使用。
\end{enumerate}

\textbf{例 M2-090：空腔零边界迫使内部势场全零}

一个接地连通导体围成完全封闭的空腔，腔内为真空且没有电荷，空腔壁电势处处为零。假设腔内电势在内部及边界满足最大值原理所需的正则条件。利用调和函数的最大值原理确定空腔内电势和电场，并说明该结论依赖哪些边界与源条件。

\par\noindent\begin{minipage}{\linewidth}
\textbf{解答}

\begin{center}
\includegraphics[width=\linewidth,height=.40\textheight,keepaspectratio]{\detokenize{assets/module-02/examples/M2-090.pdf}}
\par\smallskip{\small 图 M2-090：空腔零边界迫使内部势场全零}
\end{center}
\end{minipage}\par

图 M2-090
的封闭边界处处为零势，内部又无电荷。调和函数的最大值和最小值都只能取边界值，因而两者均为零；势处处相同，电场也随负梯度全为零。

若空腔内部无电荷，腔壁全部接地，则边界上
\(\varphi=0\)。最大值原理给出内部最大值和最小值都为零，因此

\[
\varphi=0,
\qquad
\mathbf E=0.
\]

\subsection{五、按边界形状选择求解方法}\label{m2:ux4e94ux6309ux8fb9ux754cux5f62ux72b6ux9009ux62e9ux6c42ux89e3ux65b9ux6cd5}

\par\noindent\begin{minipage}{\linewidth}
\begin{center}
\includegraphics[width=\linewidth,height=.40\textheight,keepaspectratio]{\detokenize{assets/module-02/original/tikz-m2-topic17-v8-07.png}}
\par\smallskip{\small 图示：根据源、对称性和边界几何选择静电求解方法}
\end{center}
\end{minipage}\par

{\def\LTcaptype{none} % do not increment counter
\begin{longtable}[]{@{}
  >{\raggedright\arraybackslash}p{(\linewidth - 4\tabcolsep) * \real{0.3333}}
  >{\raggedright\arraybackslash}p{(\linewidth - 4\tabcolsep) * \real{0.3333}}
  >{\raggedright\arraybackslash}p{(\linewidth - 4\tabcolsep) * \real{0.3333}}@{}}
\toprule\noalign{}
\begin{minipage}[b]{\linewidth}\raggedright
问题特征
\end{minipage} & \begin{minipage}[b]{\linewidth}\raggedright
优先方法
\end{minipage} & \begin{minipage}[b]{\linewidth}\raggedright
使用前提与检查
\end{minipage} \\
\midrule\noalign{}
\endhead
\bottomrule\noalign{}
\endlastfoot
电荷分布已知、边界影响可忽略 & 库仑积分与叠加 &
积分收敛，源区描述完整 \\
平面、柱面或球面对称 & 高斯定理 & 场强大小在高斯面上可由对称性确定 \\
接地平面、球面等特殊导体 & 电像法 &
像电荷在物理区域外，候选解满足全部边界 \\
边界与正交坐标面吻合 & 分离变量 & 正确筛选本征模态并满足端点条件 \\
一般区域与一般边界 & 格林函数 & 格林函数的边界类型与原问题一致 \\
局域源的远场 & 多极展开 & 场点在包围源区的球外，且截断误差可控 \\
平衡点或对称中心附近 & 泰勒展开 & 展开点光滑，必要时继续检查高阶项 \\
几何或材料非常复杂 & 数值方法 & 网格、边界截断、收敛性和误差验证 \\
\end{longtable}
}

方法可以组合。例如先用电像法处理导体边界，再对远区结果做多极展开；或先用分离变量得到级数，再用最大值原理和唯一性定理检查结果。

无论候选解怎样得到，都必须逐项核对：

\begin{enumerate}
\def\labelenumi{\arabic{enumi}.}
\tightlist
\item
  在每个子区域内满足正确的泊松方程或拉普拉斯方程；
\item
  点电荷奇点的位置和强度正确；
\item
  界面两侧满足电势连续与法向跳跃条件；
\item
  所有有限边界条件都满足；
\item
  无穷远行为或周期条件正确；
\item
  最后再引用唯一性定理确认。
\end{enumerate}

\subsection{六、边值问题的边界条件检查}\label{m2:ux516dux8fb9ux503cux95eeux9898ux7684ux8fb9ux754cux6761ux4ef6ux68c0ux67e5}

\par\noindent\begin{minipage}{\linewidth}
\begin{center}
\includegraphics[width=\linewidth,height=.40\textheight,keepaspectratio]{\detokenize{assets/module-02/original/tikz-m2-topic17-v8-08.png}}
\par\smallskip{\small 图示：静电边值题的区域、边界条件和唯一性检查}
\end{center}
\end{minipage}\par

\begin{enumerate}
\def\labelenumi{\arabic{enumi}.}
\tightlist
\item
  \textbf{方程必须带区域。}
  ``这里满足拉普拉斯方程''必须同时说明该区域不含哪些电荷或奇点。
\item
  \textbf{无源不等于无场。} 外部源和边界数据仍可产生非零调和场。
\item
  \textbf{第二类数据要固定法向。} 改变法向会改变 \(g\)
  和跳跃公式的符号。
\item
  \textbf{纯第二类问题先查相容条件。} 相容条件不满足时，不应继续求解。
\item
  \textbf{分段混合条件不等于第三类条件。}
  前者在不同边界段给不同数据，后者在同一点给线性组合。
\item
  \textbf{导体表面电势为常数，但常数未必已知。} 接地导体才是
  \(\varphi=0\)；孤立导体还需总电荷条件确定常数。
\item
  \textbf{对称面可能对应两种条件。} 偶对称通常给
  \(\partial_n\varphi=0\)，奇对称通常给 \(\varphi=0\)。
\item
  \textbf{电势连续有前提。}
  理想偶极层可造成电势跳跃；普通面电荷造成的是法向导数跳跃。
\item
  \textbf{最大值原理只用于调和的光滑区域。}
  不能把点电荷奇点包含在内部后仍直接使用。
\item
  \textbf{外部问题要写无穷远条件。} 无穷远也是边界的一部分。
\item
  \textbf{构造解不是猜完即止。} 控制方程、边界、奇点、渐近行为缺一不可。
\item
  \textbf{电势与势能要区分。}
  电势的局域曲率不等于任意带电粒子的势能稳定性，后者还要乘以粒子电荷。
\end{enumerate}

\subsection{七、练习与解答}\label{m2:ux4e03ux7ec3ux4e60ux4e0eux89e3ux7b54-2}

\subsubsection{作业一：两等量正电荷的远区、近中心与中垂线场强}\label{m2:ux4f5cux4e1aux4e00ux4e24ux7b49ux91cfux6b63ux7535ux8377ux7684ux8fdcux533aux8fd1ux4e2dux5fc3ux4e0eux4e2dux5782ux7ebfux573aux5f3a}

\textbf{题号：M2-091}

\textbf{层级：} 复赛　　\textbf{方法：} 边值问题、尺度分析

\(Q>0,\ell>0\)；球坐标从 \(+x\)
轴量极角，\(r=\sqrt{x^2+y^2+z^2}\)。近中心展开要求
\(r\ll\ell\)，远区展开要求 \(r\gg\ell\)，不能将两种近似混用。

两个正点电荷 \(+Q\) 分别位于 \(x=\ell\) 和
\(x=-\ell\)，整个系统处于真空，无穷远电势取零。

\begin{enumerate}
\def\labelenumi{\arabic{enumi}.}
\tightlist
\item
  在 \(r\gg\ell\)
  处求电势至四极修正，并求给定等势面在三维空间中的近似形态；
\item
  在原点附近求电势的二阶展开，讨论三维等势面和 \(z=0\) 截面等势线；
\item
  求 \(y\) 轴上 \(|E_y|\) 的最大值位置与最大值。
\end{enumerate}

\par\noindent\begin{minipage}{\linewidth}
\textbf{解答}

\begin{center}
\includegraphics[width=\linewidth,height=.40\textheight,keepaspectratio]{\detokenize{assets/module-02/examples/M2-091.pdf}}
\par\smallskip{\small 图 M2-091：两等量正电荷的远区、近中心与中垂线场强}
\end{center}
\end{minipage}\par

图 M2-091
的两正电荷关于中心对称，但远区、中心附近和中垂线上的小参数不同。远区按
\(\ell/r\) 展开，中心附近按 \(r/\ell\) 展开；在 \(y\)
轴求极值时应先使用该轴上的精确场，而不是把任一局部展开外推到全轴。

任意场点 \((x,y,z)\) 的精确电势为

\[
\boxed{
\varphi(x,y,z)
=kQ\left[
\frac1{\sqrt{(x-\ell)^2+y^2+z^2}}
+\frac1{\sqrt{(x+\ell)^2+y^2+z^2}}
\right]
},
\]

其中

\[
k=\frac1{4\pi\varepsilon_0}.
\]

\paragraph{（1）远区等势面}\label{m2:ux8fdcux533aux7b49ux52bfux9762}

取球坐标极轴为 \(x\) 轴：

\[
x=r\cos\theta.
\]

两个电荷关于原点反演对称，所以所有奇数阶多极项消失。远区展开为

\[
\boxed{
\varphi(r,\theta)
=2kQ\left[
\frac1r
+\frac{\ell^2}{r^3}P_2(\cos\theta)
+O\!\left(\frac{\ell^4}{r^5}\right)
\right]
}.
\]

单极项对应总电荷 \(2Q\)；偶极项为零；下一项是四极修正。

设所考察的等势面满足

\[
\varphi=\frac{2kQ}{R},
\qquad
R\gg\ell.
\]

令

\[
r(\theta)=R+\delta r(\theta),
\qquad
\delta r=O\!\left(\frac{\ell^2}{R}\right).
\]

展开

\[
\frac1r
\simeq
\frac1R-\frac{\delta r}{R^2},
\]

并在四极项中令 \(r\simeq R\)，得到

\[
\frac1R-\frac{\delta r}{R^2}
+\frac{\ell^2}{R^3}P_2(\cos\theta)
=\frac1R.
\]

因此

\[
\boxed{
r(\theta)
=R+\frac{\ell^2}{R}P_2(\cos\theta)
+O\!\left(\frac{\ell^4}{R^3}\right)
}.
\]

\par\noindent\begin{minipage}{\linewidth}
\begin{center}
\includegraphics[width=\linewidth,height=.40\textheight,keepaspectratio]{\detokenize{assets/module-02/original/tikz-m2-topic17-v8-09.png}}
\par\smallskip{\small 图示：两等量正电荷远区等势面的子午截面与四极形变}
\end{center}
\end{minipage}\par

在电荷轴方向 \(\theta=0,\pi\)，\(P_2=1\)：

\[
r_{\mathrm{\text{轴}}}
=R+\frac{\ell^2}{R}.
\]

在赤道方向 \(\theta=\pi/2\)，\(P_2=-1/2\)：

\[
r_{\mathrm{\text{赤道}}}
=R-\frac{\ell^2}{2R}.
\]

所以三维等势面是绕 \(x\) 轴旋转的近球面，并沿电荷轴轻微伸长。它保持
\(x\to-x\) 和任意绕 \(x\) 轴旋转的对称性，不是左右不对称的卵形面。

\paragraph{（2）原点附近的局域等势面}\label{m2:ux539fux70b9ux9644ux8fd1ux7684ux5c40ux57dfux7b49ux52bfux9762}

记

\[
s^2=x^2+y^2+z^2,
\qquad
s\ll\ell.
\]

对第一项展开：

\[
\begin{aligned}
\frac1{\sqrt{(x-\ell)^2+y^2+z^2}}
&=\frac1\ell
\left(
1-\frac{2x}{\ell}+\frac{s^2}{\ell^2}
\right)^{-1/2}\\
&\simeq\frac1\ell
\left[
1+\frac{x}{\ell}
+\frac{2x^2-y^2-z^2}{2\ell^2}
\right].
\end{aligned}
\]

另一项只需令 \(x\to-x\)。两项相加，线性项消失：

\[
\boxed{
\varphi(x,y,z)
=\frac{kQ}{\ell}
\left[
2+\frac{2x^2-y^2-z^2}{\ell^2}
+O\!\left(\frac{s^4}{\ell^4}\right)
\right]
}.
\]

相对原点电势

\[
\varphi_O=\frac{2kQ}{\ell}
\]

的偏差为

\[
\Delta\varphi
=\frac{kQ}{\ell^3}(2x^2-y^2-z^2).
\]

因此局域等势面满足

\[
\boxed{2x^2-y^2-z^2=C}.
\]

\begin{itemize}
\tightlist
\item
  \(C>0\) 时是沿 \(x\) 轴张开的双叶双曲面；
\item
  \(C<0\) 时是绕 \(x\) 轴的一叶双曲面；
\item
  \(C=0\) 时退化为圆锥面 \(y^2+z^2=2x^2\)。
\end{itemize}

在 \(z=0\) 截面内，等势线为

\[
\boxed{2x^2-y^2=C},
\]

且 \(C=0\) 时的两条渐近线是

\[
y=\pm\sqrt2x.
\]

\par\noindent\begin{minipage}{\linewidth}
\begin{center}
\includegraphics[width=\linewidth,height=.40\textheight,keepaspectratio]{\detokenize{assets/module-02/original/tikz-m2-topic17-v8-10.png}}
\par\smallskip{\small 图示：两等量正电荷原点附近的双曲等势线和中垂线场强}
\end{center}
\end{minipage}\par

原点满足
\(\nabla\varphi=0\)，但不同方向的二阶曲率异号，所以它是电势鞍点，不是局部极大值或极小值。若讨论带电粒子的势能，还必须使用
\(U=q\varphi\)，不能把电势方向直接称为任意粒子的稳定或不稳定方向。

\paragraph{\texorpdfstring{（3）\(y\)
轴上的场强极值}{（3）y 轴上的场强极值}}\label{m2:y-ux8f74ux4e0aux7684ux573aux5f3aux6781ux503c}

令场点为 \((0,H,0)\)。两个电荷的 \(x\) 分量互相抵消，\(y\) 分量相加：

\[
\boxed{
E_y(H)
=\frac{2kQH}{(H^2+\ell^2)^{3/2}}
}.
\]

求导：

\[
\frac{dE_y}{dH}
=2kQ
\frac{\ell^2-2H^2}
{(H^2+\ell^2)^{5/2}}.
\]

在 \(H>0\) 半轴，极大值满足

\[
\ell^2-2H^2=0,
\]

所以

\[
\boxed{H=\frac\ell{\sqrt2}}.
\]

由奇对称性，\(H=-\ell/\sqrt2\) 处 \(E_y\)
取最小值，二者的场强大小相同。因此

\[
\boxed{
|H|=\frac\ell{\sqrt2}
}
\]

是 \(|E_y|\) 的两个最大值位置。代回原式：

\[
\boxed{
|E_y|_{\max}
=\frac{4kQ}{3\sqrt3\,\ell^2}
}.
\]

\subsubsection{作业二：均匀带电平板层中的第一类边值问题}\label{m2:ux4f5cux4e1aux4e8cux5747ux5300ux5e26ux7535ux5e73ux677fux5c42ux4e2dux7684ux7b2cux4e00ux7c7bux8fb9ux503cux95eeux9898}

\textbf{题号：M2-092}

\textbf{层级：} 复赛　　\textbf{方法：} 边值问题

这是与 \(y,z\) 无关的真空一维平板模型，\(L>0\)；\(\rho_0,V_0\)
均可带符号，电场定义为 \(E_x=-d\varphi/dx\)，只求开区间内的场。

区域 \(0<x<L\) 内均匀分布体电荷 \(\rho_0\)，两端边界电势为

\[
\varphi(0)=0,
\qquad
\varphi(L)=V_0.
\]

\begin{enumerate}
\def\labelenumi{\arabic{enumi}.}
\tightlist
\item
  求区域内电势和电场；
\item
  求电场为零的位置，并说明该点何时位于区域内部；
\item
  检查泊松方程和两端边界条件。
\end{enumerate}

\par\noindent\begin{minipage}{\linewidth}
\textbf{解答}

\begin{center}
\includegraphics[width=\linewidth,height=.40\textheight,keepaspectratio]{\detokenize{assets/module-02/examples/M2-092.pdf}}
\par\smallskip{\small 图 M2-092：均匀带电平板层中的第一类边值问题}
\end{center}
\end{minipage}\par

图 M2-092
画出非零端电势差及非零体源的一种示例，均匀体源决定势曲线的固定曲率。二次通解先通过两端边界，再对其求负导数；场零点对应曲线水平切线，最后还要检查该横坐标是否落在
\(0<x<L\) 内。

问题只有 \(x\) 方向变化，泊松方程化为

\[
\frac{d^2\varphi}{dx^2}
=-\frac{\rho_0}{\varepsilon_0}.
\]

积分两次：

\[
\varphi(x)
=-\frac{\rho_0}{2\varepsilon_0}x^2
+Ax+B.
\]

由 \(\varphi(0)=0\) 得

\[
B=0.
\]

由 \(\varphi(L)=V_0\) 得

\[
-\frac{\rho_0L^2}{2\varepsilon_0}
+AL=V_0,
\]

所以

\[
A=\frac{V_0}{L}
+\frac{\rho_0L}{2\varepsilon_0}.
\]

最终

\[
\boxed{
\varphi(x)
=-\frac{\rho_0x^2}{2\varepsilon_0}
+\left(
\frac{V_0}{L}
+\frac{\rho_0L}{2\varepsilon_0}
\right)x
}.
\]

电场为

\[
\boxed{
E_x(x)
=-\frac{d\varphi}{dx}
=\frac{\rho_0x}{\varepsilon_0}
-\frac{V_0}{L}
-\frac{\rho_0L}{2\varepsilon_0}
}.
\]

当 \(\rho_0\ne0\) 时，\(E_x=0\) 的位置为

\[
\boxed{
x_0=\frac L2
+\frac{\varepsilon_0V_0}{\rho_0L}
}.
\]

只有满足

\[
0<x_0<L
\]

时，该零场点才位于区域内部；否则电场在整个区域内不改变方向。 若
\(\rho_0=0\)，则 \(E_x=-V_0/L\)：\(V_0\ne0\)
时区域内没有零场点，\(V_0=0\) 时整个区域都是零场。

对结果求二阶导数：

\[
\frac{d^2\varphi}{dx^2}
=-\frac{\rho_0}{\varepsilon_0},
\]

并直接代入
\(x=0,L\)，可验证控制方程和两端电势全部满足。第一类问题的唯一性保证这就是唯一解。

\subsubsection{作业三：球形区域中纯第二类问题的相容条件}\label{m2:ux4f5cux4e1aux4e09ux7403ux5f62ux533aux57dfux4e2dux7eafux7b2cux4e8cux7c7bux95eeux9898ux7684ux76f8ux5bb9ux6761ux4ef6}

\textbf{题号：M2-093}

\textbf{层级：} 复赛　　\textbf{方法：} 边值问题

球心取为原点，\(R>0\)，介质为真空；要求球心电势有限，球面上所有法向导数均相对于球体外法向定义。

半径为 \(R\) 的球形区域内均匀分布体电荷 \(\rho_0\)。在整个球面上指定

\[
\left.\frac{\partial\varphi}{\partial r}\right|_{r=R}=g,
\]

其中外法向为 \(\mathbf e_r\)。

\begin{enumerate}
\def\labelenumi{\arabic{enumi}.}
\tightlist
\item
  由相容条件确定 \(g\) 必须满足的关系；
\item
  在该关系满足时求球内电势；
\item
  说明解中仍有何种不确定性。
\end{enumerate}

\par\noindent\begin{minipage}{\linewidth}
\textbf{解答}

\begin{center}
\includegraphics[width=\linewidth,height=.40\textheight,keepaspectratio]{\detokenize{assets/module-02/examples/M2-093.pdf}}
\par\smallskip{\small 图 M2-093：球形区域中纯第二类问题的相容条件}
\end{center}
\end{minipage}\par

图 M2-093 的球面上外法向就是 \(\mathbf e_r\)，所以常量 \(g\)
乘球面积应等于负的总体电荷除以
\(\varepsilon_0\)。先用这一积分条件筛选数据，再积分径向方程；相容条件不会顺带固定电势的加法常数。

相容条件为

\[
\oint_{r=R}g\,dS
=-\frac1{\varepsilon_0}
\int_{r<R}\rho_0\,dV.
\]

两边分别为

\[
4\pi R^2g
\]

与

\[
-\frac{\rho_0}{\varepsilon_0}
\frac{4\pi R^3}{3}.
\]

因此必须有

\[
\boxed{
g=-\frac{\rho_0R}{3\varepsilon_0}
}.
\]

若 \(g\) 不满足此式，则问题无解。

当相容条件满足时，由球对称性

\[
\frac1{r^2}\frac d{dr}
\left(r^2\frac{d\varphi}{dr}\right)
=-\frac{\rho_0}{\varepsilon_0}.
\]

积分一次：

\[
r^2\frac{d\varphi}{dr}
=-\frac{\rho_0r^3}{3\varepsilon_0}+C_1.
\]

为使 \(r=0\) 处电势和电场有限，必须取

\[
C_1=0.
\]

所以

\[
\frac{d\varphi}{dr}
=-\frac{\rho_0r}{3\varepsilon_0},
\]

再积分得

\[
\boxed{
\varphi(r)
=-\frac{\rho_0r^2}{6\varepsilon_0}+C_0
}.
\]

在 \(r=R\) 处，

\[
\left.\frac{d\varphi}{dr}\right|_R
=-\frac{\rho_0R}{3\varepsilon_0}=g,
\]

边界条件满足。常数 \(C_0\) 不能由纯第二类条件确定；但

\[
\mathbf E
=-\nabla\varphi
=\frac{\rho_0r}{3\varepsilon_0}\mathbf e_r
\]

完全唯一。

\subsubsection{作业四：用唯一性定理检验接地平面的候选电势}\label{m2:ux4f5cux4e1aux56dbux7528ux552fux4e00ux6027ux5b9aux7406ux68c0ux9a8cux63a5ux5730ux5e73ux9762ux7684ux5019ux9009ux7535ux52bf}

\textbf{题号：M2-094}

\textbf{层级：} 复赛　　\textbf{方法：} 电像法、边值问题

\(a>0\)，\(k=1/(4\pi\varepsilon_0)\)；物理域不包括像源位置，点电荷自身位置的奇性需按源项处理，导体表面场取真空侧极限。

真空半空间 \(z>0\) 中，点电荷 \(q\) 位于 \((0,0,a)\)，平面 \(z=0\)
接地。柱坐标径向坐标记为 \(s=\sqrt{x^2+y^2}\)。给出候选电势

\[
\varphi(s,z)
=kq\left[
\frac1{\sqrt{s^2+(z-a)^2}}
-\frac1{\sqrt{s^2+(z+a)^2}}
\right].
\]

\begin{enumerate}
\def\labelenumi{\arabic{enumi}.}
\tightlist
\item
  验证该候选式在物理区域内满足控制方程、边界条件和无穷远条件；
\item
  说明为什么可以由唯一性定理认定它是物理解；
\item
  求导体表面感应电荷密度 \(\sigma(s)\)，并验证总感应电荷。
\end{enumerate}

\par\noindent\begin{minipage}{\linewidth}
\textbf{解答}

\begin{center}
\includegraphics[width=\linewidth,height=.40\textheight,keepaspectratio]{\detokenize{assets/module-02/examples/M2-094.pdf}}
\par\smallskip{\small 图 M2-094：用唯一性定理检验接地平面的候选电势}
\end{center}
\end{minipage}\par

图 M2-094
中平面点到真实源与镜像源的距离相同，所以候选势在整个接地面上为零。在上半空间只保留原来的奇点，两个源的组合又满足无穷远条件；确认唯一性后，从真空侧沿
\(+z\) 求场并积分面电荷。

第一项在 \((0,0,a)\) 处具有真实点电荷 \(q\) 的库仑奇点。第二项的奇点位于
\((0,0,-a)\)，在物理区域 \(z>0\) 之外。因此除真实点电荷外，候选电势在
\(z>0\) 满足

\[
\nabla^2\varphi=0.
\]

在接地平面 \(z=0\) 上，两项距离相等：

\[
\varphi(s,0)=0.
\]

当 \(s^2+z^2\to\infty\) 时，两项都趋于零，而且差值衰减更快，所以

\[
\varphi\to0.
\]

候选式具有正确的真实奇点、满足接地边界和无穷远条件。第二项的虚拟奇点位于物理区域外，不违反物理区域内的泊松方程。因此由第一类外部边值问题的唯一性，该候选式就是
\(z>0\) 中的唯一物理解。

导体外法向取 \(+\mathbf e_z\)。表面电荷密度满足

\[
\sigma=\varepsilon_0E_z(0^+)
=-\varepsilon_0
\left.\frac{\partial\varphi}{\partial z}
\right|_{z=0^+}.
\]

求导得到

\[
\left.\frac{\partial\varphi}{\partial z}
\right|_{z=0^+}
=\frac{2kqa}{(s^2+a^2)^{3/2}}.
\]

所以

\[
\boxed{
\sigma(s)
=-\frac{qa}{2\pi(s^2+a^2)^{3/2}}
}.
\]

总感应电荷为

\[
\begin{aligned}
Q_{\mathrm{\text{感}}}
&=\int_0^{\infty}\sigma(s)\,2\pi s\,ds\\
&=-qa\int_0^{\infty}
\frac{s\,ds}{(s^2+a^2)^{3/2}}\\
&=-q.
\end{aligned}
\]

因此

\[
\boxed{Q_{\mathrm{\text{感}}}=-q}.
\]

\subsubsection{作业五：球内第一类边值问题的直接构造}\label{m2:ux4f5cux4e1aux4e94ux7403ux5185ux7b2cux4e00ux7c7bux8fb9ux503cux95eeux9898ux7684ux76f4ux63a5ux6784ux9020}

\textbf{题号：M2-095}

\textbf{层级：} 复赛　　\textbf{方法：} 边值问题

采用标准球坐标，\(\theta\) 从 \(+z\)
量起，\(R>0\)；边界数据为给定的球面狄利克雷图样，\(V_0,V_1\)
为常量，球内为无源真空且不允许原点奇点。

半径为 \(R\) 的无源球形区域内满足拉普拉斯方程。球面边界电势给定为

\[
\varphi(R,\theta)=V_0+V_1\cos\theta.
\]

\begin{enumerate}
\def\labelenumi{\arabic{enumi}.}
\tightlist
\item
  构造球内有限的电势；
\item
  求球内电场；
\item
  用唯一性定理说明答案的唯一性；
\item
  验证球心电势等于边界电势的球面平均值。
\end{enumerate}

\par\noindent\begin{minipage}{\linewidth}
\textbf{解答}

\begin{center}
\includegraphics[width=\linewidth,height=.40\textheight,keepaspectratio]{\detokenize{assets/module-02/examples/M2-095.pdf}}
\par\smallskip{\small 图 M2-095：球内第一类边值问题的直接构造}
\end{center}
\end{minipage}\par

图 M2-095 的球面图样含常量和一次余弦两个角模。球内正则性分别选择 \(r^0\)
与
\(r^1\)，后一项其实是线性坐标势，因而产生匀强场；球心只保留边界的平均值，额外零边界调和项由唯一性排除。

注意

\[
r\cos\theta=z
\]

是直角坐标的一次函数，满足

\[
\nabla^2z=0.
\]

因此可直接构造

\[
\boxed{
\varphi(r,\theta)
=V_0+V_1\frac rR\cos\theta
}.
\]

它在 \(r=0\) 处有限，在 \(r=R\) 时满足给定边界：

\[
\varphi(R,\theta)=V_0+V_1\cos\theta.
\]

又因

\[
\varphi=V_0+\frac{V_1}{R}z,
\]

所以球内电场是匀强场：

\[
\boxed{
\mathbf E
=-\nabla\varphi
=-\frac{V_1}{R}\mathbf e_z
}.
\]

候选电势在球内调和、在球心有限并满足完整第一类边界条件。由唯一性定理，它是球内唯一解。

球心电势为

\[
\varphi(0)=V_0.
\]

边界球面平均值为

\[
\begin{aligned}
\langle\varphi\rangle_S
&=\frac1{4\pi}
\int_0^{2\pi}\int_0^{\pi}
(V_0+V_1\cos\theta)
\sin\theta\,d\theta\,d\phi\\
&=V_0,
\end{aligned}
\]

因为 \(\cos\theta\) 在球面上的平均值为零。故

\[
\boxed{
\varphi(0)=\langle\varphi\rangle_S
},
\]

与调和函数的平均值性质一致。

\section{专题十七附录一：拉普拉斯方程的常用解法（基础必讲与拔高选讲）}\label{m2:ux4e13ux9898ux5341ux4e03ux9644ux5f55ux4e00ux62c9ux666eux62c9ux65afux65b9ux7a0bux7684ux5e38ux7528ux89e3ux6cd5ux57faux7840ux5fc5ux8bb2ux4e0eux62d4ux9ad8ux9009ux8bb2}

静电边值问题的困难，通常不在于写出

\[
\nabla^2\varphi=0,
\]

而在于同时处理求解区域、源、边界数据、奇点与无穷远条件。方法选择若与这些条件不匹配，即使计算过程很长，也很难得到有效结果。

本附录比较分离变量法、电像法、保角变换法、格林函数法和有限差分法。五种方法分别把边值问题表示为本征模态、辅助源、标准区域、单位点源响应或离散线性方程组。

\subsection{一、先确定边值问题，再选择方法}\label{m2:ux4e00ux5148ux786eux5b9aux8fb9ux503cux95eeux9898ux518dux9009ux62e9ux65b9ux6cd5}

设求解区域为 \(\Omega\)，边界为 \(\partial\Omega\)，外法向单位矢量为
\(\mathbf n\)。静电势满足泊松方程

\[
\boxed{
\nabla^2\varphi=-\frac{\rho}{\varepsilon_0}
}.
\]

只有在所讨论的整个区域内均有 \(\rho=0\) 时，才能写成拉普拉斯方程

\[
\boxed{
\nabla^2\varphi=0
}.
\]

点电荷、线电荷或面电荷造成的奇点不能被一句``区域内无源''掩盖。正确做法是把奇点从光滑求解域中挖去，或直接保留泊松方程的分布形式。

常见边界条件有三类：

\begin{enumerate}
\def\labelenumi{\arabic{enumi}.}
\tightlist
\item
  \textbf{第一类边界条件}：在 \(\partial\Omega\) 上给定 \(\varphi=f\)；
\item
  \textbf{第二类边界条件}：在 \(\partial\Omega\) 上给定
  \(\partial\varphi/\partial n=g\)；
\item
  \textbf{混合边界条件}：边界的不同部分分别给定电势或法向导数。
\end{enumerate}

第一类边值问题在通常条件下有唯一解。纯第二类边值问题的解只确定到一个加法常量，而且数据必须满足兼容条件。对泊松方程积分并使用高斯定理，得到

\[
\boxed{
\oint_{\partial\Omega}g\,dS
=-\frac1{\varepsilon_0}\int_{\Omega}\rho\,dV
}.
\]

若无源区域采用纯第二类边界条件，则必须有
\(\oint_{\partial\Omega}g\,dS=0\)。

\textbf{例 M2-096：零体源与同号法向数据不能相容}

设有界无源真空区域 \(V\)
的整个闭合边界面积非零、分片光滑，法向一律指向区域外。给定纯第二类数据
\(\partial\varphi/\partial n=c\)，\(c\) 为常量。利用区域内
\(\nabla^2\varphi=0\)
求可解的必要相容条件，并解释即使该条件满足，电势为何仍允许整体加常数。

\par\noindent\begin{minipage}{\linewidth}
\textbf{解答}

\begin{center}
\includegraphics[width=\linewidth,height=.40\textheight,keepaspectratio]{\detokenize{assets/module-02/examples/M2-096.pdf}}
\par\smallskip{\small 图 M2-096：零体源与同号法向数据不能相容}
\end{center}
\end{minipage}\par

图 M2-096
的同号边界斜率若为非零常量，其面积分不可能为零。无源条件给出的相容性因此迫使
\(c=0\)；即使这一条件满足，纯第二类数据仍不能固定整体电势零点。

在一个无源封闭区域的整个边界上要求
\(\partial\varphi/\partial n=c\)。兼容条件给出

\[
c\,|\partial\Omega|=0,
\]

故除非
\(c=0\)，否则该边值问题根本无解；这不是求解方法不够强，而是边界数据彼此矛盾。

\par\noindent\begin{minipage}{\linewidth}
\begin{center}
\includegraphics[width=\linewidth,height=.40\textheight,keepaspectratio]{\detokenize{assets/module-02/original/tikz-m2-topic17-v8-11.png}}
\par\smallskip{\small 图示：依据几何、源和边界数据选择拉普拉斯方程求解方法}
\end{center}
\end{minipage}\par

选择方法时可按下表判断：

{\def\LTcaptype{none} % do not increment counter
\begin{longtable}[]{@{}
  >{\raggedright\arraybackslash}p{(\linewidth - 6\tabcolsep) * \real{0.2500}}
  >{\raggedright\arraybackslash}p{(\linewidth - 6\tabcolsep) * \real{0.2500}}
  >{\raggedright\arraybackslash}p{(\linewidth - 6\tabcolsep) * \real{0.2500}}
  >{\raggedright\arraybackslash}p{(\linewidth - 6\tabcolsep) * \real{0.2500}}@{}}
\toprule\noalign{}
\begin{minipage}[b]{\linewidth}\raggedright
方法
\end{minipage} & \begin{minipage}[b]{\linewidth}\raggedright
最有利的结构
\end{minipage} & \begin{minipage}[b]{\linewidth}\raggedright
主要输出
\end{minipage} & \begin{minipage}[b]{\linewidth}\raggedright
首要检查
\end{minipage} \\
\midrule\noalign{}
\endhead
\bottomrule\noalign{}
\endlastfoot
分离变量法 & 边界贴合某套正交坐标曲面 & 本征函数级数 &
是否遗漏零模、奇异模态与收敛条件 \\
电像法 & 少量区域外辅助源可满足高对称边界 & 库仑势的有限或可控叠加 &
像源是否全部位于真实求解域外 \\
保角变换法 & 二维区域可解析地映到标准区域 & 标准区域解的复映射 &
映射是否一一、解析且保持边界对应 \\
格林函数法 & 几何的单位点源响应已知或可构造 & 体源积分与边界积分 &
格林函数的符号、边界类型与法向方向 \\
有限差分法 & 几何复杂，解析展开不再经济 & 网格上的近似解 &
截断误差、残差、网格与外边界收敛 \\
\end{longtable}
}

\subsection{二、分离变量法：把边界数据分解为本征模态}\label{m2:ux4e8cux5206ux79bbux53d8ux91cfux6cd5ux628aux8fb9ux754cux6570ux636eux5206ux89e3ux4e3aux672cux5f81ux6a21ux6001}

\subsubsection{1. 最小推导}\label{m2:ux6700ux5c0fux63a8ux5bfc}

以二维条带 \(0<x<a\)、\(y>0\) 为例。无源区满足

\[
\frac{\partial^2\varphi}{\partial x^2}
+\frac{\partial^2\varphi}{\partial y^2}=0.
\]

先寻找乘积形式的非零特解

\[
\varphi(x,y)=X(x)Y(y).
\]

代入并除以 \(XY\)：

\[
\frac{X''}{X}+\frac{Y''}{Y}=0.
\]

第一项只依赖 \(x\)，第二项只依赖 \(y\)。若等式对区域内所有 \((x,y)\)
成立，两项只能分别等于常量。取分离常数为 \(-k^2\)，得到

\[
X''+k^2X=0,
\qquad
Y''-k^2Y=0.
\]

若两侧边界接地，

\[
X(0)=X(a)=0,
\]

则只有离散本征值

\[
k_n=\frac{n\pi}{a},
\qquad n=1,2,\cdots,
\]

以及本征函数

\[
X_n(x)=\sin\frac{n\pi x}{a}.
\]

再要求 \(y\to\infty\) 时电势有界并趋于零，必须舍去增长指数项，于是

\[
Y_n(y)=\exp\left(-\frac{n\pi y}{a}\right).
\]

线性叠加得到

\[
\boxed{
\varphi(x,y)
=\sum_{n=1}^{\infty}A_n
\sin\frac{n\pi x}{a}
\exp\left(-\frac{n\pi y}{a}\right)
}.
\]

若底边给定 \(\varphi(x,0)=f(x)\)，利用正弦函数的正交性可得

\[
\boxed{
A_n=\frac2a\int_0^a
f(x)\sin\frac{n\pi x}{a}\,dx
}.
\]

分离变量法先求一族乘积特解，再用完备的本征函数系叠加出满足一般边界数据的解；它并不假定任意解本身都是单个乘积。

\par\noindent\begin{minipage}{\linewidth}
\begin{center}
\includegraphics[width=\linewidth,height=.40\textheight,keepaspectratio]{\detokenize{assets/module-02/original/tikz-m2-topic17-v8-12.png}}
\par\smallskip{\small 图示：分离变量法把边界函数分解为正交本征模态}
\end{center}
\end{minipage}\par

\textbf{例 M2-097：边界正弦型只激发一个模态}

在与第三坐标无关的真空半无限条带 \(0<x<a,y>0\)
内求无源电势，\(a>0\)。两侧边 \(x=0,a\) 接地，底边 \(y=0\) 上给定
\(\varphi(x,0)=V_0\sin(\pi x/a)\)，且 \(y\to\infty\)
时电势衰减到零。用分离变量求
\(\varphi(x,y)\)，说明底边数据为何只激发一个横向正弦模态。

\par\noindent\begin{minipage}{\linewidth}
\textbf{解答}

\begin{center}
\includegraphics[width=\linewidth,height=.40\textheight,keepaspectratio]{\detokenize{assets/module-02/examples/M2-097.pdf}}
\par\smallskip{\small 图 M2-097：边界正弦型只激发一个模态}
\end{center}
\end{minipage}\par

图 M2-097
的底边只给出第一正弦模态，侧壁零势使角向基底已被固定。正交投影消去其他系数，纵向只保留衰减分支，故指数长度由同一个本征值
\(\pi/a\) 决定。

若

\[
f(x)=V_0\sin\frac{\pi x}{a},
\]

它本身就是第一本征模态，因此所有 \(n\ne1\) 的系数都为零：

\[
\boxed{
\varphi(x,y)=V_0\sin\frac{\pi x}{a}e^{-\pi y/a}
}.
\]

边界数据越光滑，其高阶模态系数通常衰减得越快；边界若有跳变，则需要更多高阶模态才能描述尖锐变化。

\subsubsection{2.
坐标系由边界几何决定}\label{m2:ux5750ux6807ux7cfbux7531ux8fb9ux754cux51e0ux4f55ux51b3ux5b9a}

分离变量是否高效，取决于边界能否写成某个正交坐标的坐标面：

\begin{itemize}
\tightlist
\item
  矩形、平行平面优先使用直角坐标；
\item
  圆柱、圆环、圆锥角域优先考虑柱坐标；
\item
  球面、同心球壳、圆锥面优先考虑球坐标。
\end{itemize}

\par\noindent\begin{minipage}{\linewidth}
\begin{center}
\includegraphics[width=\linewidth,height=.40\textheight,keepaspectratio]{\detokenize{assets/module-02/original/tikz-m2-topic17-v8-13.png}}
\par\smallskip{\small 图示：直角坐标、柱坐标和球坐标分别适配不同边界曲面}
\end{center}
\end{minipage}\par

柱坐标和球坐标中的变量分离会产生幂函数、贝塞尔函数、勒让德函数或球谐函数；它们是相应几何下拉普拉斯算符的本征函数。

\textbf{例 M2-098：真实解域筛选正幂与负幂}

采用以原点为中心、\(+z\)
为极轴的球坐标，电势轴对称且在无源真空区域满足拉普拉斯方程。一般的每阶勒让德模态含
\(r^{\ell}\) 与 \(r^{-\ell-1}\)
两类径向项。分别考虑包含球心且电势有限的球内域，以及球外域且电势在无穷远趋零两种独立边值问题，说明各自应保留哪些径向项，并写出相应级数形式。

\par\noindent\begin{minipage}{\linewidth}
\textbf{解答}

\begin{center}
\includegraphics[width=\linewidth,height=.40\textheight,keepaspectratio]{\detokenize{assets/module-02/examples/M2-098.pdf}}
\par\smallskip{\small 图 M2-098：真实解域筛选正幂与负幂}
\end{center}
\end{minipage}\par

图 M2-098
的球面把两个不同解域分开：内部要求球心有限，外部要求无穷远衰减。对每个
\(\ell\)
独立筛去不允许的径向分支，再分别匹配球面系数，不能把两种区域的正负幂同时任意叠加。

球内轴对称、在球心有限的调和势只能保留

\[
\varphi(r,\theta)
=\sum_{\ell=0}^{\infty}A_{\ell} r^{\ell}
P_{\ell}(\cos\theta),
\]

而外部无穷远衰减的解只能保留

\[
\varphi(r,\theta)
=\sum_{\ell=0}^{\infty}
\frac{B_{\ell}}{r^{\ell+1}}P_{\ell}(\cos\theta).
\]

``保留正幂还是负幂''取决于解域和允许的奇点，不能脱离物理区域机械记忆。

\subsubsection{3.
分离变量法的失败边界}\label{m2:ux5206ux79bbux53d8ux91cfux6cd5ux7684ux5931ux8d25ux8fb9ux754c}

以下情形即使形式上仍可展开，也常常不再省力：

\begin{enumerate}
\def\labelenumi{\arabic{enumi}.}
\tightlist
\item
  边界斜穿所选坐标面，使不同模态在边界条件中耦合；
\item
  区域包含多个不共心的曲面，找不到同时贴合全部边界的简单坐标系；
\item
  边界有尖角、狭缝或细小尺度，需要大量高阶模态；
\item
  只写乘积特解，却没有证明所用本征函数系足以展开边界数据；
\item
  忽略 \(k=0\) 的零模，导致常量项、对数项或线性项丢失。
\end{enumerate}

\subsection{三、电像法：用区域外辅助源实现边界条件}\label{m2:ux4e09ux7535ux50cfux6cd5ux7528ux533aux57dfux5916ux8f85ux52a9ux6e90ux5b9eux73b0ux8fb9ux754cux6761ux4ef6}

电像法适用于少量虚拟点源、线源或多极源能够精确满足边界条件的情形。它并没有改变真实区域内的控制方程，而是用辅助源构造一个容易验证的候选解，再借助唯一性定理确认。

\par\noindent\begin{minipage}{\linewidth}
\begin{center}
\includegraphics[width=\linewidth,height=.40\textheight,keepaspectratio]{\detokenize{assets/module-02/original/tikz-m2-topic17-v8-14.png}}
\par\smallskip{\small 图示：电像法中的真实源、辅助区域与三项有效性检查}
\end{center}
\end{minipage}\par

一个像源构造必须同时满足：

\begin{enumerate}
\def\labelenumi{\arabic{enumi}.}
\tightlist
\item
  所有像源均位于真实求解区域之外；
\item
  真实区域内原有源、奇点及其强度保持不变；
\item
  所有导体边界条件和无穷远条件同时满足。
\end{enumerate}

满足三项后，才可以用唯一性定理把候选势认定为真实区域内的解。像电荷不是真实导体内部的电荷分布，也不能用它直接描述导体内部的物理场。

\subsubsection{例：接地无限平面（M2-099）}\label{m2:ux4f8bux63a5ux5730ux65e0ux9650ux5e73ux9762m2-099}

真空半空间 \(z>0\) 中一个电荷量为 \(q\) 的点电荷固定于
\((0,0,d)\)，\(d>0\)；无限大平面 \(z=0\)
为接地导体边界，电势在无穷远衰减。构造像源并求真实半空间内的电势，逐一检验真实点源奇性、无源处的方程、整张平面的零势和无穷远条件，以唯一性定理说明构造成立。

\par\noindent\begin{minipage}{\linewidth}
\textbf{解答}

\begin{center}
\includegraphics[width=\linewidth,height=.40\textheight,keepaspectratio]{\detokenize{assets/module-02/examples/M2-099.pdf}}
\par\smallskip{\small 图 M2-099：平面像电荷逐点抵消边界电势}
\end{center}
\end{minipage}\par

图 M2-099 的灰色区域容纳负的镜像电荷，真实区域在平面上方。平面任一点
\(P\)
的两段距离相同，直接验证整段边界零势；域内奇点和远端衰减也一致，因而候选势由唯一性确定。

真实点电荷 \(q\) 位于 \((0,0,d)\)，物理区域为 \(z>0\)，平面 \(z=0\)
接地。在辅助区域 \(z<0\) 放置像电荷 \(-q\)，位置为
\((0,0,-d)\)。候选势为

\[
\varphi(\mathbf r)
=\frac{q}{4\pi\varepsilon_0}
\left(
\frac1{|\mathbf r-d\mathbf e_z|}
-\frac1{|\mathbf r+d\mathbf e_z|}
\right).
\]

在 \(z=0\) 上，两段距离相等，故 \(\varphi=0\)；在 \(z>0\)
内只有真实电荷处保留正确奇点；无穷远处电势趋于零。因此该势是物理区域内的唯一解。

电像法的强项是把一个边值问题化成库仑势叠加，弱项是适用几何非常有限。接地平面、球面以及少数可经反演或重复反射处理的几何最典型；一般不规则导体不能凭视觉对称随意放置像电荷。具体计算将在专题二十中展开。

\subsection{四、保角变换法：在二维中改造边界几何（拔高选讲）}\label{m2:ux56dbux4fddux89d2ux53d8ux6362ux6cd5ux5728ux4e8cux7ef4ux4e2dux6539ux9020ux8fb9ux754cux51e0ux4f55ux62d4ux9ad8ux9009ux8bb2}

设问题沿 \(z\) 方向平移不变，电势只依赖 \((x,y)\)。令

\[
z=x+iy,
\qquad
w=f(z)=u+iv,
\]

其中 \(f\) 在所讨论区域内解析且 \(f'(z)\ne0\)。若 \(U(u,v)\) 在 \(w\)
平面调和，则复合函数

\[
\varphi(x,y)=U(u(x,y),v(x,y))
\]

满足

\[
\boxed{
\nabla_z^2\varphi
=|f'(z)|^2\nabla_w^2U
=0
}.
\]

因此，可以先把原区域映到半平面、条带、圆盘等标准区域，在标准区域求解，再把结果映回。

\par\noindent\begin{minipage}{\linewidth}
\begin{center}
\includegraphics[width=\linewidth,height=.40\textheight,keepaspectratio]{\detokenize{assets/module-02/original/tikz-m2-topic17-v8-15.png}}
\par\smallskip{\small 图示：保角变换把二维劈形区域映到半平面或条带}
\end{center}
\end{minipage}\par

保角变换保持局部夹角，但不保持长度和电场大小。由复合函数求导可知，映回原平面后电场强度会乘上与
\(|f'(z)|\)
有关的尺度因子；映射导数很大或趋于零处，场强可能显著增强或出现奇异行为。

\subsubsection{例：两边电势固定的无限劈形（M2-100）}\label{m2:ux4f8bux4e24ux8fb9ux7535ux52bfux56faux5b9aux7684ux65e0ux9650ux5288ux5f62m2-100}

在二维极坐标中，真空无源域为无限劈形 \(r>0,0<\theta<\alpha\)，其中
\(0<\alpha<2\pi\)。两条径向电极分别保持
\(\varphi(r,0)=0\)、\(\varphi(r,\alpha)=V_0\)。采用不引入额外径向尺度、势只依赖角度的局部模型，求电势与电场，说明顶点附近的场强行为；写出可将此劈形映至半平面的幂映射。

\par\noindent\begin{minipage}{\linewidth}
\textbf{解答}

\begin{center}
\includegraphics[width=\linewidth,height=.40\textheight,keepaspectratio]{\detokenize{assets/module-02/examples/M2-100.pdf}}
\par\smallskip{\small 图 M2-100：不同电势的劈形边界产生径向奇性}
\end{center}
\end{minipage}\par

图 M2-100
的两条径向边界具有不同电势，红弧沿电势减小的方向。没有额外径向尺度时取势只依赖极角，积分两次角方程得到线性角势；还原场时必须保留极坐标的
\(1/r\) 因子。

考虑二维劈形

\[
0<\theta<\alpha,
\]

两条边分别保持

\[
\varphi(r,0)=0,
\qquad
\varphi(r,\alpha)=V_0.
\]

由于边界数据与 \(r\) 无关，可尝试
\(\varphi=\varphi(\theta)\)。极坐标拉普拉斯方程化为

\[
\frac1{r^2}\frac{d^2\varphi}{d\theta^2}=0,
\]

故

\[
\boxed{
\varphi(r,\theta)=\frac{V_0}{\alpha}\theta
},
\]

相应电场为

\[
\boxed{
\mathbf E
=-\frac1r\frac{\partial\varphi}{\partial\theta}\mathbf e_{\theta}
=-\frac{V_0}{\alpha r}\mathbf e_{\theta}
}.
\]

映射 \(w=z^{\pi/\alpha}\) 可把劈形送到半平面。场强随 \(1/r\)
增大，说明顶点处存在奇异性；其根源是两条边在顶点处要求不同电势，边界数据在那里不连续。

保角变换只直接保护二维拉普拉斯方程的结构。三维一般变换不具有这一性质，不能把``保角''当作任意维度下的边界变形工具。

\subsection{五、格林函数法：把几何与源分开（拔高选讲）}\label{m2:ux4e94ux683cux6797ux51fdux6570ux6cd5ux628aux51e0ux4f55ux4e0eux6e90ux5206ux5f00ux62d4ux9ad8ux9009ux8bb2}

格林函数是在给定区域和给定齐次边界条件下，单位点源产生的响应。对第一类边值问题，定义狄利克雷格林函数
\(G_D(\mathbf r,\mathbf r')\)：

\[
\boxed{
\nabla'^2G_D(\mathbf r,\mathbf r')
=-\delta^{(3)}(\mathbf r-\mathbf r'),
\qquad
G_D|_{\partial\Omega}=0
}.
\]

这里微分作用于源点坐标 \(\mathbf r'\)，而 \(\mathbf r\) 是固定的观测点。

\par\noindent\begin{minipage}{\linewidth}
\begin{center}
\includegraphics[width=\linewidth,height=.40\textheight,keepaspectratio]{\detokenize{assets/module-02/original/tikz-m2-topic17-v8-16.png}}
\par\smallskip{\small 图示：狄利克雷格林函数是指定边界下的单位点源响应}
\end{center}
\end{minipage}\par

\subsubsection{1.
表示公式的最小推导}\label{m2:ux8868ux793aux516cux5f0fux7684ux6700ux5c0fux63a8ux5bfc}

对 \(\varphi(\mathbf r')\) 与 \(G_D(\mathbf r,\mathbf r')\)
使用格林第二恒等式：

\[
\int_{\Omega}
\left(
\varphi\nabla'^2G_D
-G_D\nabla'^2\varphi
\right)dV'
=
\oint_{\partial\Omega}
\left(
\varphi\frac{\partial G_D}{\partial n'}
-G_D\frac{\partial\varphi}{\partial n'}
\right)dS'.
\]

代入

\[
\nabla'^2\varphi=-\frac{\rho}{\varepsilon_0},
\qquad
\nabla'^2G_D=-\delta^{(3)}(\mathbf r-\mathbf r'),
\]

并利用 \(G_D=0\) 以及边界上 \(\varphi=f\)，可得

\[
-\varphi(\mathbf r)
+\frac1{\varepsilon_0}\int_{\Omega} G_D\rho\,dV'
=\oint_{\partial\Omega}
f(\mathbf r')\frac{\partial G_D}{\partial n'}\,dS'.
\]

所以

\[
\boxed{
\varphi(\mathbf r)
=\frac1{\varepsilon_0}\int_{\Omega}
G_D(\mathbf r,\mathbf r')\rho(\mathbf r')\,dV'
-\oint_{\partial\Omega}
f(\mathbf r')
\frac{\partial G_D}{\partial n'}\,dS'
}.
\]

第一项叠加区域内的体源，第二项把非齐次边界数据传递到内部。格林函数一旦构造完成，同一几何下更换
\(\rho\) 或 \(f\) 时无需重新求解偏微分方程。

\subsubsection{2.
自由空间格林函数}\label{m2:ux81eaux7531ux7a7aux95f4ux683cux6797ux51fdux6570}

三维自由空间中

\[
\boxed{
G_0(\mathbf r,\mathbf r')
=\frac1{4\pi|\mathbf r-\mathbf r'|}
},
\]

满足

\[
\nabla'^2G_0=-\delta^{(3)}(\mathbf r-\mathbf r').
\]

若电荷局域分布且无穷远边界项消失，则

\[
\varphi(\mathbf r)
=\frac1{4\pi\varepsilon_0}
\int
\frac{\rho(\mathbf r')}{|\mathbf r-\mathbf r'|}
\,dV',
\]

这正是库仑势的连续叠加形式。

\textbf{例 M2-101：点源响应从三维源积分中筛出}

真空自由空间中，电荷密度为
\(\rho(\mathbf r')=q\delta^{(3)}(\mathbf r'-\mathbf r_0)\)，即点电荷固定在给定位置
\(\mathbf r_0\)。取无穷远零势，场点
\(\mathbf r\ne\mathbf r_0\)。将该密度代入自由空间电势积分
\(\varphi(\mathbf r)=\int\rho(\mathbf r')/[4\pi\varepsilon_0|\mathbf r-\mathbf r'|]\,dV'\)，利用三维筛选性质求势，区分积分源坐标与场点坐标。

\par\noindent\begin{minipage}{\linewidth}
\textbf{解答}

\begin{center}
\includegraphics[width=\linewidth,height=.40\textheight,keepaspectratio]{\detokenize{assets/module-02/examples/M2-101.pdf}}
\par\smallskip{\small 图 M2-101：点源响应从三维源积分中筛出}
\end{center}
\end{minipage}\par

图 M2-101 用有向线段连接源点 \(\mathbf r_0\) 和场点
\(\mathbf r\)。三维δ积分只筛掉源变量，保留下来的核依赖两点之差；边界存在时还应再构造满足该边界的响应，不能把自由空间核直接用于任意容器。

对点电荷

\[
\rho(\mathbf r')=q\delta^{(3)}(\mathbf r'-\mathbf r_0),
\]

积分立即给出

\[
\varphi(\mathbf r)
=\frac{q}{4\pi\varepsilon_0|\mathbf r-\mathbf r_0|}.
\]

电像法可视为某些高对称边界下构造格林函数的捷径。例如接地平面的格林函数就是``真实单位源的自由空间响应减去镜像单位源的自由空间响应''。

对纯第二类边界条件，格林函数必须改用相应的诺伊曼型构造，并处理兼容条件与加法常量。不能把狄利克雷表示式直接换一个边界符号继续使用。

\subsection{六、有限差分法：把连续问题变成离散方程（拔高选讲）}\label{m2:ux516dux6709ux9650ux5deeux5206ux6cd5ux628aux8fdeux7eedux95eeux9898ux53d8ux6210ux79bbux6563ux65b9ux7a0bux62d4ux9ad8ux9009ux8bb2}

解析方法失效时，可以在区域内建立网格。以间距均为 \(h\)
的二维正方形网格为例，中心点记为 \(C\)，东西南北相邻点记为 \(E,W,S,N\)。

在中心点作泰勒展开：

\[
\varphi_E+\varphi_W
=2\varphi_C
+h^2\varphi_{xx}
+\frac{h^4}{12}\varphi_{xxxx}
+O(h^6),
\]

\[
\varphi_N+\varphi_S
=2\varphi_C
+h^2\varphi_{yy}
+\frac{h^4}{12}\varphi_{yyyy}
+O(h^6).
\]

两式相加并除以 \(h^2\)：

\[
\boxed{
\frac{
\varphi_E+\varphi_W+\varphi_N+\varphi_S-4\varphi_C
}{h^2}
=\nabla^2\varphi
+\frac{h^2}{12}
(\varphi_{xxxx}+\varphi_{yyyy})
+O(h^4)
}.
\]

因此五点差分对拉普拉斯算符的局部截断误差为 \(O(h^2)\)。

\par\noindent\begin{minipage}{\linewidth}
\begin{center}
\includegraphics[width=\linewidth,height=.40\textheight,keepaspectratio]{\detokenize{assets/module-02/original/tikz-m2-topic17-v8-17.png}}
\par\smallskip{\small 图示：二维五点差分格式及其泊松方程更新关系}
\end{center}
\end{minipage}\par

代入泊松方程 \(\nabla^2\varphi=-\rho/\varepsilon_0\)，得到中心点关系

\[
\boxed{
\varphi_C
=\frac{
\varphi_E+\varphi_W+\varphi_N+\varphi_S
}{4}
+\frac{\rho_Ch^2}{4\varepsilon_0}
+O(h^4)
}.
\]

数值计算中把余项舍去，才得到要实际求解的离散方程。

\subsubsection{1.
四点平均不是连续平均值定理本身}\label{m2:ux56dbux70b9ux5e73ux5747ux4e0dux662fux8fdeux7eedux5e73ux5747ux503cux5b9aux7406ux672cux8eab}

当 \(\rho_C=0\) 时，离散方程写成

\[
\varphi_C
=\frac{
\varphi_E+\varphi_W+\varphi_N+\varphi_S
}{4}.
\]

这对\textbf{离散差分解}是精确的代数关系，对一般连续调和函数则只是带截断误差的近似。连续调和函数真正的平均值定理是：中心值严格等于任意完全包含于无源区域内的圆周平均值；三维时则等于球面平均值。

二维形式为

\[
\boxed{
\varphi(\mathbf r_0)
=\frac1{2\pi}
\int_0^{2\pi}
\varphi(\mathbf r_0+R\mathbf e_r)\,d\theta
}.
\]

\textbf{例 M2-102：五点平均是差分近似}

在二维均匀正方形网格上取间距 \(h>0\)，中心点为原点，四个最近邻为
\((\pm h,0),(0,\pm h)\)。给定调和多项式
\(\varphi=x^4-6x^2y^2+y^4\)。比较中心势与四邻点电势的算术平均值，判断连续调和函数是否严格满足这一离散平均关系，并说明五点差分式的近似性质。

\par\noindent\begin{minipage}{\linewidth}
\textbf{解答}

\begin{center}
\includegraphics[width=\linewidth,height=.40\textheight,keepaspectratio]{\detokenize{assets/module-02/examples/M2-102.pdf}}
\par\smallskip{\small 图 M2-102：五点平均是差分近似}
\end{center}
\end{minipage}\par

图 M2-102 的四个邻点都位于坐标轴上，函数值同为
\(h^4\)，红色中心却为零。四点平均与中心的差恰好是有限步长误差；这不否定连续调和函数的球面平均性质，而是说明有限差分平均只在离散方程中精确成立。

取调和函数

\[
\varphi(x,y)=x^4-6x^2y^2+y^4.
\]

在原点，四个轴向网格点 \((\pm h,0),(0,\pm h)\) 的函数值都等于
\(h^4\)，四点平均为 \(h^4\)，却有
\(\varphi(0,0)=0\)。因此``连续调和函数严格等于四个相邻点的平均值''是错误命题。

\subsubsection{2.
迭代与误差检查}\label{m2:ux8fedux4ee3ux4e0eux8befux5deeux68c0ux67e5}

将边界网格点固定后，内部节点组成线性方程组。常见迭代方式包括：

\begin{itemize}
\tightlist
\item
  雅可比迭代：新一轮全部使用上一轮邻点值；
\item
  高斯---赛德尔迭代：一旦得到新值，后续节点立即使用；
\item
  超松弛迭代：在更新量上乘适当松弛因子以加速收敛。
\end{itemize}

无源五点格式的离散残差可定义为

\[
R_{i,j}
=\frac{
\varphi_{i+1,j}+\varphi_{i-1,j}
+\varphi_{i,j+1}+\varphi_{i,j-1}
-4\varphi_{i,j}
}{h^2}.
\]

泊松问题则应检查

\[
R_{i,j}
=\frac{
\varphi_{i+1,j}+\varphi_{i-1,j}
+\varphi_{i,j+1}+\varphi_{i,j-1}
-4\varphi_{i,j}
}{h^2}
+\frac{\rho_{i,j}}{\varepsilon_0}.
\]

残差小只说明离散方程求得较准，不保证离散方程已经足够接近连续问题。可靠结果至少还要检查：

\begin{enumerate}
\def\labelenumi{\arabic{enumi}.}
\tightlist
\item
  网格加密后结果是否稳定；
\item
  曲边界的离散近似是否收敛；
\item
  人工截断的远边界向外移动后结果是否稳定；
\item
  电势梯度、电荷密度等派生量是否也收敛；
\item
  是否满足对称性、量纲、能量正定性和整体通量关系。
\end{enumerate}

\subsection{七、方法的配合与结果检查}\label{m2:ux4e03ux65b9ux6cd5ux7684ux914dux5408ux4e0eux7ed3ux679cux68c0ux67e5}

实际问题常常需要合用多种方法：

\begin{itemize}
\tightlist
\item
  先用唯一性定理确定需要检查的条件，再用电像法构造候选解；
\item
  先用保角变换化简二维边界，再在标准区域使用电像法；
\item
  用分离变量法构造某个几何的格林函数，再对任意源积分；
\item
  用解析解处理奇点附近的主导行为，用数值法计算其余光滑部分；
\item
  用解析极限、对称性和守恒关系校验数值结果。
\end{itemize}

\par\noindent\begin{minipage}{\linewidth}
\begin{center}
\includegraphics[width=\linewidth,height=.40\textheight,keepaspectratio]{\detokenize{assets/module-02/original/tikz-m2-topic17-v8-18.png}}
\par\smallskip{\small 图示：解析法、构造法和数值法的配合与检查}
\end{center}
\end{minipage}\par

无论采用哪种方法，最终都应逐项核对：

\begin{enumerate}
\def\labelenumi{\arabic{enumi}.}
\tightlist
\item
  \textbf{方程}：在每个光滑子区域内是否满足拉普拉斯方程或泊松方程；
\item
  \textbf{真实源}：点源、线源、面源的奇点强度和位置是否正确；
\item
  \textbf{边界}：电势、法向导数、界面跃变与无穷远条件是否全部满足；
\item
  \textbf{定义域}：候选解是否在真实区域内引入了额外奇点；
\item
  \textbf{唯一性}：当前边界类型是否足以唯一确定解；
\item
  \textbf{极限}：参数趋于零、无穷大或恢复高对称情形时是否回到已知结果；
\item
  \textbf{数值误差}：若为近似解，误差指标是否随网格或截断参数改善。
\end{enumerate}

\subsection{八、练习与解答}\label{m2:ux516bux7ec3ux4e60ux4e0eux89e3ux7b54-4}

\subsubsection{作业一：方法选择}\label{m2:ux4f5cux4e1aux4e00ux65b9ux6cd5ux9009ux62e9}

\textbf{题号：M2-103}

\textbf{层级：} 基础　　\textbf{方法：} 边值问题

以下各题都要求建立求解思路，不要求算出未知函数的数值解。对每一种静电边值问题，选择较合适的首选方法并说明依据：

\begin{enumerate}
\def\labelenumi{(\arabic{enumi})}
\item
  二维矩形 \(0<x<a,0<y<b\) 内无电荷，四边电势函数已给定。
\item
  真空中一个点电荷位于无限大接地导体平面一侧，要求该半空间的势。
\item
  二维尖角电极形成的无源劈形，其边界可由幂映射送到半平面。
\item
  一个固定几何容器的边界条件不变，但需要反复计算不同已知体电荷分布的势。
\item
  导体边界高度不规则，缺少分离变量或像源构造所需的明显对称性。
\end{enumerate}

\par\noindent\begin{minipage}{\linewidth}
\textbf{解答}

\begin{center}
\includegraphics[width=\linewidth,height=.40\textheight,keepaspectratio]{\detokenize{assets/module-02/examples/M2-103.pdf}}
\par\smallskip{\small 图 M2-103：方法选择}
\end{center}
\end{minipage}\par

图 M2-103
把边界几何与方法的适用条件对应起来。矩形适合分离变量，平面对称点源可用电像，二维角域可作保角映射；重复换源时优先构造格林响应，无明显结构的边界则进入数值离散并作收敛检查。

\begin{enumerate}
\def\labelenumi{\arabic{enumi}.}
\tightlist
\item
  首选分离变量法。矩形边界贴合直角坐标面，正弦、余弦本征函数可展开一般边界数据。
\item
  首选电像法。一个关于平面对称、异号的像电荷即可精确满足接地条件。
\item
  首选保角变换法。问题为二维，且已知解析映射能把复杂边界变成标准边界。
\item
  首选格林函数法。先求该固定容器的格林函数，之后更换 \(\rho\)
  时只需重新计算体积分。
\item
  首选数值方法。若解析坐标、像源和二维映射均无明显优势，离散化通常更直接。
\end{enumerate}

这些是``首选方法''，不是唯一方法。例如接地平面的格林函数本身就可由像电荷构造；矩形区域也可以用数值法求解。

\subsubsection{作业二：条带中的单模边界}\label{m2:ux4f5cux4e1aux4e8cux6761ux5e26ux4e2dux7684ux5355ux6a21ux8fb9ux754c}

\textbf{题号：M2-104}

\textbf{层级：} 基础　　\textbf{方法：} 边值问题

\(a>0\)，求解域为与第三坐标无关的无源真空条带，\(V_0\)
为常量；要求解在内部有限，并同时满足两侧、底边及远端条件。

在区域 \(0<x<a\)、\(y>0\) 中求解

\[
\nabla^2\varphi=0,
\]

边界条件为

\[
\varphi(0,y)=\varphi(a,y)=0,
\]

\[
\varphi(x,0)=V_0\sin\frac{\pi x}{a},
\qquad
\lim_{y\to\infty}\varphi(x,y)=0.
\]

并直接验证所得结果。

\par\noindent\begin{minipage}{\linewidth}
\textbf{解答}

\begin{center}
\includegraphics[width=\linewidth,height=.40\textheight,keepaspectratio]{\detokenize{assets/module-02/examples/M2-104.pdf}}
\par\smallskip{\small 图 M2-104：条带中的单模边界}
\end{center}
\end{minipage}\par

图 M2-104
的侧边和远端条件分别选出正弦横向基底与负指数纵向分支。底边只含一个本征模，投影系数直接读取
\(V_0\)；验证时域内二阶导数相消，两侧正弦为零，底边和无穷远也逐项满足。

两侧接地给出本征函数 \(\sin(n\pi x/a)\)，无穷远衰减给出因子
\(e^{-n\pi y/a}\)。底边只含 \(n=1\) 模态，因此

\[
\boxed{
\varphi(x,y)
=V_0\sin\frac{\pi x}{a}
e^{-\pi y/a}
}.
\]

记 \(k=\pi/a\)，则

\[
\frac{\partial^2\varphi}{\partial x^2}
=-k^2\varphi,
\qquad
\frac{\partial^2\varphi}{\partial y^2}
=k^2\varphi,
\]

所以 \(\nabla^2\varphi=0\)。在 \(x=0,a\) 时正弦因子为零；在 \(y=0\)
时恢复给定底边电势；当 \(y\to\infty\)
时指数因子趋于零。方程与全部边界条件均满足，由唯一性定理解唯一。

\subsubsection{作业三：劈形电容场}\label{m2:ux4f5cux4e1aux4e09ux5288ux5f62ux7535ux5bb9ux573a}

\textbf{题号：M2-105}

\textbf{层级：} 基础　　\textbf{方法：} 边值问题

采用二维极坐标 \((r,\theta)\)，\(r>0\)，\(0<\alpha<2\pi\)，\(\alpha\)
用弧度表示，角度从零势电极起沿域内增大方向量取。只讨论给定的无额外径向尺度局部模型。

二维无源区域为 \(0<\theta<\alpha\)。两条径向导体边界分别保持电势 \(0\)
和 \(V_0\)，并假定远离顶点的局部模型不含额外径向尺度。

\begin{enumerate}
\def\labelenumi{\arabic{enumi}.}
\tightlist
\item
  求电势与电场；
\item
  说明顶点附近场强的行为；
\item
  写出可把劈形映到半平面的幂函数。
\end{enumerate}

\par\noindent\begin{minipage}{\linewidth}
\textbf{解答}

\begin{center}
\includegraphics[width=\linewidth,height=.40\textheight,keepaspectratio]{\detokenize{assets/module-02/examples/M2-105.pdf}}
\par\smallskip{\small 图 M2-105：劈形电容场}
\end{center}
\end{minipage}\par

图 M2-105
的两条径向导体保持不同电势，局部模型没有指定新的径向边界。先解只含角度的势，再通过
\(-r^{-1}\partial_{\theta}\varphi\) 得场；顶点的 \(1/r\)
奇性来自边界电势在那里发生跳变。

边界数据与 \(r\) 无关，可取 \(\varphi=\varphi(\theta)\)。由

\[
\nabla^2\varphi
=\frac1r\frac{\partial}{\partial r}
\left(r\frac{\partial\varphi}{\partial r}\right)
+\frac1{r^2}\frac{\partial^2\varphi}{\partial\theta^2}=0
\]

得到

\[
\frac{d^2\varphi}{d\theta^2}=0.
\]

故 \(\varphi=A\theta+B\)。代入两条边界条件：

\[
B=0,
\qquad
A=\frac{V_0}{\alpha}.
\]

因此

\[
\boxed{
\varphi=\frac{V_0}{\alpha}\theta
},
\]

\[
\boxed{
\mathbf E
=-\nabla\varphi
=-\frac{V_0}{\alpha r}\mathbf e_{\theta}
}.
\]

场强大小为 \(|V_0|/(\alpha r)\)。当 \(V_0\ne0\) 时，\(r\to0\) 的场强按
\(1/r\) 发散，与顶点处边界电势跳变相联系；\(V_0=0\)
时全域电势、电场均为零。实际电极有有限曲率半径时，理想尖角模型在足够小尺度失效。

所需映射可取

\[
\boxed{
w=z^{\pi/\alpha}
}.
\]

因为 \(0<\arg z<\alpha\) 被映为 \(0<\arg w<\pi\)，即上半平面。

\subsubsection{作业四：自由空间格林函数的归一化}\label{m2:ux4f5cux4e1aux56dbux81eaux7531ux7a7aux95f4ux683cux6797ux51fdux6570ux7684ux5f52ux4e00ux5316}

\textbf{题号：M2-106}

\textbf{层级：} 决赛　　\textbf{方法：} 边值问题

这是三维真空自由空间问题，格林函数采用 \(\nabla'^2G_0=-\delta^{(3)}\)
的归一化，源点奇性按分布理解，要求势在无穷远衰减。

证明

\[
G_0(\mathbf r,\mathbf r')
=\frac1{4\pi|\mathbf r-\mathbf r'|}
\]

满足

\[
\nabla'^2G_0=-\delta^{(3)}(\mathbf r-\mathbf r'),
\]

并由此求点电荷 \(q\) 位于 \(\mathbf r_0\) 时的电势。

\par\noindent\begin{minipage}{\linewidth}
\textbf{解答}

\begin{center}
\includegraphics[width=\linewidth,height=.40\textheight,keepaspectratio]{\detokenize{assets/module-02/examples/M2-106.pdf}}
\par\smallskip{\small 图 M2-106：自由空间格林函数的归一化}
\end{center}
\end{minipage}\par

图 M2-106 的小球包围核的奇点，向外法向下 \(\nabla^{\prime}G_0\) 的通量为
\(-1\)。球外普通拉普拉斯为零，有限通量只集中于源点，因此得到负的δ归一化；将核乘
\(q/\varepsilon_0\) 才还原点电荷势。

令

\[
R=|\mathbf r-\mathbf r'|.
\]

在 \(R\ne0\) 处，直接计算有
\(\nabla'^2(1/R)=0\)。为确定奇点强度，以源点为中心取半径 \(a\)
的小球。对源点坐标 \(\mathbf r'\) 而言，\(R\) 沿小球外法向增大，故

\[
\frac{\partial G_0}{\partial n'}
=-\frac1{4\pi a^2}.
\]

于是

\[
\oint_{S_a}
\frac{\partial G_0}{\partial n'}\,dS'
=-\frac1{4\pi a^2}\cdot4\pi a^2
=-1.
\]

由散度定理，这说明 \(\nabla'^2G_0\) 在分布意义下的积分为 \(-1\)，因此

\[
\nabla'^2G_0
=-\delta^{(3)}(\mathbf r-\mathbf r').
\]

对

\[
\rho(\mathbf r')
=q\delta^{(3)}(\mathbf r'-\mathbf r_0),
\]

自由空间表示式给出

\[
\begin{aligned}
\varphi(\mathbf r)
&=\frac1{\varepsilon_0}
\int G_0(\mathbf r,\mathbf r')
q\delta^{(3)}(\mathbf r'-\mathbf r_0)\,dV'\\
&=\boxed{
\frac{q}{4\pi\varepsilon_0
|\mathbf r-\mathbf r_0|}
}.
\end{aligned}
\]

\subsubsection{作业五：泊松方程的五点差分}\label{m2:ux4f5cux4e1aux4e94ux6ccaux677eux65b9ux7a0bux7684ux4e94ux70b9ux5deeux5206}

\textbf{题号：M2-107}

\textbf{层级：} 基础　　\textbf{方法：} 边值问题

采用与第三坐标无关的真空泊松方程
\(\partial_x^2\varphi+\partial_y^2\varphi=-\rho/\varepsilon_0\)。四个邻点分别距中心
\(h>0\)，只能在场足够光滑的网格区域使用普通五点近似。

在二维均匀正方形网格上，设中心点附近电荷密度可近似为常量
\(\rho_C\)，网格间距为 \(h\)。

\begin{enumerate}
\def\labelenumi{\arabic{enumi}.}
\tightlist
\item
  写出中心点电势的五点差分更新式；
\item
  若四个邻点电势均为零且 \(\rho_C>0\)，判断中心点电势的符号；
\item
  说明为什么迭代残差趋于零后仍需加密网格。
\end{enumerate}

\par\noindent\begin{minipage}{\linewidth}
\textbf{解答}

\begin{center}
\includegraphics[width=\linewidth,height=.40\textheight,keepaspectratio]{\detokenize{assets/module-02/examples/M2-107.pdf}}
\par\smallskip{\small 图 M2-107：泊松方程的五点差分}
\end{center}
\end{minipage}\par

图 M2-107
的四个邻点和中心构成同一均匀五点模板。把二阶差分代入泊松方程后，正密度项使中心势高于四邻点平均；迭代残差只测离散方程是否解准，改变网格尺度才检验连续问题的近似误差。

五点差分近似为

\[
\frac{
\varphi_E+\varphi_W+\varphi_N+\varphi_S-4\varphi_C
}{h^2}
\approx-\frac{\rho_C}{\varepsilon_0}.
\]

整理得

\[
\boxed{
\varphi_C
\approx
\frac{
\varphi_E+\varphi_W+\varphi_N+\varphi_S
}{4}
+\frac{\rho_Ch^2}{4\varepsilon_0}
}.
\]

若四个邻点均为零且 \(\rho_C>0\)，则

\[
\boxed{
\varphi_C\approx
\frac{\rho_Ch^2}{4\varepsilon_0}>0
}.
\]

这与正电荷倾向于产生正电势一致。

残差趋于零只表示当前网格上的离散线性方程已经解得充分准确；五点格式相对于连续泊松方程仍有
\(O(h^2)\) 的截断误差，曲边界和人工远边界还会带来额外误差。因此必须比较
\(h,h/2\) 等不同网格下的结果，确认目标物理量随网格加密趋于稳定。

处理边值问题时，先辨认几何和边界数据，再选择求解方法并检查结果。专题十八讨论直角坐标和球坐标，专题十九转入柱坐标，重点比较本征模态、端点条件和级数筛选。

\section{专题十八：直角坐标和球坐标中的分离变量法（基础必讲与拔高选讲）}\label{m2:ux4e13ux9898ux5341ux516bux76f4ux89d2ux5750ux6807ux548cux7403ux5750ux6807ux4e2dux7684ux5206ux79bbux53d8ux91cfux6cd5ux57faux7840ux5fc5ux8bb2ux4e0eux62d4ux9ad8ux9009ux8bb2}

\begin{quote}
阅读准备：二阶方程、端点条件与正交展开 令 \(\varphi=X(x)Y(y)\)
后，每个因子满足二阶常微分方程，因此各有两个线性独立解；有限性、周期性和边界值负责筛掉不允许的解，并把分离常数限制为离散值。正交性保证不同模态的系数可以分别求出。

\textbf{例 M2-108：正交投影只留下匹配的本征模}
\end{quote}

在有限区间 \(0<y<b\)，\(b>0\)，选用满足两端齐次边界的正弦基
\(Y_n(y)=\sin(n\pi y/b)\)，\(n=1,2,\ldots\)。将给定边界函数
\(f(y)=\sin(2\pi y/b)\) 展成
\(\sum_{n\ge1}c_nY_n(y)\)，利用正交投影求非零系数，并说明为什么其他阶不被激发。

\par\noindent\begin{minipage}{\linewidth}
\textbf{解答}

\begin{center}
\includegraphics[width=\linewidth,height=.40\textheight,keepaspectratio]{\detokenize{assets/module-02/examples/M2-108.pdf}}
\par\smallskip{\small 图 M2-108：正交投影只留下匹配的本征模}
\end{center}
\end{minipage}\par

图 M2-108 的边界曲线在 \(0<y<b\)
内有正、负两瓣，正好是第二正弦模态。与每个 \(Y_n\) 作正交投影时，只有
\(n=2\) 的积分非零；没有引入任何其他模态去近似这条已经匹配的函数。

在 \(0<y<b\) 且两端为零时，\(Y_n=\sin(n\pi y/b)\)；若边界数据正好是
\(\sin(2\pi y/b)\)，其余模态系数都为零。

\par\noindent\begin{minipage}{\linewidth}
\begin{center}
\includegraphics[width=\linewidth,height=.40\textheight,keepaspectratio]{\detokenize{assets/module-02/theory/tikz-prereq-m2-topic18-v7-01-separation-modes.pdf}}
\par\smallskip{\small 图示：分离变量把边界图样拆成独立模态并逐模匹配}
\end{center}
\end{minipage}\par

分离变量法用满足齐次边界条件的正交本征函数展开边界数据，各模态再独立延伸到区域内部。直角坐标中的正弦级数和球坐标中的勒让德多项式形式不同，但选项原则相同：坐标系适应几何边界，本征函数满足齐次条件，展开系数由非齐次数据确定，端点条件筛选径向或纵向解。

\par\noindent\begin{minipage}{\linewidth}
\begin{center}
\includegraphics[width=\linewidth,height=.40\textheight,keepaspectratio]{\detokenize{assets/module-02/original/tikz-m2-topic18-separation-v2-01-workflow.png}}
\par\smallskip{\small 图示：分离变量法从几何、齐次边界、正交投影到模态叠加的完整路径}
\end{center}
\end{minipage}\par

\subsection{一、使用分离变量法前必须检查的条件}\label{m2:ux4e00ux4f7fux7528ux5206ux79bbux53d8ux91cfux6cd5ux524dux5fc5ux987bux68c0ux67e5ux7684ux6761ux4ef6}

以拉普拉斯方程为例，分离变量法最适合下列问题：

\begin{enumerate}
\def\labelenumi{\arabic{enumi}.}
\tightlist
\item
  求解区域的边界与某个正交坐标面重合；
\item
  控制方程线性且系数在所选坐标中可分离；
\item
  大部分边界条件已经齐次化；
\item
  剩余边界数据可在所得本征函数系中展开；
\item
  原点、无穷远或其他端点处的正则性条件明确。
\end{enumerate}

若不同边界都非齐次，通常先利用线性叠加把原问题拆成若干个``只有一段边界非齐次''的子问题。若控制方程含源，可先寻找一个特解，再对剩余齐次方程使用分离变量法。

\textbf{例 M2-109：先减去满足部分边界的简单解}

在二维真空矩形 \(0<x<a,0<y<b\) 内，电势满足拉普拉斯方程。左右边分别给定
\(\varphi(0,y)=0\)、\(\varphi(a,y)=V_0\)，上下边为已知函数
\(f_0(x),f_b(x)\)，端点与相邻边电势相容。构造一个简单的辅助调和势，使减去它后的新未知量在左右边满足齐次条件；写出新未知量的方程和四边数据，说明它为何适合正弦展开。

\par\noindent\begin{minipage}{\linewidth}
\textbf{解答}

\begin{center}
\includegraphics[width=\linewidth,height=.40\textheight,keepaspectratio]{\detokenize{assets/module-02/examples/M2-109.pdf}}
\par\smallskip{\small 图 M2-109：先减去满足部分边界的简单解}
\end{center}
\end{minipage}\par

图 M2-109 的竖直虚线表示线性基准势的等势线，左右边值已由 \(V_0x/a\)
匹配。将它从原势中减去后，新未知量在左右边为零，上下边的剩余数据也必须同样相减，然后才可作正弦展开。

在矩形内若左右两边分别固定为 \(0,V_0\)，上下两边还有附加数据，可先取

\[
\varphi_0(x)=\frac{V_0x}{a}.
\]

\(\varphi_0\) 已满足拉普拉斯方程和左右边界。令
\(u=\varphi-\varphi_0\)，则图 M2-109 中的矩形内

\[
\nabla^2u=0,\qquad
u(0,y)=u(a,y)=0,
\]

\[
u(x,0)=f_0(x)-\frac{V_0x}{a},\qquad
u(x,b)=f_b(x)-\frac{V_0x}{a}.
\]

左右边界已齐次，因此可以沿 \(x\)
选正弦基；上下边界仍须按上面变换后的数据分别展开。

\subsection{二、半无限条带：边界模态向内部指数衰减}\label{m2:ux4e8cux534aux65e0ux9650ux6761ux5e26ux8fb9ux754cux6a21ux6001ux5411ux5185ux90e8ux6307ux6570ux8870ux51cf}

考虑区域

\[
0<x<a,
\qquad y>0,
\]

满足

\[
\nabla^2\varphi=0,
\]

边界条件为

\[
\varphi(0,y)=\varphi(a,y)=0,
\qquad
\varphi(x,0)=f(x),
\qquad
\varphi(x,y)\to0\quad(y\to\infty).
\]

\subsubsection{1.
分离常数与端点筛选}\label{m2:ux5206ux79bbux5e38ux6570ux4e0eux7aefux70b9ux7b5bux9009}

令

\[
\varphi(x,y)=X(x)Y(y).
\]

代入 \(\varphi_{xx}+\varphi_{yy}=0\) 并除以 \(XY\)：

\[
\frac{X''}{X}=-\frac{Y''}{Y}=-k^2.
\]

于是

\[
X''+k^2X=0,
\qquad
Y''-k^2Y=0.
\]

左右边界给出 \(X(0)=X(a)=0\)，非零解只在

\[
\boxed{k_n=\frac{n\pi}{a}},
\qquad n=1,2,3,\ldots
\]

时存在，且

\[
X_n(x)=\sin\frac{n\pi x}{a}.
\]

\(Y_n\) 的一般解为
\(A_ne^{k_ny}+B_ne^{-k_ny}\)。无穷远衰减要求舍去增长项，故

\[
Y_n(y)=e^{-k_ny}.
\]

叠加全部模态：

\[
\boxed{
\varphi(x,y)
=\sum_{n=1}^{\infty}
b_n\sin\frac{n\pi x}{a}
e^{-n\pi y/a}
}.
\]

利用正交关系

\[
\int_0^a
\sin\frac{n\pi x}{a}
\sin\frac{m\pi x}{a}
\,dx
=\frac a2\delta_{mn},
\]

由底边数据得到

\[
\boxed{
b_n=\frac2a\int_0^a
f(x)\sin\frac{n\pi x}{a}\,dx
}.
\]

\par\noindent\begin{minipage}{\linewidth}
\begin{center}
\includegraphics[width=\linewidth,height=.40\textheight,keepaspectratio]{\detokenize{assets/module-02/original/tikz-m2-topic18-separation-v2-02-strip-modes.png}}
\par\smallskip{\small 图示：半无限条带中高阶边界模态比低阶模态衰减得更快}
\end{center}
\end{minipage}\par

每个模态的纵向衰减长度为

\[
\ell_n=\frac1{k_n}=\frac{a}{n\pi}.
\]

所以离边界足够远时，最低的非零模态占主导；边界上的细小空间结构对应较大的
\(n\)，只能影响靠近边界的区域。

\textbf{例 M2-110：第二正弦模态衰减得更快}

真空无源半无限条带为
\(0<x<a,y>0\)，\(a>0\)。两竖直侧边电势均为零，底边为
\(V_0\sin(2\pi x/a)\)，电势在 \(y\to\infty\)
时趋零。求电势的单模解，并从沿 \(+y\)
方向的衰减因子说明第二正弦模态与第一模态相比谁衰减更快。

\par\noindent\begin{minipage}{\linewidth}
\textbf{解答}

\begin{center}
\includegraphics[width=\linewidth,height=.40\textheight,keepaspectratio]{\detokenize{assets/module-02/examples/M2-110.pdf}}
\par\smallskip{\small 图 M2-110：第二正弦模态衰减得更快}
\end{center}
\end{minipage}\par

图 M2-110 的底边由两个符号相反的模态瓣组成，侧边仍为零势。正交性留下
\(n=2\)，对应纵向衰减指数为
\(2\pi/a\)；这条模态的衰减长度是第一模态的一半，并非只在近边界成立。

若

\[
f(x)=V_0\sin\frac{2\pi x}{a},
\]

正交性使得只有 \(b_2=V_0\)，于是

\[
\boxed{
\varphi=V_0\sin\frac{2\pi x}{a}e^{-2\pi y/a}
}.
\]

这不是近似，而是该边值问题的精确解。

\subsection{三、有限矩形：双曲函数连接两条水平边界}\label{m2:ux4e09ux6709ux9650ux77e9ux5f62ux53ccux66f2ux51fdux6570ux8fdeux63a5ux4e24ux6761ux6c34ux5e73ux8fb9ux754c}

考虑

\[
0<x<a,
\qquad 0<y<b,
\]

其中左、右、下三边接地，上边给定

\[
\varphi(x,b)=f(x).
\]

左右边界仍选择正弦本征函数。下边 \(y=0\) 要求纵向函数为零，所以应取
\(\sinh(k_ny)\)。归一化到上边后，解为

\[
\boxed{
\varphi(x,y)
=\sum_{n=1}^{\infty}
b_n\sin\frac{n\pi x}{a}
\frac{\sinh(n\pi y/a)}
{\sinh(n\pi b/a)}
},
\]

其中

\[
\boxed{
b_n=\frac2a\int_0^a
f(x)\sin\frac{n\pi x}{a}\,dx
}.
\]

\par\noindent\begin{minipage}{\linewidth}
\begin{center}
\includegraphics[width=\linewidth,height=.40\textheight,keepaspectratio]{\detokenize{assets/module-02/original/tikz-m2-topic18-separation-v2-03-rectangle.png}}
\par\smallskip{\small 图示：有限矩形中的每个正弦模态由双曲正弦函数从上边界传播到内部}
\end{center}
\end{minipage}\par

写成比值形式而不是把系数留作两个独立常量，可直接看出：

\[
\varphi_n(x,0)=0,
\qquad
\varphi_n(x,b)=b_n\sin\frac{n\pi x}{a}.
\]

\textbf{例 M2-111：有限矩形中的纵向因子是双曲正弦}

在二维真空矩形 \(0<x<a,0<y<b\)
内无电荷，\(a,b>0\)。左、右、下三边接地，上边 \(y=b\) 的电势为
\(V_0\sin(\pi x/a)\)。求内部满足拉普拉斯方程的电势，要求同时满足四条边界；说明有限高度及底边零势如何决定纵向因子，而不是直接沿用半无限条带的指数形式。

\par\noindent\begin{minipage}{\linewidth}
\textbf{解答}

\begin{center}
\includegraphics[width=\linewidth,height=.40\textheight,keepaspectratio]{\detokenize{assets/module-02/examples/M2-111.pdf}}
\par\smallskip{\small 图 M2-111：有限矩形中的纵向因子是双曲正弦}
\end{center}
\end{minipage}\par

图 M2-111
的有限顶边要求把纵向函数归一化为顶边取1、底边取0。满足这两个端值的是双曲正弦比值，而不是半无限区域的单个指数；它与横向正弦相乘后同时满足四边条件。

若

\[
f(x)=V_0\sin\frac{\pi x}{a},
\]

则

\[
\boxed{
\varphi(x,y)
=V_0\sin\frac{\pi x}{a}
\frac{\sinh(\pi y/a)}{\sinh(\pi b/a)}
}.
\]

若矩形角点处相邻两边给出的电势不一致，边界函数在角点不连续，内部级数仍可收敛到调和解，但角点附近电场往往具有奇异性；不能要求级数在不连续角点同时等于两个不同的数值。

\subsection{四、球坐标中的轴对称分离：勒让德模态}\label{m2:ux56dbux7403ux5750ux6807ux4e2dux7684ux8f74ux5bf9ux79f0ux5206ux79bbux52d2ux8ba9ux5fb7ux6a21ux6001}

若边界和电荷分布绕 \(z\) 轴旋转不变，电势与方位角无关：

\[
\varphi=\varphi(r,\theta).
\]

球坐标拉普拉斯方程化为

\[
\boxed{
\frac1{r^2}\frac{\partial}{\partial r}
\left(r^2\frac{\partial\varphi}{\partial r}\right)
+\frac1{r^2\sin\theta}
\frac{\partial}{\partial\theta}
\left(\sin\theta\frac{\partial\varphi}{\partial\theta}\right)
=0
}.
\]

令 \(\varphi=R(r)\Theta(\theta)\)，分离后得到

\[
\frac{d}{dr}\left(r^2\frac{dR}{dr}\right)
-L(L+1)R=0,
\]

\[
\frac1{\sin\theta}\frac{d}{d\theta}
\left(\sin\theta\frac{d\Theta}{d\theta}\right)
+L(L+1)\Theta=0.
\]

在极轴 \(\theta=0,\pi\) 上有限且单值的角向解为

\[
\Theta_L(\theta)=P_L(\cos\theta),
\qquad L=0,1,2,\ldots,
\]

径向解为

\[
R_L(r)=A_Lr^L+B_Lr^{-(L+1)}.
\]

因此轴对称无源区域的通解为

\[
\boxed{
\varphi(r,\theta)
=\sum_{L=0}^{\infty}
\left(A_Lr^L+\frac{B_L}{r^{L+1}}\right)
P_L(\cos\theta)
}.
\]

\par\noindent\begin{minipage}{\linewidth}
\begin{center}
\includegraphics[width=\linewidth,height=.40\textheight,keepaspectratio]{\detokenize{assets/module-02/original/tikz-m2-topic18-separation-v3-04-sphere-modes.png}}
\par\smallskip{\small 图示：球坐标分离后，角向勒让德模态与两类径向幂次必须配合端点条件筛选}
\end{center}
\end{minipage}\par

\textbf{例 M2-112：前三阶勒让德模态的角向结构}

在轴对称球坐标展开中，\(\theta\) 从对称轴正向量起，勒让德多项式按
\(P_L(1)=1\) 归一化。写出
\(P_0(\cos\theta)\)、\(P_1(\cos\theta)\)、\(P_2(\cos\theta)\)
的具体表达式，比较它们关于赤道反射的对称性，并说明这三阶分别对应球对称、偶极型和四极型的哪种角分布。

\par\noindent\begin{minipage}{\linewidth}
\textbf{解答}

\begin{center}
\includegraphics[width=\linewidth,height=.40\textheight,keepaspectratio]{\detokenize{assets/module-02/examples/M2-112.pdf}}
\par\smallskip{\small 图 M2-112：前三阶勒让德模态的角向结构}
\end{center}
\end{minipage}\par

图 M2-112 的横坐标是
\(\mu=\cos\theta\)，三条曲线分别给出零阶、一次与二次勒让德函数。球面边界按这些角模分解后，每阶再配相同阶数的径向解，不能仅凭曲线看起来相似而混配幂次。

\[
P_0(\cos\theta)=1,
\qquad
P_1(\cos\theta)=\cos\theta,
\qquad
P_2(\cos\theta)=\frac12(3\cos^2\theta-1).
\]

\(L=0\) 表示球对称模态，\(L=1\) 表示偶极型角分布，\(L=2\)
表示四极型角分布。这里的 \(L\) 与专题十一多极展开的阶数完全一致。

\subsection{五、球内与球外狄利克雷问题}\label{m2:ux4e94ux7403ux5185ux4e0eux7403ux5916ux72c4ux5229ux514bux96f7ux95eeux9898}

设半径为 \(R\) 的球面给定轴对称电势

\[
\varphi(R,\theta)=f(\theta).
\]

令 \(\mu=\cos\theta\)。勒让德多项式满足正交关系

\[
\int_{-1}^{1}P_L(\mu)P_M(\mu)\,d\mu
=\frac{2}{2L+1}\delta_{LM}.
\]

因此边界函数可展开为

\[
f(\theta)=\sum_{L=0}^{\infty}f_LP_L(\cos\theta),
\]

其中

\[
\boxed{
f_L=\frac{2L+1}{2}
\int_{-1}^{1}f(\arccos\mu)P_L(\mu)\,d\mu
}.
\]

\subsubsection{1. 球内问题}\label{m2:ux7403ux5185ux95eeux9898}

求解区域为 \(0\le r<R\) 且原点无真实奇点。\(r^{-(L+1)}\)
在原点发散，必须舍去。因此

\[
\boxed{
\varphi_{\rm \text{内}}(r,\theta)
=\sum_{L=0}^{\infty}
f_L\left(\frac rR\right)^L
P_L(\cos\theta)
}.
\]

\subsubsection{2. 球外问题}\label{m2:ux7403ux5916ux95eeux9898}

求解区域为 \(r>R\) 且要求 \(\varphi\to0\)。\(r^L\)
在无穷远不衰减，必须舍去。因此

\[
\boxed{
\varphi_{\rm \text{外}}(r,\theta)
=\sum_{L=0}^{\infty}
f_L\left(\frac Rr\right)^{L+1}
P_L(\cos\theta)
}.
\]

\par\noindent\begin{minipage}{\linewidth}
\begin{center}
\includegraphics[width=\linewidth,height=.40\textheight,keepaspectratio]{\detokenize{assets/module-02/original/tikz-m2-topic18-separation-v5-05-inside-outside.png}}
\par\smallskip{\small 图示：同一球面边界模态在球内按正幂传播，在球外按负幂衰减}
\end{center}
\end{minipage}\par

球内高阶模态离球面越远衰减越快；球外高阶模态随距离增加衰减更快。因此球心附近主要看最低非零内模态，远场主要看最低非零多极矩。

\textbf{例 M2-113：同一球面图样有不同内外径向延伸}

半径为 \(R>0\) 的球面上给定轴对称电势
\(\varphi(R,\theta)=V_0+V_1\cos\theta\)，\(V_0,V_1\) 为常量，\(\theta\)
从 \(+z\)
轴量起。分别作两个独立的无源真空边值问题：球内要求球心有限，球外要求无穷远衰减。求两种解的径向延伸，解释同一球面角向数据为何在内、外域产生不同径向幂次。

\par\noindent\begin{minipage}{\linewidth}
\textbf{解答}

\begin{center}
\includegraphics[width=\linewidth,height=.40\textheight,keepaspectratio]{\detokenize{assets/module-02/examples/M2-113.pdf}}
\par\smallskip{\small 图 M2-113：同一球面图样有不同内外径向延伸}
\end{center}
\end{minipage}\par

图 M2-113
的着色表示扣除平均值后的偶极角分布。常量部分在球内保持常量、在球外按
\(1/r\) 衰减；一次余弦部分在内外分别配 \(r\) 和 \(r^{-2}\)，两式在同一个
\(r=R\) 上恢复完全相同的数据。

因边界只含 \(L=0,1\)，球内解为

\[
\boxed{
\varphi_{\rm \text{内}}
=V_0+V_1\frac rR\cos\theta
},
\]

球外衰减解为

\[
\boxed{
\varphi_{\rm \text{外}}
=V_0\frac Rr
+V_1\left(\frac Rr\right)^2\cos\theta
}.
\]

第二项具有偶极电势的角向和径向形式。

\subsection{六、非轴对称球谐展开（拔高选讲）}\label{m2:ux516dux975eux8f74ux5bf9ux79f0ux7403ux8c10ux5c55ux5f00ux62d4ux9ad8ux9009ux8bb2}

若电势依赖方位角 \(\phi\)，角向本征函数不再只是
\(P_L(\cos\theta)\)，而是球谐函数 \(Y_{Lm}(\theta,\phi)\)：

\[
L=0,1,2,\ldots,
\qquad
m=-L,-L+1,\ldots,L.
\]

它们在单位球面上满足

\[
\int Y_{Lm}^*Y_{L'm'}\,d\Omega
=\delta_{LL'}\delta_{mm'}.
\]

若

\[
f(\theta,\phi)
=\sum_{L,m}f_{Lm}Y_{Lm}(\theta,\phi),
\]

则

\[
f_{Lm}=\int fY_{Lm}^*\,d\Omega.
\]

球内、球外解分别为

\[
\boxed{
\varphi_{\rm \text{内}}
=\sum_{L,m}f_{Lm}
\left(\frac rR\right)^L
Y_{Lm}
},
\]

\[
\boxed{
\varphi_{\rm \text{外}}
=\sum_{L,m}f_{Lm}
\left(\frac Rr\right)^{L+1}
Y_{Lm}
}.
\]

使用复球谐时，实电势要求系数满足相应共轭关系；竞赛手算中也常改用
\(\cos m\phi\)、\(\sin m\phi\) 形式的实球谐。

\textbf{例 M2-114：横向匀强场也是偶极阶实球谐}

真空中给定沿 \(+x\) 方向的匀强静电场 \(\mathbf E=E_0\mathbf e_x\)。取
\(\varphi(0)=0\)，采用 \(x=r\sin\theta\cos\phi\)
的标准球坐标约定。将对应电势写为球坐标形式，判断其球谐阶数及是否轴对称，并说明它作为球心有限调和势应属于哪一类径向项。

\par\noindent\begin{minipage}{\linewidth}
\textbf{解答}

\begin{center}
\includegraphics[width=\linewidth,height=.40\textheight,keepaspectratio]{\detokenize{assets/module-02/examples/M2-114.pdf}}
\par\smallskip{\small 图 M2-114：横向匀强场也是偶极阶实球谐}
\end{center}
\end{minipage}\par

图 M2-114 的电场沿横向 \(x\)，电势因此是 \(-E_0x\)。代入球坐标的
\(x=r\sin\theta\cos\phi\)，即可识别正则的一阶实球谐；它不是轴对称模态，但与沿
\(z\) 的偶极角模同属 \(l=1\)。

电势

\[
\varphi=-E_0x=-E_0r\sin\theta\cos\phi
\]

是 \(L=1\) 的非轴对称实球谐模态。它在原点有限，故属于球内允许的 \(r^1\)
项。

\subsection{七、正交投影、模态衰减与物理图像}\label{m2:ux4e03ux6b63ux4ea4ux6295ux5f71ux6a21ux6001ux8870ux51cfux4e0eux7269ux7406ux56feux50cf}

\par\noindent\begin{minipage}{\linewidth}
\begin{center}
\includegraphics[width=\linewidth,height=.40\textheight,keepaspectratio]{\detokenize{assets/module-02/original/tikz-m2-topic18-separation-v2-06-projection-decay.png}}
\par\smallskip{\small 图示：边界函数先投影成正交模态，各模态再按自身尺度向区域内部传播}
\end{center}
\end{minipage}\par

直角坐标和球坐标下的展开都可写成

\[
\text{边界函数}
=\sum_n(\text{投影系数})\times(\text{边界本征函数}),
\]

\[
\text{区域内解}
=\sum_n(\text{同一投影系数})
\times(\text{端点允许的传播因子})
\times(\text{本征函数}).
\]

正交性保证每个边界模态的系数可独立求出；线性性保证各模态可以直接叠加；唯一性定理保证满足方程和全部边界条件的级数就是物理解。

\subsection{八、分离变量法的使用条件}\label{m2:ux516bux5206ux79bbux53d8ux91cfux6cd5ux7684ux4f7fux7528ux6761ux4ef6}

\begin{enumerate}
\def\labelenumi{\arabic{enumi}.}
\tightlist
\item
  \textbf{先齐次化边界。}
  本征函数由齐次边界决定，非齐次边界只负责确定系数。
\item
  \textbf{分离常数的符号由边界决定。}
  两端为零的有限区间需要正弦模态，对应
  \(X''+k^2X=0\)；不能预先机械规定正负号。
\item
  \textbf{端点行为决定项的去留。}
  球内无源才可舍去负幂；球外若有给定匀强背景场，增长的 \(rP_1\)
  项反而必须保留。
\item
  \textbf{边界展开要使用正确权重。} 正弦系数的积分变量是
  \(dx\)；勒让德系数写成 \(d\mu\)，等价的 \(\theta\) 积分要带
  \(\sin\theta\)。
\item
  \textbf{角点不连续要单独解释。}
  级数在光滑边界点收敛到给定值，在跳跃点通常收敛到左右极限的平均；电场可能在角点奇异。
\item
  \textbf{不要把单个分离解当作通解。} 一般边界必须叠加全部允许模态。
\item
  \textbf{最后逐项验证。}
  检查偏微分方程、每段边界、原点或无穷远条件、量纲和已知极限。
\end{enumerate}

\subsection{九、练习与解答}\label{m2:ux4e5dux7ec3ux4e60ux4e0eux89e3ux7b54-2}

\subsubsection{作业一：半无限条带中的单模场}\label{m2:ux4f5cux4e1aux4e00ux534aux65e0ux9650ux6761ux5e26ux4e2dux7684ux5355ux6a21ux573a}

\textbf{题号：M2-115}

\textbf{层级：} 复赛　　\textbf{方法：} 边值问题

\(a>0\)，\(n\) 为正整数，\(V_0\) 为常量；域内为无源真空，电场按
\(\mathbf E=-\nabla\varphi\) 定义，模态衰减长度沿 \(+y\) 测量。

在 \(0<x<a,y>0\) 中，边界条件为

\[
\varphi(0,y)=\varphi(a,y)=0,
\]

\[
\varphi(x,0)=V_0\sin\frac{n\pi x}{a},
\qquad
\varphi\to0\quad(y\to\infty).
\]

求电势和电场，并给出该模态的衰减长度。

\par\noindent\begin{minipage}{\linewidth}
\textbf{解答}

\begin{center}
\includegraphics[width=\linewidth,height=.40\textheight,keepaspectratio]{\detokenize{assets/module-02/examples/M2-115.pdf}}
\par\smallskip{\small 图 M2-115：半无限条带中的单模场}
\end{center}
\end{minipage}\par

图 M2-115 的条带宽度决定横向本征值 \(n\pi/a\)，底边图样只留下对应的
\(n\)。纵向衰减用同一本征值，求场时还须分别微分正弦与指数；衰减长度因此为
\(a/(n\pi)\)。

只有第 \(n\) 个正弦模态存在：

\[
\boxed{
\varphi
=V_0\sin(k_nx)e^{-k_ny}
},
\qquad
k_n=\frac{n\pi}{a}.
\]

于是

\[
\boxed{
E_x=-k_nV_0\cos(k_nx)e^{-k_ny}
},
\]

\[
\boxed{
E_y=k_nV_0\sin(k_nx)e^{-k_ny}
}.
\]

直接计算得 \(\varphi_{xx}+\varphi_{yy}=0\)。衰减长度为

\[
\boxed{\ell_n=\frac a{n\pi}}.
\]

\subsubsection{作业二：上边等势、其余三边接地的矩形}\label{m2:ux4f5cux4e1aux4e8cux4e0aux8fb9ux7b49ux52bfux5176ux4f59ux4e09ux8fb9ux63a5ux5730ux7684ux77e9ux5f62}

\textbf{题号：M2-116}

\textbf{层级：} 复赛　　\textbf{方法：} 分离变量、边值问题

\(a,b>0\)；域内为二维无源真空，\(V_0\)
为常量。四角处按相邻开边界的侧极限讨论，不额外要求两种不相容角点电势同时成立。

在 \(0<x<a,0<y<b\) 中，左、右、下三边电势为零，上边除两个角点外保持
\(V_0\)。求内部电势，并说明角点处的情况。

\par\noindent\begin{minipage}{\linewidth}
\textbf{解答}

\begin{center}
\includegraphics[width=\linewidth,height=.40\textheight,keepaspectratio]{\detokenize{assets/module-02/examples/M2-116.pdf}}
\par\smallskip{\small 图 M2-116：上边等势、其余三边接地的矩形}
\end{center}
\end{minipage}\par

图 M2-116
的上边电势为常量而其余三边为零，上边数据投影后只含奇数正弦项。每项的纵向双曲因子在底边为零、顶边为1；两个角点的数据不连续，不能要求级数在那里同时等于相邻两边的不同极限。

把常量 \(V_0\) 展开成正弦级数：

\[
b_n=\frac{2V_0}{a}\int_0^a
\sin\frac{n\pi x}{a}\,dx
=\frac{2V_0}{n\pi}[1-(-1)^n].
\]

所以偶数 \(n\) 的系数为零，奇数 \(n\) 时 \(b_n=4V_0/(n\pi)\)。因此

\[
\boxed{
\varphi(x,y)
=\frac{4V_0}{\pi}
\sum_{\substack{n=1\\n\text{ 为奇数}}}^{\infty}
\frac1n
\sin\frac{n\pi x}{a}
\frac{\sinh(n\pi y/a)}
{\sinh(n\pi b/a)}
}.
\]

上边在 \(0<x<a\) 恢复
\(V_0\)，左右边恢复零。两个上角点同时属于电势为零的侧边和电势为 \(V_0\)
的上边，理想边界数据不连续；正弦级数在两个端点逐项为零；从域内沿不同方向趋近角点一般无唯一极限，角点附近场强可奇异，但区域内部的电势仍由该级数唯一确定。

\subsubsection{\texorpdfstring{作业三：球面电势为 \(V_0\cos^2\theta\)
的球内问题}{作业三：球面电势为 V\_0\textbackslash cos\^{}2\textbackslash theta 的球内问题}}\label{m2:ux4f5cux4e1aux4e09ux7403ux9762ux7535ux52bfux4e3a-v_0cos2theta-ux7684ux7403ux5185ux95eeux9898}

\textbf{题号：M2-117}

\textbf{层级：} 复赛　　\textbf{方法：} 边值问题

\(R>0\)，\(\theta\) 从 \(+z\) 轴量起，\(V_0\)
为常量；球面为给定电势边界，球内为无源真空，要求原点正则。

半径为 \(R\) 的球面给定

\[
\varphi(R,\theta)=V_0\cos^2\theta.
\]

球内无电荷，求球内电势和球心电势。

\par\noindent\begin{minipage}{\linewidth}
\textbf{解答}

\begin{center}
\includegraphics[width=\linewidth,height=.40\textheight,keepaspectratio]{\detokenize{assets/module-02/examples/M2-117.pdf}}
\par\smallskip{\small 图 M2-117：球面电势为 \(V_0\cos^2\theta\) 的球内问题}
\end{center}
\end{minipage}\par

图 M2-117 的平方余弦边界不是纯二阶项，它含一个常量平均值。先用
\(\cos^2\theta=1/3+2P_2/3\)
分解，再把两阶分别向球心延伸；二阶项在球心消失，因此中心电势为
\(V_0/3\)。

由

\[
P_2(\cos\theta)=\frac12(3\cos^2\theta-1)
\]

得到

\[
\cos^2\theta
=\frac13+\frac23P_2(\cos\theta).
\]

球内保留 \(r^L\) 项，因此

\[
\boxed{
\varphi_{\rm \text{内}}(r,\theta)
=V_0\left[
\frac13
+\frac23\left(\frac rR\right)^2
P_2(\cos\theta)
\right]
}.
\]

在球心 \(r=0\)，所有 \(L\ge1\) 模态消失，故

\[
\boxed{\varphi(0)=\frac{V_0}{3}}.
\]

这也等于边界电势的球面平均值，因为 \(\langle\cos^2\theta\rangle=1/3\)。

\subsubsection{作业四：球外的偶极型边界}\label{m2:ux4f5cux4e1aux56dbux7403ux5916ux7684ux5076ux6781ux578bux8fb9ux754c}

\textbf{题号：M2-118}

\textbf{层级：} 复赛　　\textbf{方法：} 多极展开、边值问题

\(R>0\)，\(\theta\) 从 \(+z\) 轴量起，\(V_1\) 为常量；只求 \(r>R\)
的真空外域，允许源位于被排除的球内，偶极矩方向应相对于同一极轴判断。

半径为 \(R\) 的球面电势为

\[
\varphi(R,\theta)=V_1\cos\theta,
\]

球外无电荷且无穷远电势为零。

\begin{enumerate}
\def\labelenumi{\arabic{enumi}.}
\tightlist
\item
  求球外电势；
\item
  与理想偶极电势比较，求等效偶极矩。
\end{enumerate}

\par\noindent\begin{minipage}{\linewidth}
\textbf{解答}

\begin{center}
\includegraphics[width=\linewidth,height=.40\textheight,keepaspectratio]{\detokenize{assets/module-02/examples/M2-118.pdf}}
\par\smallskip{\small 图 M2-118：球外的偶极型边界}
\end{center}
\end{minipage}\par

图 M2-118 的球面数据只有一次余弦角模，外域无穷远条件筛出 \(r^{-2}\)
分支。把归一化系数固定在 \(r=R\)
后，再与标准偶极势比较，等效偶极矩应沿图中轴线，其正负由 \(V_1\) 决定。

边界只含 \(L=1\) 模态，球外保留 \(r^{-2}\) 项：

\[
\boxed{
\varphi_{\rm \text{外}}(r,\theta)
=V_1\left(\frac Rr\right)^2\cos\theta
}.
\]

理想偶极电势为

\[
\varphi_{\rm dip}
=\frac{p\cos\theta}{4\pi\varepsilon_0r^2}.
\]

比较系数得到

\[
\boxed{p=4\pi\varepsilon_0V_1R^2},
\]

方向沿 \(+z\) 轴；若 \(V_1<0\)，方向反向。

\subsubsection{作业五：球面单一勒让德模态的内外衰减}\label{m2:ux4f5cux4e1aux4e94ux7403ux9762ux5355ux4e00ux52d2ux8ba9ux5fb7ux6a21ux6001ux7684ux5185ux5916ux8870ux51cf}

\textbf{题号：M2-119}

\textbf{层级：} 复赛　　\textbf{方法：} 边值问题

\(R>0\)，\(L\) 为非负整数，\(P_L\) 为归一化的勒让德多项式，\(V_L\)
为常量。两种解分别在无源真空球内、球外独立求解，比较的是同一角模的径向振幅。

球面电势为

\[
\varphi(R,\theta)=V_LP_L(\cos\theta),
\]

分别求球内正则解与球外衰减解，并比较在 \(r=R/2\) 与 \(r=2R\)
处的模态振幅。

\par\noindent\begin{minipage}{\linewidth}
\textbf{解答}

\begin{center}
\includegraphics[width=\linewidth,height=.40\textheight,keepaspectratio]{\detokenize{assets/module-02/examples/M2-119.pdf}}
\par\smallskip{\small 图 M2-119：球面单一勒让德模态的内外衰减}
\end{center}
\end{minipage}\par

图 M2-119
以同一个球面角模作为内外问题的边界值。内部正幂在半径减半时给振幅因子
\(2^{-L}\)，外部负幂在半径加倍时给
\(2^{-L-1}\)；二者方向图样相同，径向衰减规律却不同。

由正交性，边界只激发第 \(L\) 阶模态。球内解为

\[
\boxed{
\varphi_{\rm \text{内}}
=V_L\left(\frac rR\right)^LP_L(\cos\theta)
},
\]

球外解为

\[
\boxed{
\varphi_{\rm \text{外}}
=V_L\left(\frac Rr\right)^{L+1}P_L(\cos\theta)
}.
\]

在 \(r=R/2\)，球内振幅相对球面为

\[
\boxed{2^{-L}};
\]

在 \(r=2R\)，球外振幅相对球面为

\[
\boxed{2^{-(L+1)}}.
\]

阶数越高，离开边界后的衰减越快。这正是球心附近只保留低阶内模态、远场只保留低阶多极项的定量原因。

\section{专题十九：柱坐标中的分离变量法（基础必讲）}\label{m2:ux4e13ux9898ux5341ux4e5dux67f1ux5750ux6807ux4e2dux7684ux5206ux79bbux53d8ux91cfux6cd5ux57faux7840ux5fc5ux8bb2}

当电荷分布、材料分布和边界条件沿某一直线方向完全不变时，三维静电问题可化为二维截面问题。取该方向为
\(z\) 轴，则

\[
\frac{\partial\varphi}{\partial z}=0,
\qquad
\varphi=\varphi(\rho,\phi).
\]

这里 \(\rho\) 是到轴线的距离，\(\phi\) 是方位角，电势仍记为
\(\varphi\)。这不是把三维图形粗略投影到平面，而是每个横截面都具有相同的源和边界数据，因此唯一解也与
\(z\) 无关。

\par\noindent\begin{minipage}{\linewidth}
\begin{center}
\includegraphics[width=\linewidth,height=.40\textheight,keepaspectratio]{\detokenize{assets/module-02/original/tikz-m2-topic19-v8-01.png}}
\par\smallskip{\small 图示：轴向平移不变的三维圆柱问题化为二维极坐标截面}
\end{center}
\end{minipage}\par

求解域是否包含原点或延伸到无穷远，决定径向项的取舍；边界函数的傅里叶模态决定角向项。求出电势后，还要计算电场、面电荷密度和二维多极矩，并检查法向符号。

\subsection{一、二维柱坐标中的常用解法}\label{m2:ux4e00ux4e8cux7ef4ux67f1ux5750ux6807ux4e2dux7684ux5e38ux7528ux89e3ux6cd5}

\subsubsection{1.
拉普拉斯方程与电场分量}\label{m2:ux62c9ux666eux62c9ux65afux65b9ux7a0bux4e0eux7535ux573aux5206ux91cf}

在无源区域中，

\[
\boxed{
\frac1\rho\frac{\partial}{\partial\rho}
\left(
\rho\frac{\partial\varphi}{\partial\rho}
\right)
+\frac1{\rho^2}
\frac{\partial^2\varphi}{\partial\phi^2}
=0
}.
\]

由 \(\mathbf E=-\nabla\varphi\)，二维极坐标中的电场分量为

\[
\boxed{
E_{\rho}=-\frac{\partial\varphi}{\partial\rho},
\qquad
E_{\phi}=-\frac1\rho\frac{\partial\varphi}{\partial\phi}
}.
\]

其中 \(E_{\phi}\) 前的 \(1/\rho\) 不能遗漏。它来自方位角方向的弧长微元
\(dl_{\phi}=\rho\,d\phi\)。

\textbf{例 M2-120：柱坐标分量合成同一个直角矢量}

以 \(z\) 轴为柱坐标轴，给定电势
\(\varphi(\rho,\phi,z)=-E_0\rho\cos\phi\)，\(E_0\)
为常量，\(x=\rho\cos\phi\)、\(y=\rho\sin\phi\)。在 \(\rho>0\) 处求
\(E_{\rho},E_{\phi},E_z\)，再将局部基矢展开为直角坐标基，合成为同一电场矢量；说明分量随方位角变化是否等于实际电场方向随位置改变。

\par\noindent\begin{minipage}{\linewidth}
\textbf{解答}

\begin{center}
\includegraphics[width=\linewidth,height=.40\textheight,keepaspectratio]{\detokenize{assets/module-02/examples/M2-120.pdf}}
\par\smallskip{\small 图 M2-120：柱坐标分量合成同一个直角矢量}
\end{center}
\end{minipage}\par

图 M2-120
的绿色与橙色箭头是点处的正径向、正方位基矢，蓝色场向始终向右。分别求两分量时保留柱坐标的
\(1/\rho\)，再用基矢的直角展开合成，即可验证所有角依赖相消。

若

\[
\varphi=-E_0\rho\cos\phi=-E_0x,
\]

则

\[
E_{\rho}=E_0\cos\phi,
\qquad
E_{\phi}=-E_0\sin\phi.
\]

再利用

\[
\mathbf e_x=\cos\phi\,\mathbf e_{\rho}-\sin\phi\,\mathbf e_{\phi},
\]

便有 \(\mathbf E=E_0\mathbf e_x\)，确实是沿 \(x\) 方向的匀强电场。

\subsubsection{2.
分离变量给出的模态}\label{m2:ux5206ux79bbux53d8ux91cfux7ed9ux51faux7684ux6a21ux6001}

令

\[
\varphi(\rho,\phi)=R(\rho)\Phi(\phi).
\]

分离后得到

\[
\Phi''+m^2\Phi=0,
\]

\[
\rho^2R''+\rho R'-m^2R=0.
\]

电势必须满足单值性

\[
\Phi(\phi+2\pi)=\Phi(\phi),
\]

所以角向模态只能取

\[
m=0,1,2,\ldots.
\]

当 \(m\ge1\) 时，

\[
\Phi_m=C_m\cos m\phi+D_m\sin m\phi,
\]

\[
R_m=A_m\rho^m+B_m\rho^{-m}.
\]

当 \(m=0\) 时，角向函数为常量，径向解为

\[
R_0=A_0+B_0\ln\rho.
\]

\par\noindent\begin{minipage}{\linewidth}
\begin{center}
\includegraphics[width=\linewidth,height=.40\textheight,keepaspectratio]{\detokenize{assets/module-02/original/tikz-m2-topic19-v8-02.png}}
\par\smallskip{\small 图示：柱坐标分离变量中角向单值性与径向端点条件的作用}
\end{center}
\end{minipage}\par

叠加全部独立模态，得到二维柱坐标下无源区域的通解

\[
\boxed{
\begin{aligned}
\varphi(\rho,\phi)
={}&A_0+B_0\ln\rho\\
&+\sum_{m=1}^{\infty}
\left[
\rho^m(a_m\cos m\phi+b_m\sin m\phi)
+\rho^{-m}(c_m\cos m\phi+d_m\sin m\phi)
\right].
\end{aligned}
}
\]

这里的``通解''只表示局部形式完备，并不表示所有项在某个具体物理区域中都允许存在。

\subsubsection{3.
先按解域筛项}\label{m2:ux5148ux6309ux89e3ux57dfux7b5bux9879}

\par\noindent\begin{minipage}{\linewidth}
\begin{center}
\includegraphics[width=\linewidth,height=.40\textheight,keepaspectratio]{\detokenize{assets/module-02/original/tikz-m2-topic19-v8-03.png}}
\par\smallskip{\small 图示：圆盘内域、圆外区域和圆环区域允许保留不同径向项}
\end{center}
\end{minipage}\par

三类典型解域的筛项规则如下。

{\def\LTcaptype{none} % do not increment counter
\begin{longtable}[]{@{}
  >{\raggedright\arraybackslash}p{(\linewidth - 4\tabcolsep) * \real{0.3333}}
  >{\raggedright\arraybackslash}p{(\linewidth - 4\tabcolsep) * \real{0.3333}}
  >{\raggedright\arraybackslash}p{(\linewidth - 4\tabcolsep) * \real{0.3333}}@{}}
\toprule\noalign{}
\begin{minipage}[b]{\linewidth}\raggedright
解域
\end{minipage} & \begin{minipage}[b]{\linewidth}\raggedright
端点要求
\end{minipage} & \begin{minipage}[b]{\linewidth}\raggedright
通常允许的径向项
\end{minipage} \\
\midrule\noalign{}
\endhead
\bottomrule\noalign{}
\endlastfoot
圆盘 \(0\le\rho<R\) & 原点无真实源时电势有限 & \(1,\rho^m\) \\
圆外 \(\rho>R\) & 远场已指定 & 已知背景项加衰减修正 \(\rho^{-m}\) \\
圆环 \(a<\rho<b\) & 不含原点，也不到无穷远 &
\(1,\ln\rho,\rho^m,\rho^{-m}\) 均可能出现 \\
\end{longtable}
}

上表中的筛项必须结合真实源理解：

\begin{itemize}
\tightlist
\item
  若原点存在均匀线电荷，则 \(\ln\rho\) 的发散是物理奇点，不能舍去；
\item
  若原点存在二维线偶极子，则 \(\rho^{-1}\cos\phi\)
  是真实源项，不能舍去；
\item
  若圆外区域处在匀强背景电场中，则 \(\rho\cos\phi\)
  虽在无穷远增长，却是给定背景，必须保留；
\item
  感应电荷产生的修正项不能比给定远场增长得更快。
\end{itemize}

\begin{quote}
\begin{quote}
筛项的正确顺序是：先确定真实求解区域，再标出允许的真实奇点和给定远场，最后才讨论``发散项''是否应舍去。只背``正幂在内、负幂在外''会在含源问题中立即出错。
\end{quote}
\end{quote}

\subsubsection{4.
三个基础模态的物理意义}\label{m2:ux4e09ux4e2aux57faux7840ux6a21ux6001ux7684ux7269ux7406ux610fux4e49}

\textbf{零阶对数模态。} 无限长线电荷密度为 \(\lambda\) 时，

\[
\boxed{
\varphi(\rho)
=-\frac{\lambda}{2\pi\varepsilon_0}
\ln\frac{\rho}{\rho_0}
},
\]

其中 \(\rho_0\) 只确定电势参考零点。对应电场为

\[
\mathbf E
=\frac{\lambda}{2\pi\varepsilon_0\rho}\mathbf e_{\rho}.
\]

\textbf{一次正幂模态。}

\[
-E_0\rho\cos\phi=-E_0x
\]

对应沿 \(+x\) 方向的匀强电场。

\textbf{一次负幂模态。} 理想二维线偶极矩为

\[
\mathbf p_{\ell}=p_{\ell}\mathbf e_x,
\]

其远场电势为

\[
\boxed{
\varphi_{\rm dip}
=\frac{p_{\ell}}{2\pi\varepsilon_0\rho}\cos\phi
}.
\]

\(p_{\ell}\) 是``每单位轴向长度的三维偶极矩''，量纲为
\(\mathrm C\)。若用两根相距 \(h\) 的异号线电荷构造，则

\[
p_{\ell}=\lambda h.
\]

\(m=2\) 对应二维四极型角分布，更高的 \(m\) 描述更细的角向起伏。

\subsection{二、圆周边界数据的傅里叶传播}\label{m2:ux4e8cux5706ux5468ux8fb9ux754cux6570ux636eux7684ux5085ux91ccux53f6ux4f20ux64ad}

在圆周 \(\rho=R\) 上给定电势

\[
f(\phi)
=F_0+\sum_{m=1}^{\infty}
\left(
F_m\cos m\phi+G_m\sin m\phi
\right).
\]

三角函数正交性保证不同 \(m\)
的边界模态彼此独立，不会在拉普拉斯方程中相互混合。

\subsubsection{1.
圆盘内的正则解}\label{m2:ux5706ux76d8ux5185ux7684ux6b63ux5219ux89e3}

若圆盘内无源且原点处有限，则

\[
\boxed{
\varphi_{\rm in}(\rho,\phi)
=F_0+\sum_{m=1}^{\infty}
\left(\frac{\rho}{R}\right)^m
\left(
F_m\cos m\phi+G_m\sin m\phi
\right)
}.
\]

每个非零角模向圆心传播时都按 \((\rho/R)^m\) 衰减。\(m\)
越大，边界上的精细起伏越难深入圆盘内部。

\textbf{例 M2-121：纯二阶圆周边界在圆心没有匀强场}

在与 \(z\) 无关的二维无源真空圆盘 \(0\le\rho<R\) 内求电势，\(R>0\)，极角
\(\phi\) 从 \(+x\) 轴起算。圆周给定
\(\varphi(R,\phi)=V_0\cos2\phi\)，要求电势在圆心正则且关于 \(\phi\)
单值。求圆盘内电势及圆心电场，说明圆心处应如何取极坐标表达式的极限。

\par\noindent\begin{minipage}{\linewidth}
\textbf{解答}

\begin{center}
\includegraphics[width=\linewidth,height=.40\textheight,keepaspectratio]{\detokenize{assets/module-02/examples/M2-121.pdf}}
\par\smallskip{\small 图 M2-121：纯二阶圆周边界在圆心没有匀强场}
\end{center}
\end{minipage}\par

图 M2-121 的圆周正负分区呈二阶角向对称，正则径向因子必须是
\(\rho^2\)。改写为 \(x^2-y^2\)
后，对两个坐标求导都在原点消失；边界非零并不要求中心有匀强电场。

若

\[
\varphi(R,\phi)=V_0\cos2\phi,
\]

则圆盘内唯一正则解为

\[
\boxed{
\varphi(\rho,\phi)
=V_0\left(\frac{\rho}{R}\right)^2\cos2\phi
}.
\]

圆心处 \(\varphi=0\)，且场也趋于零，说明纯四极边界不会在圆心产生匀强场。

\subsubsection{2.
圆外的衰减解}\label{m2:ux5706ux5916ux7684ux8870ux51cfux89e3}

若圆外无源，并要求 \(\rho\to\infty\) 时 \(\varphi\to0\)，则只允许
\(m\ge1\) 的负次幂模态：

\[
\boxed{
\varphi_{\rm out}(\rho,\phi)
=\sum_{m=1}^{\infty}
\left(\frac{R}{\rho}\right)^m
\left(
F_m\cos m\phi+G_m\sin m\phi
\right)
}.
\]

\par\noindent\begin{minipage}{\linewidth}
\begin{center}
\includegraphics[width=\linewidth,height=.40\textheight,keepaspectratio]{\detokenize{assets/module-02/original/tikz-m2-topic19-v8-04-margin12-v11.png}}
\par\smallskip{\small 图示：圆周上的傅里叶模态分别按幂律向圆盘内或圆外延拓}
\end{center}
\end{minipage}\par

这里没有 \(F_0\) 项。原因是二维圆外区域中，不存在同时满足

\[
\varphi(R)=F_0\ne0,
\qquad
\varphi(\infty)=0
\]

的有界轴对称调和函数：\(m=0\) 解只能是
\(A_0+B_0\ln\rho\)。这反映了无限长孤立圆柱相对于无穷远的电容定义需要格外小心。

\textbf{例 M2-122：有限外边界允许二维常值内边界}

两根无限长同轴圆柱之间为无源真空，横截面是
\(R<\rho<b\)，\(0<R<b\)。内、外圆柱面分别保持常电势 \(V_0\) 与零，系统与
\(z\)、方位角无关。求圆环域内电势，核对两个边界，并说明为何本题零角模的径向解应包含对数。

\par\noindent\begin{minipage}{\linewidth}
\textbf{解答}

\begin{center}
\includegraphics[width=\linewidth,height=.40\textheight,keepaspectratio]{\detokenize{assets/module-02/examples/M2-122.pdf}}
\par\smallskip{\small 图 M2-122：有限外边界允许二维常值内边界}
\end{center}
\end{minipage}\par

图 M2-122 的内圆柱保持
\(V_0\)、有限外圆柱接地。零角模的两项为常量和对数，由两圆上的电势确定系数；无穷远极限会同时涉及电荷和电势约束，不能在两种固定量之间直接交换。

若希望半径 \(R\) 的圆周保持常量电势
\(V_0\)，通常必须引入有限外边界。例如在 \(R<\rho<b\)
的同轴圆环区域内规定

\[
\varphi(R)=V_0,
\qquad
\varphi(b)=0,
\]

则

\[
\boxed{
\varphi(\rho)
=V_0\frac{\ln(b/\rho)}{\ln(b/R)}
}.
\]

当 \(b\to\infty\) 时，在固定电荷与固定电势之间不能简单互换极限。

\subsection{三、圆柱边值问题的求解步骤}\label{m2:ux4e09ux5706ux67f1ux8fb9ux503cux95eeux9898ux7684ux6c42ux89e3ux6b65ux9aa4}

处理柱坐标分离变量题时，可以按以下顺序进行：

\begin{enumerate}
\def\labelenumi{\arabic{enumi}.}
\tightlist
\item
  \textbf{确定解域}：圆盘内、圆外还是圆环；
\item
  \textbf{标出真实源}：原点是否存在对数、偶极或更高多极奇点；
\item
  \textbf{写出远场}：匀强场、多极背景或趋于零；
\item
  \textbf{利用对称性}：判断保留余弦项、正弦项还是两者都要；
\item
  \textbf{按端点行为筛径向项}；
\item
  \textbf{把边界函数展开为傅里叶级数并逐模匹配}；
\item
  \textbf{由电势求电场和表面电荷}；
\item
  \textbf{检查方程、边界、切向场、总电荷和极限行为}。
\end{enumerate}

下面两个模型展示``外部背景''和``内部真实奇点''如何对同一套通解作出相反筛选。

\subsection{四、接地圆柱置于匀强外电场}\label{m2:ux56dbux63a5ux5730ux5706ux67f1ux7f6eux4e8eux5300ux5f3aux5916ux7535ux573a}

半径为 \(R\) 的无限长接地导体圆柱，轴线沿 \(z\) 方向。外加匀强电场为

\[
\mathbf E_0=E_0\mathbf e_x.
\]

求圆柱外部 \(\rho\ge R\) 的电势、电场和表面电荷密度。

\subsubsection{1. 定解条件}\label{m2:ux5b9aux89e3ux6761ux4ef6}

真实求解域为
\(\rho>R\)，其中无自由电荷，所以满足拉普拉斯方程。边界条件为

\[
\varphi(R,\phi)=0,
\]

\[
\varphi(\rho,\phi)
\sim-E_0\rho\cos\phi,
\qquad
\rho\to\infty.
\]

系统关于 \(x\) 轴镜面对称，因此电势是 \(\phi\) 的偶函数，只保留余弦项。

\par\noindent\begin{minipage}{\linewidth}
\begin{center}
\includegraphics[width=\linewidth,height=.40\textheight,keepaspectratio]{\detokenize{assets/module-02/original/tikz-m2-topic19-v8-05-margin12-v11.png}}
\par\smallskip{\small 图示：接地圆柱在沿正x方向匀强外场中的近场畸变、远场条件和感应电荷符号}
\end{center}
\end{minipage}\par

\subsubsection{2.
远场先锁定角模}\label{m2:ux8fdcux573aux5148ux9501ux5b9aux89d2ux6a21}

远场只有

\[
-E_0\rho\cos\phi
\]

这一 \(m=1\)
正次幂模态。导体的影响在远处必须衰减，因此感应修正只能取同一角模的负次幂：

\[
\varphi(\rho,\phi)
=-E_0\rho\cos\phi
+\frac{C}{\rho}\cos\phi.
\]

不出现 \(\ln\rho\)，因为它会改变规定的远场，并对应额外净线电荷。也不出现
\(m\ge2\) 项，因为圆周边界条件中没有相应角模。

\subsubsection{3.
匹配接地边界}\label{m2:ux5339ux914dux63a5ux5730ux8fb9ux754c}

在 \(\rho=R\) 上，

\[
-E_0R\cos\phi+\frac{C}{R}\cos\phi=0.
\]

因此

\[
C=E_0R^2.
\]

最终电势为

\[
\boxed{
\varphi(\rho,\phi)
=-E_0
\left(
\rho-\frac{R^2}{\rho}
\right)\cos\phi,
\qquad
\rho\ge R
}.
\]

\par\noindent\begin{minipage}{\linewidth}
\begin{center}
\includegraphics[width=\linewidth,height=.40\textheight,keepaspectratio]{\detokenize{assets/module-02/original/tikz-m2-topic19-v8-06.png}}
\par\smallskip{\small 图示：外部圆柱解由匀强背景模态和同角模的衰减感应项组成}
\end{center}
\end{minipage}\par

\subsubsection{4. 电场分量}\label{m2:ux7535ux573aux5206ux91cf}

径向分量为

\[
\begin{aligned}
E_{\rho}
&=-\frac{\partial\varphi}{\partial\rho}\\
&=\boxed{
E_0
\left(
1+\frac{R^2}{\rho^2}
\right)\cos\phi
}.
\end{aligned}
\]

方位角分量为

\[
\begin{aligned}
E_{\phi}
&=-\frac1\rho\frac{\partial\varphi}{\partial\phi}\\
&=\boxed{
-E_0
\left(
1-\frac{R^2}{\rho^2}
\right)\sin\phi
}.
\end{aligned}
\]

在导体表面

\[
E_{\phi}(R,\phi)=0,
\]

与静电平衡时导体表面切向场为零一致。

\subsubsection{5.
面电荷密度与二维偶极矩}\label{m2:ux9762ux7535ux8377ux5bc6ux5ea6ux4e0eux4e8cux7ef4ux5076ux6781ux77e9}

圆柱外表面从导体指向真空的法向为
\(+\mathbf e_{\rho}\)。导体内部电场为零，所以

\[
\sigma(\phi)
=\varepsilon_0E_{\rho}(R^+,\phi).
\]

由

\[
E_{\rho}(R^+,\phi)=2E_0\cos\phi,
\]

得到

\[
\boxed{
\sigma(\phi)
=2\varepsilon_0E_0\cos\phi
}.
\]

\(\phi=0\) 一侧带正电，\(\phi=\pi\) 一侧带负电。每单位长度的总感应电荷为

\[
\int_0^{2\pi}\sigma(\phi)R\,d\phi=0.
\]

感应修正势

\[
\varphi_{\rm ind}
=\frac{E_0R^2}{\rho}\cos\phi
\]

与二维线偶极势

\[
\frac{p_{\ell}}{2\pi\varepsilon_0\rho}\cos\phi
\]

比较，得到圆柱的感应二维偶极矩

\[
\boxed{
p_{\ell}=2\pi\varepsilon_0E_0R^2
}.
\]

结果同时通过了五项检查：真实区域内调和、圆柱面接地、表面切向场为零、远处恢复匀强场、总感应电荷为零。

\subsection{五、接地圆柱壳内的中心线偶极子}\label{m2:ux4e94ux63a5ux5730ux5706ux67f1ux58f3ux5185ux7684ux4e2dux5fc3ux7ebfux5076ux6781ux5b50}

半径为 \(R\)
的无限长接地圆柱壳内部处于真空。在轴线上放置理想二维线偶极子

\[
\mathbf p_{\ell}=p_{\ell}\mathbf e_x.
\]

求内腔 \(0<\rho<R\) 中的电势、诱导电场和内表面感应电荷密度。

\subsubsection{1.
真实奇点决定必须保留的项}\label{m2:ux771fux5b9eux5947ux70b9ux51b3ux5b9aux5fc5ux987bux4fddux7559ux7684ux9879}

除原点外，内腔满足拉普拉斯方程。原点处的真实偶极奇点要求

\[
\varphi(\rho,\phi)
\sim
\frac{p_{\ell}}{2\pi\varepsilon_0\rho}\cos\phi,
\qquad
\rho\to0.
\]

壳上感应电荷在原点附近产生的电势必须光滑。与 \(m=1\)
奇异项具有相同角向结构的最低阶正则项是 \(\rho\cos\phi\)，所以候选势写成

\[
\varphi(\rho,\phi)
=\frac{p_{\ell}}{2\pi\varepsilon_0\rho}\cos\phi
+A\rho\cos\phi.
\]

\par\noindent\begin{minipage}{\linewidth}
\begin{center}
\includegraphics[width=\linewidth,height=.40\textheight,keepaspectratio]{\detokenize{assets/module-02/original/tikz-m2-topic19-v8-07.png}}
\par\smallskip{\small 图示：中心线偶极子的奇异势由圆柱内壁产生的同角模正则势抵消}
\end{center}
\end{minipage}\par

\subsubsection{2.
接地边界确定正则项}\label{m2:ux63a5ux5730ux8fb9ux754cux786eux5b9aux6b63ux5219ux9879}

在 \(\rho=R\) 上要求 \(\varphi=0\)：

\[
\frac{p_{\ell}}{2\pi\varepsilon_0R}\cos\phi
+AR\cos\phi=0.
\]

因此

\[
A=-\frac{p_{\ell}}{2\pi\varepsilon_0R^2}.
\]

内腔电势为

\[
\boxed{
\varphi(\rho,\phi)
=\frac{p_{\ell}}{2\pi\varepsilon_0}
\left(
\frac1\rho-\frac{\rho}{R^2}
\right)\cos\phi
}.
\]

第一项是原点真实线偶极子的奇异势，第二项是内壁感应电荷产生的正则势。

\subsubsection{3.
诱导电场方向}\label{m2:ux8bf1ux5bfcux7535ux573aux65b9ux5411}

感应势单独写成

\[
\varphi_{\rm ind}
=-\frac{p_{\ell}}{2\pi\varepsilon_0R^2}\rho\cos\phi
=-\frac{p_{\ell}}{2\pi\varepsilon_0R^2}x.
\]

因此

\[
\boxed{
\mathbf E_{\rm ind}
=-\nabla\varphi_{\rm ind}
=\frac{p_{\ell}}{2\pi\varepsilon_0R^2}\mathbf e_x
}.
\]

诱导电场沿 \(+x\) 方向，与 \(\mathbf p_{\ell}\)
同向。右侧靠近正线电荷的内壁感应负电，左侧靠近负线电荷的内壁感应正电；两侧感应电荷在中心产生的电场都沿
\(+x\) 方向。

\subsubsection{4.
内壁面电荷密度}\label{m2:ux5185ux58c1ux9762ux7535ux8377ux5bc6ux5ea6}

总电场分量为

\[
\boxed{
E_{\rho}
=\frac{p_{\ell}}{2\pi\varepsilon_0}
\left(
\frac1{\rho^2}+\frac1{R^2}
\right)\cos\phi
},
\]

\[
\boxed{
E_{\phi}
=\frac{p_{\ell}}{2\pi\varepsilon_0}
\left(
\frac1{\rho^2}-\frac1{R^2}
\right)\sin\phi
}.
\]

在 \(\rho=R\) 上，

\[
E_{\phi}(R^-,\phi)=0,
\]

再次满足导体表面切向场为零。

对圆柱壳内表面，从导体指向真空内腔的法向是

\[
\mathbf n=-\mathbf e_{\rho}.
\]

因此不能直接写 \(\sigma=\varepsilon_0E_{\rho}\)，而应写

\[
\sigma_{\rm in}
=\varepsilon_0\mathbf E(R^-)\cdot\mathbf n
=-\varepsilon_0E_{\rho}(R^-,\phi).
\]

由于

\[
E_{\rho}(R^-,\phi)
=\frac{p_{\ell}}{\pi\varepsilon_0R^2}\cos\phi,
\]

得到

\[
\boxed{
\sigma_{\rm in}(\phi)
=-\frac{p_{\ell}}{\pi R^2}\cos\phi
}.
\]

\par\noindent\begin{minipage}{\linewidth}
\begin{center}
\includegraphics[width=\linewidth,height=.40\textheight,keepaspectratio]{\detokenize{assets/module-02/original/tikz-m2-topic19-v8-08.png}}
\par\smallskip{\small 图示：圆柱外表面和圆柱壳内表面从导体指向真空的法向相反}
\end{center}
\end{minipage}\par

在 \(\phi=0\) 处
\(\sigma_{\rm in}<0\)，与内壁靠近正线电荷的一侧感应负电一致。总感应电荷为

\[
\int_0^{2\pi}\sigma_{\rm in}R\,d\phi=0,
\]

与中心偶极子的净线电荷为零相符。

\subsection{六、两类圆柱题的对照}\label{m2:ux516dux4e24ux7c7bux5706ux67f1ux9898ux7684ux5bf9ux7167}

{\def\LTcaptype{none} % do not increment counter
\begin{longtable}[]{@{}
  >{\raggedright\arraybackslash}p{(\linewidth - 4\tabcolsep) * \real{0.3333}}
  >{\raggedright\arraybackslash}p{(\linewidth - 4\tabcolsep) * \real{0.3333}}
  >{\raggedright\arraybackslash}p{(\linewidth - 4\tabcolsep) * \real{0.3333}}@{}}
\toprule\noalign{}
\begin{minipage}[b]{\linewidth}\raggedright
项目
\end{minipage} & \begin{minipage}[b]{\linewidth}\raggedright
圆柱外的匀强背景
\end{minipage} & \begin{minipage}[b]{\linewidth}\raggedright
圆柱壳内的中心偶极子
\end{minipage} \\
\midrule\noalign{}
\endhead
\bottomrule\noalign{}
\endlastfoot
真实求解域 & \(\rho>R\) & \(0<\rho<R\) \\
必须保留的已知项 & \(-E_0\rho\cos\phi\) &
\(p_{\ell}\cos\phi/(2\pi\varepsilon_0\rho)\) \\
该项的来源 & 无穷远背景 & 原点真实奇点 \\
导体产生的修正 & \(+E_0R^2\cos\phi/\rho\) &
\(-p_{\ell}\rho\cos\phi/(2\pi\varepsilon_0R^2)\) \\
表面法向 & \(+\mathbf e_{\rho}\) & \(-\mathbf e_{\rho}\) \\
总感应电荷 & \(0\) & \(0\) \\
\end{longtable}
}

两题都只含 \(m=1\)
余弦模态，但筛项方向正好相反。决定正幂或负幂去留的不是``导体在里面还是外面''这一句话，而是已知背景、真实奇点和求解域端点行为。

\subsection{七、柱坐标解的端点条件}\label{m2:ux4e03ux67f1ux5750ux6807ux89e3ux7684ux7aefux70b9ux6761ux4ef6}

\begin{enumerate}
\def\labelenumi{\arabic{enumi}.}
\item
  \textbf{把三维点偶极公式用于二维线偶极。}\\
  二维线偶极势按 \(1/\rho\) 衰减，系数为
  \(p_{\ell}/(2\pi\varepsilon_0)\)；三维点偶极势按 \(1/r^2\) 衰减。
\item
  \textbf{看到 \(\rho^{-m}\) 发散就一律删除。}\\
  若原点处有真实多极源，该发散正是必须保留的物理奇点。
\item
  \textbf{认为圆外问题只能保留负次幂。}\\
  已知远场可以包含
  \(\rho^m\)；需要衰减的是导体造成的修正，而不是给定背景本身。
\item
  \textbf{遗漏 \(m=0\) 的 \(\ln\rho\)。}\\
  同轴圆柱、电容器和线电荷问题往往由对数项主导。
\item
  \textbf{遗漏 \(E_{\phi}\) 中的 \(1/\rho\)。}
\item
  \textbf{内壁仍取 \(+\mathbf e_{\rho}\) 为导体外法向。}\\
  圆柱壳内表面从导体指向真空内腔的法向是 \(-\mathbf e_{\rho}\)。
\item
  \textbf{把感应势的符号直接当作感应场方向。}\\
  若 \(\varphi_{\rm ind}=-Cx\)，则
  \(\mathbf E_{\rm ind}=+C\mathbf e_x\)。
\item
  \textbf{在二维圆外区域随意同时规定非零常量边界势和无穷远零势。}\\
  轴对称解只有 \(A+B\ln\rho\)，必须检查这两个条件是否相容。
\end{enumerate}

\subsection{八、练习与解答}\label{m2:ux516bux7ec3ux4e60ux4e0eux89e3ux7b54-5}

\subsubsection{作业一：圆盘内的一般低阶边界}\label{m2:ux4f5cux4e1aux4e00ux5706ux76d8ux5185ux7684ux4e00ux822cux4f4eux9636ux8fb9ux754c}

\textbf{题号：M2-123}

\textbf{层级：} 复赛　　\textbf{方法：} 分离变量、边值问题

在与第三坐标无关的真空圆盘 \(0\le\rho<R\) 内没有电荷或原点奇性。采用
\(x=\rho\cos\phi,y=\rho\sin\phi\)，圆周电势为
\(V_0+V_1\cos\phi+V_2\sin2\phi\)，其中 \(R>0\)，\(V_0,V_1,V_2\)
为已知常量。求整个圆盘内有限、单值的电势，并用直角坐标或正则极限求圆心电场。

\par\noindent\begin{minipage}{\linewidth}
\textbf{解答}

\begin{center}
\includegraphics[width=\linewidth,height=.40\textheight,keepaspectratio]{\detokenize{assets/module-02/examples/M2-123.pdf}}
\par\smallskip{\small 图 M2-123：圆盘内的一般低阶边界}
\end{center}
\end{minipage}\par

图 M2-123
的圆周数据分成零阶、一次余弦和二次正弦三部分。每部分向圆心延伸时分别配
\(1,\rho,\rho^2\)；求中心场时只有一次部分保留非零梯度，二次部分只从局部场梯度进入。

边界包含 \(m=0\)、\(m=1\) 余弦和 \(m=2\)
正弦模态。圆盘内要求原点有限，因此各模态按 \((\rho/R)^m\) 向内延拓：

\[
\boxed{
\varphi(\rho,\phi)
=V_0
+V_1\frac{\rho}{R}\cos\phi
+V_2\left(\frac{\rho}{R}\right)^2\sin2\phi
}.
\]

利用

\[
\rho\cos\phi=x,
\qquad
\rho^2\sin2\phi=2xy,
\]

可写成

\[
\varphi(x,y)
=V_0+\frac{V_1}{R}x+\frac{2V_2}{R^2}xy.
\]

于是

\[
\mathbf E
=-\nabla\varphi
=-\left(
\frac{V_1}{R}+\frac{2V_2}{R^2}y
\right)\mathbf e_x
-\frac{2V_2}{R^2}x\,\mathbf e_y.
\]

在圆心

\[
\boxed{
\mathbf E(0)
=-\frac{V_1}{R}\mathbf e_x
}.
\]

只有 \(m=1\) 边界模态能在圆心留下非零匀强场；\(m\ge2\)
模态在圆心的场为零。

\subsubsection{作业二：接地圆柱在匀强外场中}\label{m2:ux4f5cux4e1aux4e8cux63a5ux5730ux5706ux67f1ux5728ux5300ux5f3aux5916ux573aux4e2d}

\textbf{题号：M2-124}

\textbf{层级：} 复赛　　\textbf{方法：} 分离变量、边值问题

无限长接地导体圆柱半径为 \(R>0\)，轴线为 \(z\)
轴，柱外为真空。远处施加匀强场
\(\mathbf E_{\infty}=E_0\mathbf e_x\)，\(E_0>0\)；以外加势
\(-E_0\rho\cos\phi\)
为参考，要求感应附加势在无穷远衰减。求柱外电势、柱面的感应自由面电荷密度，以及单位轴向长度的等效诱导二维偶极矩，并说明它的方向。

\par\noindent\begin{minipage}{\linewidth}
\textbf{解答}

\begin{center}
\includegraphics[width=\linewidth,height=.40\textheight,keepaspectratio]{\detokenize{assets/module-02/examples/M2-124.pdf}}
\par\smallskip{\small 图 M2-124：接地圆柱在匀强外场中}
\end{center}
\end{minipage}\par

图 M2-124 的导体圆周是零势边界，远处是沿 \(x\)
的匀强背景。分别保留增长的一阶背景与衰减的一阶反应项，使它们在圆周抵消；取外法向导数再得到有正负分区、总量为零的感应面电荷。

无穷远背景锁定 \(m=1\) 正次幂项，感应修正必须衰减且具有相同角模：

\[
\varphi
=-E_0\rho\cos\phi+\frac{C}{\rho}\cos\phi.
\]

接地边界给出

\[
-E_0R+\frac{C}{R}=0,
\]

故

\[
C=E_0R^2.
\]

电势为

\[
\boxed{
\varphi
=-E_0
\left(
\rho-\frac{R^2}{\rho}
\right)\cos\phi
}.
\]

径向电场为

\[
E_{\rho}
=E_0
\left(
1+\frac{R^2}{\rho^2}
\right)\cos\phi.
\]

外表面法向为 \(+\mathbf e_{\rho}\)，所以

\[
\boxed{
\sigma(\phi)
=\varepsilon_0E_{\rho}(R^+,\phi)
=2\varepsilon_0E_0\cos\phi
}.
\]

感应势与

\[
\varphi_{\rm dip}
=\frac{p_{\ell}}{2\pi\varepsilon_0\rho}\cos\phi
\]

比较，得到

\[
\boxed{
p_{\ell}=2\pi\varepsilon_0E_0R^2
}.
\]

并且

\[
\int_0^{2\pi}\sigma R\,d\phi=0.
\]

\subsubsection{作业三：接地圆柱壳内的中心线偶极子}\label{m2:ux4f5cux4e1aux4e09ux63a5ux5730ux5706ux67f1ux58f3ux5185ux7684ux4e2dux5fc3ux7ebfux5076ux6781ux5b50}

\textbf{题号：M2-125}

\textbf{层级：} 复赛　　\textbf{方法：} 分离变量、多极展开

半径为 \(R>0\) 的无限长接地导体圆柱壳内为真空，轴线为 \(z\)
轴。在轴线上置一理想线电偶极源，单位轴向长度的偶极矩为
\(\mathbf p_{\ell}=p_{\ell}\mathbf e_x\)。它在未放置导体时的二维源势为
\(p_{\ell}\cos\phi/(2\pi\varepsilon_0\rho)\)。求 \(0<\rho<R\)
内的电势、内壁感应面电荷密度，以及轴线上由感应电荷产生的有限电场；最后一项不包含线偶极源自身的奇异场。

\par\noindent\begin{minipage}{\linewidth}
\textbf{解答}

\begin{center}
\includegraphics[width=\linewidth,height=.40\textheight,keepaspectratio]{\detokenize{assets/module-02/examples/M2-125.pdf}}
\par\smallskip{\small 图 M2-125：接地圆柱壳内的中心线偶极子}
\end{center}
\end{minipage}\par

图 M2-125
的中心线偶极产生允许的轴线奇点，不能因为在原点发散就把真实源项删去。再补一个正则的一阶调和项，使圆柱内壁电势为零；内壁密度和中心附近感应场都应从这一补项及边界导数读取。

真实偶极奇点给出

\[
\frac{p_{\ell}}{2\pi\varepsilon_0\rho}\cos\phi.
\]

内壁感应电荷在原点附近产生正则势，且只需同一 \(m=1\) 角模：

\[
\varphi
=\frac{p_{\ell}}{2\pi\varepsilon_0\rho}\cos\phi
+A\rho\cos\phi.
\]

由 \(\varphi(R,\phi)=0\)，

\[
A=-\frac{p_{\ell}}{2\pi\varepsilon_0R^2}.
\]

所以

\[
\boxed{
\varphi
=\frac{p_{\ell}}{2\pi\varepsilon_0}
\left(
\frac1\rho-\frac{\rho}{R^2}
\right)\cos\phi
}.
\]

感应势为

\[
\varphi_{\rm ind}
=-\frac{p_{\ell}}{2\pi\varepsilon_0R^2}x,
\]

故

\[
\boxed{
\mathbf E_{\rm ind}
=\frac{p_{\ell}}{2\pi\varepsilon_0R^2}\mathbf e_x
}.
\]

方向沿 \(+x\)，与偶极矩同向。

壁面内侧的径向场为

\[
E_{\rho}(R^-,\phi)
=\frac{p_{\ell}}{\pi\varepsilon_0R^2}\cos\phi.
\]

内表面从导体指向真空的法向为 \(-\mathbf e_{\rho}\)，因此

\[
\boxed{
\sigma_{\rm in}(\phi)
=-\varepsilon_0E_{\rho}(R^-,\phi)
=-\frac{p_{\ell}}{\pi R^2}\cos\phi
}.
\]

右侧感应负电、左侧感应正电，且总感应电荷为零。

\subsubsection{作业四：同轴圆柱电容器}\label{m2:ux4f5cux4e1aux56dbux540cux8f74ux5706ux67f1ux7535ux5bb9ux5668}

\textbf{题号：M2-126}

\textbf{层级：} 复赛　　\textbf{方法：} 分离变量、边值问题

无限长同轴导体圆柱的半径分别为
\(a,b\)，\(0<a<b\)，两者之间为真空且没有体电荷。内、外圆柱分别保持常电势
\(V_a,V_b\)。求 \(a<\rho<b\)
的电势和电场，并由单位长度电荷与电势差的关系求单位长度电容；写明径向正方向及电势差的次序。

\par\noindent\begin{minipage}{\linewidth}
\textbf{解答}

\begin{center}
\includegraphics[width=\linewidth,height=.40\textheight,keepaspectratio]{\detokenize{assets/module-02/examples/M2-126.pdf}}
\par\smallskip{\small 图 M2-126：同轴圆柱电容器}
\end{center}
\end{minipage}\par

图 M2-126 的同轴几何消去方位角依赖，径向通解为
\(A+B\ln\rho\)。两圆分别给定
\(V_a,V_b\)，电场取其差的符号；把单位长度电荷除以电势差时，公用的对数半径比决定
\(C^{\prime}\)。

系统轴对称，只保留 \(m=0\) 模态：

\[
\varphi=A+B\ln\rho.
\]

由

\[
A+B\ln a=V_a,
\]

\[
A+B\ln b=V_b,
\]

相减得到

\[
B=\frac{V_b-V_a}{\ln(b/a)}.
\]

因此

\[
\boxed{
\varphi(\rho)
=V_a
+(V_b-V_a)
\frac{\ln(\rho/a)}{\ln(b/a)}
}.
\]

电场为

\[
\boxed{
\mathbf E
=-\frac{d\varphi}{d\rho}\mathbf e_{\rho}
=\frac{V_a-V_b}{\rho\ln(b/a)}\mathbf e_{\rho}
}.
\]

设内导体每单位长度电荷为 \(\lambda\)。由高斯定理

\[
E_{\rho}(\rho)\,2\pi\rho
=\frac{\lambda}{\varepsilon_0},
\]

所以

\[
\lambda
=\frac{2\pi\varepsilon_0(V_a-V_b)}{\ln(b/a)}.
\]

每单位长度电容为

\[
\boxed{
C'
=\frac{\lambda}{V_a-V_b}
=\frac{2\pi\varepsilon_0}{\ln(b/a)}
}.
\]

\subsubsection{作业五：任意阶外部背景模态}\label{m2:ux4f5cux4e1aux4e94ux4efbux610fux9636ux5916ux90e8ux80ccux666fux6a21ux6001}

\textbf{题号：M2-127}

\textbf{层级：} 复赛　　\textbf{方法：} 分离变量、边值问题

无限长接地导体圆柱半径为 \(R>0\)，柱外为与 \(z\)
无关的无源真空。给定远处背景势
\(\varphi_{\infty}=-K\rho^m\cos m\phi\)，其中 \(m\ge1\) 为整数、\(K\)
为有适当量纲的常量。要求感应附加势随 \(\rho\to\infty\)
衰减。求柱外电势和柱面自由面电荷密度，说明角向周期性与远处条件如何筛选径向项。

\par\noindent\begin{minipage}{\linewidth}
\textbf{解答}

\begin{center}
\includegraphics[width=\linewidth,height=.40\textheight,keepaspectratio]{\detokenize{assets/module-02/examples/M2-127.pdf}}
\par\smallskip{\small 图 M2-127：任意阶外部背景模态}
\end{center}
\end{minipage}\par

图 M2-127 的外部背景保留增长的 \(m\)
阶项，导体反应则选同阶衰减项。令两者在 \(\rho=R\)
恰好抵消，随后沿径向求导得到表面密度；背景阶数大于1时，不能把它误称为远处匀强场。

背景只有 \(m\) 阶余弦模态。导体产生的衰减修正必须具有相同角模，因此

\[
\varphi
=-K\rho^m\cos m\phi
+C\rho^{-m}\cos m\phi.
\]

由 \(\varphi(R,\phi)=0\)，

\[
-KR^m+CR^{-m}=0,
\]

所以

\[
C=KR^{2m}.
\]

圆柱外电势为

\[
\boxed{
\varphi
=-K
\left(
\rho^m-\frac{R^{2m}}{\rho^m}
\right)\cos m\phi
}.
\]

径向电场为

\[
\begin{aligned}
E_{\rho}
&=-\frac{\partial\varphi}{\partial\rho}\\
&=Km
\left(
\rho^{m-1}
+R^{2m}\rho^{-m-1}
\right)\cos m\phi.
\end{aligned}
\]

在 \(\rho=R\) 上，

\[
E_{\rho}(R^+,\phi)
=2mKR^{m-1}\cos m\phi.
\]

故表面电荷密度为

\[
\boxed{
\sigma(\phi)
=2m\varepsilon_0KR^{m-1}\cos m\phi
}.
\]

当 \(m=1\)、\(K=E_0\) 时，立即退化为匀强外场中的接地圆柱结果。

角模由源和边界确定，径向幂次由解域端点确定，系数由圆周边界逐模匹配。下一附录从变量分离、周期性和欧拉方程出发，推导这里使用的二维柱坐标通解。

\section{专题十九附录：二维柱坐标拉普拉斯方程通解的推导（拔高选讲）}\label{m2:ux4e13ux9898ux5341ux4e5dux9644ux5f55ux4e8cux7ef4ux67f1ux5750ux6807ux62c9ux666eux62c9ux65afux65b9ux7a0bux901aux89e3ux7684ux63a8ux5bfcux62d4ux9ad8ux9009ux8bb2}

本附录严格推导二维柱坐标下拉普拉斯方程的通解，并说明这套通解成立所需的几何条件。

设电荷分布、介质和边界条件沿 \(z\)
方向完全平移不变，且相应的定解条件足以保证解的唯一性，则

\[
\frac{\partial\varphi}{\partial z}=0.
\]

无源区域中的电势只依赖 \((\rho,\phi)\)，满足

\[
\boxed{
\frac1\rho\frac{\partial}{\partial\rho}
\left(
\rho\frac{\partial\varphi}{\partial\rho}
\right)
+\frac1{\rho^2}
\frac{\partial^2\varphi}{\partial\phi^2}
=0
}.
\]

严格推导需要依次回答四个问题：

\begin{enumerate}
\def\labelenumi{\arabic{enumi}.}
\tightlist
\item
  分离常数为何非负；
\item
  完整圆周中角向本征值为何是整数平方；
\item
  \(m=0\) 时为何出现 \(\ln\rho\)；
\item
  一个乘积解为何最终能扩展为四组独立系数。
\end{enumerate}

\subsection{一、变量分离与一般分离常数}\label{m2:ux4e00ux53d8ux91cfux5206ux79bbux4e0eux4e00ux822cux5206ux79bbux5e38ux6570}

先寻找非零乘积形式的特解

\[
\varphi(\rho,\phi)=R(\rho)\Phi(\phi).
\]

代入拉普拉斯方程：

\[
\frac{\Phi}{\rho}
\frac{d}{d\rho}
\left(
\rho\frac{dR}{d\rho}
\right)
+\frac{R}{\rho^2}
\frac{d^2\Phi}{d\phi^2}
=0.
\]

在 \(R\Phi\ne0\) 的局部区间内，两边乘以 \(\rho^2/(R\Phi)\)：

\[
\frac{\rho}{R}
\frac{d}{d\rho}
\left(
\rho\frac{dR}{d\rho}
\right)
+\frac{\Phi''}{\Phi}
=0.
\]

第一项只依赖 \(\rho\)，第二项只依赖 \(\phi\)。若等式对所有
\((\rho,\phi)\) 成立，则两项必须分别为常量。暂时不预设常量的符号，定义

\[
\boxed{
-\frac{\Phi''}{\Phi}
=\frac{\rho}{R}
\frac{d}{d\rho}
\left(
\rho\frac{dR}{d\rho}
\right)
=\lambda
}.
\]

于是得到角向本征方程

\[
\boxed{
-\Phi''=\lambda\Phi
},
\]

以及径向方程

\[
\boxed{
\rho^2R''+\rho R'-\lambda R=0
}.
\]

\par\noindent\begin{minipage}{\linewidth}
\begin{center}
\includegraphics[width=\linewidth,height=.40\textheight,keepaspectratio]{\detokenize{assets/module-02/original/tikz-m2-topic19-v8-09.png}}
\par\smallskip{\small 图示：二维柱坐标拉普拉斯方程分离为角向本征问题和径向响应问题}
\end{center}
\end{minipage}\par

在 \(\Phi\) 或 \(R\)
的孤立零点处，上述除法形式不能直接使用，但分离所得常微分方程的解可由连续性延拓过去，因此不影响最终本征函数。

\subsection{二、完整圆周中的角向本征谱}\label{m2:ux4e8cux5b8cux6574ux5706ux5468ux4e2dux7684ux89d2ux5411ux672cux5f81ux8c31}

本节首先假定求解域包含完整圆周，即 \(\phi\) 与 \(\phi+2\pi\)
表示同一空间方向。电势是单值标量，因此

\[
\Phi(\phi+2\pi)=\Phi(\phi).
\]

对光滑本征函数，其导数也必须首尾相接：

\[
\Phi'(\phi+2\pi)=\Phi'(\phi).
\]

\subsubsection{1.
分离常数必定非负}\label{m2:ux5206ux79bbux5e38ux6570ux5fc5ux5b9aux975eux8d1f}

从

\[
-\Phi''=\lambda\Phi
\]

出发，设 \(\Phi\) 为实函数。两边乘以 \(\Phi\) 并在一个周期上积分：

\[
-\int_0^{2\pi}\Phi\Phi''\,d\phi
=\lambda\int_0^{2\pi}\Phi^2\,d\phi.
\]

对左边分部积分：

\[
-\int_0^{2\pi}\Phi\Phi''\,d\phi
=-\left[\Phi\Phi'\right]_0^{2\pi}
+\int_0^{2\pi}(\Phi')^2\,d\phi.
\]

周期边界使边界项为零，所以

\[
\boxed{
\lambda
=\frac{\displaystyle\int_0^{2\pi}(\Phi')^2\,d\phi}
{\displaystyle\int_0^{2\pi}\Phi^2\,d\phi}
\ge0
}.
\]

只要 \(\Phi\) 不是恒零函数，分母为正。由此严格排除 \(\lambda<0\)
的非平凡周期解。

\par\noindent\begin{minipage}{\linewidth}
\begin{center}
\includegraphics[width=\linewidth,height=.40\textheight,keepaspectratio]{\detokenize{assets/module-02/original/tikz-m2-topic19-v8-10.png}}
\par\smallskip{\small 图示：负分离常数、零分离常数和正分离常数在周期条件下的不同结局}
\end{center}
\end{minipage}\par

\subsubsection{2. 零本征值}\label{m2:ux96f6ux672cux5f81ux503c}

当 \(\lambda=0\) 时，

\[
\Phi''=0,
\]

故

\[
\Phi=A+B\phi.
\]

周期性要求

\[
A+B(\phi+2\pi)=A+B\phi,
\]

所以 \(B=0\)。零模的角向函数只能是常量。

\textbf{例 M2-128：零角模对应对数径向解}

在二维圆环 \(a<\rho<b\) 中，\(0<a<b\)，无源电势与方位角无关，满足
\(\rho^{-1}(\rho\varphi')'=0\)。直接求这个零角模的径向通解，并与非零角模
\(m\ge1\) 的两支幂函数作比较；解释为何不能把非零阶公式简单令 \(m=0\)
后仍当作两个独立解。

\par\noindent\begin{minipage}{\linewidth}
\textbf{解答}

\begin{center}
\includegraphics[width=\linewidth,height=.40\textheight,keepaspectratio]{\detokenize{assets/module-02/examples/M2-128.pdf}}
\par\smallskip{\small 图 M2-128：零角模对应对数径向解}
\end{center}
\end{minipage}\par

图 M2-128 的两条同心圆不随方位变化，因而属于 \(m=0\)。径向方程积分一次给
\(\rho\varphi^{\prime}\)
为常数，再积分才得到对数；非零角模的两种幂在这里合并，需单独处理零模。

同轴圆柱电容器的电势与 \(\phi\) 无关，正是 \(m=0\)
角模；其径向部分不是幂函数，而是下一节得到的 \(A+B\ln\rho\)。

\subsubsection{3.
正本征值与整数谱}\label{m2:ux6b63ux672cux5f81ux503cux4e0eux6574ux6570ux8c31}

当 \(\lambda>0\) 时，记

\[
\lambda=\nu^2,
\qquad \nu>0.
\]

角向通解为

\[
\Phi(\phi)
=A\cos\nu\phi+B\sin\nu\phi.
\]

在 \(\phi=0\) 与 \(\phi=2\pi\) 处匹配函数和导数，可写成

\[
\begin{pmatrix}
\cos2\pi\nu&\sin2\pi\nu\\
-\sin2\pi\nu&\cos2\pi\nu
\end{pmatrix}
\binom AB
=\binom AB.
\]

若 \((A,B)\ne(0,0)\)，系数矩阵减去单位矩阵后的行列式必须为零：

\[
\begin{aligned}
0
&=
\begin{vmatrix}
\cos2\pi\nu-1&\sin2\pi\nu\\
-\sin2\pi\nu&\cos2\pi\nu-1
\end{vmatrix}\\
&=(\cos2\pi\nu-1)^2+\sin^22\pi\nu\\
&=2(1-\cos2\pi\nu).
\end{aligned}
\]

因此

\[
\cos2\pi\nu=1,
\]

即

\[
\nu\in\mathbb Z.
\]

正负整数产生相同的实函数空间，所以只需取

\[
\boxed{
m=1,2,3,\ldots
}.
\]

对应角向本征函数为

\[
\boxed{
\cos m\phi,\qquad \sin m\phi
}.
\]

\par\noindent\begin{minipage}{\linewidth}
\begin{center}
\includegraphics[width=\linewidth,height=.40\textheight,keepaspectratio]{\detokenize{assets/module-02/original/tikz-m2-topic19-v8-11.png}}
\par\smallskip{\small 图示：方位角转过二π回到同一点使角向本征值成为整数}
\end{center}
\end{minipage}\par

\subsubsection{4. 正交性}\label{m2:ux6b63ux4ea4ux6027}

不同角模在 \([0,2\pi]\) 上正交：

\[
\int_0^{2\pi}
\cos m\phi\cos n\phi\,d\phi
=\pi\delta_{mn},
\qquad m,n\ge1,
\]

\[
\int_0^{2\pi}
\sin m\phi\sin n\phi\,d\phi
=\pi\delta_{mn},
\]

\[
\int_0^{2\pi}
\sin m\phi\cos n\phi\,d\phi
=0.
\]

零模满足

\[
\int_0^{2\pi}1\,d\phi=2\pi,
\]

并且与所有非零角模正交。正交性使圆周边界上的每个傅里叶模态可以独立匹配。

\subsection{三、径向方程与对数坐标}\label{m2:ux4e09ux5f84ux5411ux65b9ux7a0bux4e0eux5bf9ux6570ux5750ux6807}

对 \(m\ge0\)，径向方程为

\[
\rho^2R''+\rho R'-m^2R=0.
\]

令

\[
s=\ln\rho,
\qquad
\rho=e^s.
\]

把 \(R(\rho)\) 看成 \(R(s)\)，则

\[
\frac{dR}{d\rho}
=\frac1\rho\frac{dR}{ds},
\]

\[
\frac{d^2R}{d\rho^2}
=\frac1{\rho^2}
\left(
\frac{d^2R}{ds^2}-\frac{dR}{ds}
\right).
\]

因此

\[
\rho^2R''+\rho R'
=\frac{d^2R}{ds^2}.
\]

径向方程化为常系数方程

\[
\boxed{
\frac{d^2R}{ds^2}-m^2R=0
}.
\]

\par\noindent\begin{minipage}{\linewidth}
\begin{center}
\includegraphics[width=\linewidth,height=.40\textheight,keepaspectratio]{\detokenize{assets/module-02/original/tikz-m2-topic19-v8-12.png}}
\par\smallskip{\small 图示：对数坐标把径向半轴映成实轴并把欧拉型方程化为常系数方程}
\end{center}
\end{minipage}\par

\subsubsection{1. 非零角模}\label{m2:ux975eux96f6ux89d2ux6a21}

当 \(m\ge1\) 时，

\[
R=Ae^{ms}+Be^{-ms}.
\]

代回 \(s=\ln\rho\)：

\[
\boxed{
R_m(\rho)
=A_m\rho^m+B_m\rho^{-m}
}.
\]

\subsubsection{2. 零角模}\label{m2:ux96f6ux89d2ux6a21}

当 \(m=0\) 时，

\[
\frac{d^2R}{ds^2}=0,
\]

所以

\[
R=A+Bs.
\]

恢复 \(\rho\)：

\[
\boxed{
R_0(\rho)
=A_0+B_0\ln\rho
}.
\]

这比``重根时补一个对数解''的口诀更直接：\(\ln\rho\)
只是对数坐标中线性函数 \(s\) 的原变量表达。

\subsubsection{3. 端点行为}\label{m2:ux7aefux70b9ux884cux4e3a}

\par\noindent\begin{minipage}{\linewidth}
\begin{center}
\includegraphics[width=\linewidth,height=.40\textheight,keepaspectratio]{\detokenize{assets/module-02/original/tikz-m2-topic19-v8-13.png}}
\par\smallskip{\small 图示：正次幂、负次幂和对数零模在原点与无穷远处的行为}
\end{center}
\end{minipage}\par

{\def\LTcaptype{none} % do not increment counter
\begin{longtable}[]{@{}llll@{}}
\toprule\noalign{}
径向函数 & \(\rho\to0\) & \(\rho\to\infty\) & 常见物理来源 \\
\midrule\noalign{}
\endhead
\bottomrule\noalign{}
\endlastfoot
\(1\) & 有限 & 有限 & 零阶参考电势 \\
\(\ln\rho\) & 发散 & 发散 & 轴上线电荷或圆环区域零模 \\
\(\rho^m\) & 趋于零 & 增长 & 外部边界或给定远场 \\
\(\rho^{-m}\) & 发散 & 趋于零 & 原点附近二维多极源或外域衰减修正 \\
\end{longtable}
}

``发散''本身不是删除某项的充分理由。若发散点位于真实求解域之外，或者它就是题目给定的真实源，该项必须保留。

\subsubsection{4.
对数项与线电荷奇点}\label{m2:ux5bf9ux6570ux9879ux4e0eux7ebfux7535ux8377ux5947ux70b9}

在 \(\rho\ne0\) 处，

\[
\nabla_{\perp}^2\ln\rho=0.
\]

但以原点为中心取半径 \(a\) 的小圆，外法向导数为

\[
\frac{\partial\ln\rho}{\partial n}
=\frac1a.
\]

因此

\[
\oint_{\rho=a}
\nabla\ln\rho\cdot\mathbf n\,dl
=\frac1a(2\pi a)
=2\pi.
\]

所以在分布意义下

\[
\boxed{
\nabla_{\perp}^2\ln\rho
=2\pi\delta^{(2)}(\boldsymbol\rho)
}.
\]

无限长线电荷的电势

\[
\varphi
=-\frac{\lambda_{\ell}}{2\pi\varepsilon_0}
\ln\frac{\rho}{\rho_0}
\]

满足

\[
\nabla_{\perp}^2\varphi
=-\frac{\lambda_{\ell}}{\varepsilon_0}
\delta^{(2)}(\boldsymbol\rho),
\]

正好对应三维泊松方程在每单位轴向长度上的二维形式。

\subsection{四、为何每个非零角模有四个独立系数}\label{m2:ux56dbux4e3aux4f55ux6bcfux4e2aux975eux96f6ux89d2ux6a21ux6709ux56dbux4e2aux72ecux7acbux7cfbux6570}

若只写一个乘积

\[
\left(A\rho^m+B\rho^{-m}\right)
\left(C\cos m\phi+D\sin m\phi\right),
\]

展开后四个系数为 \(AC,AD,BC,BD\)，它们彼此受同一组 \(A,B,C,D\)
的乘积关系约束。这只是一个分离特解族，还不是该本征值下最一般的线性组合。

拉普拉斯方程是线性的。对同一个 \(m\ge1\)，以下四个函数分别都是解：

\[
\rho^m\cos m\phi,
\qquad
\rho^m\sin m\phi,
\]

\[
\rho^{-m}\cos m\phi,
\qquad
\rho^{-m}\sin m\phi.
\]

它们线性独立，因此必须允许四个独立系数：

\[
\boxed{
\begin{aligned}
\varphi_m(\rho,\phi)
={}&
\rho^m
\left(
a_m\cos m\phi+b_m\sin m\phi
\right)\\
&+\rho^{-m}
\left(
c_m\cos m\phi+d_m\sin m\phi
\right).
\end{aligned}
}
\]

\par\noindent\begin{minipage}{\linewidth}
\begin{center}
\includegraphics[width=\linewidth,height=.40\textheight,keepaspectratio]{\detokenize{assets/module-02/original/tikz-m2-topic19-v8-14.png}}
\par\smallskip{\small 图示：每个非零角模由两个径向解和两个角向解组成四个独立基函数}
\end{center}
\end{minipage}\par

再加入 \(m=0\) 的两个径向解，得到完整圆周上的二维柱坐标通解：

\[
\boxed{
\begin{aligned}
\varphi(\rho,\phi)
={}&A_0+B_0\ln\rho\\
&+\sum_{m=1}^{\infty}
\left[
\rho^m(a_m\cos m\phi+b_m\sin m\phi)
+\rho^{-m}(c_m\cos m\phi+d_m\sin m\phi)
\right].
\end{aligned}
}
\]

这套展开适用于不跨越源和材料界面的连通无源区域。若区域被面电荷、介质界面或其他奇点分割，应在每个光滑子区域分别展开，再用界面条件连接。

\subsection{五、傅里叶完备性与边界系数}\label{m2:ux4e94ux5085ux91ccux53f6ux5b8cux5907ux6027ux4e0eux8fb9ux754cux7cfbux6570}

设圆周 \(\rho=R\) 上的第一类边界数据为 \(f(\phi)\)。若 \(f\)
分段光滑，它可展开为

\[
f(\phi)
=F_0+\sum_{m=1}^{\infty}
\left(
F_m\cos m\phi+G_m\sin m\phi
\right),
\]

其中

\[
\boxed{
F_0
=\frac1{2\pi}
\int_0^{2\pi}f(\phi)\,d\phi
},
\]

\[
\boxed{
F_m
=\frac1\pi
\int_0^{2\pi}
f(\phi)\cos m\phi\,d\phi
},
\]

\[
\boxed{
G_m
=\frac1\pi
\int_0^{2\pi}
f(\phi)\sin m\phi\,d\phi
}.
\]

\par\noindent\begin{minipage}{\linewidth}
\begin{center}
\includegraphics[width=\linewidth,height=.40\textheight,keepaspectratio]{\detokenize{assets/module-02/original/tikz-m2-topic19-v8-15.png}}
\par\smallskip{\small 图示：圆周边界函数通过正交投影得到各个傅里叶模态系数}
\end{center}
\end{minipage}\par

\subsubsection{1.
圆盘内的正则延拓}\label{m2:ux5706ux76d8ux5185ux7684ux6b63ux5219ux5ef6ux62d3}

若 \(0\le\rho<R\) 内无源且原点有限，只保留 \(1\) 和 \(\rho^m\)：

\[
\boxed{
\varphi(\rho,\phi)
=F_0+\sum_{m=1}^{\infty}
\left(\frac{\rho}{R}\right)^m
\left(
F_m\cos m\phi+G_m\sin m\phi
\right)
}.
\]

当 \(\rho<R\) 时，因子 \((\rho/R)^m\)
使高阶项快速衰减，因此即使边界函数有尖锐变化，内部任一点的级数通常比边界级数收敛得更好。

\textbf{例 M2-129：三阶边界与三次径向幂配对}

二维无源真空圆盘半径为 \(R>0\)，圆心为原点，\(\phi\) 为从 \(+x\)
轴起算的方位角。圆周电势为
\(\varphi(R,\phi)=V_0+V_3\cos3\phi\)，要求盘内无奇性且电势单值。求盘内电势，并指出常值部分和三阶角模各自的径向形式。

\par\noindent\begin{minipage}{\linewidth}
\textbf{解答}

\begin{center}
\includegraphics[width=\linewidth,height=.40\textheight,keepaspectratio]{\detokenize{assets/module-02/examples/M2-129.pdf}}
\par\smallskip{\small 图 M2-129：三阶边界与三次径向幂配对}
\end{center}
\end{minipage}\par

图 M2-129
的三组正负扇区表示三阶角模，另有未着色表示的常量平均值。正则性分别留下常量与
\(\rho^3\)，在圆周匹配系数后，向圆心靠近时三阶偏差迅速消失。

若

\[
f(\phi)=V_0+V_3\cos3\phi,
\]

则

\[
\boxed{
\varphi(\rho,\phi)
=V_0+V_3
\left(\frac{\rho}{R}\right)^3
\cos3\phi
}.
\]

\subsubsection{2.
圆盘的泊松积分公式}\label{m2:ux5706ux76d8ux7684ux6ccaux677eux79efux5206ux516cux5f0f}

令

\[
q=\frac{\rho}{R}<1.
\]

利用恒等式

\[
1+2\sum_{m=1}^{\infty}q^m\cos m\theta
=\frac{1-q^2}{1-2q\cos\theta+q^2},
\]

可把傅里叶级数写成积分形式：

\[
\boxed{
\varphi(\rho,\phi)
=\frac1{2\pi}
\int_0^{2\pi}
\frac{1-q^2}
{1-2q\cos(\phi-\phi')+q^2}
f(\phi')\,d\phi'
}.
\]

积分核

\[
P_q(\theta)
=\frac{1-q^2}{1-2q\cos\theta+q^2}
\]

称为圆盘的泊松核。它为正，且满足

\[
\frac1{2\pi}
\int_0^{2\pi}P_q(\theta)\,d\theta=1.
\]

因此圆盘内部电势是边界电势的正权平均，这与调和函数的最大值原理一致。

若 \(f\) 在某点连续，则 \(\rho\to R^-\) 时内部解趋于该点的边界值；若
\(f\)
在某点发生跳变，傅里叶级数在该点趋于左右极限的平均值。实际导体边界的尖角或电势跳变处往往伴随场强奇异，需要单独分析。

\subsection{六、适用范围：完整圆周与劈形不能混用}\label{m2:ux516dux9002ux7528ux8303ux56f4ux5b8cux6574ux5706ux5468ux4e0eux5288ux5f62ux4e0dux80fdux6df7ux7528}

整数 \(m\) 来自 \(2\pi\)
周期条件，只适用于完整绕轴一周后回到同一点的角域。

若求解域是劈形

\[
0<\phi<\alpha,
\]

两条射线 \(\phi=0\) 与 \(\phi=\alpha\) 是不同的物理边界，不需要满足
\(2\pi\) 周期性。

\par\noindent\begin{minipage}{\linewidth}
\begin{center}
\includegraphics[width=\linewidth,height=.40\textheight,keepaspectratio]{\detokenize{assets/module-02/original/tikz-m2-topic19-v8-16.png}}
\par\smallskip{\small 图示：完整圆周由周期条件产生整数角模，接地劈形由两侧边界产生另一套本征值}
\end{center}
\end{minipage}\par

\subsubsection{1.
两侧接地的劈形}\label{m2:ux4e24ux4fa7ux63a5ux5730ux7684ux5288ux5f62}

若

\[
\Phi(0)=\Phi(\alpha)=0,
\]

角向本征函数为

\[
\boxed{
\Phi_n(\phi)
=\sin\frac{n\pi\phi}{\alpha},
\qquad n=1,2,\ldots
},
\]

对应

\[
\boxed{
\nu_n=\frac{n\pi}{\alpha}
}.
\]

径向方程随之变为

\[
\rho^2R''+\rho R'-\nu_n^2R=0,
\]

所以

\[
R=A\rho^{\nu_n}+B\rho^{-\nu_n}.
\]

当 \(\alpha\) 不是 \(\pi\) 的特殊分数时，\(\nu_n\)
一般不是整数。这并不违反单值性，因为劈形区域并未包含完整圆周。

\textbf{例 M2-130：劈形可以容许非整数角向阶数}

真空无源劈形在二维极坐标中为
\(0<r<a\)、\(0<\phi<2\pi/3\)，两条径向边接地，要求电势在顶点有限；外圆弧上的非零数据可按角向本征模展开。求满足两条齐次角边界的最低非零角向阶数，并解释这一劈形内的角函数为何不必满足整个平面上的
\(2\pi\) 周期条件。

\par\noindent\begin{minipage}{\linewidth}
\textbf{解答}

\begin{center}
\includegraphics[width=\linewidth,height=.40\textheight,keepaspectratio]{\detokenize{assets/module-02/examples/M2-130.pdf}}
\par\smallskip{\small 图 M2-130：劈形可以容许非整数角向阶数}
\end{center}
\end{minipage}\par

图 M2-130 的两条径向边界相隔
\(120^{\circ}\)，两边零势使最低角因子在两端各有一个零点。由角跨度得到
\(\nu=3/2\)，这一非整数分支在扇区内单值，但不能不加支割地延伸为全圆周的周期势。

若

\[
\alpha=\frac{2\pi}{3},
\]

最低角向本征值为

\[
\nu_1=\frac{\pi}{\alpha}=\frac32.
\]

相应的正则调和函数

\[
\rho^{3/2}\sin\frac{3\phi}{2}
\]

在该劈形内部完全合法，但它不能作为完整圆周上的单值势函数。

\subsubsection{2.
边界电势不同时的零型解}\label{m2:ux8fb9ux754cux7535ux52bfux4e0dux540cux65f6ux7684ux96f6ux578bux89e3}

若两条劈形边界分别保持

\[
\varphi(\rho,0)=0,
\qquad
\varphi(\rho,\alpha)=V_0,
\]

可取与 \(\rho\) 无关的解

\[
\boxed{
\varphi(\rho,\phi)
=\frac{V_0}{\alpha}\phi
}.
\]

在完整圆周中，线性角项 \(C\phi\) 会破坏单值性；在劈形中，\(\phi=0\) 和
\(\phi=\alpha\) 是不同边界，因此该解合法。

\subsection{七、存在轴向变化时：普通与修正贝塞尔方程}\label{m2:ux4e03ux5b58ux5728ux8f74ux5411ux53d8ux5316ux65f6ux666eux901aux4e0eux4feeux6b63ux8d1dux585eux5c14ux65b9ux7a0b}

若 \(\partial\varphi/\partial z\ne0\)，三维拉普拉斯方程为

\[
\frac1\rho\frac{\partial}{\partial\rho}
\left(
\rho\frac{\partial\varphi}{\partial\rho}
\right)
+\frac1{\rho^2}
\frac{\partial^2\varphi}{\partial\phi^2}
+\frac{\partial^2\varphi}{\partial z^2}
=0.
\]

令

\[
\varphi=R(\rho)\Phi(\phi)Z(z),
\]

角向仍取 \(m\) 阶模态。

\subsubsection{1.
轴向为振荡函数}\label{m2:ux8f74ux5411ux4e3aux632fux8361ux51fdux6570}

若

\[
Z''+k^2Z=0,
\]

即 \(Z\) 为正弦、余弦或 \(e^{ikz}\)，径向方程为

\[
\boxed{
\rho^2R''+\rho R'
-\left(
m^2+k^2\rho^2
\right)R
=0
}.
\]

这是修正贝塞尔方程，解为

\[
I_m(k\rho),
\qquad
K_m(k\rho).
\]

\subsubsection{2.
轴向为指数函数}\label{m2:ux8f74ux5411ux4e3aux6307ux6570ux51fdux6570}

若

\[
Z''-k^2Z=0,
\]

径向方程为

\[
\boxed{
\rho^2R''+\rho R'
+\left(
k^2\rho^2-m^2
\right)R
=0
}.
\]

这是普通贝塞尔方程，解为

\[
J_m(k\rho),
\qquad
Y_m(k\rho).
\]

因此，``存在 \(z\)
方向变化就出现贝塞尔函数''只能作为总提示；具体使用普通贝塞尔函数还是修正贝塞尔函数，取决于轴向分离常数的符号。

\subsection{八、推导结果的检查}\label{m2:ux516bux63a8ux5bfcux7ed3ux679cux7684ux68c0ux67e5}

得到候选展开后，应逐项检查：

\begin{enumerate}
\def\labelenumi{\arabic{enumi}.}
\tightlist
\item
  每个基函数在所讨论的无源子区域内是否满足拉普拉斯方程；
\item
  完整圆周问题是否满足 \(2\pi\) 周期与单值性；
\item
  劈形问题是否误用了整数角模；
\item
  原点或无穷远处的发散是否与真实源、远场和解域一致；
\item
  \(m=0\) 的常数项和对数项是否遗漏；
\item
  同一 \(m\) 的四个独立系数是否被错误绑定；
\item
  边界函数的傅里叶级数是否在跳变点需要取平均值；
\item
  若含 \(z\) 变化，径向方程的贝塞尔类型是否与分离常数符号一致。
\end{enumerate}

\subsection{九、练习与解答}\label{m2:ux4e5dux7ec3ux4e60ux4e0eux89e3ux7b54-3}

\subsubsection{作业一：用积分证明角向本征值非负}\label{m2:ux4f5cux4e1aux4e00ux7528ux79efux5206ux8bc1ux660eux89d2ux5411ux672cux5f81ux503cux975eux8d1f}

\textbf{题号：M2-131}

\textbf{层级：} 决赛　　\textbf{方法：} 边值问题

设非零实函数 \(\Phi\) 在 \([0,2\pi]\) 上二阶连续可微，满足角向本征方程
\(-\Phi''=\lambda\Phi\)，以及周期条件
\(\Phi(0)=\Phi(2\pi)\)、\(\Phi'(0)=\Phi'(2\pi)\)。通过分部积分证明实本征值
\(\lambda\ge0\)；讨论 \(\lambda=0\)
时的函数形式，并说明两个周期条件在证明中各起什么作用。

\par\noindent\begin{minipage}{\linewidth}
\textbf{解答}

\begin{center}
\includegraphics[width=\linewidth,height=.40\textheight,keepaspectratio]{\detokenize{assets/module-02/examples/M2-131.pdf}}
\par\smallskip{\small 图 M2-131：用积分证明角向本征值非负}
\end{center}
\end{minipage}\par

图 M2-131
的角坐标在一整圈闭合，函数值和一阶导数都周期相接。将本征方程乘原函数并分部积分，首尾边界项成对抵消，留下梯度平方积分；非零函数的模方积分为正，因此本征值非负。

两边乘以 \(\Phi\) 并积分：

\[
-\int_0^{2\pi}\Phi\Phi''\,d\phi
=\lambda\int_0^{2\pi}\Phi^2\,d\phi.
\]

分部积分得

\[
-\left[\Phi\Phi'\right]_0^{2\pi}
+\int_0^{2\pi}(\Phi')^2\,d\phi
=\lambda\int_0^{2\pi}\Phi^2\,d\phi.
\]

周期边界使第一项为零，因此

\[
\lambda
=\frac{\displaystyle\int_0^{2\pi}(\Phi')^2\,d\phi}
{\displaystyle\int_0^{2\pi}\Phi^2\,d\phi}
\ge0.
\]

若 \(\lambda=0\)，则

\[
\int_0^{2\pi}(\Phi')^2\,d\phi=0.
\]

被积函数非负，所以必须有 \(\Phi'=0\)，即

\[
\boxed{
\Phi=\text{常量}
}.
\]

\subsubsection{作业二：圆盘内外的三阶模态}\label{m2:ux4f5cux4e1aux4e8cux5706ux76d8ux5185ux5916ux7684ux4e09ux9636ux6a21ux6001}

\textbf{题号：M2-132}

\textbf{层级：} 决赛　　\textbf{方法：} 边值问题

半径为 \(R>0\) 的圆周上给定二维电势
\(V_0\cos3\phi\)。分别求两个独立的真空无源问题：圆盘内电势在原点正则；圆外
\(\rho>R\)
的电势在无穷远衰减。要求两种解均单值，分别核对其在同一圆周上的边界值，说明为何内外径向幂次不同。

\par\noindent\begin{minipage}{\linewidth}
\textbf{解答}

\begin{center}
\includegraphics[width=\linewidth,height=.40\textheight,keepaspectratio]{\detokenize{assets/module-02/examples/M2-132.pdf}}
\par\smallskip{\small 图 M2-132：圆盘内外的三阶模态}
\end{center}
\end{minipage}\par

图 M2-132 的同一圆周三阶图样在内外有不同径向延伸。内域用 \(\rho^3\)
避免中心奇点，外域用 \(\rho^{-3}\) 满足远端衰减，再把两者分别归一到
\(\rho=R\)，即可逐点检查边界一致。

边界只含 \(m=3\) 余弦模态。

圆盘内必须保留正次幂并按边界归一化：

\[
\boxed{
\varphi_{\rm in}
=V_0
\left(
\frac{\rho}{R}
\right)^3
\cos3\phi
}.
\]

圆外必须保留衰减的负次幂：

\[
\boxed{
\varphi_{\rm out}
=V_0
\left(
\frac{R}{\rho}
\right)^3
\cos3\phi
}.
\]

两式分别由 \(\rho^3\cos3\phi\) 和 \(\rho^{-3}\cos3\phi\)
构成，均在各自无源域内调和。在 \(\rho=R\) 上，

\[
\varphi_{\rm in}(R,\phi)
=\varphi_{\rm out}(R,\phi)
=V_0\cos3\phi.
\]

\subsubsection{作业三：非整数角模的劈形}\label{m2:ux4f5cux4e1aux4e09ux975eux6574ux6570ux89d2ux6a21ux7684ux5288ux5f62}

\textbf{题号：M2-133}

\textbf{层级：} 决赛　　\textbf{方法：} 边值问题

二维真空无源域为 \(0<\rho<a\)、\(0<\phi<2\pi/3\)，\(a>0\)。两条径向边
\(\phi=0,2\pi/3\) 接地，圆弧 \(\rho=a\) 上给定
\(\varphi=V_0\sin(3\phi/2)\)，顶点处要求电势有限。用分离变量求域内电势，逐一核对三条边界和顶点条件。

\par\noindent\begin{minipage}{\linewidth}
\textbf{解答}

\begin{center}
\includegraphics[width=\linewidth,height=.40\textheight,keepaspectratio]{\detokenize{assets/module-02/examples/M2-133.pdf}}
\par\smallskip{\small 图 M2-133：非整数角模的劈形}
\end{center}
\end{minipage}\par

图 M2-133 的扇区角宽为
\(2\pi/3\)，弧边数据正好匹配最低非整数正弦角模。径向正则性选
\(\rho^{3/2}\)，再用外弧半径 \(a\)
归一；该模态在两条径向接地边上均为零。

劈形角为

\[
\alpha=\frac{2\pi}{3}.
\]

接地侧边的角向本征值为

\[
\nu_n=\frac{n\pi}{\alpha}
=\frac{3n}{2}.
\]

边界数据只含最低模态 \(n=1\)：

\[
\Phi_1(\phi)
=\sin\frac{3\phi}{2}.
\]

原点有限要求保留正次幂 \(\rho^{3/2}\)，故

\[
\varphi
=C\rho^{3/2}
\sin\frac{3\phi}{2}.
\]

由 \(\rho=a\) 处的边界条件，

\[
Ca^{3/2}=V_0.
\]

所以

\[
\boxed{
\varphi(\rho,\phi)
=V_0
\left(
\frac{\rho}{a}
\right)^{3/2}
\sin\frac{3\phi}{2}
}.
\]

在 \(\phi=0\) 和 \(2\pi/3\) 上正弦因子为零，在 \(\rho=a\)
上恢复给定圆弧边界。指数 \(3/2\) 不是整数，但在该劈形域内电势单值。

\subsubsection{作业四：对数势的线源强度}\label{m2:ux4f5cux4e1aux56dbux5bf9ux6570ux52bfux7684ux7ebfux6e90ux5f3aux5ea6}

\textbf{题号：M2-134}

\textbf{层级：} 决赛　　\textbf{方法：} 边值问题

无限长线源沿 \(z\) 轴分布，真空中考察与 \(z\)、方位角无关的势
\(\varphi(\rho)=B\ln(\rho/\rho_*)\)，其中 \(\rho_*>0\)
是任意参考长度。先在 \(\rho>0\)
处计算普通拉普拉斯，再取半径趋零的同轴高斯圆柱，求穿过单位长度侧面的电场通量；由此解释轴线上分布源的含义，并求它代表线电荷密度
\(\lambda_{\ell}\) 时的 \(B\)。

\par\noindent\begin{minipage}{\linewidth}
\textbf{解答}

\begin{center}
\includegraphics[width=\linewidth,height=.40\textheight,keepaspectratio]{\detokenize{assets/module-02/examples/M2-134.pdf}}
\par\smallskip{\small 图 M2-134：对数势的线源强度}
\end{center}
\end{minipage}\par

图 M2-134
的小圆围住被排除的轴线，普通拉普拉斯在圆外为零。由负梯度得到径向场后，圆周通量仍是与半径无关的常量；把这一常量与单位长度线电荷对应，才固定对数系数的符号。

在 \(\rho\ne0\) 处，

\[
\nabla_{\perp}^2\varphi
=\frac1\rho\frac{d}{d\rho}
\left(
\rho\frac{d}{d\rho}B\ln\rho
\right).
\]

因为

\[
\frac{d}{d\rho}B\ln\rho=\frac{B}{\rho},
\]

所以

\[
\rho\frac{d\varphi}{d\rho}=B,
\]

从而

\[
\nabla_{\perp}^2\varphi=0,
\qquad \rho\ne0.
\]

小圆上的梯度通量为

\[
\oint\nabla\varphi\cdot\mathbf n\,dl
=\frac{B}{a}(2\pi a)
=2\pi B.
\]

因此

\[
\boxed{
\nabla_{\perp}^2\varphi
=2\pi B\delta^{(2)}(\boldsymbol\rho)
}.
\]

线电荷对应二维泊松方程

\[
\nabla_{\perp}^2\varphi
=-\frac{\lambda_{\ell}}{\varepsilon_0}
\delta^{(2)}(\boldsymbol\rho).
\]

比较系数：

\[
2\pi B=-\frac{\lambda_{\ell}}{\varepsilon_0}.
\]

所以

\[
\boxed{
B=-\frac{\lambda_{\ell}}{2\pi\varepsilon_0}
}.
\]

\subsubsection{作业五：轴向分离常数与贝塞尔类型}\label{m2:ux4f5cux4e1aux4e94ux8f74ux5411ux5206ux79bbux5e38ux6570ux4e0eux8d1dux585eux5c14ux7c7bux578b}

\textbf{题号：M2-135}

\textbf{层级：} 决赛　　\textbf{方法：} 边值问题

在三维无源真空区采用柱坐标，设分离形式
\(\varphi(\rho,\phi,z)=R(\rho)\cos(m\phi)Z(z)\)，\(m\ge0\)
为整数，\(k>0\)。分别对 \(Z''+k^2Z=0\) 与 \(Z''-k^2Z=0\)
两种情况推导径向方程，判断径向解属于普通贝塞尔函数还是修正贝塞尔函数；说明函数族的选择由哪个分离常数符号决定。

\par\noindent\begin{minipage}{\linewidth}
\textbf{解答}

\begin{center}
\includegraphics[width=\linewidth,height=.40\textheight,keepaspectratio]{\detokenize{assets/module-02/examples/M2-135.pdf}}
\par\smallskip{\small 图 M2-135：轴向分离常数与贝塞尔类型}
\end{center}
\end{minipage}\par

图 M2-135 将纵向振荡与纵向指数两种情形分开。把各自的
\(Z^{\prime\prime}/Z\) 代回原拉普拉斯方程，径向方程中的 \(k^2\rho^2\)
项随之反号，所以普通与修正贝塞尔类型不能只按函数名称记忆。

代入三维拉普拉斯方程并除以 \(R\cos m\phi\,Z\)：

\[
\frac1R
\frac1\rho\frac{d}{d\rho}
\left(
\rho\frac{dR}{d\rho}
\right)
-\frac{m^2}{\rho^2}
+\frac{Z''}{Z}
=0.
\]

第一种情况下，

\[
\frac{Z''}{Z}=-k^2.
\]

因此

\[
\frac1\rho\frac{d}{d\rho}
\left(
\rho\frac{dR}{d\rho}
\right)
-\frac{m^2}{\rho^2}R
-k^2R
=0.
\]

乘以 \(\rho^2\)：

\[
\boxed{
\rho^2R''+\rho R'
-\left(
m^2+k^2\rho^2
\right)R
=0
}.
\]

这是修正贝塞尔方程，解为

\[
\boxed{
R=A I_m(k\rho)+B K_m(k\rho)
}.
\]

第二种情况下，

\[
\frac{Z''}{Z}=k^2.
\]

同理得到

\[
\boxed{
\rho^2R''+\rho R'
+\left(
k^2\rho^2-m^2
\right)R
=0
}.
\]

这是普通贝塞尔方程，解为

\[
\boxed{
R=A J_m(k\rho)+B Y_m(k\rho)
}.
\]

具体保留哪一个径向函数，还要继续根据原点、无穷远和圆柱边界条件判断。

本附录完成了二维柱坐标通解从分离常数、角向谱、径向方程到傅里叶完备性的严格推导。下一专题将转向电像法，利用唯一性定理和区域外辅助源直接构造静电边值问题的候选解。

\section{专题二十：电像法应用（基础必讲与拔高选讲）}\label{m2:ux4e13ux9898ux4e8cux5341ux7535ux50cfux6cd5ux5e94ux7528ux57faux7840ux5fc5ux8bb2ux4e0eux62d4ux9ad8ux9009ux8bb2}

电像法在\textbf{真实求解区域之外}布置辅助源，以构造满足原方程和全部边界条件的候选电势。像电荷只是数学辅助源；候选解通过边界检查后，再由唯一性定理确认。

像源必须位于真实区域之外，并使候选势满足区域内的方程、全部边界条件和无穷远条件。由候选势可以进一步求出电场、表面电荷、作用力与能量。接地平面是基本例；半球凸起、导体劈形和球面问题需要多次反射、反演或总电荷约束。第一次阅读可先完成作业一，再处理选讲题。

\subsection{一、唯一性定理与电像法的合法性（基础必讲）}\label{m2:ux4e00ux552fux4e00ux6027ux5b9aux7406ux4e0eux7535ux50cfux6cd5ux7684ux5408ux6cd5ux6027ux57faux7840ux5fc5ux8bb2}

\subsubsection{1.
本专题调用的唯一性结论}\label{m2:ux672cux4e13ux9898ux8c03ux7528ux7684ux552fux4e00ux6027ux7ed3ux8bba}

专题十五已经严格证明：在连通区域内给定电荷分布与边界电势后，泊松方程的解唯一。专题十七进一步整理了第一类、第二类、混合边界以及无界区域的调用条件。电像法最常用的是第一类边值问题：接地导体表面满足
\(\varphi=0\)，固定电势导体表面满足 \(\varphi=V_0\)。

因此，电像法无需再次证明唯一性；真正需要做的是对候选电势完成以下检验：

\begin{enumerate}
\def\labelenumi{\arabic{enumi}.}
\tightlist
\item
  在真实区域内，它的真实源、奇点位置和奇点强度与原问题完全相同；
\item
  所有像源均位于真实区域外，不改变真实区域内的泊松方程；
\item
  候选电势在\textbf{整段边界}上满足给定电势；
\item
  对外部问题，还满足题目规定的无穷远渐近条件。
\end{enumerate}

四项全部成立后，唯一性定理才说明候选式就是物理解。若问题给定的是总电荷而不是导体电势，还必须额外加入总电荷约束；此时不能只核对``表面等势''。

\subsubsection{紧随例：为什么只验一个边界点不够（M2-136）}\label{m2:ux7d27ux968fux4f8bux4e3aux4ec0ux4e48ux53eaux9a8cux4e00ux4e2aux8fb9ux754cux70b9ux4e0dux591fm2-136}

在三维直角坐标中，一段接地导体边界的形状为抛物柱面 \(z=cx^2\)，其中
\(c>0\)，且边界包含 \(x\ne0\) 的点。有人提出候选调和函数
\(\varphi(x,y,z)=z\)，理由是它满足拉普拉斯方程且在原点为零。把候选势代入整张曲面，判断它是否满足接地条件，并说明仅检验边界上一个点为什么不足。

\par\noindent\begin{minipage}{\linewidth}
\textbf{解答}

\begin{center}
\includegraphics[width=\linewidth,height=.40\textheight,keepaspectratio]{\detokenize{assets/module-02/examples/M2-136.pdf}}
\par\smallskip{\small 图 M2-136：在一个点满足零势还不够}
\end{center}
\end{minipage}\par

图 M2-136 的接地曲面只在 \(O\) 与 \(z=0\) 相交。候选函数 \(\varphi=z\)
在域内调和，也在 \(O\) 为零，但在曲面其他点 \(P\) 取
\(cx_P^2\ne0\)；所以单个边界点不能代替整段边界检验。

考虑过原点的弯曲接地导体表面 \(z=cx^2\)（\(c>0\)）。函数 \(\varphi=z\)
在原点满足 \(\varphi=0\)，但在该表面上取值 \(cx^2\)，除 \(x=0\)
外并非零势。边界条件必须在\textbf{整段边界上逐点成立}，而不是在一个特殊点成立。

\subsubsection{2.
电像法的三项检验}\label{m2:ux7535ux50cfux6cd5ux7684ux4e09ux9879ux68c0ux9a8c}

\par\noindent\begin{minipage}{\linewidth}
\begin{center}
\includegraphics[width=\linewidth,height=.40\textheight,keepaspectratio]{\detokenize{assets/module-02/original/tikz-m2-topic20-v8-01.png}}
\par\smallskip{\small 图示：电像法的真实区域、辅助区域与合法性检验}
\end{center}
\end{minipage}\par

一个像源构造必须同时满足：

\begin{enumerate}
\def\labelenumi{\arabic{enumi}.}
\tightlist
\item
  所有像源都位于真实求解区域之外；
\item
  真实区域内的真实电荷、奇点和介质分布没有被改变；
\item
  每一段导体边界以及无穷远条件都满足原题要求。
\end{enumerate}

只有三项全部成立，才能援引唯一性定理。

像电荷只是求解工具。由像电荷算出的电势和电场只在真实区域内有效，导体内部不能按这个辅助电势作物理解释。

\subsubsection{3.
从像源回到真实物理量}\label{m2:ux4eceux50cfux6e90ux56deux5230ux771fux5b9eux7269ux7406ux91cf}

设导体表面外侧单位法向量 \(\mathbf n\) 从导体指向真空，则表面电荷密度为

\[
\boxed{
\sigma=\varepsilon_0\mathbf E_{\text{外}}\cdot\mathbf n
=-\varepsilon_0\left.\frac{\partial\varphi}{\partial n}\right|_{\text{外}}
}.
\]

真实点电荷 \(q\) 所受的力，可以直接用像电荷在真实电荷位置产生的场计算：

\[
\boxed{
\mathbf F=q\mathbf E_{\text{像}}(\mathbf r_q)
}.
\]

这里必须排除真实电荷自己的奇异场。

导体感应电荷与真实电荷的相互作用能则为

\[
\boxed{
W=\frac12q\varphi_{\text{感应}}(\mathbf r_q)
}.
\]

若有多个真实点电荷，感应部分的能量推广为

\[
W_{\text{感应}}
=\frac12\sum_i q_i\varphi_{\text{感应}}(\mathbf r_i).
\]

\par\noindent\begin{minipage}{\linewidth}
\begin{center}
\includegraphics[width=\linewidth,height=.40\textheight,keepaspectratio]{\detokenize{assets/module-02/original/tikz-m2-topic20-v8-02.png}}
\par\smallskip{\small 图示：像电荷计算作用力与相互作用能时的区别}
\end{center}
\end{minipage}\par

能量中出现
\(1/2\)，是因为真实电荷从无穷远移入时，导体上的感应电荷从零逐渐建立；不能把像电荷误当成位置和电量始终固定的真实电荷。

\subsubsection{紧随例：力与能量是否相容（M2-137）}\label{m2:ux7d27ux968fux4f8bux529bux4e0eux80fdux91cfux662fux5426ux76f8ux5bb9m2-137}

真空中一个点电荷 \(q\) 位于无限大接地导体平面上方，垂直距离为
\(d>0\)。以增大 \(d\)
为力的正方向，取点电荷远离平面时的有限诱导相互作用能为零，不计点电荷自身自能。先写出该能量随
\(d\)
的关系，再由势能导数求受力；说明直接把真实电荷与像电荷的库仑配对能当作体系势能会出现什么问题。

\par\noindent\begin{minipage}{\linewidth}
\textbf{解答}

\begin{center}
\includegraphics[width=\linewidth,height=.40\textheight,keepaspectratio]{\detokenize{assets/module-02/examples/M2-137.pdf}}
\par\smallskip{\small 图 M2-137：像电荷力与感应能必须一致}
\end{center}
\end{minipage}\par

图 M2-137
的像源随真实电荷高度一同反射移动。系统感应能先取通常两点相互作用表达式的一半，再对真实高度
\(d\) 求导；这使能量法的吸引力恰好与直接用像场算出的力一致。

若某一单参数构型的能量为 \(W(d)\)，则沿 \(d\) 增大的方向有

\[
F_d=-\frac{\mathrm dW}{\mathrm dd}.
\]

后面的接地平面模型将给出

\[
W(d)=-\frac{q^2}{16\pi\varepsilon_0d},
\]

于是

\[
F_d=-\frac{\mathrm dW}{\mathrm dd}
=-\frac{q^2}{16\pi\varepsilon_0d^2},
\]

与直接计算像电荷作用力完全一致。

\subsection{二、点电荷与无限大接地平面（基础必讲）}\label{m2:ux4e8cux70b9ux7535ux8377ux4e0eux65e0ux9650ux5927ux63a5ux5730ux5e73ux9762ux57faux7840ux5fc5ux8bb2}

\subsubsection{1.
像电荷构造与电势}\label{m2:ux50cfux7535ux8377ux6784ux9020ux4e0eux7535ux52bf}

取接地平面为 \(z=0\)，真实区域为 \(z>0\)。把真实电荷 \(q\) 放在

\[
\mathbf r_q=(0,0,d),
\qquad d>0.
\]

在平面下方的镜像点

\[
\mathbf r'=(0,0,-d)
\]

放置像电荷 \(-q\)。用柱坐标横向距离
\(s=\sqrt{x^2+y^2}\)，真实区域内的候选电势为

\[
\boxed{
\varphi(s,z)
=\frac{q}{4\pi\varepsilon_0}
\left[
\frac{1}{\sqrt{s^2+(z-d)^2}}
-\frac{1}{\sqrt{s^2+(z+d)^2}}
\right]
},
\qquad z>0.
\]

在 \(z=0\) 上，两项分母相等，故 \(\varphi=0\)；在真实区域中只有 \(q\)
的奇点，像电荷位于区域外；无穷远处
\(\varphi\to0\)。因此该候选解就是唯一物理解。

\par\noindent\begin{minipage}{\linewidth}
\begin{center}
\includegraphics[width=\linewidth,height=.40\textheight,keepaspectratio]{\detokenize{assets/module-02/original/tikz-m2-topic20-v8-03.png}}
\par\smallskip{\small 图示：接地平面的真实电荷、像电荷和表面电荷密度}
\end{center}
\end{minipage}\par

\subsubsection{紧随例：等势面的直观判断（M2-138）}\label{m2:ux7d27ux968fux4f8bux7b49ux52bfux9762ux7684ux76f4ux89c2ux5224ux65adm2-138}

真空半空间 \(z>0\) 中一点电荷 \(q\) 固定于 \((0,0,d)\)，\(d>0\)，平面
\(z=0\) 是无限大接地导体。拟在导体一侧的虚拟区域放置一个像电荷 \(-q\) 于
\((0,0,-d)\)。对边界任意点 \((x,y,0)\)
比较到两源的距离，证明整张平面的电势为零，并说明这一检验与只核对原点有何区别。

\par\noindent\begin{minipage}{\linewidth}
\textbf{解答}

\begin{center}
\includegraphics[width=\linewidth,height=.40\textheight,keepaspectratio]{\detokenize{assets/module-02/examples/M2-138.pdf}}
\par\smallskip{\small 图 M2-138：整平面上的等距关系保证零势}
\end{center}
\end{minipage}\par

图 M2-138 的边界点 \(P\)
可在整个平面任意移动，两段距离始终相等。两项势的系数等大反号，所以抵消逐点成立；若只验证电荷正下方的一个点，就没有证明这一整面几何关系。

平面上任一点到 \(q\) 与 \(-q\)
的距离相同，而两电荷异号，所以两项电势逐点抵消。这个判断对平面上的每一点成立，而不只对电荷正下方成立。

\subsubsection{2.
表面电荷密度}\label{m2:ux8868ux9762ux7535ux8377ux5bc6ux5ea6}

导体指向真空的法向量为
\(\mathbf n=\mathbf e_z\)。先求平面外侧的法向电场：

\[
E_z(s,0^+)
=-\left.\frac{\partial\varphi}{\partial z}\right|_{z=0^+}
=-\frac{qd}{2\pi\varepsilon_0(s^2+d^2)^{3/2}}.
\]

所以

\[
\boxed{
\sigma(s)
=-\frac{qd}{2\pi(s^2+d^2)^{3/2}}
}.
\]

当 \(q>0\) 时，\(\sigma<0\)，且 \(s=0\)
处绝对值最大。横向距离增大时，感应电荷密度按 \(s^{-3}\) 衰减。

总感应电荷为

\[
\begin{aligned}
Q_{\text{感应}}
&=\int_0^{\infty}\sigma(s)\,2\pi s\,\mathrm ds\\
&=-qd\int_0^{\infty}\frac{s\,\mathrm ds}{(s^2+d^2)^{3/2}}\\
&=-q.
\end{aligned}
\]

因此

\[
\boxed{Q_{\text{感应}}=-q}.
\]

这里总感应电荷恰好等于像电荷，但真实分布是连续的
\(\sigma(s)\)，并不是平面下方真的存在一个点电荷。

\subsubsection{3.
作用力与能量}\label{m2:ux4f5cux7528ux529bux4e0eux80fdux91cf}

真实电荷与像电荷相距 \(2d\)，故真实电荷所受法向力为

\[
\boxed{
\mathbf F
=-\frac{q^2}{16\pi\varepsilon_0d^2}\mathbf e_z
}.
\]

负号表示力指向导体平面。无论 \(q\) 为正还是负，接地平面对它都是吸引的。

真实电荷位置的感应电势为

\[
\varphi_{\text{感应}}(\mathbf r_q)
=\frac{1}{4\pi\varepsilon_0}\frac{-q}{2d}.
\]

所以

\[
\boxed{
W
=\frac12q\varphi_{\text{感应}}(\mathbf r_q)
=-\frac{q^2}{16\pi\varepsilon_0d}
}.
\]

\subsubsection{\texorpdfstring{紧随例：把电荷从 \(d\) 移到
\(2d\)（M2-139）}{紧随例：把电荷从 d 移到 2d（M2-139）}}\label{m2:ux7d27ux968fux4f8bux628aux7535ux8377ux4ece-d-ux79fbux5230-2dm2-139}

点电荷 \(q\)
位于真空中无限大接地导体平面上方。保持导体接地，把点电荷沿法向从距离
\(d\) 缓慢移动到距离
\(2d\)，\(d>0\)；忽略动能变化及辐射。求外力所做的功，并用该过程中诱导相互作用能的变化核对符号；能量零点取距离无穷远处，不计点电荷自能。

\par\noindent\begin{minipage}{\linewidth}
\textbf{解答}

\begin{center}
\includegraphics[width=\linewidth,height=.40\textheight,keepaspectratio]{\detokenize{assets/module-02/examples/M2-139.pdf}}
\par\smallskip{\small 图 M2-139：把电荷拉离接地平面需做正功}
\end{center}
\end{minipage}\par

图 M2-139 中把真实电荷从距平面 \(d\) 移到
\(2d\)，辅助像源也随构型变化。系统能量由较负的值升到较接近零的值，因此外力做正功；不能把像源当作固定真实电荷，用普通两点势能不带
\(1/2\) 地计算。

外力缓慢移动电荷所做的功等于势能增量：

\[
W_{\text{外}}
=W(2d)-W(d)
=\frac{q^2}{32\pi\varepsilon_0d}>0.
\]

也就是说，要把电荷拉离接地平面，外力必须做正功。

\subsection{三、接地平面上的半球形凸起（拔高选讲）}\label{m2:ux4e09ux63a5ux5730ux5e73ux9762ux4e0aux7684ux534aux7403ux5f62ux51f8ux8d77ux62d4ux9ad8ux9009ux8bb2}

\subsubsection{1. 几何与边界}\label{m2:ux51e0ux4f55ux4e0eux8fb9ux754c}

在接地平面 \(z=0\) 上有半径为 \(R\)
的导体半球，球心位于平面内的原点。真实电荷 \(Q\) 位于对称轴上的
\(z=d\)，其中
\(d>R\)。真实区域是上半空间中半球外部的部分，平面和半球表面都要求

\[
\varphi=0.
\]

\subsubsection{2. 四电荷构造}\label{m2:ux56dbux7535ux8377ux6784ux9020}

记

\[
a=\frac{R^2}{d}.
\]

除真实电荷外，取三个像电荷：

\[
\begin{array}{c|c|c}
\text{电荷}&\text{电量}&z\text{ 坐标}\\ \hline
q_1&-QR/d&+a\\
q_2&-Q&-d\\
q_3&+QR/d&-a
\end{array}
\]

\par\noindent\begin{minipage}{\linewidth}
\begin{center}
\includegraphics[width=\linewidth,height=.40\textheight,keepaspectratio]{\detokenize{assets/module-02/original/tikz-m2-topic20-v8-04.png}}
\par\smallskip{\small 图示：接地平面半球凸起的三个像电荷及两类边界检验}
\end{center}
\end{minipage}\par

真实区域中的电势为

\[
\boxed{
\begin{aligned}
\varphi(\mathbf r)=\frac{1}{4\pi\varepsilon_0}\Bigg[
&\frac{Q}{|\mathbf r-d\mathbf e_z|}
-\frac{QR/d}{|\mathbf r-a\mathbf e_z|}\\
&-\frac{Q}{|\mathbf r+d\mathbf e_z|}
+\frac{QR/d}{|\mathbf r+a\mathbf e_z|}
\Bigg].
\end{aligned}
}
\]

验证分成两组：

\begin{enumerate}
\def\labelenumi{\arabic{enumi}.}
\tightlist
\item
  \(Q\) 与 \(q_2=-Q\) 关于平面反对称，\(q_1\) 与 \(q_3=-q_1\)
  也关于平面反对称，因此 \(z=0\) 上电势为零；
\item
  \(Q\) 与 \(q_1\) 构成完整球面的反演对，\(q_2\) 与 \(q_3\)
  也构成完整球面的反演对，因此 \(r=R\) 上电势为零。
\end{enumerate}

三个像电荷都位于真实区域之外：\(q_1\) 在半球导体内部，\(q_2\) 与 \(q_3\)
在平面下方。由唯一性定理，上式就是物理解。

\subsubsection{紧随例：为什么不能只放平面像电荷（M2-140）}\label{m2:ux7d27ux968fux4f8bux4e3aux4ec0ux4e48ux4e0dux80fdux53eaux653eux5e73ux9762ux50cfux7535ux8377m2-140}

接地导体由无限平面 \(z=0\) 及其上半径为 \(R>0\)
的半球形凸起组成，半球球心在原点，凸起处于 \(z\ge0\)。真实点电荷 \(Q\)
位于对称轴上
\((0,0,d)\)，\(d>R\)，外部为真空，远处势衰减。说明为何仅作平面反射或仅作球面反演都不足，并构造满足两部分零势边界的三个像电荷，标明位置与电荷量。

\par\noindent\begin{minipage}{\linewidth}
\textbf{解答}

\begin{center}
\includegraphics[width=\linewidth,height=.40\textheight,keepaspectratio]{\detokenize{assets/module-02/examples/M2-140.pdf}}
\par\smallskip{\small 图 M2-140：半球和平面要同时成为零势面}
\end{center}
\end{minipage}\par

图 M2-140
的三枚像源都位于灰色非物理解域内。关于平面的两对反射负责平面零势，关于完整球面的两对反演负责半球零势；只取其中一枚像源只能满足一部分边界，不能援引唯一性。

若只在 \(z=-d\) 放置
\(-Q\)，平面边界确实成为零势，但半球面上一般不再等势。反过来，若只放球面反演像电荷
\(-QR/d\)，完整球面为零势，平面边界却一般不为零。两个边界必须同时检查，因此需要三枚像电荷。

\subsection{四、两接地半平面形成的导体劈形（拔高选讲）}\label{m2:ux56dbux4e24ux63a5ux5730ux534aux5e73ux9762ux5f62ux6210ux7684ux5bfcux4f53ux5288ux5f62ux62d4ux9ad8ux9009ux8bb2}

设两接地半平面沿同一直线相交，横截面中的真实区域为

\[
0<\theta<\alpha.
\]

真实点电荷位于 \((\rho_0,\theta_0,z_0)\)，其中
\(0<\theta_0<\alpha\)。所有像电荷都保持相同的 \(\rho_0\) 与
\(z_0\)，只改变角坐标和符号。

\subsubsection{\texorpdfstring{1. \(\alpha=\pi/n\)
时的有限闭合}{1. \textbackslash alpha=\textbackslash pi/n 时的有限闭合}}\label{m2:alphapin-ux65f6ux7684ux6709ux9650ux95edux5408}

若

\[
\alpha=\frac{\pi}{n},
\qquad n\in\mathbb N,
\]

则可取两组共 \(2n\) 个电荷：

\[
\boxed{
\begin{aligned}
+q:&\quad \theta_k^+=\theta_0+2k\alpha,\\
-q:&\quad \theta_k^-=-\theta_0+2k\alpha,
\end{aligned}
\qquad k=0,1,\ldots,n-1.
}
\]

其中 \(k=0\) 的正电荷是真实电荷，其余 \(2n-1\) 个均为像电荷。

在 \(\theta=0\) 边界，两组角坐标互为镜像且电量异号；在 \(\theta=\alpha\)
边界，集合在反射 \(\theta\mapsto2\alpha-\theta\)
下仍逐项异号配对。因此两边界上的电势都为零。

\par\noindent\begin{minipage}{\linewidth}
\begin{center}
\includegraphics[width=\linewidth,height=.40\textheight,keepaspectratio]{\detokenize{assets/module-02/original/tikz-m2-topic20-v8-05.png}}
\par\smallskip{\small 图示：六十度接地劈形的六电荷有限闭合构造}
\end{center}
\end{minipage}\par

\subsubsection{\texorpdfstring{紧随例：\(60^{\circ}\)
劈形（M2-141）}{紧随例：60\^{}\{\textbackslash circ\} 劈形（M2-141）}}\label{m2:ux7d27ux968fux4f8b60circ-ux5288ux5f62m2-141}

两块无限接地导体半平面相交于 \(z\) 轴，真实真空区域的横截面为
\(\rho>0\)、\(0<\theta<\pi/3\)。点电荷 \(q\) 位于
\((\rho_0,\theta_0,z_0)\)，\(\rho_0>0\)、\(0<\theta_0<\pi/3\)，远处势衰减。用有限反射电像构造满足两面接地的解，列出除真实电荷之外各像电荷的符号及方位角；所有源均保持相同的
\(\rho_0,z_0\)。

\par\noindent\begin{minipage}{\linewidth}
\textbf{解答}

\begin{center}
\includegraphics[width=\linewidth,height=.40\textheight,keepaspectratio]{\detokenize{assets/module-02/examples/M2-141.pdf}}
\par\smallskip{\small 图 M2-141：60度接地劈形的完整像源列}
\end{center}
\end{minipage}\par

图 M2-141 的绿色 \(60^{\circ}\)
扇区只含一枚真实正电荷。每次反射改变符号，两边交替反射后形成相隔
\(120^{\circ}\)
的三枚正源及三枚负源；闭合一周共六枚，扣除真实源才得到五枚像电荷。

当 \(\alpha=\pi/3\) 时，\(n=3\)。六个电荷的角坐标为

\[
\begin{array}{c|c}
\text{电量}&\text{角坐标}\\ \hline
+q&\theta_0,\ \theta_0+120^{\circ},\ \theta_0+240^{\circ}\\
-q&-\theta_0,\ 120^{\circ}-\theta_0,\ 240^{\circ}-\theta_0
\end{array}
\]

除 \(\theta_0\) 处的真实电荷外，其余五个电荷都在 \(0<\theta<60^{\circ}\)
之外，故构造合法。

\subsubsection{\texorpdfstring{2. \(72^{\circ}\)
劈形为何失败}{2. 72\^{}\{\textbackslash circ\} 劈形为何失败}}\label{m2:circ-ux5288ux5f62ux4e3aux4f55ux5931ux8d25}

当

\[
\alpha=72^{\circ}=\frac{2\pi}{5}
\]

时，连续反射的轨道虽然会有限重复，但这个 \(72^{\circ}\)
扇区不是反射群的基本扇区。像电荷角坐标集合中会出现

\[
\theta=72^{\circ}-\theta_0
\]

处的 \(-q\)。因为

\[
0<\theta_0<72^{\circ},
\]

必有

\[
0<72^{\circ}-\theta_0<72^{\circ}.
\]

这枚像电荷落入真实区域，改变了区域内的泊松方程，因此构造非法。

\par\noindent\begin{minipage}{\linewidth}
\begin{center}
\includegraphics[width=\linewidth,height=.40\textheight,keepaspectratio]{\detokenize{assets/module-02/original/tikz-m2-topic20-v8-06.png}}
\par\smallskip{\small 图示：七十二度劈形的反射轨道把异号像电荷带回真实区域}
\end{center}
\end{minipage}\par

对一般位置的单个点电荷，标准有限点像构造若要求所有像源都在真实区域之外，则劈形角应为

\[
\boxed{\alpha=\frac{\pi}{n}}.
\]

其他角度不能简单理解为``再多反射几次就行''。根据具体角度，通常要改用分离变量、格林函数、保角变换或无限级数。

\subsubsection{\texorpdfstring{紧随例：\(90^{\circ}\)
劈形需要几个像电荷（M2-142）}{紧随例：90\^{}\{\textbackslash circ\} 劈形需要几个像电荷（M2-142）}}\label{m2:ux7d27ux968fux4f8b90circ-ux5288ux5f62ux9700ux8981ux51e0ux4e2aux50cfux7535ux8377m2-142}

两块无限接地导体半平面相交于 \(z\) 轴，真空求解域为
\(\rho>0\)、\(0<\theta<\pi/2\)。点电荷 \(q\) 固定于
\((\rho_0,\theta_0,z_0)\)，\(\rho_0>0\)、\(0<\theta_0<\pi/2\)，要求远处势衰减。通过逐次反射列出全部像电荷的符号和方位角，并说明反射群为何在直角情形闭合为有限组；像源的径向与轴向坐标和真实源相同。

\par\noindent\begin{minipage}{\linewidth}
\textbf{解答}

\begin{center}
\includegraphics[width=\linewidth,height=.40\textheight,keepaspectratio]{\detokenize{assets/module-02/examples/M2-142.pdf}}
\par\smallskip{\small 图 M2-142：90度接地劈形的完整像源列}
\end{center}
\end{minipage}\par

图 M2-142
的真实域是直角扇区，关于每一条边反射一次各产生一个负像源，连续关于两边反射的像源又变回正号。四个角位置全部列出后，真实域中只有原来那枚正电荷，其余三枚都在域外。

此时 \(\alpha=\pi/2\)，故 \(n=2\)，总电荷数为
\(2n=4\)。除真实电荷外需要三枚像电荷，角坐标分别为

\[
-\theta_0,\qquad
\pi-\theta_0,\qquad
\pi+\theta_0,
\]

符号依次为 \(-q,-q,+q\)。

\subsection{五、点电荷与接地导体球（拔高选讲）}\label{m2:ux4e94ux70b9ux7535ux8377ux4e0eux63a5ux5730ux5bfcux4f53ux7403ux62d4ux9ad8ux9009ux8bb2}

\subsubsection{1.
反演点与像电荷}\label{m2:ux53cdux6f14ux70b9ux4e0eux50cfux7535ux8377}

设半径为 \(R\) 的接地导体球以原点为球心，真实电荷 \(q\) 位于

\[
\mathbf r_q=d\mathbf e_z,
\qquad d>R.
\]

在球内同一轴线上放置像电荷

\[
\boxed{
q'=-q\frac{R}{d},
\qquad
\mathbf r'=\frac{R^2}{d}\mathbf e_z
}.
\]

\par\noindent\begin{minipage}{\linewidth}
\begin{center}
\includegraphics[width=\linewidth,height=.40\textheight,keepaspectratio]{\detokenize{assets/module-02/original/tikz-m2-topic20-v8-07-margin12-v11.png}}
\par\smallskip{\small 图示：接地导体球的反演点、距离比和像电荷}
\end{center}
\end{minipage}\par

\subsubsection{2.
距离比的证明}\label{m2:ux8dddux79bbux6bd4ux7684ux8bc1ux660e}

在球面上任取一点 \(P\)，有 \(|\mathbf r_P|=R\)。记 \(\theta\) 为
\(\mathbf r_P\) 与 \(\mathbf e_z\) 的夹角，则

\[
|\mathbf r_P-\mathbf r_q|^2
=R^2+d^2-2Rd\cos\theta.
\]

另一方面，

\[
\begin{aligned}
|\mathbf r_P-\mathbf r'|^2
&=R^2+\frac{R^4}{d^2}
-2\frac{R^3}{d}\cos\theta\\
&=\frac{R^2}{d^2}
\left(R^2+d^2-2Rd\cos\theta\right).
\end{aligned}
\]

因此

\[
\boxed{
|\mathbf r_P-\mathbf r'|
=\frac{R}{d}|\mathbf r_P-\mathbf r_q|
}.
\]

球面上的电势于是满足

\[
\begin{aligned}
\varphi(P)
&=\frac{1}{4\pi\varepsilon_0}
\left(
\frac{q}{|\mathbf r_P-\mathbf r_q|}
+\frac{q'}{|\mathbf r_P-\mathbf r'|}
\right)\\
&=\frac{1}{4\pi\varepsilon_0|\mathbf r_P-\mathbf r_q|}
\left(q+\frac{d}{R}q'\right)=0.
\end{aligned}
\]

球外真实区域中的电势为

\[
\boxed{
\varphi(\mathbf r)
=\frac{1}{4\pi\varepsilon_0}
\left[
\frac{q}{|\mathbf r-d\mathbf e_z|}
-\frac{qR/d}{|\mathbf r-(R^2/d)\mathbf e_z|}
\right]
},
\qquad r\ge R.
\]

\subsubsection{紧随例：为什么像电荷不能放在球心（M2-143）}\label{m2:ux7d27ux968fux4f8bux4e3aux4ec0ux4e48ux50cfux7535ux8377ux4e0dux80fdux653eux5728ux7403ux5fc3m2-143}

真空中半径为 \(R>0\) 的接地导体球以原点为球心，外部点电荷 \(q\ne0\) 位于
\(z=d\)，\(d>R\)。有人拟仅在球心放一个像电荷，使某一球面点的电势抵消。比较不同球面点到真实源的距离，解释球心像为何不能满足整球面零势，并指出正确构造需要同时确定哪些像源参数。

\par\noindent\begin{minipage}{\linewidth}
\textbf{解答}

\begin{center}
\includegraphics[width=\linewidth,height=.40\textheight,keepaspectratio]{\detokenize{assets/module-02/examples/M2-143.pdf}}
\par\smallskip{\small 图 M2-143：中心像源不能抵消角向变化}
\end{center}
\end{minipage}\par

图 M2-143 的反演像位于球心与真实电荷之间，距球心
\(R^2/d\)。球面各点到两偏心源的距离成固定比例，才可抵消变化的边界势；球心源只能给常量，不能消去真实外部源带来的角向变化。

球心点电荷在球面上产生常量电势，只能整体抬高或降低球面电势，不能抵消外部偏心电荷在球面上的角向变化。接地球需要偏心像电荷；中心像电荷将在孤立球问题中用来调节净电荷。

\subsubsection{3.
表面电荷密度与总感应电荷}\label{m2:ux8868ux9762ux7535ux8377ux5bc6ux5ea6ux4e0eux603bux611fux5e94ux7535ux8377}

取 \(\theta=0\) 指向真实电荷。由球面外侧法向场得到

\[
\boxed{
\sigma(\theta)
=-\frac{q}{4\pi R}
\frac{d^2-R^2}
{\left(d^2+R^2-2dR\cos\theta\right)^{3/2}}
}.
\]

当 \(q>0\) 时，整个球面上 \(\sigma<0\)，最近点 \(\theta=0\)
处绝对值最大。积分可得

\[
\boxed{
Q_{\text{感应}}
=\int_S\sigma\,\mathrm dS
=-\frac{qR}{d}
}.
\]

总感应电荷等于偏心像电荷 \(q'\)，但表面电荷密度并不均匀。

\par\noindent\begin{minipage}{\linewidth}
\begin{center}
\includegraphics[width=\linewidth,height=.40\textheight,keepaspectratio]{\detokenize{assets/module-02/original/tikz-m2-topic20-v8-08-margin12-v11.png}}
\par\smallskip{\small 图示：接地导体球的表面电荷、总感应电荷、作用力与能量}
\end{center}
\end{minipage}\par

\subsubsection{4.
作用力与能量}\label{m2:ux4f5cux7528ux529bux4e0eux80fdux91cf-1}

真实电荷到像电荷的距离为

\[
d-\frac{R^2}{d}=\frac{d^2-R^2}{d}.
\]

取 \(+\mathbf e_z\) 为远离球心的方向，真实电荷所受力为

\[
\boxed{
\mathbf F
=-\frac{q^2Rd}{4\pi\varepsilon_0(d^2-R^2)^2}\mathbf e_z
}.
\]

相互作用能为

\[
\boxed{
W
=\frac12\frac{1}{4\pi\varepsilon_0}
\frac{qq'}{d-R^2/d}
=-\frac{q^2R}{8\pi\varepsilon_0(d^2-R^2)}
}.
\]

两式满足 \(F_z=-\mathrm dW/\mathrm dd\)。

\subsubsection{紧随例：近平面极限（M2-144）}\label{m2:ux7d27ux968fux4f8bux8fd1ux5e73ux9762ux6781ux9650m2-144}

真空中点电荷 \(q\) 位于半径 \(R>0\) 的接地导体球外，对称轴距离为
\(d=R+h\)，其中 \(h>0\)。取无穷远零势，诱导相互作用能不计点电荷自能且以
\(d\to\infty\) 为零。求点电荷受力及该能量，并在 \(h\ll R\)
时取主阶极限，与距接地平面 \(h\) 的情形比较，明确应保持的最近表面距离。

\par\noindent\begin{minipage}{\linewidth}
\textbf{解答}

\begin{center}
\includegraphics[width=\linewidth,height=.40\textheight,keepaspectratio]{\detokenize{assets/module-02/examples/M2-144.pdf}}
\par\smallskip{\small 图 M2-144：球面局部在小间距下趋于平面}
\end{center}
\end{minipage}\par

图 M2-144 的最近间距是 \(h=d-R\)，而不是球心距
\(d\)。在力与能量的精确分母中把 \(d^2-R^2\) 展为
\(2Rh\)，再保留最低阶，就回到接地平面公式；这一极限要求间距远小于曲率半径。

令电荷到球面的最近距离为

\[
h=d-R,
\qquad h\ll R.
\]

此时 \(d^2-R^2\simeq2Rh\)，故

\[
F_z\simeq-\frac{q^2}{16\pi\varepsilon_0h^2},
\qquad
W\simeq-\frac{q^2}{16\pi\varepsilon_0h}.
\]

这正是点电荷距接地平面为 \(h\)
时的结果。局部曲率半径远大于间距时，球面在电荷附近近似为平面。

\subsection{六、孤立导体球：等势条件与净电荷约束（拔高选讲）}\label{m2:ux516dux5b64ux7acbux5bfcux4f53ux7403ux7b49ux52bfux6761ux4ef6ux4e0eux51c0ux7535ux8377ux7ea6ux675fux62d4ux9ad8ux9009ux8bb2}

偏心像电荷

\[
q_1=-q\frac{R}{d},
\qquad
r_1=\frac{R^2}{d}
\]

负责消除球面电势的角向变化。若还要满足给定净电荷，就在球心加入中心像电荷
\(q_0\)。中心电荷在整个球面上产生相同电势，不会破坏等势条件。

\par\noindent\begin{minipage}{\linewidth}
\begin{center}
\includegraphics[width=\linewidth,height=.40\textheight,keepaspectratio]{\detokenize{assets/module-02/original/tikz-m2-topic20-v8-09.png}}
\par\smallskip{\small 图示：接地球、孤立不带电球和孤立带电球的像电荷比较}
\end{center}
\end{minipage}\par

\subsubsection{1.
孤立不带电导体球}\label{m2:ux5b64ux7acbux4e0dux5e26ux7535ux5bfcux4f53ux7403}

净电荷条件为

\[
q_1+q_0=0,
\]

所以

\[
\boxed{q_0=+q\frac{R}{d}}.
\]

球面电势为

\[
\boxed{
U=\frac{1}{4\pi\varepsilon_0}\frac{q_0}{R}
=\frac{q}{4\pi\varepsilon_0d}
}.
\]

球面电荷密度可写成接地球的非均匀部分与中心像电荷对应的均匀部分之和：

\[
\boxed{
\sigma_{\text{中性}}(\theta)
=-\frac{q}{4\pi R}
\frac{d^2-R^2}
{(d^2+R^2-2dR\cos\theta)^{3/2}}
+\frac{q}{4\pi Rd}
}.
\]

第一项积分为 \(-qR/d\)，第二项积分为 \(+qR/d\)，故总电荷为零。

真实电荷所受力为两个像电荷作用力之和：

\[
\boxed{
F_z
=\frac{q^2R}{4\pi\varepsilon_0}
\left[
\frac{1}{d^3}
-\frac{d}{(d^2-R^2)^2}
\right]
}.
\]

把括号通分，

\[
\frac{1}{d^3}-\frac{d}{(d^2-R^2)^2}
=-\frac{R^2(2d^2-R^2)}{d^3(d^2-R^2)^2}<0.
\]

所以

\[
\boxed{F_z<0}.
\]

孤立不带电导体球对外部点电荷总是吸引。其物理原因是同号感应电荷被推向远侧，异号感应电荷集中在近侧，近侧作用占优。

\subsubsection{紧随例：远距离极限（M2-145）}\label{m2:ux7d27ux968fux4f8bux8fdcux8dddux79bbux6781ux9650m2-145}

半径为 \(R>0\)
的孤立中性导体球位于真空中，球心在原点，球的总自由电荷固定为零而不接地。点电荷
\(q\ne0\) 位于距球心 \(d>R\)
处，取无穷远零势。利用适当像源求点电荷受力的远区主阶
\(d\gg R\)，判断方向，并与同半径接地球的远区受力作比较，解释总电荷约束的作用。

\par\noindent\begin{minipage}{\linewidth}
\textbf{解答}

\begin{center}
\includegraphics[width=\linewidth,height=.40\textheight,keepaspectratio]{\detokenize{assets/module-02/examples/M2-145.pdf}}
\par\smallskip{\small 图 M2-145：孤立中性球远场只剩诱导偶极}
\end{center}
\end{minipage}\par

图 M2-145
在偏心反演像之外又加球心补偿源，以保持球净电荷为零。远处单极贡献因此消失，主要相互作用来自诱导偶极，力按
\(d^{-5}\) 衰减；接地球允许交换净电荷，故不能使用同一远区结论。

当 \(d\gg R\) 时，

\[
F_z\simeq
-\frac{2q^2R^3}{4\pi\varepsilon_0d^5}.
\]

接地球在远处保留净感应电荷，力按 \(d^{-3}\)
衰减；孤立中性球只能产生偶极极化，力按 \(d^{-5}\)
衰减。两种边界条件不能混用。

\subsubsection{2.
孤立带电导体球}\label{m2:ux5b64ux7acbux5e26ux7535ux5bfcux4f53ux7403}

若导体球给定净电荷 \(Q_0\)，则

\[
q_1+q_0=Q_0,
\]

从而

\[
\boxed{
q_0=Q_0+q\frac{R}{d}
}.
\]

球面电势为

\[
\boxed{
U=\frac{1}{4\pi\varepsilon_0}
\left(
\frac{Q_0}{R}+\frac{q}{d}
\right)
}.
\]

表面电荷密度为

\[
\boxed{
\sigma(\theta)
=-\frac{q}{4\pi R}
\frac{d^2-R^2}
{(d^2+R^2-2dR\cos\theta)^{3/2}}
+\frac{Q_0+qR/d}{4\pi R^2}
}.
\]

真实电荷所受力为

\[
\boxed{
F_z
=\frac{q}{4\pi\varepsilon_0}
\left[
\frac{Q_0}{d^2}
+\frac{qR}{d^3}
-\frac{qRd}{(d^2-R^2)^2}
\right]
}.
\]

前两项来自中心像电荷，最后一项来自偏心像电荷。即使 \(q\) 与 \(Q_0\)
同号，只要电荷足够靠近球面，最后一项仍按 \((d-R)^{-2}\)
发散并占主导，作用力最终指向球面。

\subsubsection{紧随例：同号电荷何时仍相吸（M2-146）}\label{m2:ux7d27ux968fux4f8bux540cux53f7ux7535ux8377ux4f55ux65f6ux4ecdux76f8ux5438m2-146}

半径为 \(R>0\) 的孤立导体球在真空中带固定总电荷 \(Q_0\)，外部点电荷
\(q>0\) 位于距球心 \(d>R\)
的轴线上。球不接地，取无穷远电势为零，以远离球心为受力正向。求点电荷受力由吸引变为排斥时的临界
\(Q_0\)，并分析 \(d\to R^+\)
时临界条件的变化；必须同时满足球面等势和球的总电荷约束。

\par\noindent\begin{minipage}{\linewidth}
\textbf{解答}

\begin{center}
\includegraphics[width=\linewidth,height=.40\textheight,keepaspectratio]{\detokenize{assets/module-02/examples/M2-146.pdf}}
\par\smallskip{\small 图 M2-146：同号净电荷仍可能被极化吸引}
\end{center}
\end{minipage}\par

图 M2-146
的中心补偿源要同时承担给定净电荷和抵消反演像总量的任务。由真实电荷处的总像场令力为零，得到随
\(d\) 变化的临界
\(Q_0\)；越接近球面，极化吸引越强，因此有限同号净电荷也可能对应吸引。

设 \(q>0\)，并规定 \(F_z>0\) 为斥力。令 \(F_z=0\)，得到临界球电荷

\[
\boxed{
Q_{0,\text{临界}}
=\frac{qR^3(2d^2-R^2)}{d(d^2-R^2)^2}
}.
\]

于是

\[
Q_0>Q_{0,\text{临界}}
\Rightarrow F_z>0,
\]

\[
Q_0<Q_{0,\text{临界}}
\Rightarrow F_z<0.
\]

由于 \(d\to R^+\) 时
\(Q_{0,\text{临界}}\to\infty\)，任何有限的同号球电荷都无法阻止足够近处的极化吸引。

\subsection{七、电像法的使用条件与检查}\label{m2:ux4e03ux7535ux50cfux6cd5ux7684ux4f7fux7528ux6761ux4ef6ux4e0eux68c0ux67e5}

\par\noindent\begin{minipage}{\linewidth}
\begin{center}
\includegraphics[width=\linewidth,height=.40\textheight,keepaspectratio]{\detokenize{assets/module-02/original/tikz-m2-topic20-v8-10.png}}
\par\smallskip{\small 图示：电像法的像源构造和真实物理量计算}
\end{center}
\end{minipage}\par

\subsubsection{1.
必须区分的三类条件}\label{m2:ux5fc5ux987bux533aux5206ux7684ux4e09ux7c7bux6761ux4ef6}

\begin{enumerate}
\def\labelenumi{\arabic{enumi}.}
\tightlist
\item
  \textbf{接地导体：}
  导体电势给定，导体可与大地交换电荷，净电荷通常未知；
\item
  \textbf{孤立中性导体：} 导体电势是未知常量，净电荷固定为零；
\item
  \textbf{孤立带电导体：} 导体电势是未知常量，净电荷固定为 \(Q_0\)。
\end{enumerate}

接地条件与``净电荷为零''不是同一件事。

\subsubsection{2.
力、能量和表面电荷不能混用规则}\label{m2:ux529bux80fdux91cfux548cux8868ux9762ux7535ux8377ux4e0dux80fdux6df7ux7528ux89c4ux5219}

\begin{enumerate}
\def\labelenumi{\arabic{enumi}.}
\tightlist
\item
  力用 \(q\mathbf E_{\text{像}}\) 直接计算，不带 \(1/2\)；
\item
  感应相互作用能用 \(\frac12q\varphi_{\text{感应}}\)；
\item
  真实表面电荷必须由导体外侧法向场恢复，不能把像电荷位置当作真实电荷位置；
\item
  像电荷电势只在真实求解区域内使用。
\end{enumerate}

\subsubsection{3.
有限像电荷法的范围}\label{m2:ux6709ux9650ux50cfux7535ux8377ux6cd5ux7684ux8303ux56f4}

平面、球面以及某些夹角为 \(\pi/n\)
的劈形具有特殊的反射或反演结构，因而能用有限个点像电荷处理。一般曲面、多导体、多球体系或非整分劈形往往会产生无限像电荷序列，或者根本不适合点像表示。

选择电像法时依次检查：

\begin{enumerate}
\def\labelenumi{\arabic{enumi}.}
\tightlist
\item
  边界几何是否允许反射或反演后保持简单源形式；
\item
  像源是否全部位于真实区域之外；
\item
  所有边界能否同时满足；
\item
  是否还有固定净电荷、固定电势或无穷远条件需要补充。
\end{enumerate}

\subsection{八、练习与解答}\label{m2:ux516bux7ec3ux4e60ux4e0eux89e3ux7b54-6}

\subsubsection{作业一：接地平面的电势、表面电荷与力（基础必做）}\label{m2:ux4f5cux4e1aux4e00ux63a5ux5730ux5e73ux9762ux7684ux7535ux52bfux8868ux9762ux7535ux8377ux4e0eux529bux57faux7840ux5fc5ux505a}

\textbf{题号：M2-147}

\textbf{层级：} 基础　　\textbf{方法：} 电像法

真空中无限大接地导体占据 \(z\le0\)，点电荷 \(q\) 固定于
\((0,0,d)\)，\(d>0\)，无穷远电势为零。用电像法求 \(z>0\)
的电势、平面上的感应面电荷密度及其总量、真实电荷受到的力。再求与导体感应有关的有限相互作用能，以
\(d\to\infty\)
为零，不包含真实点电荷自身发散的自能，并核对力与势能导数。

\par\noindent\begin{minipage}{\linewidth}
\textbf{解答}

\begin{center}
\includegraphics[width=\linewidth,height=.40\textheight,keepaspectratio]{\detokenize{assets/module-02/examples/M2-147.pdf}}
\par\smallskip{\small 图 M2-147：接地平面的电势、表面电荷与力（基础必做）}
\end{center}
\end{minipage}\par

图 M2-147
的反射点保证接地面处处零势。对真实半空间电势求导得到表面密度，再积分核对总感应电荷为
\(-q\)；力直接用像场，感应能则保留一半系数，两者应由高度导数互相检验。

在平面下方对称位置放置像电荷 \(-q\)。以平面内距垂足的距离为 \(s\)，则

\[
\varphi(s,z)
=\frac{q}{4\pi\varepsilon_0}
\left[
\frac{1}{\sqrt{s^2+(z-d)^2}}
-\frac{1}{\sqrt{s^2+(z+d)^2}}
\right].
\]

平面外侧法向场为

\[
E_z(s,0^+)
=-\frac{qd}{2\pi\varepsilon_0(s^2+d^2)^{3/2}},
\]

故

\[
\boxed{
\sigma(s)=-\frac{qd}{2\pi(s^2+d^2)^{3/2}}
}.
\]

积分得

\[
Q_{\text{感应}}
=\int_0^{\infty}\sigma(s)2\pi s\,\mathrm ds
=-q.
\]

真实电荷与像电荷相距 \(2d\)，所以

\[
\boxed{
\mathbf F=-\frac{q^2}{16\pi\varepsilon_0d^2}\mathbf e_z
}.
\]

感应电势为 \(-q/(8\pi\varepsilon_0d)\)，因此

\[
\boxed{
W=\frac12q\varphi_{\text{感应}}
=-\frac{q^2}{16\pi\varepsilon_0d}
}.
\]

最后检验 \(-\mathrm dW/\mathrm dd=F_z\)，结果相符。

\subsubsection{作业二：接地平面上的半球凸起（拔高选做）}\label{m2:ux4f5cux4e1aux4e8cux63a5ux5730ux5e73ux9762ux4e0aux7684ux534aux7403ux51f8ux8d77ux62d4ux9ad8ux9009ux505a}

\textbf{题号：M2-148}

\textbf{层级：} 复赛　　\textbf{方法：} 电像法

无限大接地导体平面 \(z=0\) 上连有一半径为 \(R>0\)
的导体半球凸起，半球中心为原点、位于
\(z\ge0\)。真空真实域是导体外部的上半空间，真实点电荷 \(Q\) 位于
\((0,0,d)\)，\(d>R\)。构造全部像电荷，给出位置、电荷量及所在虚拟区域；分别在平面部分和半球部分任取边界点证明总势为零，并检验远处衰减条件。

\par\noindent\begin{minipage}{\linewidth}
\textbf{解答}

\begin{center}
\includegraphics[width=\linewidth,height=.40\textheight,keepaspectratio]{\detokenize{assets/module-02/examples/M2-148.pdf}}
\par\smallskip{\small 图 M2-148：接地平面上的半球凸起（拔高选做）}
\end{center}
\end{minipage}\par

图 M2-148
的平面边界与半球边界需要同时匹配，完整像源列含平面反射、球面反演及反演源的反射。验证时分别在平面和球面上用等距或比例距离成对抵消，并确认三枚辅助源都在真实区域之外。

令 \(a=R^2/d\)。三枚像电荷为

\[
q_1=-\frac{QR}{d},\quad z_1=a;
\]

\[
q_2=-Q,\quad z_2=-d;
\]

\[
q_3=+\frac{QR}{d},\quad z_3=-a.
\]

先检验平面：\(Q\) 与 \(q_2\) 关于 \(z=0\) 等距异号，\(q_1\) 与 \(q_3\)
也关于 \(z=0\) 等距异号，所以平面上总电势为零。

再检验球面：\(Q\) 与 \(q_1\)
是外部电荷与球内反演像电荷的一对，故它们在完整球面上总势为零；\(q_2\) 与
\(q_3\) 同样构成另一组反演对，故它们在完整球面上总势也为零。

所有像电荷均位于真实区域外，且无穷远处电势趋零。由唯一性定理，这一四电荷体系在真实区域中给出的电势就是物理解。

\subsubsection{\texorpdfstring{作业三：\(60^{\circ}\)
接地劈形（拔高选做）}{作业三：60\^{}\{\textbackslash circ\} 接地劈形（拔高选做）}}\label{m2:ux4f5cux4e1aux4e0960circ-ux63a5ux5730ux5288ux5f62ux62d4ux9ad8ux9009ux505a}

\textbf{题号：M2-149}

\textbf{层级：} 复赛　　\textbf{方法：} 电像法

无限接地导体劈形的真空开角为 \(\alpha=\pi/3\)，两边界半平面交于 \(z\)
轴。点电荷 \(q\) 位于柱坐标
\((\rho_0,\theta_0,z_0)\)，\(\rho_0>0\)、\(0<\theta_0<\alpha\)，远处势衰减。构造有限组电像，列出五个像源的电荷量和角位置，说明其径向、轴向坐标，并用成对反射核对两张接地边界。

\par\noindent\begin{minipage}{\linewidth}
\textbf{解答}

\begin{center}
\includegraphics[width=\linewidth,height=.40\textheight,keepaspectratio]{\detokenize{assets/module-02/examples/M2-149.pdf}}
\par\smallskip{\small 图 M2-149：\(60^{\circ}\) 接地劈形（拔高选做）}
\end{center}
\end{minipage}\par

图 M2-149 的六个角位置由同一半径和两组旋转角确定，只有 \(\theta_0\)
位于真实扇区。分别对每个接地边界寻找异号镜像配对，既检验零势，也能核对是否少算了闭合反射链的某一枚电荷。

这里

\[
\alpha=60^{\circ}=\frac{\pi}{3},
\qquad n=3.
\]

用统一公式

\[
\theta_k^+=\theta_0+2k\alpha,
\qquad
\theta_k^-=-\theta_0+2k\alpha,
\qquad k=0,1,2.
\]

正电荷位于

\[
\theta_0,\qquad
\theta_0+120^{\circ},\qquad
\theta_0+240^{\circ};
\]

负电荷位于

\[
-\theta_0,\qquad
120^{\circ}-\theta_0,\qquad
240^{\circ}-\theta_0.
\]

共得到六个电荷，其中 \(\theta_0\) 处的 \(+q\)
是原电荷，其余五个是像电荷。它们在两条边界的反射下逐项异号配对，故两边界电势均为零；五枚像电荷均不在
\(0<\theta<60^{\circ}\) 内，构造合法。

\subsubsection{作业四：接地导体球的反演构造（拔高选做）}\label{m2:ux4f5cux4e1aux56dbux63a5ux5730ux5bfcux4f53ux7403ux7684ux53cdux6f14ux6784ux9020ux62d4ux9ad8ux9009ux505a}

\textbf{题号：M2-150}

\textbf{层级：} 复赛　　\textbf{方法：} 电像法

半径为 \(R>0\) 的接地导体球球心在原点，真空中一点电荷 \(q\) 位于
\(z=d\)，\(d>R\)，要求无穷远势为零。设球内轴线上有一个待定像电荷，从整球面等势条件推导其位置与电荷量；据此求真实点电荷受到的力和诱导相互作用能，能量以
\(d\to\infty\) 为零且不计真实点电荷自能。

\par\noindent\begin{minipage}{\linewidth}
\textbf{解答}

\begin{center}
\includegraphics[width=\linewidth,height=.40\textheight,keepaspectratio]{\detokenize{assets/module-02/examples/M2-150.pdf}}
\par\smallskip{\small 图 M2-150：接地导体球的反演构造（拔高选做）}
\end{center}
\end{minipage}\par

图 M2-150
的真实源、球心和反演像共线。球面距离比例先给像位置，再给电荷系数；真实电荷处的像场决定力，感应势乘
\(q/2\) 决定能量，整个过程只在球外真实域解释为物理解。

设像电荷位于同一轴线上，距球心 \(a\)，电量为 \(q'\)。要使球面任一点
\(P\) 上电势为零，需要两距离之比为常数。取

\[
a=\frac{R^2}{d},
\]

则对 \(|\mathbf r_P|=R\)，有

\[
|\mathbf r_P-a\mathbf e_z|
=\frac{R}{d}|\mathbf r_P-d\mathbf e_z|.
\]

因此

\[
\frac{q}{|\mathbf r_P-d\mathbf e_z|}
+\frac{q'}{|\mathbf r_P-a\mathbf e_z|}=0
\]

要求

\[
q+\frac{d}{R}q'=0,
\]

故

\[
\boxed{
a=\frac{R^2}{d},
\qquad
q'=-q\frac{R}{d}
}.
\]

两电荷间距为 \((d^2-R^2)/d\)，所以

\[
\boxed{
F_z
=\frac{1}{4\pi\varepsilon_0}
\frac{qq'}{(d-a)^2}
=-\frac{q^2Rd}{4\pi\varepsilon_0(d^2-R^2)^2}
}.
\]

感应相互作用能为

\[
\boxed{
W
=\frac12\frac{1}{4\pi\varepsilon_0}\frac{qq'}{d-a}
=-\frac{q^2R}{8\pi\varepsilon_0(d^2-R^2)}
}.
\]

\subsubsection{作业五：证明孤立中性球总是吸引点电荷（拔高选做）}\label{m2:ux4f5cux4e1aux4e94ux8bc1ux660eux5b64ux7acbux4e2dux6027ux7403ux603bux662fux5438ux5f15ux70b9ux7535ux8377ux62d4ux9ad8ux9009ux505a}

\textbf{题号：M2-151}

\textbf{层级：} 复赛　　\textbf{方法：} 电像法

真空中半径为 \(R>0\)
的孤立中性导体球不接地，总自由电荷始终为零。外部点电荷 \(q\ne0\)
位于距球心 \(d>R\)
的轴线上，取无穷远零势。构造满足球面等势及总电荷为零的像源组，求真实点电荷所受径向力，并证明它在全部
\(d>R\) 范围内均指向球心。

\par\noindent\begin{minipage}{\linewidth}
\textbf{解答}

\begin{center}
\includegraphics[width=\linewidth,height=.40\textheight,keepaspectratio]{\detokenize{assets/module-02/examples/M2-151.pdf}}
\par\smallskip{\small 图 M2-151：证明孤立中性球总是吸引点电荷（拔高选做）}
\end{center}
\end{minipage}\par

图 M2-151
的偏心反演像为负，球心补偿源为正，两者总和为零。两项在真实电荷处产生相反方向的力，但偏心像更近；将精确结果通分后，其分子和分母都可判号，得到始终指向球心的净力。

偏心像电荷为

\[
q_1=-q\frac{R}{d},
\qquad
r_1=\frac{R^2}{d}.
\]

为使导体球净电荷为零，在球心加入

\[
q_0=+q\frac{R}{d}.
\]

取远离球心方向为正，真实电荷所受力为

\[
F_z
=\frac{q^2R}{4\pi\varepsilon_0}
\left[
\frac1{d^3}-\frac{d}{(d^2-R^2)^2}
\right].
\]

通分得

\[
F_z
=-\frac{q^2R^3(2d^2-R^2)}
{4\pi\varepsilon_0d^3(d^2-R^2)^2}.
\]

因为 \(d>R\)，分母为正且 \(2d^2-R^2>0\)，所以

\[
\boxed{F_z<0}.
\]

因此无论 \(q\) 的符号如何，孤立不带电导体球对它总是吸引。

\subsubsection{作业六：孤立带电球的吸引与排斥分界（拔高选做）}\label{m2:ux4f5cux4e1aux516dux5b64ux7acbux5e26ux7535ux7403ux7684ux5438ux5f15ux4e0eux6392ux65a5ux5206ux754cux62d4ux9ad8ux9009ux505a}

\textbf{题号：M2-152}

\textbf{层级：} 复赛　　\textbf{方法：} 电像法

真空中半径为 \(R>0\) 的孤立导体球带固定总自由电荷 \(Q_0\)，外部点电荷
\(q>0\) 位于距球心 \(d>R\)
的轴线上。球面等势但电势数值未预先指定，取无穷远零势。构造满足这些条件的像源，求点电荷受力，并给出力为零的临界
\(Q_0\) 以及临界两侧的吸引、排斥判据。

\par\noindent\begin{minipage}{\linewidth}
\textbf{解答}

\begin{center}
\includegraphics[width=\linewidth,height=.40\textheight,keepaspectratio]{\detokenize{assets/module-02/examples/M2-152.pdf}}
\par\smallskip{\small 图 M2-152：孤立带电球的吸引与排斥分界（拔高选做）}
\end{center}
\end{minipage}\par

图 M2-152 的中心源除了补偿反演像，还要补入指定的球净电荷
\(Q_0\)。以远离球心为正方向，把两项像力相加后令总力为零，得到临界量；判断吸引还是排斥时保持同一方向约定。

偏心像电荷仍为

\[
q_1=-q\frac{R}{d},
\qquad
r_1=\frac{R^2}{d}.
\]

球心像电荷由净电荷条件确定：

\[
q_0=Q_0+q\frac{R}{d}.
\]

作用力为

\[
F_z
=\frac{q}{4\pi\varepsilon_0}
\left[
\frac{Q_0}{d^2}
+\frac{qR}{d^3}
-\frac{qRd}{(d^2-R^2)^2}
\right].
\]

令 \(F_z=0\)，得到

\[
\frac{Q_0}{d^2}
=\frac{qRd}{(d^2-R^2)^2}
-\frac{qR}{d^3}.
\]

乘以 \(d^2\) 并通分：

\[
\begin{aligned}
Q_{0,\text{临界}}
&=qR\left[
\frac{d^3}{(d^2-R^2)^2}-\frac1d
\right]\\
&=\boxed{
\frac{qR^3(2d^2-R^2)}{d(d^2-R^2)^2}
}.
\end{aligned}
\]

因此

\[
Q_0>Q_{0,\text{临界}}
\quad\Longrightarrow\quad
F_z>0,
\]

表现为排斥；反之表现为吸引。特别地，当 \(d\to R^+\)
时，临界电荷发散，说明足够靠近导体表面时，局部极化吸引总会占优势。

上述构造涵盖接地平面、半球凸起、导体劈形和三类导体球。像源求得后，仍需由候选势计算真实的表面电荷、作用力与能量。遇到新的几何边界时，先确定真实区域和边界类型，再判断有限个像源能否同时满足全部条件。

\section{模块二附录一：勒让德方程和贝塞尔方程的推导（拔高选讲）}\label{m2:ux6a21ux5757ux4e8cux9644ux5f55ux4e00ux52d2ux8ba9ux5fb7ux65b9ux7a0bux548cux8d1dux585eux5c14ux65b9ux7a0bux7684ux63a8ux5bfcux62d4ux9ad8ux9009ux8bb2}

\begin{quote}
阅读准备：特殊函数来自几何与边界
特殊函数并不是额外假设。把拉普拉斯方程写到球坐标或柱坐标并分离变量后，角向或径向方程自然成为勒让德方程或贝塞尔方程。端点有限性、原点正则性和边界条件共同选择允许的函数与本征值。

\textbf{例 M2-153：轴线是否在域内决定能否保留奇异分支}
\end{quote}

柱坐标分离变量得到普通贝塞尔径向方程，考虑固定整数阶 \(m\ge0\)、波数
\(k>0\) 时的分支 \(Y_m(k\rho)\)。分别比较包含轴线的实心无源域
\(0\le\rho<a\) 与不含轴线的圆环域
\(a<\rho<b\)，\(0<a<b\)。在要求真实域内电势正则的前提下，判断是否能因
\(\rho=0\) 的奇性而舍去这一分支，说明圆环情形还需检查哪些条件。

\par\noindent\begin{minipage}{\linewidth}
\textbf{解答}

\begin{center}
\includegraphics[width=\linewidth,height=.40\textheight,keepaspectratio]{\detokenize{assets/module-02/examples/M2-153.pdf}}
\par\smallskip{\small 图 M2-153：轴线是否在域内决定能否保留奇异分支}
\end{center}
\end{minipage}\par

图 M2-153 的左域包含轴线，右域把轴线挖去。是否要求 \(Y_m\) 在 \(\rho=0\)
有限取决于该点是否属于真实区域；对于圆环，必须转而检查实际内、外圆柱边界，不能只凭域外奇性删除某个分支。

柱内包含轴线时，第二类贝塞尔函数 \(Y_m(kr)\) 在 \(r=0\)
发散，通常必须舍去；若求解区域是空心圆环，轴线不在区域内，就不能仅凭这一理由舍去它。

球面和圆柱面边界会在变量分离后产生勒让德函数和贝塞尔函数。坐标系决定微分算子，周期性限制角向整数，端点正则性和边界条件再筛选离散本征值。

以下推导分别涉及方程来源、本征值取值、不同函数的适用区域和正交展开：

\begin{enumerate}
\def\labelenumi{\arabic{enumi}.}
\tightlist
\item
  勒让德方程和贝塞尔方程怎样从拉普拉斯方程中产生；
\item
  为什么 \(l\)、\(m\) 取整数，径向本征值却由函数零点决定；
\item
  第一类、第二类和修正函数分别在什么区域内可用；
\item
  如何利用正交性把任意边界函数展开为本征模。
\end{enumerate}

\par\noindent\begin{minipage}{\linewidth}
\begin{center}
\includegraphics[width=\linewidth,height=.40\textheight,keepaspectratio]{\detokenize{assets/module-02/original/tikz-m2-topic20-v8-11.png}}
\par\smallskip{\small 图示：勒让德函数与贝塞尔函数由坐标和边界条件共同筛选}
\end{center}
\end{minipage}\par

\subsection{一、球坐标分离变量与勒让德方程}\label{m2:ux4e00ux7403ux5750ux6807ux5206ux79bbux53d8ux91cfux4e0eux52d2ux8ba9ux5fb7ux65b9ux7a0b}

\subsubsection{1.
轴对称拉普拉斯方程}\label{m2:ux8f74ux5bf9ux79f0ux62c9ux666eux62c9ux65afux65b9ux7a0b}

球坐标中的拉普拉斯方程为

\[
\frac{1}{r^2}\frac{\partial}{\partial r}
\left(r^2\frac{\partial\varphi}{\partial r}\right)
+\frac{1}{r^2\sin\theta}\frac{\partial}{\partial\theta}
\left(\sin\theta\frac{\partial\varphi}{\partial\theta}\right)
+\frac{1}{r^2\sin^2\theta}
\frac{\partial^2\varphi}{\partial\phi^2}=0.
\]

先考虑轴对称情形

\[
\frac{\partial\varphi}{\partial\phi}=0.
\]

令

\[
\varphi(r,\theta)=R(r)\Theta(\theta),
\]

代入并乘以 \(r^2/(R\Theta)\)：

\[
\frac1R\frac{\mathrm d}{\mathrm dr}
\left(r^2\frac{\mathrm dR}{\mathrm dr}\right)
=-\frac1{\Theta\sin\theta}
\frac{\mathrm d}{\mathrm d\theta}
\left(\sin\theta\frac{\mathrm d\Theta}{\mathrm d\theta}\right).
\]

左边只含 \(r\)，右边只含
\(\theta\)，因此二者必须等于同一个常数。把这个常数记为 \(l(l+1)\)，得到

\[
\boxed{
\frac{\mathrm d}{\mathrm dr}\left(r^2R'\right)-l(l+1)R=0
}
\]

和

\[
\boxed{
\frac1{\sin\theta}\frac{\mathrm d}{\mathrm d\theta}
\left(\sin\theta\,\Theta'\right)
+l(l+1)\Theta=0
}.
\]

\par\noindent\begin{minipage}{\linewidth}
\begin{center}
\includegraphics[width=\linewidth,height=.40\textheight,keepaspectratio]{\detokenize{assets/module-02/original/tikz-m2-topic20-v8-12.png}}
\par\smallskip{\small 图示：球坐标分离后径向方程与角向方程共享本征值}
\end{center}
\end{minipage}\par

径向方程是欧拉方程。设 \(R=r^s\)，则

\[
s(s-1)+2s-l(l+1)=0,
\]

即

\[
(s-l)(s+l+1)=0.
\]

所以

\[
\boxed{
R(r)=Ar^l+Br^{-l-1}
}.
\]

包含原点的区域要求电势在 \(r=0\) 有限，通常舍去
\(r^{-l-1}\)；延伸到无穷远且要求电势衰减的区域则舍去 \(r^l\)。

\subsubsection{紧随例：偶极势的径向幂次（M2-154）}\label{m2:ux7d27ux968fux4f8bux5076ux6781ux52bfux7684ux5f84ux5411ux5e42ux6b21m2-154}

三维真空无源球外域中，某个电势分量的角向依赖为
\(\cos\theta\)，\(\theta\) 从 \(+z\)
轴量起。要求该分量在无穷远衰减，源均位于被排除的内域。用球坐标分离变量确定其径向幂次，并说明它对应哪一阶多极项；不要把势的衰减幂次与场强的衰减幂次混同。

\par\noindent\begin{minipage}{\linewidth}
\textbf{解答}

\begin{center}
\includegraphics[width=\linewidth,height=.40\textheight,keepaspectratio]{\detokenize{assets/module-02/examples/M2-154.pdf}}
\par\smallskip{\small 图 M2-154：偶极角模与负二次径向幂配对}
\end{center}
\end{minipage}\par

图 M2-154 的正负两极角分布对应
\(P_1(\cos\theta)\)，所以角阶已固定为1。球外衰减再筛出
\(r^{-l-1}\)，得到偶极的 \(r^{-2}\)
电势；不能只根据图样是一阶，就在外域选增长的 \(r\) 分支。

电偶极子的远场角向因子为 \(\cos\theta=P_1(\cos\theta)\)，对应
\(l=1\)。无穷远处必须衰减，所以径向函数应取

\[
R\propto r^{-l-1}=r^{-2},
\]

这就得到偶极势 \(\varphi\propto\cos\theta/r^2\)。

\subsubsection{2.
从角向方程到勒让德方程}\label{m2:ux4eceux89d2ux5411ux65b9ux7a0bux5230ux52d2ux8ba9ux5fb7ux65b9ux7a0b}

令

\[
x=\cos\theta,
\qquad -1\le x\le1.
\]

利用

\[
\frac{\mathrm d}{\mathrm d\theta}
=-\sin\theta\frac{\mathrm d}{\mathrm dx},
\qquad
\sin^2\theta=1-x^2,
\]

角向方程化为

\[
\boxed{
\frac{\mathrm d}{\mathrm dx}
\left[(1-x^2)\frac{\mathrm d\Theta}{\mathrm dx}\right]
+l(l+1)\Theta=0
}.
\]

展开即

\[
\boxed{
(1-x^2)\Theta''-2x\Theta'+l(l+1)\Theta=0
}.
\]

这就是勒让德方程。

\subsection{\texorpdfstring{二、端点正则性为什么把 \(l\)
筛成整数}{二、端点正则性为什么把 l 筛成整数}}\label{m2:ux4e8cux7aefux70b9ux6b63ux5219ux6027ux4e3aux4ec0ux4e48ux628a-l-ux7b5bux6210ux6574ux6570}

\subsubsection{1. 幂级数递推}\label{m2:ux5e42ux7ea7ux6570ux9012ux63a8}

在常点 \(x=0\) 附近设

\[
\Theta(x)=\sum_{n=0}^{\infty}a_nx^n.
\]

代入勒让德方程，比较 \(x^n\) 的系数：

\[
(n+2)(n+1)a_{n+2}
+\left[l(l+1)-n(n+1)\right]a_n=0.
\]

因此

\[
\boxed{
a_{n+2}
=\frac{n(n+1)-l(l+1)}
{(n+1)(n+2)}a_n
}.
\]

偶数系数由 \(a_0\) 决定，奇数系数由 \(a_1\) 决定，两条递推链彼此独立。

若 \(l\) 是非负整数，则与 \(l\) 同奇偶性的递推链在 \(n=l\) 处出现

\[
n(n+1)-l(l+1)=0,
\]

因而截断成 \(l\)
次多项式。另一条递推链不截断，构成第二个线性无关解的一部分。

\subsubsection{2.
只在一个端点有限还不够}\label{m2:ux53eaux5728ux4e00ux4e2aux7aefux70b9ux6709ux9650ux8fd8ux4e0dux591f}

勒让德方程在 \(x=\pm1\) 处是正则奇点。对一般非整数参数
\(\nu\)，第一类勒让德函数 \(P_{\nu}(x)\) 可在 \(x=1\) 处有限，却通常在
\(x=-1\) 出现对数奇性。要求完整球面上的电势在南、北两极都有限，才得到

\[
\boxed{
l=0,1,2,\ldots
}.
\]

对应的正则本征函数是勒让德多项式 \(P_l(x)\)。第二类解 \(Q_l(x)\)
在极轴端点发散；只有物理区域避开相应端点时，才可能保留它。

\par\noindent\begin{minipage}{\linewidth}
\begin{center}
\includegraphics[width=\linewidth,height=.40\textheight,keepaspectratio]{\detokenize{assets/module-02/theory/unified7-legendre-regularity.pdf}}
\par\smallskip{\small 图示：极轴两端同时正则把勒让德阶数筛成非负整数}
\end{center}
\end{minipage}\par

这里应避免一个常见误解：并不是``所有不截断级数都在两个端点同时发散''，而是完整球面要求同一个解在两个端点同时正则，这一条件才把本征值离散化。

\subsubsection{紧随例：区域不包含极轴时可以出现什么（M2-155）}\label{m2:ux7d27ux968fux4f8bux533aux57dfux4e0dux5305ux542bux6781ux8f74ux65f6ux53efux4ee5ux51faux73b0ux4ec0ux4e48m2-155}

球坐标分离变量得到角向勒让德方程。现在求解域只含
\(\theta_1<\theta<\theta_2\)，其中
\(0<\theta_1<\theta_2<\pi\)，南北两极均不属于真实域。讨论能否仅因第二类勒让德函数在极轴上的奇性而一律排除它，说明本域的正则性与实际角边界条件应如何决定允许的线性组合。

\par\noindent\begin{minipage}{\linewidth}
\textbf{解答}

\begin{center}
\includegraphics[width=\linewidth,height=.40\textheight,keepaspectratio]{\detokenize{assets/module-02/examples/M2-155.pdf}}
\par\smallskip{\small 图 M2-155：端点正则性只在真实区域内要求}
\end{center}
\end{minipage}\par

图 M2-155
的蓝色角带避开了两极端点。分离后的角函数只需在这段真实角域内有限并满足其边界，极轴处的奇性位于域外；因此是否允许第二类或非整数阶函数，应由实际角带重新判断。

若只研究一个避开 \(\theta=0,\pi\)
的角环区域，端点奇性不在物理区域内，第二类勒让德函数或非整数阶函数可能成为合法解。是否舍去某个函数，必须由真实区域决定。

\subsection{三、勒让德多项式、正交性与边界展开}\label{m2:ux4e09ux52d2ux8ba9ux5fb7ux591aux9879ux5f0fux6b63ux4ea4ux6027ux4e0eux8fb9ux754cux5c55ux5f00}

\subsubsection{1.
罗德里格斯公式}\label{m2:ux7f57ux5fb7ux91ccux683cux65afux516cux5f0f}

取归一化

\[
P_l(1)=1,
\]

勒让德多项式可写成

\[
\boxed{
P_l(x)
=\frac{1}{2^ll!}
\frac{\mathrm d^l}{\mathrm dx^l}(x^2-1)^l
}.
\]

前四阶为

\[
P_0=1,
\qquad
P_1=x,
\]

\[
P_2=\frac12(3x^2-1),
\qquad
P_3=\frac12(5x^3-3x).
\]

它们满足奇偶性

\[
\boxed{
P_l(-x)=(-1)^lP_l(x)
}
\]

以及三项递推

\[
\boxed{
(l+1)P_{l+1}
=(2l+1)xP_l-lP_{l-1}
}.
\]

\par\noindent\begin{minipage}{\linewidth}
\begin{center}
\includegraphics[width=\linewidth,height=.40\textheight,keepaspectratio]{\detokenize{assets/module-02/original/tikz-m2-topic20-v8-14.png}}
\par\smallskip{\small 图示：前四阶勒让德多项式的奇偶性与零点}
\end{center}
\end{minipage}\par

\subsubsection{2.
正交性的来源}\label{m2:ux6b63ux4ea4ux6027ux7684ux6765ux6e90}

勒让德方程可写成自伴随形式

\[
-\frac{\mathrm d}{\mathrm dx}
\left[(1-x^2)P_l'\right]
=l(l+1)P_l.
\]

分别写出 \(l\) 与 \(l'\) 的方程，交叉相乘后相减并积分：

\[
\begin{aligned}
&\left[l(l+1)-l'(l'+1)\right]
\int_{-1}^{1}P_lP_{l'}\,\mathrm dx\\
&\quad=
\left[
(1-x^2)
\left(P_{l'}P_l'-P_lP_{l'}'\right)
\right]_{-1}^{1}.
\end{aligned}
\]

正则解及其导数在端点有限，而 \(1-x^2\) 在端点为零，所以右端为零。若
\(l\ne l'\)，便有

\[
\int_{-1}^{1}P_l(x)P_{l'}(x)\,\mathrm dx=0.
\]

结合归一化积分，

\[
\boxed{
\int_{-1}^{1}P_lP_{l'}\,\mathrm dx
=\frac{2}{2l+1}\delta_{ll'}
}.
\]

因此，足够光滑的轴对称球面边界函数可展开为

\[
f(\cos\theta)
=\sum_{l=0}^{\infty}A_lP_l(\cos\theta),
\]

其中

\[
\boxed{
A_l=\frac{2l+1}{2}
\int_{-1}^{1}f(x)P_l(x)\,\mathrm dx
}.
\]

\subsubsection{\texorpdfstring{紧随例：球面边界为
\(\cos^2\theta\)（M2-156）}{紧随例：球面边界为 \textbackslash cos\^{}2\textbackslash theta（M2-156）}}\label{m2:ux7d27ux968fux4f8bux7403ux9762ux8fb9ux754cux4e3a-cos2thetam2-156}

球心为原点、半径为 \(R>0\) 的真空球内无电荷，球面给定轴对称电势
\(\varphi(R,\theta)=V_0\cos^2\theta\)，\(\theta\) 从 \(+z\)
量起。要求球心电势有限。先将边界图样分解为适当勒让德模，再求球内电势，核对球面值及原点正则性。

\par\noindent\begin{minipage}{\linewidth}
\textbf{解答}

\begin{center}
\includegraphics[width=\linewidth,height=.40\textheight,keepaspectratio]{\detokenize{assets/module-02/examples/M2-156.pdf}}
\par\smallskip{\small 图 M2-156：平方余弦拆成零阶与二阶}
\end{center}
\end{minipage}\par

图 M2-156
的着色仅代表二阶偏差，完整平方余弦边界还包含常量平均值。两部分分别配球内径向幂
\(1\) 和 \(r^2\)，归一到球面后相加，球心只留下常量部分。

由

\[
P_2(x)=\frac12(3x^2-1)
\]

得

\[
\cos^2\theta
=\frac13+\frac23P_2(\cos\theta).
\]

若半径为 \(R\) 的球内无电荷，边界电势为
\(V_0\cos^2\theta\)，且原点处有限，则

\[
\boxed{
\varphi(r,\theta)
=V_0\left[
\frac13+\frac23\left(\frac rR\right)^2
P_2(\cos\theta)
\right]
}.
\]

每个角向阶数只能与相同 \(l\) 的径向幂次配对。

\subsection{四、非轴对称问题与伴随勒让德函数}\label{m2:ux56dbux975eux8f74ux5bf9ux79f0ux95eeux9898ux4e0eux4f34ux968fux52d2ux8ba9ux5fb7ux51fdux6570}

若保留方位角依赖，令

\[
\varphi(r,\theta,\phi)
=R(r)\Theta(\theta)\Phi(\phi).
\]

方位角方程为

\[
\Phi''+m^2\Phi=0.
\]

单值性

\[
\Phi(\phi+2\pi)=\Phi(\phi)
\]

要求

\[
\boxed{m\in\mathbb Z}.
\]

取 \(\Phi=e^{im\phi}\)，再令 \(x=\cos\theta\)，角向方程变为

\[
\boxed{
(1-x^2)\Theta''-2x\Theta'
+\left[
l(l+1)-\frac{m^2}{1-x^2}
\right]\Theta=0
}.
\]

对整数 \(l\ge0\)、非负整数
\(m=0,1,\ldots,l\)，正则伴随勒让德函数可定义为

\[
\boxed{
P_l^m(x)=(-1)^m(1-x^2)^{m/2}
\frac{\mathrm d^mP_l(x)}{\mathrm dx^m}
},\qquad 0\le m\le l.
\]

负阶由正阶定义，而不使用``负阶导数''：

\[
\boxed{P_l^{-m}(x)=(-1)^m\frac{(l-m)!}{(l+m)!}P_l^m(x)},
\qquad m=0,1,\ldots,l.
\]

例如 \(P_1^1(x)=-\sqrt{1-x^2}\)，所以
\(P_1^{-1}(x)=\sqrt{1-x^2}/2\)。完整球谐的阶数范围仍是
\(-l\le m\le l\)。

完整角向本征函数是球谐函数

\[
Y_{lm}(\theta,\phi)
=N_{lm}P_l^m(\cos\theta)e^{im\phi}.
\]

其中 \(m\) 决定方位方向的周期和节点面数，\(l-|m|\) 决定纬向节点数。

\par\noindent\begin{minipage}{\linewidth}
\begin{center}
\includegraphics[width=\linewidth,height=.40\textheight,keepaspectratio]{\detokenize{assets/module-02/theory/unified7-associated-legendre.pdf}}
\par\smallskip{\small 图示：球谐函数中 l 与 m 对角向节点的不同控制}
\end{center}
\end{minipage}\par

\subsubsection{紧随例：三个偶极角向模（M2-157）}\label{m2:ux7d27ux968fux4f8bux4e09ux4e2aux5076ux6781ux89d2ux5411ux6a21m2-157}

采用标准球坐标
\(x=r\sin\theta\cos\phi\)、\(y=r\sin\theta\sin\phi\)、\(z=r\cos\theta\)。分别写出沿
\(z,x,y\)
三个方向的实偶极角函数，说明这些函数属于同一球谐阶数的理由，并指出何者对所选
\(z\) 轴轴对称。只比较角结构，不要求给定偶极矩数值。

\par\noindent\begin{minipage}{\linewidth}
\textbf{解答}

\begin{center}
\includegraphics[width=\linewidth,height=.40\textheight,keepaspectratio]{\detokenize{assets/module-02/examples/M2-157.pdf}}
\par\smallskip{\small 图 M2-157：三个实偶极模只是取向不同}
\end{center}
\end{minipage}\par

图 M2-157 的三条坐标方向分别对应三个实偶极角模。用单位径向矢量在
\(x,y,z\) 三轴的投影读取它们，可直接看出三者都属于
\(l=1\)，只是方位因子的实线性组合不同。

三个实偶极角向函数可取

\[
\cos\theta,\qquad
\sin\theta\cos\phi,\qquad
\sin\theta\sin\phi.
\]

它们分别对应沿 \(z\)、\(x\)、\(y\) 方向的偶极子。三者都有
\(l=1\)，差别只在 \(m\) 的线性组合。

\subsection{五、柱坐标分离变量与贝塞尔方程}\label{m2:ux4e94ux67f1ux5750ux6807ux5206ux79bbux53d8ux91cfux4e0eux8d1dux585eux5c14ux65b9ux7a0b}

\subsubsection{1.
分离常数的符号}\label{m2:ux5206ux79bbux5e38ux6570ux7684ux7b26ux53f7}

柱坐标拉普拉斯方程为

\[
\frac1\rho\frac{\partial}{\partial\rho}
\left(\rho\frac{\partial\varphi}{\partial\rho}\right)
+\frac1{\rho^2}\frac{\partial^2\varphi}{\partial\phi^2}
+\frac{\partial^2\varphi}{\partial z^2}=0.
\]

令

\[
\varphi=R(\rho)e^{im\phi}Z(z),
\qquad m\in\mathbb Z.
\]

若取

\[
Z''=+k^2Z,
\]

则径向方程为

\[
\boxed{
\rho^2R''+\rho R'
+(k^2\rho^2-m^2)R=0
},
\]

其解为

\[
R=A J_m(k\rho)+B Y_m(k\rho).
\]

若取

\[
Z''=-k^2Z,
\]

则径向方程为

\[
\boxed{
\rho^2R''+\rho R'
-(k^2\rho^2+m^2)R=0
},
\]

其解为

\[
R=A I_m(k\rho)+B K_m(k\rho).
\]

\par\noindent\begin{minipage}{\linewidth}
\begin{center}
\includegraphics[width=\linewidth,height=.40\textheight,keepaspectratio]{\detokenize{assets/module-02/original/tikz-m2-topic20-v8-16.png}}
\par\smallskip{\small 图示：柱坐标分离常数符号决定贝塞尔函数类型}
\end{center}
\end{minipage}\par

指数衰减函数 \(Z=e^{-kz}\) 满足 \(Z''=k^2Z\)，所以半无限圆柱内沿 \(z\)
衰减的静电势通常与普通贝塞尔函数 \(J_m\)
配对。不能仅凭``出现指数衰减''就误选修正贝塞尔函数。

\subsubsection{2.
弗罗贝尼乌斯级数}\label{m2:ux5f17ux7f57ux8d1dux5c3cux4e4cux65afux7ea7ux6570}

令

\[
x=k\rho,
\]

普通贝塞尔方程化为

\[
x^2y''+xy'+(x^2-m^2)y=0.
\]

\(x=0\) 是正则奇点。设

\[
y=\sum_{n=0}^{\infty}a_nx^{n+s},
\qquad a_0\ne0.
\]

最低次幂给出指标方程

\[
\boxed{s^2-m^2=0},
\]

即

\[
s=\pm m.
\]

比较一般次幂，得到

\[
\boxed{
a_n
=-\frac{a_{n-2}}{(n+s)^2-m^2}
},
\qquad n\ge2.
\]

取 \(s=m\ge0\)，由 \(x^{m+1}\) 的系数可知
\(a_1=0\)，故只剩偶数项。按标准归一化得到

\[
\boxed{
J_m(x)
=\sum_{p=0}^{\infty}
\frac{(-1)^p}{p!\,\Gamma(p+m+1)}
\left(\frac x2\right)^{2p+m}
}.
\]

当 \(m\) 为非负整数，

\[
\Gamma(p+m+1)=(p+m)!.
\]

\subsubsection{3.
第二类与修正贝塞尔函数}\label{m2:ux7b2cux4e8cux7c7bux4e0eux4feeux6b63ux8d1dux585eux5c14ux51fdux6570}

整数阶时，\(J_{-m}\) 与 \(J_m\) 不独立，第二个独立解定义为
\(Y_m\)。原点附近，

\[
J_m(x)\sim\frac1{m!}\left(\frac x2\right)^m,
\]

\[
Y_0(x)\sim\frac2\pi\ln x,
\]

\[
Y_m(x)\sim-\frac{(m-1)!}{\pi}
\left(\frac2x\right)^m,
\qquad m\ge1.
\]

因此，若真实区域包含轴线 \(\rho=0\)，有限电势通常要求舍去 \(Y_m\)。

修正贝塞尔函数的远场行为为

\[
I_m(x)\sim\frac{e^x}{\sqrt{2\pi x}},
\qquad
K_m(x)\sim\sqrt{\frac{\pi}{2x}}e^{-x}.
\]

所以外部无界区域要求径向衰减时，通常保留
\(K_m\)；内部包含轴线时则通常保留 \(I_m\)。

\subsubsection{紧随例：区域是否包含轴线决定函数选择（M2-158）}\label{m2:ux7d27ux968fux4f8bux533aux57dfux662fux5426ux5305ux542bux8f74ux7ebfux51b3ux5b9aux51fdux6570ux9009ux62e9m2-158}

柱坐标分离变量在不同分离常数下会给出普通贝塞尔分支 \(J_m,Y_m\)
或修正分支
\(I_m,K_m\)，\(m\ge0\)、\(k>0\)。比较轴线属于真实域且电势正则的实心圆柱内域，与不含轴线、要求径向无穷远衰减的圆柱外域
\(\rho>a>0\)。说明 \(Y_m(k\rho)\)、\(K_m(k\rho)\)
在两种区域中是否可因轴线奇性被排除，以及修正方程的外域衰减条件如何筛选分支。

\par\noindent\begin{minipage}{\linewidth}
\textbf{解答}

\begin{center}
\includegraphics[width=\linewidth,height=.40\textheight,keepaspectratio]{\detokenize{assets/module-02/examples/M2-158.pdf}}
\par\smallskip{\small 图 M2-158：外域衰减和轴线正则是两项不同条件}
\end{center}
\end{minipage}\par

图 M2-158
的实心域要求轴线正则，外域则根本不包含轴线。对于给定分离常数，先排除真实域内不允许的奇点，再检查远端是否增长；因此某些在内域须舍去的分支，在外域恰好满足衰减。

在实心圆柱内部，\(\rho=0\) 属于真实区域，\(Y_m\) 与 \(K_m\)
的轴线奇性通常不允许。对于圆柱外部区域
\(\rho>a\)，轴线不属于真实区域，若还要求无穷远衰减，则 \(K_m\)
反而是自然选择。

\subsection{六、贝塞尔零点、正交性与圆柱边界}\label{m2:ux516dux8d1dux585eux5c14ux96f6ux70b9ux6b63ux4ea4ux6027ux4e0eux5706ux67f1ux8fb9ux754c}

\subsubsection{\texorpdfstring{1. 外圆边界使 \(k\)
离散}{1. 外圆边界使 k 离散}}\label{m2:ux5916ux5706ux8fb9ux754cux4f7f-k-ux79bbux6563}

若圆柱侧壁 \(\rho=a\) 接地，且包含轴线，则径向函数取

\[
R=J_m(k\rho).
\]

边界条件

\[
R(a)=0
\]

要求

\[
J_m(ka)=0.
\]

记 \(J_m\) 的第 \(n\) 个正零点为 \(j_{mn}\)，则

\[
\boxed{
k_{mn}=\frac{j_{mn}}{a}
}.
\]

\par\noindent\begin{minipage}{\linewidth}
\begin{center}
\includegraphics[width=\linewidth,height=.40\textheight,keepaspectratio]{\detokenize{assets/module-02/original/tikz-m2-topic20-v8-17.png}}
\par\smallskip{\small 图示：轴线有界性和外圆零势条件依次筛选贝塞尔本征模}
\end{center}
\end{minipage}\par

角向整数 \(m\) 来自 \(2\pi\) 周期性，径向序号 \(n\)
来自外边界零点，两者来源不同。

\subsubsection{2. 径向正交性}\label{m2:ux5f84ux5411ux6b63ux4ea4ux6027}

固定 \(m\)，令

\[
R_n(\rho)=J_m\left(\frac{j_{mn}}a\rho\right).
\]

径向方程可写成

\[
-\frac{\mathrm d}{\mathrm d\rho}
\left(\rho R_n'\right)
+\frac{m^2}{\rho}R_n
=k_{mn}^2\rho R_n.
\]

这是权函数为 \(\rho\) 的自伴随本征值问题。因此

\[
\boxed{
\int_0^a
\rho R_n(\rho)R_{n'}(\rho)\,\mathrm d\rho
=\frac{a^2}{2}
\left[J_{m+1}(j_{mn})\right]^2
\delta_{nn'}
}.
\]

径向边界函数可按这些本征模投影。若

\[
f(\rho)
=\sum_{n=1}^{\infty}A_n
J_m\left(\frac{j_{mn}}a\rho\right),
\]

则

\[
\boxed{
A_n
=\frac{2}
{a^2[J_{m+1}(j_{mn})]^2}
\int_0^a
\rho f(\rho)
J_m\left(\frac{j_{mn}}a\rho\right)
\mathrm d\rho
}.
\]

\par\noindent\begin{minipage}{\linewidth}
\begin{center}
\includegraphics[width=\linewidth,height=.40\textheight,keepaspectratio]{\detokenize{assets/module-02/original/tikz-m2-topic20-v8-18.png}}
\par\smallskip{\small 图示：勒让德模与贝塞尔模都通过自伴随正交性投影边界数据}
\end{center}
\end{minipage}\par

\subsubsection{紧随例：半无限接地圆筒（M2-159）}\label{m2:ux7d27ux968fux4f8bux534aux65e0ux9650ux63a5ux5730ux5706ux7b52m2-159}

半无限圆柱的真空无源内域为 \(0\le\rho<a,z>0\)，\(a>0\)，问题轴对称。侧壁
\(\rho=a\) 接地，底面 \(z=0\) 给定足够规则的电势
\(f(\rho)\)，要求轴线正则且 \(z\to\infty\) 时势衰减。用 \(J_0\)
的正零点组织分离变量级数，写出径向模态、轴向因子及由底面数据确定系数的正交投影关系。

\par\noindent\begin{minipage}{\linewidth}
\textbf{解答}

\begin{center}
\includegraphics[width=\linewidth,height=.40\textheight,keepaspectratio]{\detokenize{assets/module-02/examples/M2-159.pdf}}
\par\smallskip{\small 图 M2-159：圆筒本征模同时满足三个边界要求}
\end{center}
\end{minipage}\par

图 M2-159
的两侧是接地圆筒壁，底面是给定的非齐次数据，顶部表示半无限延伸。\(J_0\)
保证轴线有限，其零点固定侧壁条件；每个径向零点还决定同一模态沿 \(z\)
的指数衰减率。

设区域为

\[
0\le\rho<a,\qquad z>0,
\]

侧壁 \(\rho=a\) 接地，轴对称底面电势为 \(\varphi(\rho,0)=f(\rho)\)，且
\(z\to\infty\) 时电势衰减。解为

\[
\boxed{
\varphi(\rho,z)
=\sum_{n=1}^{\infty}A_n
J_0\left(\frac{j_{0n}}a\rho\right)
e^{-j_{0n}z/a}
},
\]

其中

\[
A_n
=\frac{2}
{a^2[J_1(j_{0n})]^2}
\int_0^a
\rho f(\rho)
J_0\left(\frac{j_{0n}}a\rho\right)
\mathrm d\rho.
\]

这里 \(e^{-j_{0n}z/a}\) 与 \(J_0(j_{0n}\rho/a)\)
的组合同时满足远端衰减、轴线有限和侧壁零势。

\subsection{七、特殊函数与边值方法的选择}\label{m2:ux4e03ux7279ux6b8aux51fdux6570ux4e0eux8fb9ux503cux65b9ux6cd5ux7684ux9009ux62e9}

{\def\LTcaptype{none} % do not increment counter
\begin{longtable}[]{@{}llll@{}}
\toprule\noalign{}
真实区域与边界 & 通常保留的解 & 通常舍去的解 & 主要原因 \\
\midrule\noalign{}
\endhead
\bottomrule\noalign{}
\endlastfoot
完整球面角域 & \(P_l,\ P_l^m\) & \(Q_l,\ Q_l^m\) & 两极正则 \\
包含球心的内部区域 & \(r^l\) & \(r^{-l-1}\) & 原点有限 \\
延伸到无穷远的球外区域 & \(r^{-l-1}\) & \(r^l\) & 远场衰减 \\
包含圆柱轴线 & \(J_m\) 或 \(I_m\) & \(Y_m\) 或 \(K_m\) & 轴线有限 \\
不含轴线且要求径向衰减 & \(K_m\) & \(I_m\) & 无穷远衰减 \\
接地圆柱侧壁 & \(J_m(j_{mn}\rho/a)\) & 任意连续 \(k\) & 外边界零点 \\
\end{longtable}
}

函数名称本身不能决定取舍。先确定真实区域是否包含奇点，再检查有限性、无穷远条件和导体边界。

\subsection{八、练习与解答}\label{m2:ux516bux7ec3ux4e60ux4e0eux89e3ux7b54-7}

\subsubsection{作业一：勒让德递推与前三阶多项式}\label{m2:ux4f5cux4e1aux4e00ux52d2ux8ba9ux5fb7ux9012ux63a8ux4e0eux524dux4e09ux9636ux591aux9879ux5f0f}

\textbf{题号：M2-160}

\textbf{层级：} 决赛　　\textbf{方法：} 边值问题

对轴对称勒让德方程
\((1-x^2)y''-2xy'+\ell(\ell+1)y=0\)，\(-1\le x\le1\)，要求解为端点有限的多项式，并按
\(P_{\ell}(1)=1\)
归一化。用幂级数代入推导系数递推关系，说明级数截断与非负整数阶的联系，再具体求
\(P_0,P_1,P_2,P_3\)。

\par\noindent\begin{minipage}{\linewidth}
\textbf{解答}

\begin{center}
\includegraphics[width=\linewidth,height=.40\textheight,keepaspectratio]{\detokenize{assets/module-02/examples/M2-160.pdf}}
\par\smallskip{\small 图 M2-160：勒让德递推与前三阶多项式}
\end{center}
\end{minipage}\par

图 M2-160 左侧把偶、奇系数分成两条递推链，右侧列出归一化后的
\(P_0,P_1,P_2,P_3\)。幂级数递推决定每一支的系数比例，整数阶时在对应幂次终止，再用
\(P_l(1)=1\) 固定常数；因此偶阶多项式只含偶次幂，奇阶只含奇次幂。

递推关系为

\[
a_{n+2}
=\frac{n(n+1)-l(l+1)}
{(n+1)(n+2)}a_n.
\]

\(l=0\) 时偶链在 \(n=0\) 截断，取 \(P_0(1)=1\)：

\[
P_0=1.
\]

\(l=1\) 时奇链在 \(n=1\) 截断：

\[
P_1=x.
\]

\(l=2\) 时由 \(n=0\)，

\[
a_2=-3a_0.
\]

故 \(\Theta=a_0(1-3x^2)\)。令 \(P_2(1)=1\)，得到

\[
P_2=\frac12(3x^2-1).
\]

\(l=3\) 时由 \(n=1\)，

\[
a_3=-\frac53a_1.
\]

归一化后

\[
P_3=\frac12(5x^3-3x).
\]

\subsubsection{作业二：球内轴对称边值问题}\label{m2:ux4f5cux4e1aux4e8cux7403ux5185ux8f74ux5bf9ux79f0ux8fb9ux503cux95eeux9898}

\textbf{题号：M2-161}

\textbf{层级：} 复赛　　\textbf{方法：} 边值问题

真空球形区域 \(0\le r<R\) 内没有电荷或点奇性，\(R>0\)。球面电势给定为
\(V_0(1+3\cos\theta+2\cos^2\theta)\)，\(\theta\) 从 \(+z\)
轴量起。将边界值按归一化勒让德多项式展开，求球心正则的球内电势，并逐项核对边界数据。

\par\noindent\begin{minipage}{\linewidth}
\textbf{解答}

\begin{center}
\includegraphics[width=\linewidth,height=.40\textheight,keepaspectratio]{\detokenize{assets/module-02/examples/M2-161.pdf}}
\par\smallskip{\small 图 M2-161：球内轴对称边值问题}
\end{center}
\end{minipage}\par

图 M2-161 的边界先把 \(\cos^2\theta\)
换成零阶和二阶模态，再与已有一次项合并。球内分别配
\(1,r,r^2\)，各阶独立在 \(r=R\)
归一；不要把球面上的二次角函数直接当作整个内部的纯二阶势。

利用

\[
\cos\theta=P_1(\cos\theta),
\]

\[
\cos^2\theta
=\frac13+\frac23P_2(\cos\theta),
\]

边界函数化为

\[
\varphi(R,\theta)
=V_0\left[
\frac53+3P_1(\cos\theta)
+\frac43P_2(\cos\theta)
\right].
\]

球内包含原点，所以第 \(l\) 阶角向模与 \((r/R)^l\) 配对：

\[
\boxed{
\varphi(r,\theta)
=V_0\left[
\frac53
+3\frac rR P_1(\cos\theta)
+\frac43\left(\frac rR\right)^2
P_2(\cos\theta)
\right]
}.
\]

在 \(r=R\) 处立即恢复原边界函数，在 \(r=0\) 处保持有限。

\subsubsection{作业三：贝塞尔函数的首三项}\label{m2:ux4f5cux4e1aux4e09ux8d1dux585eux5c14ux51fdux6570ux7684ux9996ux4e09ux9879}

\textbf{题号：M2-162}

\textbf{层级：} 决赛　　\textbf{方法：} 边值问题

考虑第一类贝塞尔方程 \(x^2y''+xy'+(x^2-m^2)y=0\)。对 \(m=0,1\) 分别在
\(x=0\) 附近作 Frobenius 展开，选取正则分支，使用 \(J_0(0)=1\) 与
\(J_1(x)\sim x/2\) 的标准归一化。推导系数递推式，写出 \(J_0,J_1\)
的前三个非零项，并说明第二独立分支为何不满足实心域轴线正则条件。

\par\noindent\begin{minipage}{\linewidth}
\textbf{解答}

\begin{center}
\includegraphics[width=\linewidth,height=.40\textheight,keepaspectratio]{\detokenize{assets/module-02/examples/M2-162.pdf}}
\par\smallskip{\small 图 M2-162：贝塞尔函数的首三项}
\end{center}
\end{minipage}\par

图 M2-162
的级数在原点附近分别从常量与一次项开始。将弗罗贝尼乌斯递推按相邻偶次步进展开，\(J_0\)
保留偶幂、\(J_1\) 保留奇幂；截断曲线只用于局部检验，不能替代大参数渐近。

在图 M2-162 的轴线邻域令

\[
y=x^s\sum_{n=0}^{\infty}a_nx^n,\qquad a_0\ne0.
\]

代入方程，最低次项给指标方程 \(s^2-m^2=0\)。正则分支取
\(s=m\)；逐次比较得

\[
[(n+m)^2-m^2]a_n+a_{n-2}=0,\qquad
a_n=-\frac{a_{n-2}}{n(n+2m)},\quad n\ge2,\qquad a_1=0.
\]

对 \(m=0\)，取 \(a_0=1\)，有 \(a_2=-1/4\)、\(a_4=1/64\)；对 \(m=1\)，取
\(a_0=1/2\)，有 \(a_2=-1/16\)、\(a_4=1/384\)。第二独立分支在 \(m=0\)
时含 \(Y_0(x)\sim(2/\pi)\ln x\)，在 \(m=1\) 时含
\(Y_1(x)\sim-2/(\pi x)\)，均在轴线上发散，不能用于要求正则的实心域。由递推归纳得到标准级数

\[
J_m(x)
=\sum_{p=0}^{\infty}
\frac{(-1)^p}{p!(p+m)!}
\left(\frac x2\right)^{2p+m}.
\]

取 \(m=0\)：

\[
\boxed{
J_0(x)
=1-\frac{x^2}{4}
+\frac{x^4}{64}
+O(x^6)
}.
\]

取 \(m=1\)：

\[
\boxed{
J_1(x)
=\frac x2-\frac{x^3}{16}
+\frac{x^5}{384}
+O(x^7)
}.
\]

因此 \(J_0(0)=1\)，而 \(J_1(x)\sim x/2\)，两者在轴线上均有限。

\subsubsection{作业四：半无限圆筒的单一本征模}\label{m2:ux4f5cux4e1aux56dbux534aux65e0ux9650ux5706ux7b52ux7684ux5355ux4e00ux672cux5f81ux6a21}

\textbf{题号：M2-163}

\textbf{层级：} 复赛　　\textbf{方法：} 边值问题

轴对称真空无源域为半无限圆柱
\(0\le\rho<a,z>0\)，侧壁接地，轴线要求正则，远端 \(z\to\infty\)
的势趋零。底面给定 \(\varphi(\rho,0)=V_0J_0(j_{02}\rho/a)\)，其中
\(a>0\)，\(j_{02}\) 是 \(J_0\)
的第二个正零点。求域内电势，解释底面正交投影为何仅留下一个径向模态，并检验侧壁及远端条件。

\par\noindent\begin{minipage}{\linewidth}
\textbf{解答}

\begin{center}
\includegraphics[width=\linewidth,height=.40\textheight,keepaspectratio]{\detokenize{assets/module-02/examples/M2-163.pdf}}
\par\smallskip{\small 图 M2-163：半无限圆筒的单一本征模}
\end{center}
\end{minipage}\par

图 M2-163 的底面数据已经是第二径向本征函数。按带权 \(\rho\)
的贝塞尔正交关系投影，其他径向系数均为零；同一零点 \(j_{02}\)
还确定轴向衰减率，而非另选第一零点。

侧壁条件要求径向模为

\[
J_0\left(\frac{j_{0n}}a\rho\right).
\]

底面函数恰好是 \(n=2\) 的一个本征函数。由正交性，对 \(n\ne2\)，

\[
\int_0^a\rho
J_0\left(\frac{j_{02}}a\rho\right)
J_0\left(\frac{j_{0n}}a\rho\right)
\mathrm d\rho=0.
\]

所以只有 \(A_2=V_0\)，其余系数均为零。沿 \(z\) 衰减的解为

\[
\boxed{
\varphi(\rho,z)
=V_0
J_0\left(\frac{j_{02}}a\rho\right)
e^{-j_{02}z/a}
}.
\]

该解在轴线上有限，在 \(\rho=a\) 为零，并在 \(z\to\infty\) 时衰减。

在边值问题中，微分方程给出候选函数，端点正则性和周期性限制允许的阶数，边界零点确定离散本征值，正交性用于展开边界数据。

\section{模块二附录二：二维静电场的保角变换法（拔高选讲）}\label{m2:ux6a21ux5757ux4e8cux9644ux5f55ux4e8cux4e8cux7ef4ux9759ux7535ux573aux7684ux4fddux89d2ux53d8ux6362ux6cd5ux62d4ux9ad8ux9009ux8bb2}

\begin{quote}
阅读准备：复数同时表示平面上的点 写 \(z=x+iy\) 后，复函数
\(w=f(z)=u(x,y)+iv(x,y)\) 同时给出两个实函数。解析性要求柯西---黎曼关系
\(u_x=v_y\)、\(u_y=-v_x\)；它使 \(u,v\)
都满足二维拉普拉斯方程，并让局部角度在非临界点保持不变。多值函数还必须明确割线与分支。

\textbf{例 M2-164：平方映射把直角区域展成半平面}
\end{quote}

取复平面 \(z=re^{\mathrm i\theta}\) 的开直角劈形
\(r>0,0<\theta<\pi/2\)。用 \(w=z^2\)
映射此域，分别追踪两条边界射线与域内点的像，说明所选域为何一一对应到上半平面，并指出若把原域扩大到整个去心平面会发生什么。

\par\noindent\begin{minipage}{\linewidth}
\textbf{解答}

\begin{center}
\includegraphics[width=\linewidth,height=.40\textheight,keepaspectratio]{\detokenize{assets/module-02/examples/M2-164.pdf}}
\par\smallskip{\small 图 M2-164：平方映射把直角区域展成半平面}
\end{center}
\end{minipage}\par

图 M2-164 的直角扇区用极角 \(0<\theta<\pi/2\) 描述。平方映射把模长变为
\(r^2\)、把极角加倍，因此两条直边分别映到实轴两侧，扇区内部对应上半平面。

\(w=z^2\) 把极角 \(\theta\) 变成 \(2\theta\)，因此张角为 \(\pi/2\)
的区域可映成张角为 \(\pi\) 的上半平面。

二维静电问题可以利用调和函数与解析函数的关系，把复杂边界映射为较易处理的边界。本附录所说的``二维''是指系统沿
\(z\) 轴平移不变，电势只依赖 \(x,y\)；电荷与电容均按 \(z\)
方向单位长度计。

这一方法包含三层内容：

\begin{enumerate}
\def\labelenumi{\arabic{enumi}.}
\tightlist
\item
  \textbf{物理层}：无源区域满足拉普拉斯方程，导体边界是等势线；
\item
  \textbf{函数层}：电势可作为解析函数的实部，虚部给出电力线；
\item
  \textbf{几何层}：解析且局部单叶的映射保持角度，把边界形状变简单。
\end{enumerate}

\par\noindent\begin{minipage}{\linewidth}
\begin{center}
\includegraphics[width=\linewidth,height=.40\textheight,keepaspectratio]{\detokenize{assets/module-02/original/tikz-m2-topic20-v8-19.png}}
\par\smallskip{\small 图示：二维静电场中物理问题、复势与保角映射的对应关系}
\end{center}
\end{minipage}\par

\subsection{一、适用范围：先确认``二维、无源、可映射''}\label{m2:ux4e00ux9002ux7528ux8303ux56f4ux5148ux786eux8ba4ux4e8cux7ef4ux65e0ux6e90ux53efux6620ux5c04}

设考察区域为 \(D\subset\mathbb C\)，以

\[
z=x+\mathrm i y
\]

表示平面内一点。若区域内没有体电荷，且介质均匀、各向同性，则

\[
\nabla^2\varphi
=
\frac{\partial^2\varphi}{\partial x^2}
+
\frac{\partial^2\varphi}{\partial y^2}
=0.
\]

保角变换法直接适用于下列问题：

\begin{itemize}
\tightlist
\item
  场与边界沿第三个方向平移不变；
\item
  所映射的区域内部无自由电荷；
\item
  介电常数在该区域内为常数；
\item
  导体边界给定电势，或绝缘边界给定法向导数；
\item
  映射在所选区域内解析、单叶，并且除单独处理的角点外满足 \(f'(z)\ne0\)。
\end{itemize}

若区域内含点电荷、线电荷或介质分界面，方法仍可使用，但必须显式保留奇点或分段匹配，不能把整个区域直接当作光滑调和域。

\begin{quote}
先检查区域
从积分形式转为微分方程、从原区域转到像区域之前，都必须检查奇点、割线、角点和多连通结构。映射公式正确，不等于所选分支和边界对应关系正确。
\end{quote}

\subsection{二、调和共轭：电势与电力线来自同一个解析函数}\label{m2:ux4e8cux8c03ux548cux5171ux8f6dux7535ux52bfux4e0eux7535ux529bux7ebfux6765ux81eaux540cux4e00ux4e2aux89e3ux6790ux51fdux6570}

\subsubsection{1.
从拉普拉斯方程构造调和共轭}\label{m2:ux4eceux62c9ux666eux62c9ux65afux65b9ux7a0bux6784ux9020ux8c03ux548cux5171ux8f6d}

设 \(\varphi(x,y)\) 在单连通区域 \(D\) 内二阶连续并满足

\[
\varphi_{xx}+\varphi_{yy}=0.
\]

定义微分形式

\[
\mathrm d\psi=-\varphi_y\,\mathrm dx+\varphi_x\,\mathrm dy.
\]

它成为某个单值函数 \(\psi\) 的全微分，需要路径无关。计算闭合条件：

\[
\frac{\partial}{\partial x}(\varphi_x)
-
\frac{\partial}{\partial y}(-\varphi_y)
=
\varphi_{xx}+\varphi_{yy}
=0.
\]

因此，在单连通区域中可定义

\[
\psi(P)-\psi(P_0)
=
\int_{P_0}^{P}
\left(-\varphi_y\,\mathrm dx+\varphi_x\,\mathrm dy\right),
\]

而积分与路径无关。于是

\[
\boxed{
\varphi_x=\psi_y,\qquad
\varphi_y=-\psi_x
}
\]

即柯西---黎曼条件。

令

\[
W(z)=\varphi(x,y)+\mathrm i\psi(x,y),
\]

则 \(W\) 在 \(D\) 内解析，称为\textbf{复势}。其中
\(\varphi=\text{常量}\) 是等势线，\(\psi=\text{常量}\) 是电力线。

\par\noindent\begin{minipage}{\linewidth}
\begin{center}
\includegraphics[width=\linewidth,height=.40\textheight,keepaspectratio]{\detokenize{assets/module-02/original/tikz-m2-topic20-v8-20.png}}
\par\smallskip{\small 图示：调和共轭使等势线与电力线组成正交曲线网}
\end{center}
\end{minipage}\par

\subsubsection{\texorpdfstring{2. 为什么 \(\psi=\text{常量}\)
是电力线}{2. 为什么 \textbackslash psi=\textbackslash text\{常量\} 是电力线}}\label{m2:ux4e3aux4ec0ux4e48-psitextux5e38ux91cf-ux662fux7535ux529bux7ebf}

由柯西---黎曼条件，

\[
\nabla\psi=(-\varphi_y,\varphi_x).
\]

而

\[
\mathbf E=-\nabla\varphi=(-\varphi_x,-\varphi_y).
\]

两者点积为

\[
\mathbf E\cdot\nabla\psi
=
(-\varphi_x)(-\varphi_y)
+
(-\varphi_y)\varphi_x
=0.
\]

因此 \(\mathbf E\) 与 \(\nabla\psi\) 垂直，也就与曲线
\(\psi=\text{常量}\) 相切。另一方面，

\[
\nabla\varphi\cdot\nabla\psi=0,
\]

所以等势线与电力线在非临界点处正交。

\subsubsection{3. 复电场}\label{m2:ux590dux7535ux573a}

由

\[
W'(z)
=
\varphi_x+\mathrm i\psi_x
=
\varphi_x-\mathrm i\varphi_y,
\]

定义复电场

\[
\mathcal E(z)=E_x-\mathrm iE_y,
\]

便有

\[
\boxed{\mathcal E(z)=-W'(z)}.
\]

这个记号同时编码电场大小与方向：

\[
|\mathcal E|=|\mathbf E|,
\qquad
\arg\mathcal E=-\arg(E_x+\mathrm iE_y).
\]

\textbf{例 M2-165：复势的实虚部分形成正交曲线网}

在二维静电问题中取 \(z=x+\mathrm iy\)，约定复势
\(W(z)=\varphi(x,y)+\mathrm i\psi(x,y)\)，电场为
\(\mathbf E=-\nabla\varphi\)。给定解析复势
\(W(z)=z^2\)，把它写成实部与虚部，求
\(\varphi,\psi,E_x,E_y\)，确定电场零点，并说明等势线与共轭函数等值线在非零场处的关系。

\par\noindent\begin{minipage}{\linewidth}
\textbf{解答}

\begin{center}
\includegraphics[width=\linewidth,height=.40\textheight,keepaspectratio]{\detokenize{assets/module-02/examples/M2-165.pdf}}
\par\smallskip{\small 图 M2-165：复势的实虚部分形成正交曲线网}
\end{center}
\end{minipage}\par

图 M2-165 的绿色曲线是 \(\psi=2xy\) 的常量族，红箭头沿
\((-2x,2y)\)。分别对实部求负梯度、对复势求导应得到相同场；原点两种一阶导数都消失，不能在那里套用普通正交网的唯一切向。

取

\[
W(z)=z^2=(x^2-y^2)+\mathrm i(2xy).
\]

于是

\[
\varphi=x^2-y^2,\qquad
\psi=2xy,
\]

\[
\mathcal E=-2z,
\qquad
\mathbf E=(-2x,\,2y).
\]

双曲线 \(\varphi=\text{常量}\) 与 \(\psi=\text{常量}\) 彼此正交；在原点
\(W'(0)=0\)，电场为零，曲线网在此处不能用普通的``唯一切向''描述。

\subsection{三、保角变换为什么保持拉普拉斯方程}\label{m2:ux4e09ux4fddux89d2ux53d8ux6362ux4e3aux4ec0ux4e48ux4fddux6301ux62c9ux666eux62c9ux65afux65b9ux7a0b}

设

\[
w=f(z)=u(x,y)+\mathrm iv(x,y)
\]

在 \(D\) 内解析，\(\Phi(u,v)\) 是像区域中的电势。把它拉回原区域：

\[
\varphi(x,y)=\Phi(u(x,y),v(x,y)).
\]

解析性给出

\[
u_x=v_y,\qquad u_y=-v_x,
\]

并且 \(u,v\)
均为调和函数。对复合函数求二阶导数，交叉项借助柯西---黎曼条件消去，得到

\[
\boxed{
\nabla_z^2\varphi
=
|f'(z)|^2\nabla_w^2\Phi
}.
\]

所以只要

\[
\nabla_w^2\Phi=0,
\]

就有

\[
\nabla_z^2\varphi=0.
\]

这里的 \(|f'(z)|\) 是局部长度伸缩率，\(\arg f'(z)\) 是局部旋转角。只要
\(f'(z)\ne0\)，任意两条相交曲线的夹角在映射前后保持不变。

\par\noindent\begin{minipage}{\linewidth}
\begin{center}
\includegraphics[width=\linewidth,height=.40\textheight,keepaspectratio]{\detokenize{assets/module-02/original/tikz-m2-topic20-v8-21.png}}
\par\smallskip{\small 图示：解析映射在微小邻域内等效于统一伸缩与旋转}
\end{center}
\end{minipage}\par

\subsubsection{1.
电场如何变换}\label{m2:ux7535ux573aux5982ux4f55ux53d8ux6362}

设像平面中的复势为 \(W_w(w)\)，原平面中的复势为

\[
W_z(z)=W_w(f(z)).
\]

链式法则给出

\[
\mathcal E_z(z)
=
-\frac{\mathrm dW_z}{\mathrm dz}
=
-\frac{\mathrm dW_w}{\mathrm dw}f'(z)
=
\boxed{\mathcal E_w(f(z))\,f'(z)}.
\]

因此

\[
\boxed{
|\mathbf E_z|
=
|\mathbf E_w|\,|f'(z)|
}.
\]

映射越强烈地放大像平面中的小长度，拉回原平面后的场就越强。

\subsubsection{2.
线电荷密度与总电荷如何变换}\label{m2:ux7ebfux7535ux8377ux5bc6ux5ea6ux4e0eux603bux7535ux8377ux5982ux4f55ux53d8ux6362}

导体表面外侧有

\[
\sigma=\varepsilon E_n.
\]

保角映射保持法线与切线的正交关系，故

\[
\sigma_z=|f'(z)|\,\sigma_w.
\]

而弧长满足

\[
\mathrm ds_w=|f'(z)|\,\mathrm ds_z.
\]

于是对应边界段上的电荷满足

\[
\boxed{
\sigma_z\,\mathrm ds_z
=
\sigma_w\,\mathrm ds_w
}.
\]

也就是说，单位轴向长度的总电荷在一一对应的导体边界之间保持不变。若两边的电势差也相同，则二维电容同样保持不变。

\textbf{例 M2-166：映射长度与局部场强反向缩放}

同一均匀电容率的二维静电域通过线性保角映射 \(w=az\) 对应，\(a\ne0\)
为复常数，电势按 \(\varphi_z(z)=\varphi_w(az)\)
拉回。已知像平面对应点场强大小为 \(E_w\)。求原平面场强与 \(E_w\)
的尺度关系，说明 \(a\)
的相位如何影响方向；再比较对应导体边界段的弧长、面电荷密度及单位轴向长度的电荷积分。

\par\noindent\begin{minipage}{\linewidth}
\textbf{解答}

\begin{center}
\includegraphics[width=\linewidth,height=.40\textheight,keepaspectratio]{\detokenize{assets/module-02/examples/M2-166.pdf}}
\par\smallskip{\small 图 M2-166：映射长度与局部场强反向缩放}
\end{center}
\end{minipage}\par

图 M2-166
比较同一映射两侧的对应边界段。复导数的模给长度放大率，电势差不变因而局部场强作相反的尺度换算；场强与弧长相乘保持电荷积分不变，复数系数的相位还会旋转方向。

若 \(w=az\)，像平面中电场大小为 \(E_w\)，则

\[
E_z=|a|E_w.
\]

同时原平面的边界长度缩短为像平面的 \(1/|a|\)，故电荷密度增大 \(|a|\)
倍，但总电荷不变。

\subsection{四、单叶性、分支与多连通区域}\label{m2:ux56dbux5355ux53f6ux6027ux5206ux652fux4e0eux591aux8fdeux901aux533aux57df}

保角变换最常见的错误不是微分算错，而是没有说清``在哪个区域、取哪个分支''。

\subsubsection{1.
局部条件与整体条件}\label{m2:ux5c40ux90e8ux6761ux4ef6ux4e0eux6574ux4f53ux6761ux4ef6}

\begin{itemize}
\tightlist
\item
  \(f\) 解析且 \(f'(z)\ne0\)：保证局部保角；
\item
  \(f\) 在所选区域内一一对应：保证整个物理区域不会折叠到一起；
\item
  边界按预定顺序对应：保证边界条件没有被错配；
\item
  反函数必须在像区域中选定同一分支。
\end{itemize}

例如

\[
w=z^2
\]

在整个去心平面上不是一一对应，因为 \(z\) 与 \(-z\)
有相同像；但在张角小于 \(\pi\) 的扇形内可以单叶。

\subsubsection{2.
对数与幂函数的分支}\label{m2:ux5bf9ux6570ux4e0eux5e42ux51fdux6570ux7684ux5206ux652f}

\[
\operatorname{Log}z=\ln r+\mathrm i\theta
\]

只有在选定

\[
\theta_0<\theta<\theta_0+2\pi
\]

后才是单值函数。割线不是额外的物理导体，而是为了定义单值分支引入的数学边界；若原问题跨越割线，必须检查割线两侧是否应重新粘合。

同理，

\[
z^{\beta}=\exp\!\bigl(\beta\operatorname{Log}z\bigr)
\]

也依赖对数分支。其在角点处可能满足 \(f'(0)=0\)
或发散，因此角点本身要单独用渐近分析处理。

\par\noindent\begin{minipage}{\linewidth}
\begin{center}
\includegraphics[width=\linewidth,height=.40\textheight,keepaspectratio]{\detokenize{assets/module-02/original/tikz-m2-topic20-v8-22.png}}
\par\smallskip{\small 图示：同一解析式在不同区域与分支上可能对应不同几何问题}
\end{center}
\end{minipage}\par

\subsubsection{3.
多连通区域中的调和共轭}\label{m2:ux591aux8fdeux901aux533aux57dfux4e2dux7684ux8c03ux548cux5171ux8f6d}

在多连通区域中，闭形式未必是全局恰当形式。即使 \(\varphi\)
是单值调和函数，\(\psi\) 也可能绕孔一周后改变常量。

例如

\[
W(z)=\operatorname{Log}z
=
\ln r+\mathrm i\theta.
\]

电势 \(\varphi=\ln r\) 单值，电场

\[
\mathcal E=-\frac1z
\]

也单值，但 \(\psi=\theta\) 绕原点一周增加
\(2\pi\)。这并不妨碍电场的物理意义，却说明复势本身需要分支或割线。

\subsection{五、楔形区域：对数映射把角度变成直线距离}\label{m2:ux4e94ux6954ux5f62ux533aux57dfux5bf9ux6570ux6620ux5c04ux628aux89d2ux5ea6ux53d8ux6210ux76f4ux7ebfux8dddux79bb}

考虑扇形区域

\[
0<\theta<\alpha.
\]

取

\[
w=\operatorname{Log}z=\ln r+\mathrm i\theta.
\]

则两条边界

\[
\theta=0,\qquad \theta=\alpha
\]

分别映成 \(w\) 平面中的两条平行直线

\[
\operatorname{Im}w=0,\qquad
\operatorname{Im}w=\alpha.
\]

\par\noindent\begin{minipage}{\linewidth}
\begin{center}
\includegraphics[width=\linewidth,height=.40\textheight,keepaspectratio]{\detokenize{assets/module-02/original/tikz-m2-topic20-v8-23.png}}
\par\smallskip{\small 图示：对数映射把楔形区域展开为平行直线之间的条带}
\end{center}
\end{minipage}\par

设下边界电势为 \(0\)，上边界电势为
\(V\)。条带中的电势沿虚部线性变化，拉回原平面得到

\[
\boxed{
\varphi(r,\theta)=\frac{V}{\alpha}\theta
}.
\]

直接用极坐标求梯度：

\[
\mathbf E
=
-\frac{\partial\varphi}{\partial r}\mathbf e_r
-\frac1r\frac{\partial\varphi}{\partial\theta}\mathbf e_{\theta}
=
\boxed{
-\frac{V}{\alpha r}\mathbf e_{\theta}
}.
\]

场线是圆弧 \(r=\text{常量}\)，等势线是射线 \(\theta=\text{常量}\)。

若场域为 \(0<\theta<\alpha\)，则两导体朝向场域的法线分别为
\(+\mathbf e_{\theta}\) 与 \(-\mathbf e_{\theta}\)。因此

\[
\boxed{
\sigma_{\theta=0}
=
-\frac{\varepsilon V}{\alpha r},
\qquad
\sigma_{\theta=\alpha}
=
+\frac{\varepsilon V}{\alpha r}
}.
\]

电荷密度按 \(1/r\)
增大，说明若两条不同电势的导体在顶点直接相接，顶点处存在边界电势不连续导致的强奇性。

若只考察

\[
a<r<b
\]

且两条圆弧边界满足零法向场，正极板单位轴向长度上的电荷为

\[
\lambda
=
\int_a^b
\frac{\varepsilon V}{\alpha r}\,\mathrm dr
=
\frac{\varepsilon V}{\alpha}\ln\frac ba.
\]

故单位长度电容

\[
\boxed{
C'
=
\frac{\lambda}{V}
=
\frac{\varepsilon}{\alpha}\ln\frac ba
}.
\]

\textbf{例 M2-167：直角扇形边界电势线性随角度改变}

真空二维无限直角劈形为 \(r>0,0<\theta<\pi/2\)，两条径向电极分别保持
\(0\) 与常电势
\(V\)。选取只依赖角度、无额外径向尺度的局部解。用平方映射或对数复势求电势和电场，核对两边界，并指出顶点附近的场强变化及其与不同电势电极相接的关系。

\par\noindent\begin{minipage}{\linewidth}
\textbf{解答}

\begin{center}
\includegraphics[width=\linewidth,height=.40\textheight,keepaspectratio]{\detokenize{assets/module-02/examples/M2-167.pdf}}
\par\smallskip{\small 图 M2-167：直角扇形边界电势线性随角度改变}
\end{center}
\end{minipage}\par

图 M2-167 的上下两条径向边界分别对应角度0和
\(\pi/2\)。线性角势先匹配两边电势，负的角导数除以半径给场；\(w=z^2\)
的指数与把直角放大成半平面张角的倍数相同。

当 \(\alpha=\pi/2\) 时，

\[
\varphi=\frac{2V}{\pi}\theta,
\qquad
|\mathbf E|=\frac{2V}{\pi r}.
\]

把映射取为 \(w=z^2\)，直角扇形又可映成上半平面；指数 \(2=\pi/\alpha\)
正是把角度 \(\alpha\) 放大为 \(\pi\) 所需的倍数。

\subsection{六、圆环区域：对数映射与二维同轴电容}\label{m2:ux516dux5706ux73afux533aux57dfux5bf9ux6570ux6620ux5c04ux4e0eux4e8cux7ef4ux540cux8f74ux7535ux5bb9}

考虑

\[
a<r<b,
\]

内圆柱电势为 \(V\)，外圆柱电势为 \(0\)。轴对称使电势只依赖
\(r\)，拉普拉斯方程化为

\[
\frac1r\frac{\mathrm d}{\mathrm dr}
\left(
r\frac{\mathrm d\varphi}{\mathrm dr}
\right)=0.
\]

两次积分得到

\[
\varphi=A\ln r+B.
\]

代入边界条件：

\[
\boxed{
\varphi(r)
=
V\frac{\ln(b/r)}{\ln(b/a)}
}.
\]

于是

\[
\boxed{
\mathbf E(r)
=
\frac{V}{r\ln(b/a)}\mathbf e_r
}.
\]

内导体单位长度的电荷为

\[
\lambda
=
2\pi a\varepsilon E_r(a)
=
\frac{2\pi\varepsilon V}{\ln(b/a)},
\]

故

\[
\boxed{
C'
=
\frac{2\pi\varepsilon}{\ln(b/a)}
}.
\]

\par\noindent\begin{minipage}{\linewidth}
\begin{center}
\includegraphics[width=\linewidth,height=.40\textheight,keepaspectratio]{\detokenize{assets/module-02/original/tikz-m2-topic20-v8-24.png}}
\par\smallskip{\small 图示：圆环在对数映射下成为具有周期识别的直条带}
\end{center}
\end{minipage}\par

这里

\[
w=\operatorname{Log}z=\ln r+\mathrm i\theta
\]

把 \(r=a,b\) 映成两条竖直线，而 \(\theta\) 与 \(\theta+2\pi\)
对应同一点。因此像区域不是普通开口矩形，而是上下边周期识别的条带。若人为画出割线，割线两侧的电势与法向导数必须匹配。

\textbf{例 M2-168：二维同轴电容只依赖半径比}

两根很长的同轴圆柱形成均匀电容率为 \(\varepsilon\) 的电容器，内、外半径
\(0<a<b\)，有效长度 \(L\)，忽略端部效应。将横截面两半径同时乘以
\(s>0\)，求单位长度电容的变化。分别在 \(L\) 保持不变、\(L\) 也乘以 \(s\)
两种情况下讨论总电容，说明二维尺度变换与整个三维装置同比缩放的区别。

\par\noindent\begin{minipage}{\linewidth}
\textbf{解答}

\begin{center}
\includegraphics[width=\linewidth,height=.40\textheight,keepaspectratio]{\detokenize{assets/module-02/examples/M2-168.pdf}}
\par\smallskip{\small 图 M2-168：二维同轴电容只依赖半径比}
\end{center}
\end{minipage}\par

图 M2-168
的两条同心圆一起改变尺度时，径向积分的上下限也同时缩放。对数只留下半径比，故单位长度电容不变；若轴向长度也改变，仍应另乘改变后的长度，不能误称总电容也必不变。

若 \(a,b\) 同时乘以 \(s>0\)，则 \(b/a\) 不变，所以 \(C'\)
不变。这是二维电容的尺度不变性；总电容仍等于 \(C'L\)，随轴向长度 \(L\)
变化。

\subsection{七、幂映射与角点场的渐近形式}\label{m2:ux4e03ux5e42ux6620ux5c04ux4e0eux89d2ux70b9ux573aux7684ux6e10ux8fd1ux5f62ux5f0f}

考虑场域在角点附近占据

\[
0<\theta<\alpha,
\]

两条边界在角点附近同为零电势。用分离变量

\[
\varphi(r,\theta)=R(r)\Theta(\theta)
\]

代入极坐标拉普拉斯方程

\[
\frac1r\frac{\partial}{\partial r}
\left(
r\frac{\partial\varphi}{\partial r}
\right)
+
\frac1{r^2}\frac{\partial^2\varphi}{\partial\theta^2}
=0,
\]

得到

\[
\frac{r(rR')'}{R}
=
-\frac{\Theta''}{\Theta}
=
\lambda^2.
\]

角向方程和边界条件为

\[
\Theta''+\lambda^2\Theta=0,
\qquad
\Theta(0)=\Theta(\alpha)=0.
\]

非零解要求

\[
\lambda_n=\frac{n\pi}{\alpha},
\qquad
\Theta_n=\sin\frac{n\pi\theta}{\alpha}.
\]

径向方程给出

\[
R_n=A_nr^{\lambda_n}+B_nr^{-\lambda_n}.
\]

若电势在角点保持有界，应舍去 \(r^{-\lambda_n}\)。一般边界激励下第一模
\(n=1\) 占主导：

\[
\boxed{
\varphi
\sim
A r^{\pi/\alpha}
\sin\frac{\pi\theta}{\alpha}
}.
\]

因此

\[
\boxed{
|\mathbf E|
\sim
\text{常量}\times r^{\pi/\alpha-1}
}.
\]

\par\noindent\begin{minipage}{\linewidth}
\begin{center}
\includegraphics[width=\linewidth,height=.40\textheight,keepaspectratio]{\detokenize{assets/module-02/original/tikz-m2-topic20-v8-25.png}}
\par\smallskip{\small 图示：场域张角决定角点附近电场趋零、有限或发散}
\end{center}
\end{minipage}\par

由指数可立即判断：

\[
\begin{array}{c|c|c}
\text{场域张角} & \pi/\alpha-1 & r\to0\text{ 时的电场}\\
\hline
\alpha<\pi & >0 & |\mathbf E|\to0\\
\alpha=\pi & =0 & |\mathbf E|\text{ 一般趋于有限值}\\
\alpha>\pi & <0 & |\mathbf E|\to\infty
\end{array}
\]

突出导体尖端外部的场域通常满足 \(\alpha>\pi\)，所以表面电荷密度

\[
\sigma=\varepsilon E_n
\]

也会发散。这是``尖端放电''中几何场增强的局部数学来源。

幂映射

\[
w=z^{\pi/\alpha}
\]

把张角 \(\alpha\) 的楔形映成上半平面。其导数

\[
f'(z)
=
\frac{\pi}{\alpha}z^{\pi/\alpha-1}
\]

恰好给出同一个场强指数。

\begin{quote}
两类奇性不要混淆
上述角点渐近假设两条边界在角点附近具有相同电势。若两条相接边界被强行保持为不同电势，则
\(\varphi=V\theta/\alpha\)，场强按 \(1/r\)
发散；这是边界电势不连续造成的奇性，不是同势导体尖角的本征模奇性。
\end{quote}

\subsection{八、分式线性变换：圆与直线可以互相转换}\label{m2:ux516bux5206ux5f0fux7ebfux6027ux53d8ux6362ux5706ux4e0eux76f4ux7ebfux53efux4ee5ux4e92ux76f8ux8f6cux6362}

一般分式线性变换为

\[
w=\frac{az+b}{cz+d},
\qquad
ad-bc\ne0.
\]

它由平移、旋转伸缩和反演组合而成，并具有以下性质：

\begin{enumerate}
\def\labelenumi{\arabic{enumi}.}
\tightlist
\item
  广义圆（圆或直线）仍映成广义圆；
\item
  角度保持；
\item
  极点 \(z=-d/c\) 映到无穷远；
\item
  只要极点不在物理域内，映射在适当区域中可保持单叶。
\end{enumerate}

特别地，

\[
w=\frac{z-z_1}{z-z_2}
\]

把 \(z_1\) 映到 \(0\)，把 \(z_2\) 映到无穷远。选择 \(z_1,z_2\)
为一族圆的两个极限点，可把偏心圆族映成同心圆族；专题18中偏心圆柱电容正可由这一思想处理。

\textbf{例 M2-169：实轴与单位圆的边界对应}

设 \(z=x+\mathrm iy\)，\(a>0\)，考察 Möbius 变换
\(\zeta=(z-\mathrm ia)/(z+\mathrm ia)\)。求实轴的像、点 \(z=\mathrm ia\)
的像和整个上半平面 \(y>0\)
的像域，说明如何由模长判断像点位于单位圆内或外，并检查变换的极点是否在所选开域中。

\par\noindent\begin{minipage}{\linewidth}
\textbf{解答}

\begin{center}
\includegraphics[width=\linewidth,height=.40\textheight,keepaspectratio]{\detokenize{assets/module-02/examples/M2-169.pdf}}
\par\smallskip{\small 图 M2-169：实轴与单位圆的边界对应}
\end{center}
\end{minipage}\par

图 M2-169 左侧的实轴点到 \(ia\) 与 \(-ia\)
等距，因此像点模长为1。上半平面内部的 \(ia\)
映到圆心，结合解析性和边界对应即可确定像域是单位圆内部，而不是外部。

取 \(a>0\)，

\[
\zeta=\frac{z-\mathrm ia}{z+\mathrm ia}.
\]

当 \(z=x\) 位于实轴时，

\[
|\zeta|
=
\frac{|x-\mathrm ia|}{|x+\mathrm ia|}
=1.
\]

而 \(z=\mathrm ia\) 映到圆心
\(\zeta=0\)。因此上半平面映成单位圆内部，实轴映成单位圆边界。

\subsection{九、保角变换的使用条件与检查}\label{m2:ux4e5dux4fddux89d2ux53d8ux6362ux7684ux4f7fux7528ux6761ux4ef6ux4e0eux68c0ux67e5}

\par\noindent\begin{minipage}{\linewidth}
\begin{center}
\includegraphics[width=\linewidth,height=.40\textheight,keepaspectratio]{\detokenize{assets/module-02/original/tikz-m2-topic20-v8-26.png}}
\par\smallskip{\small 图示：二维静电问题中的保角映射和边界条件}
\end{center}
\end{minipage}\par

面对二维边值问题，可按以下顺序处理：

\begin{enumerate}
\def\labelenumi{\arabic{enumi}.}
\tightlist
\item
  \textbf{确认二维性}：场量是否确实与第三坐标无关；
\item
  \textbf{标出物理域}：奇点、导体、介质界面、孔洞和无穷远分别在哪里；
\item
  \textbf{选映射与分支}：写出区域范围，验证 \(f'(z)\ne0\) 和单叶性；
\item
  \textbf{逐段对应边界}：不仅画形状，还要核对边界方向和电势值；
\item
  \textbf{在像平面求解}：优先利用平行线、同心圆、上半平面等标准模型；
\item
  \textbf{拉回电势和电场}：
\end{enumerate}

\[
\varphi_z=\varphi_w\circ f,
\qquad
\mathcal E_z=\mathcal E_w(f(z))f'(z);
\]

\begin{enumerate}
\def\labelenumi{\arabic{enumi}.}
\setcounter{enumi}{6}
\tightlist
\item
  \textbf{恢复电荷与电容}：使用 \(\sigma=\varepsilon E_n\) 和边界积分；
\item
  \textbf{做极限检验}：远场、对称极限、量纲、符号和边界值必须同时正确。
\end{enumerate}

常用映射可归纳为

\[
\begin{array}{c|c}
\text{映射} & \text{主要用途}\\
\hline
w=az+b & \text{平移、旋转、统一缩放}\\
w=1/z & \text{圆与直线的转换、近远区域互换}\\
w=\operatorname{Log}z & \text{楔形或圆环展开为条带}\\
w=z^{\beta} & \text{改变角度、分析角点指数}\\
w=(az+b)/(cz+d) & \text{广义圆之间的转换}
\end{array}
\]

\subsection{十、练习与解答}\label{m2:ux5341ux7ec3ux4e60ux4e0eux89e3ux7b54-1}

\subsubsection{作业一：由电势求调和共轭与复电场}\label{m2:ux4f5cux4e1aux4e00ux7531ux7535ux52bfux6c42ux8c03ux548cux5171ux8f6dux4e0eux590dux7535ux573a}

\textbf{题号：M2-170}

\textbf{层级：} 决赛　　\textbf{方法：} 边值问题

在整个二维平面给定真空调和势 \(\varphi(x,y)=x^2-y^2+2x\)。取
\(z=x+\mathrm iy\)，按 \(W=\varphi+\mathrm i\psi\) 的约定求调和共轭
\(\psi\)（允许加实常量），写出解析复势
\(W(z)\)，求电场分量与复电场，并确定电场零点。

\par\noindent\begin{minipage}{\linewidth}
\textbf{解答}

\begin{center}
\includegraphics[width=\linewidth,height=.40\textheight,keepaspectratio]{\detokenize{assets/module-02/examples/M2-170.pdf}}
\par\smallskip{\small 图 M2-170：由电势求调和共轭与复电场}
\end{center}
\end{minipage}\par

图 M2-170
的势可整理为平移后的二次鞍形。由柯西---黎曼条件积分得到调和共轭，再将其合成解析复势；电场零点由复导数为零确定，它随线性项移到新的横坐标。

先检验

\[
\varphi_{xx}+\varphi_{yy}=2-2=0.
\]

由柯西---黎曼条件，

\[
\psi_y=\varphi_x=2x+2,
\qquad
\psi_x=-\varphi_y=2y.
\]

对第二式关于 \(x\) 积分：

\[
\psi=2xy+g(y).
\]

再由

\[
\psi_y=2x+g'(y)=2x+2
\]

得到 \(g'(y)=2\)，故可取

\[
\boxed{\psi=2xy+2y}.
\]

于是

\[
W(z)=z^2+2z+\mathrm iC,
\]

舍去无关的虚常量后，

\[
\boxed{W(z)=z^2+2z}.
\]

复电场为

\[
\mathcal E=-W'(z)=-2z-2.
\]

所以

\[
\boxed{
\mathbf E=(-2x-2,\,2y)
}.
\]

电场零点满足 \(z=-1\)，即

\[
\boxed{(x,y)=(-1,0)}.
\]

\subsubsection{作业二：有限直角扇形电容}\label{m2:ux4f5cux4e1aux4e8cux6709ux9650ux76f4ux89d2ux6247ux5f62ux7535ux5bb9}

\textbf{题号：M2-171}

\textbf{层级：} 决赛　　\textbf{方法：} 边值问题

均匀电容率为 \(\varepsilon>0\) 的二维无源域为有限四分之一圆环
\(a<r<b,0<\theta<\pi/2\)，\(0<a<b\)，沿轴向无限延伸。径向导体边
\(\theta=0,\pi/2\) 分别保持 \(0,V\)，两条圆弧 \(r=a,b\)
规定零法向电场，不把圆弧当作等势导体。求电势、电场、两条直导体面朝向场区一侧的自由面电荷密度，以及单位轴向长度电容。

\par\noindent\begin{minipage}{\linewidth}
\textbf{解答}

\begin{center}
\includegraphics[width=\linewidth,height=.40\textheight,keepaspectratio]{\detokenize{assets/module-02/examples/M2-171.pdf}}
\par\smallskip{\small 图 M2-171：有限直角扇形电容}
\end{center}
\end{minipage}\par

图 M2-171
的两段圆弧是零法向场边界，电势只依赖极角即可自动满足。两条径向导体则固定角势斜率；在直边上沿各自区域外法向取导数，面电荷符号相反，其径向积分产生对数电容。

边界条件只依赖 \(\theta\)，可设 \(\varphi=A\theta+B\)。由

\[
\varphi(r,0)=0,\qquad
\varphi\left(r,\frac{\pi}{2}\right)=V
\]

得

\[
\boxed{
\varphi(r,\theta)=\frac{2V}{\pi}\theta
}.
\]

该解与 \(r\) 无关，所以在 \(r=a,b\) 上

\[
\frac{\partial\varphi}{\partial r}=0,
\]

自动满足给定的零法向场条件。

电场为

\[
\boxed{
\mathbf E
=
-\frac{2V}{\pi r}\mathbf e_{\theta}
}.
\]

在 \(\theta=0\) 的导体上，朝向场域的法线是 \(+\mathbf e_{\theta}\)，故

\[
\boxed{
\sigma_0
=
-\frac{2\varepsilon V}{\pi r}
}.
\]

在 \(\theta=\pi/2\) 的导体上，朝向场域的法线是
\(-\mathbf e_{\theta}\)，故

\[
\boxed{
\sigma_V
=
+\frac{2\varepsilon V}{\pi r}
}.
\]

正极单位长度电荷为

\[
\lambda
=
\int_a^b\sigma_V\,\mathrm dr
=
\frac{2\varepsilon V}{\pi}\ln\frac ba.
\]

所以

\[
\boxed{
C'
=
\frac{2\varepsilon}{\pi}\ln\frac ba
}.
\]

\subsubsection{作业三：二维同轴电容的能量检验}\label{m2:ux4f5cux4e1aux4e09ux4e8cux7ef4ux540cux8f74ux7535ux5bb9ux7684ux80fdux91cfux68c0ux9a8c}

\textbf{题号：M2-172}

\textbf{层级：} 决赛　　\textbf{方法：} 边值问题、能量法

两根很长的同轴导体圆柱半径为 \(0<a<b\)，有效长度为
\(L\)，间隙充满均匀线性介质，绝对电容率为
\(\varepsilon>0\)。内柱保持电势
\(V\)、外柱接地，忽略端部效应。求场与电荷---电压关系，分别通过场能积分和单位长度电荷计算推导电容，并检查两种路径的一致性。

\par\noindent\begin{minipage}{\linewidth}
\textbf{解答}

\begin{center}
\includegraphics[width=\linewidth,height=.40\textheight,keepaspectratio]{\detokenize{assets/module-02/examples/M2-172.pdf}}
\par\smallskip{\small 图 M2-172：二维同轴电容的能量检验}
\end{center}
\end{minipage}\par

图 M2-172 的场能位于两圆柱之间，体积元每单位长度为
\(2\pi\rho d\rho\)。把平方场乘这一微元积分后，恰好留下一个对数半径比；与
\(C^{\prime}V^2/2\) 比较，必须得到和电荷法相同的结果。

电场为

\[
\mathbf E(r)
=
\frac{V}{r\ln(b/a)}\mathbf e_r.
\]

场能

\[
U
=
\frac{\varepsilon}{2}
\int_0^L\!\mathrm dz
\int_0^{2\pi}\!\mathrm d\theta
\int_a^b E^2 r\,\mathrm dr.
\]

代入：

\[
\begin{aligned}
U
&=
\frac{\varepsilon}{2}
L(2\pi)
\frac{V^2}{\ln^2(b/a)}
\int_a^b\frac{\mathrm dr}{r}\\
&=
\frac{\pi\varepsilon L V^2}{\ln(b/a)}.
\end{aligned}
\]

由

\[
U=\frac12CV^2
\]

得到

\[
\boxed{
C
=
\frac{2\pi\varepsilon L}{\ln(b/a)}
},
\qquad
\boxed{
C'=\frac CL
=
\frac{2\pi\varepsilon}{\ln(b/a)}
}.
\]

这与内导体电荷

\[
Q
=
\varepsilon E(a)(2\pi aL)
=
\frac{2\pi\varepsilon LV}{\ln(b/a)}
\]

给出的 \(C=Q/V\) 完全一致。

\subsubsection{作业四：导体尖角的场强指数}\label{m2:ux4f5cux4e1aux56dbux5bfcux4f53ux5c16ux89d2ux7684ux573aux5f3aux6307ux6570}

\textbf{题号：M2-173}

\textbf{层级：} 决赛　　\textbf{方法：} 边值问题

二维导体的一个角具有内部角 \(\beta\)，外部无源均匀介质的局部场域角为
\(\alpha=2\pi-\beta\)。相邻两导体面属于同一等势导体，取它们为零势；要求角点电势有限，并假设远处激励对最低允许角模的系数非零。推导角点附近电场大小随距离
\(r\) 的主阶幂律，分别讨论
\(\beta=\pi/2,\pi,3\pi/2\)，说明何种角形产生场增强。

\par\noindent\begin{minipage}{\linewidth}
\textbf{解答}

\begin{center}
\includegraphics[width=\linewidth,height=.40\textheight,keepaspectratio]{\detokenize{assets/module-02/examples/M2-173.pdf}}
\par\smallskip{\small 图 M2-173：导体尖角的场强指数}
\end{center}
\end{minipage}\par

图 M2-173 区分导体内角与外部场域张角，分离变量应使用后者
\(\alpha\)。两侧同为零势时，最低角模为正弦而径向幂为
\(r^{\pi/\alpha}\)；求梯度后指数减1，据此分别判断凸角、平面与凹角附近的场。

外部场域中的第一角向本征值为

\[
\lambda_1=\frac{\pi}{\alpha}
=
\frac{\pi}{2\pi-\beta}.
\]

因此

\[
\varphi
\sim
A r^{\lambda_1}
\sin(\lambda_1\theta),
\]

\[
\boxed{
|\mathbf E|
\sim
r^{\lambda_1-1}
=
r^{\frac{\pi}{2\pi-\beta}-1}
}.
\]

\begin{enumerate}
\def\labelenumi{\arabic{enumi}.}
\tightlist
\item
  当 \(\beta=\pi/2\) 时，
\end{enumerate}

\[
\alpha=\frac{3\pi}{2},
\qquad
\lambda_1=\frac23,
\]

故

\[
\boxed{|\mathbf E|\sim r^{-1/3}}.
\]

这是向场域突出的导体直角尖角，电场发散。

\begin{enumerate}
\def\labelenumi{\arabic{enumi}.}
\setcounter{enumi}{1}
\tightlist
\item
  当 \(\beta=\pi\) 时，
\end{enumerate}

\[
\alpha=\pi,\qquad \lambda_1=1,
\]

故

\[
\boxed{|\mathbf E|\sim r^0},
\]

一般趋于有限常量。

\begin{enumerate}
\def\labelenumi{\arabic{enumi}.}
\setcounter{enumi}{2}
\tightlist
\item
  当 \(\beta=3\pi/2\) 时，
\end{enumerate}

\[
\alpha=\frac{\pi}{2},
\qquad
\lambda_1=2,
\]

故

\[
\boxed{|\mathbf E|\sim r}.
\]

这是导体的凹角，场域张角小于 \(\pi\)，电场在角点趋于零。

\subsubsection{\texorpdfstring{作业五：边界电势跳变产生的 \(1/r\)
奇性}{作业五：边界电势跳变产生的 1/r 奇性}}\label{m2:ux4f5cux4e1aux4e94ux8fb9ux754cux7535ux52bfux8df3ux53d8ux4ea7ux751fux7684-1r-ux5947ux6027}

\textbf{题号：M2-174}

\textbf{层级：} 决赛　　\textbf{方法：} 边值问题

二维真空上半平面 \(y>0\) 无电荷，实轴正半轴 \(x>0\) 保持电势
\(V\)，负半轴 \(x<0\)
保持零势，原点处允许边界电势跳变。取不引入额外径向尺度的有界角向解，求电势与电场，分析
\(r\to0\)
的场强。再与``同一等势导体两面形成角点、势连续且最低角模系数非零''的情形比较，解释两种奇性的边界条件差别。

\par\noindent\begin{minipage}{\linewidth}
\textbf{解答}

\begin{center}
\includegraphics[width=\linewidth,height=.40\textheight,keepaspectratio]{\detokenize{assets/module-02/examples/M2-174.pdf}}
\par\smallskip{\small 图 M2-174：边界电势跳变产生的 \(1/r\) 奇性}
\end{center}
\end{minipage}\par

图 M2-174
的实轴两侧给不同常量边值，原点处存在跳变。角势无需径向幂即可匹配这一数据，因此梯度直接产生
\(1/r\)；它与前题两侧同势时的最低齐次角模属于不同边界条件。

在上半平面取极坐标

\[
0<\theta<\pi.
\]

正半轴对应 \(\theta=0\)，负半轴对应
\(\theta=\pi\)。满足边界条件的调和函数为

\[
\boxed{
\varphi(r,\theta)
=
V\left(1-\frac{\theta}{\pi}\right)
}.
\]

电势与 \(r\) 无关，因此

\[
\mathbf E
=
-\frac1r\frac{\partial\varphi}{\partial\theta}\mathbf e_{\theta}
=
\boxed{
\frac{V}{\pi r}\mathbf e_{\theta}
}.
\]

场强为

\[
\boxed{|\mathbf E|=\frac{|V|}{\pi r}}.
\]

当 \(V\ne0\) 时在原点按 \(1/r\) 发散；\(V=0\)
时该有界角向解恒为零，没有此奇性。

作业4中，角点两侧属于同一等势导体，允许的局部解由齐次边界条件下的第一本征模决定，指数为
\(\pi/\alpha-1\)。本题中边界电势在原点从 \(V\) 跳到
\(0\)，电势主项本身就是 \(\theta\) 的线性函数，故产生更强的 \(1/r\)
奇性。两者的边界条件不同，不能套用同一个尖角指数。

\section{模块二附录三：导体边值问题与电像法综合题（拔高选讲）}\label{m2:ux6a21ux5757ux4e8cux9644ux5f55ux4e09ux5bfcux4f53ux8fb9ux503cux95eeux9898ux4e0eux7535ux50cfux6cd5ux7efcux5408ux9898ux62d4ux9ad8ux9009ux8bb2}

本附录把唯一性定理、电像法、特殊坐标、级数渐近和能量法放在同一组导体边值问题中。所有电容均先说明电势约束：

\begin{itemize}
\tightlist
\item
  \textbf{孤立体系的自电容}：所有连通导体同为电势 \(V\)，以
  \(C=Q_{\text{总}}/V\) 定义；
\item
  \textbf{两导体间电容}：两导体电势为 \(+U,-U\)，电势差为 \(2U\)，以
  \(C=Q/(2U)\) 定义；
\item
  \textbf{柱形体系}：忽略端部效应，\(C'\) 与 \(F'\)
  分别表示单位轴向长度的电容和力。
\end{itemize}

这种区分不是符号细节。同一对金属球在``等电势''与``反相电势''下激发的是不同对称模，极限行为也不同。

\subsection{综合题1：相切等大金属球的自电容}\label{m2:ux7efcux5408ux98981ux76f8ux5207ux7b49ux5927ux91d1ux5c5eux7403ux7684ux81eaux7535ux5bb9}

\textbf{题号：M2-175}

\textbf{层级：} 复赛　　\textbf{方法：} 电像法

真空中两个半径均为 \(R>0\) 的理想金属球外切，球心距为
\(2R\)，通过接触点保持为同一导体。除两球外没有其他源，取无穷远零势，整个导体保持常电势
\(V\ne0\)。求两球总自由电荷 \(Q_{\rm tot}\)
及这一连通导体相对于无穷远的自电容；应在接触极限中处理像源或等势面，而不是把两个孤立球电容直接相加。

\par\noindent\begin{minipage}{\linewidth}
\subsubsection{解答}\label{m2:ux89e3ux7b54}

\begin{center}
\includegraphics[width=\linewidth,height=.40\textheight,keepaspectratio]{\detokenize{assets/module-02/examples/M2-175.pdf}}
\par\smallskip{\small 图 M2-175：相切等大金属球的自电容}
\end{center}
\end{minipage}\par

图 M2-175
的两球同电势且连通，接触点是两条反演链的积聚点。沿每次反演的几何距离递推位置与电量，再把单球内部的有符号电荷相加；两球同相时得到交错级数，不能使用反相双球的正项级数。

\subsubsection{1. 像电荷递推}\label{m2:ux50cfux7535ux8377ux9012ux63a8}

以相切点为原点，球心位于轴线上 \(x=\pm R\)。先在每个球心放置电荷

\[
q_0=4\pi\varepsilon_0RV.
\]

球心电荷在本球表面产生常量电势
\(V\)。另一球中的电荷对本球表面造成的非均匀电势，可用球面反演不断消去。

设某次待反演电荷量的绝对值为 \(q_n\)，它到相切点的距离为
\(r_n\)。它到另一球球心的距离为 \(R+r_n\)。半径为 \(R\)
的接地球对外部点电荷的像点距球心

\[
d'=\frac{R^2}{R+r_n},
\]

故新像点距相切点

\[
\begin{aligned}
r_{n+1}
&=R-d'\\
&=R-\frac{R^2}{R+r_n}\\
&=\frac{Rr_n}{R+r_n}.
\end{aligned}
\]

取倒数：

\[
\frac1{r_{n+1}}
=
\frac1{r_n}+\frac1R.
\]

由 \(r_0=R\) 得

\[
\boxed{
r_n=\frac{R}{n+1}
}.
\]

球面反演还给出像电荷量

\[
q_{n+1}
=
q_n\frac{R}{R+r_n}.
\]

代入 \(r_n=R/(n+1)\)：

\[
q_{n+1}
=
q_n\frac{n+1}{n+2},
\]

所以

\[
\boxed{
q_n=\frac{q_0}{n+1}
}.
\]

每次关于接地球面的反演改变电荷符号。叠加从左右两个球心出发的两条反射链后，每个球内部的有符号电荷依次为

\[
+q_0,\ -q_1,\ +q_2,\ -q_3,\ldots
\]

并逐渐聚集到相切点。

\par\noindent\begin{minipage}{\linewidth}
\begin{center}
\includegraphics[width=\linewidth,height=.40\textheight,keepaspectratio]{\detokenize{assets/module-02/process/M2-175a.pdf}}
\par\smallskip{\small 图 M2-175a：反演链在相切点两侧聚集}
\end{center}
\end{minipage}\par

图 M2-175a 的像点向相切点 \(O\)
聚集。每次球面反演同时改变位置、大小和符号；沿每个球内部读出交错级数，才可对电荷作代数求和。

\subsubsection{2. 求总电荷}\label{m2:ux6c42ux603bux7535ux8377}

单个球所含像电荷的代数和为

\[
\begin{aligned}
Q_{\text{单球}}
&=
q_0\sum_{n=0}^{\infty}
\frac{(-1)^n}{n+1}\\
&=
q_0\left(
1-\frac12+\frac13-\frac14+\cdots
\right)\\
&=
q_0\ln2.
\end{aligned}
\]

因此

\[
Q_{\text{总}}
=
2q_0\ln2
=
8\pi\varepsilon_0RV\ln2.
\]

自电容为

\[
\boxed{
C_{\text{touch}}
=
\frac{Q_{\text{总}}}{V}
=
8\pi\varepsilon_0R\ln2
}.
\]

\subsubsection{3.
用相切球坐标独立验证}\label{m2:ux7528ux76f8ux5207ux7403ux5750ux6807ux72ecux7acbux9a8cux8bc1}

在圆柱坐标 \((\rho,z,\phi)\) 中定义

\[
u=\frac{\rho}{\rho^2+z^2},
\qquad
v=\frac{z}{\rho^2+z^2}.
\]

逆变换为

\[
\rho=\frac{u}{u^2+v^2},
\qquad
z=\frac{v}{u^2+v^2}.
\]

把 \(v=c\) 代入逆变换，可得

\[
\rho^2+
\left(z-\frac1{2c}\right)^2
=
\frac1{4c^2}.
\]

所以 \(v=c\) 是一个与原点相切的球面，球心在 \(z=1/(2c)\)，半径为
\(1/(2|c|)\)。两个半径为 \(R\) 的相切球对应

\[
v=\pm v_0,
\qquad
v_0=\frac1{2R}.
\]

\par\noindent\begin{minipage}{\linewidth}
\begin{center}
\includegraphics[width=\linewidth,height=.40\textheight,keepaspectratio]{\detokenize{assets/module-02/process/M2-175b.pdf}}
\par\smallskip{\small 图 M2-175b：反演坐标把相切球面化成两条直线}
\end{center}
\end{minipage}\par

图 M2-175b 画的是变换后的 \(u,v\)
坐标域，两条竖线对应原来的两球面，真实域是两线之间的 \(u\ge0\)
半条带。物理远区应在坐标原点附近读取 \(1/r\)
系数，不能把坐标无穷远误当成物理远区。

对轴对称电势作代换

\[
\Phi(u,v)=\sqrt{u^2+v^2}\,\Psi(u,v),
\]

拉普拉斯方程化为

\[
\Psi_{uu}+\frac1u\Psi_u+\Psi_{vv}=0.
\]

由于两球同电势，解关于 \(v\) 为偶函数。满足

\[
\Phi(u,\pm v_0)=V
\]

的解可写成

\[
\Phi(u,v)
=
V\sqrt{u^2+v^2}
\int_0^{\infty}
\frac{e^{-kv_0}}{\cosh(kv_0)}
\cosh(kv)J_0(ku)\,\mathrm dk.
\]

这里使用了积分恒等式

\[
\int_0^{\infty} e^{-kv_0}J_0(ku)\,\mathrm dk
=
\frac1{\sqrt{u^2+v_0^2}}.
\]

物理无穷远 \(r\to\infty\) 对应 \((u,v)\to(0,0)\)，并且

\[
\sqrt{u^2+v^2}=\frac1r.
\]

因此远场系数为

\[
\frac{Q_{\text{总}}}{4\pi\varepsilon_0}
=
V\int_0^{\infty}
\frac{e^{-kv_0}}{\cosh(kv_0)}\,\mathrm dk.
\]

令 \(x=2kv_0\)，得到

\[
\begin{aligned}
\int_0^{\infty}
\frac{e^{-kv_0}}{\cosh(kv_0)}\,\mathrm dk
&=
\frac1{v_0}
\int_0^{\infty}
\frac{e^{-x}}{1+e^{-x}}\,\mathrm dx\\
&=
\frac{\ln2}{v_0}\\
&=
2R\ln2.
\end{aligned}
\]

故再次得到

\[
Q_{\text{总}}
=
8\pi\varepsilon_0RV\ln2.
\]

两种方法中的 \(\ln2\)
是同一个结构：像电荷法中它来自交错调和级数，相切球坐标中它来自核函数的几何级数展开。

\subsubsection{4.
极限与量纲检查}\label{m2:ux6781ux9650ux4e0eux91cfux7eb2ux68c0ux67e5}

单个半径为 \(R\) 的孤立球电容为
\(4\pi\varepsilon_0R\)。两个相距很远且同电势的球，总自电容趋于

\[
8\pi\varepsilon_0R.
\]

相切后

\[
8\pi\varepsilon_0R\ln2
<
8\pi\varepsilon_0R,
\]

因为两个球靠近时，同号电荷相互排斥，使给定电势下可容纳的总电荷减少。结果与长度
\(R\) 成正比，量纲正确。

\subsection{综合题2：两只分离等大金属球的电容}\label{m2:ux7efcux5408ux98982ux4e24ux53eaux5206ux79bbux7b49ux5927ux91d1ux5c5eux7403ux7684ux7535ux5bb9}

\textbf{题号：M2-176}

\textbf{层级：} 复赛　　\textbf{方法：} 电像法

真空中两个互不接触的导体球半径均为 \(R>0\)，球心距为
\(D=2R+\Delta\)，\(\Delta>0\)。两球分别由外电源保持电势
\(+U,-U\)，\(U>0\)，无穷远零势。由对称性两球电荷为
\(+Q,-Q\)。求二导体电容 \(C=Q/(2U)\)，给出可计算的精确表示，并分析窄隙
\(\Delta/R\ll1\) 时的主阶渐近行为；注意电压应取两球电势差。

\par\noindent\begin{minipage}{\linewidth}
\subsubsection{解答}\label{m2:ux89e3ux7b54-1}

\begin{center}
\includegraphics[width=\linewidth,height=.40\textheight,keepaspectratio]{\detokenize{assets/module-02/examples/M2-176.pdf}}
\par\smallskip{\small 图 M2-176：两只分离等大金属球的电容}
\end{center}
\end{minipage}\par

图 M2-176 的两球电势为 \(+U,-U\)，所以电压差是
\(2U\)。几何参数由球心距和半径共同固定，反相模给单球电荷
\(Q\)，最后要除以
\(2U\)；窄隙近似只能确定对数主项，常数项仍须由完整级数给出。

\subsubsection{1.
双球坐标与几何参数}\label{m2:ux53ccux7403ux5750ux6807ux4e0eux51e0ux4f55ux53c2ux6570}

定义双球坐标 \((\eta,\mu,\phi)\)：

\[
\rho
=
\frac{a\sin\mu}{\cosh\eta-\cos\mu},
\qquad
z
=
\frac{a\sinh\eta}{\cosh\eta-\cos\mu}.
\]

\(\eta=\eta_0\) 是球面。把坐标式消去 \(\mu\)，可得

\[
\rho^2+
\left(z-a\coth\eta_0\right)^2
=
a^2\operatorname{csch}^2\eta_0.
\]

所以球心与半径分别为

\[
z_c=a\coth\eta_0,
\qquad
R_{\eta}=a|\operatorname{csch}\eta_0|.
\]

两只等大球取为

\[
\eta=\pm\alpha.
\]

于是

\[
R=a\operatorname{csch}\alpha,
\qquad
D=2a\coth\alpha,
\]

从而

\[
\boxed{
\cosh\alpha=\frac{D}{2R}
=
1+\frac{\Delta}{2R}
},
\qquad
\boxed{
a=R\sinh\alpha
}.
\]

\par\noindent\begin{minipage}{\linewidth}
\begin{center}
\includegraphics[width=\linewidth,height=.40\textheight,keepaspectratio]{\detokenize{assets/module-02/process/M2-176a.pdf}}
\par\smallskip{\small 图 M2-176a：球心距与双球坐标参数相互确定}
\end{center}
\end{minipage}\par

图 M2-176a 把中心距 \(D\) 与表面窄隙 \(\Delta\) 分开标注。先由
\(D/(2R)\) 确定 \(\alpha\)，再定坐标尺度
\(a\)，后续级数中的两个参数便不再独立。

\subsubsection{2.
轴对称拉普拉斯方程的分离解}\label{m2:ux8f74ux5bf9ux79f0ux62c9ux666eux62c9ux65afux65b9ux7a0bux7684ux5206ux79bbux89e3}

在双球坐标中，轴对称调和函数可展开为

\[
\Phi(\eta,\mu)
=
\sqrt{\cosh\eta-\cos\mu}
\sum_{n=0}^{\infty}
\left[
A_ne^{(n+\frac12)\eta}
+
B_ne^{-(n+\frac12)\eta}
\right]
P_n(\cos\mu).
\]

边界上常电势的展开来自勒让德生成函数：

\[
\boxed{
\frac1{\sqrt{\cosh\alpha-\cos\mu}}
=
\sqrt2
\sum_{n=0}^{\infty}
e^{-(n+\frac12)\alpha}P_n(\cos\mu)
}.
\]

在 \(\eta=\pm\alpha\) 分别匹配 \(V_1,V_2\)，再用勒让德正交性确定
\(A_n,B_n\)。对球面法向电场积分后得到电容系数

\[
Q_1=C_{11}V_1+C_{12}V_2,
\qquad
Q_2=C_{12}V_1+C_{11}V_2,
\]

其中

\[
\boxed{
C_{11}
=
4\pi\varepsilon_0a
\sum_{n=1}^{\infty}
\operatorname{csch}\!\bigl((2n-1)\alpha\bigr)
},
\]

\[
\boxed{
C_{12}
=
-4\pi\varepsilon_0a
\sum_{n=1}^{\infty}
\operatorname{csch}(2n\alpha)
}.
\]

\(C_{12}<0\) 表示另一导体电势升高会减少本导体在给定电势下的电荷。

\subsubsection{3.
同相模与反相模}\label{m2:ux540cux76f8ux6a21ux4e0eux53cdux76f8ux6a21}

由于电容矩阵具有形式

\[
\begin{pmatrix}
C_{11}&C_{12}\\
C_{12}&C_{11}
\end{pmatrix},
\]

它的两个本征向量分别为

\[
\begin{pmatrix}1\\1\end{pmatrix},
\qquad
\begin{pmatrix}1\\-1\end{pmatrix}.
\]

\par\noindent\begin{minipage}{\linewidth}
\begin{center}
\includegraphics[width=\linewidth,height=.40\textheight,keepaspectratio]{\detokenize{assets/module-02/process/M2-176b.pdf}}
\par\smallskip{\small 图 M2-176b：电容矩阵的两个电势模态}
\end{center}
\end{minipage}\par

图 M2-176b 上下两组是不同的电势控制方式。本题采用下方反相模，单球电荷为
\(\lambda_-U\)，而两球电势差为 \(2U\)，所以球间电容还要除以2。

对本题的反相电势 \(V_1=+U,V_2=-U\)，有

\[
Q
=
(C_{11}-C_{12})U.
\]

电势差为 \(2U\)，故

\[
C_{\text{pair}}
=
\frac{C_{11}-C_{12}}2.
\]

把奇数项与偶数项合并：

\[
\boxed{
C_{\text{pair}}
=
2\pi\varepsilon_0a
\sum_{m=1}^{\infty}
\operatorname{csch}(m\alpha)
}.
\]

作为比较，若两球连通并同为电势 \(V\)，总自电容为

\[
\boxed{
C_{\text{self}}
=
8\pi\varepsilon_0a
\sum_{m=1}^{\infty}
(-1)^{m+1}\operatorname{csch}(m\alpha)
}.
\]

当 \(\alpha\to0\) 时，后一式趋于

\[
8\pi\varepsilon_0R\ln2,
\]

与综合题1连续衔接。

\subsubsection{4. 窄隙渐近}\label{m2:ux7a84ux9699ux6e10ux8fd1}

当 \(\Delta\ll R\) 时，

\[
\cosh\alpha
=
1+\frac{\Delta}{2R}.
\]

利用

\[
\cosh\alpha
=
1+\frac{\alpha^2}{2}+O(\alpha^4),
\]

得

\[
\alpha
=
\sqrt{\frac{\Delta}{R}}
\left[
1+O\!\left(\frac{\Delta}{R}\right)
\right],
\]

并且

\[
a=R\sinh\alpha
=
R\alpha\left[1+O(\alpha^2)\right].
\]

双曲余割和的欧拉---麦克劳林渐近为

\[
\sum_{m=1}^{\infty}\operatorname{csch}(m\alpha)
=
\frac1\alpha
\left[
\ln\frac2\alpha+\gamma
\right]
+
O\!\left(\alpha\ln\frac1\alpha\right),
\]

其中 \(\gamma\) 是欧拉常数。代回精确级数：

\[
\begin{aligned}
C_{\text{pair}}
&=
2\pi\varepsilon_0R
\left[
\ln\frac2\alpha+\gamma
\right]
+
O\!\left(
\varepsilon_0R\alpha^2\ln\frac1\alpha
\right)\\
&=
\boxed{
\pi\varepsilon_0R
\left[
\ln\frac{4R}{\Delta}
+2\gamma
+O\!\left(
\frac{\Delta}{R}\ln\frac{R}{\Delta}
\right)
\right]
}.
\end{aligned}
\]

\par\noindent\begin{minipage}{\linewidth}
\begin{center}
\includegraphics[width=\linewidth,height=.40\textheight,keepaspectratio]{\detokenize{assets/module-02/process/M2-176c.pdf}}
\par\smallskip{\small 图 M2-176c：窄隙中的局部厚度决定对数主项}
\end{center}
\end{minipage}\par

图 M2-176c
的两条抛物线是球面在最近点附近的局部展开。把薄环面积除以局部间隙并积分可解释对数发散；式中的常数项仍须由完整球面解确定。

局部间隙近似为

\[
h(\rho)
\simeq
\Delta+\frac{\rho^2}{R}.
\]

把每个环带近似成平行板可得

\[
\mathrm dC
\sim
\frac{\varepsilon_0\,2\pi\rho\,\mathrm d\rho}
{\Delta+\rho^2/R},
\]

其积分产生 \(\ln(R/\Delta)\)。但局部近似无法确定常数项；\(2\gamma\)
必须由完整双球级数得到。

\subsubsection{5. 极限检查}\label{m2:ux6781ux9650ux68c0ux67e5}

当 \(D\gg R\) 时，\(\alpha\gg1\)，并有

\[
a\operatorname{csch}\alpha=R.
\]

级数首项给出

\[
C_{\text{pair}}\to2\pi\varepsilon_0R.
\]

这正是两个相距很远、各自自电容为 \(4\pi\varepsilon_0R\)
的导体串联后的电容。

\subsection{综合题3：两接地平板之间带电单摆的稳定性}\label{m2:ux7efcux5408ux98983ux4e24ux63a5ux5730ux5e73ux677fux4e4bux95f4ux5e26ux7535ux5355ux6446ux7684ux7a33ux5b9aux6027}

\textbf{题号：M2-177}

\textbf{层级：} 复赛　　\textbf{方法：} 电像法、能量法

真空中两块无限大接地导体平板位于
\(x=\pm l\)，\(l>0\)，重力方向竖直且平行于平板。质量为 \(m>0\)、电荷为
\(q\) 的小球用长 \(h>0\) 的轻绝缘细线悬挂，悬点在中面
\(x=0\)。把小球近似为点电荷，忽略细线对电场的影响、空气阻力及辐射；小球在垂直于平板的竖直面内摆动，重力加速度为
\(g\)。求小球偏离中面一个微小横向位移时的线性电像力，并与重力恢复作用比较，给出竖直平衡的线性稳定判据及临界电荷量，说明临界点在线性近似中的限度。

\par\noindent\begin{minipage}{\linewidth}
\subsubsection{解答}\label{m2:ux89e3ux7b54-2}

\begin{center}
\includegraphics[width=\linewidth,height=.40\textheight,keepaspectratio]{\detokenize{assets/module-02/examples/M2-177.pdf}}
\par\smallskip{\small 图 M2-177：两接地平板之间带电单摆的稳定性}
\end{center}
\end{minipage}\par

图 M2-177
的小球偏向右侧后，右板的感应吸引增强，电像合力继续放大位移。先由两列周期像源求这一反恢复刚度，再与摆长给出的重力恢复刚度比较；临界时二次项消失，不能仅凭线性项决定更高阶运动。

\subsubsection{1. 两列像电荷}\label{m2:ux4e24ux5217ux50cfux7535ux8377}

实电荷位于 \(x=a\)，其中 \(|a|<l\)。满足两板零电势的像电荷可写成两列：

\[
\boxed{
+q:\quad x=a+4nl,
\qquad n\in\mathbb Z
},
\]

其中 \(n=0\) 是实电荷，其余为同号像电荷；

\[
\boxed{
-q:\quad x=-a+(4n+2)l,
\qquad n\in\mathbb Z
}.
\]

\par\noindent\begin{minipage}{\linewidth}
\begin{center}
\includegraphics[width=\linewidth,height=.40\textheight,keepaspectratio]{\detokenize{assets/module-02/process/M2-177a.pdf}}
\par\smallskip{\small 图 M2-177a：两列像源逐项满足两块接地板}
\end{center}
\end{minipage}\par

图 M2-177a
中左右两块板分别是两组距离配对的对称面；蓝色负像源均在板间真实区域之外。逐项配对使两板零势，同号像源的对称力则互相抵消。

在 \(x=l\) 上，任取同号像电荷 \(a+4nl\)，都有一个负列索引 \(m=-n\)
的异号像电荷

\[
-a+(2-4n)l
\]

与它关于该平面对称，因此二者电势逐项抵消。在 \(x=-l\) 上对应负列索引
\(m=-n-1\)，像点为
\(-a-(4n+2)l\)，也可逐项配对，所以两块板均为零电势。由唯一性定理，这组像电荷在板间给出的电势就是所求解。

\subsubsection{2.
小位移电像力}\label{m2:ux5c0fux4f4dux79fbux7535ux50cfux529b}

同号像电荷成对位于实电荷左右等距处，对实电荷的轴向力彼此抵消。异号像电荷位于

\[
x=\pm2(2j+1)l-a,
\qquad
j=0,1,2,\ldots
\]

附近，分别吸引实电荷向右、向左。故

\[
F_e(a)
=
\frac{kq^2}{4}
\sum_{j=0}^{\infty}
\left[
\frac1{((2j+1)l-a)^2}
-
\frac1{((2j+1)l+a)^2}
\right],
\]

其中

\[
k=\frac1{4\pi\varepsilon_0}.
\]

对 \(|a|\ll l\)，展开

\[
\frac1{(b-a)^2}
-
\frac1{(b+a)^2}
=
\frac{4a}{b^3}
+
O\!\left(\frac{a^3}{b^5}\right).
\]

于是

\[
F_e(a)
=
\frac{kq^2a}{l^3}
\sum_{j=0}^{\infty}
\frac1{(2j+1)^3}
+
O(a^3).
\]

利用

\[
\sum_{j=0}^{\infty}
\frac1{(2j+1)^3}
=
\left(1-\frac1{2^3}\right)\zeta(3)
=
\frac78\zeta(3),
\]

得到

\[
\boxed{
F_e(a)
=
\frac{7\zeta(3)kq^2}{8l^3}\,a
+
O(a^3)
}.
\]

当 \(a>0\)
时，\(F_e>0\)：离右板稍近后，右板吸引更强，电像力继续放大偏移。因此电像力是\textbf{反回复力}。

\subsubsection{3.
与重力回复力比较}\label{m2:ux4e0eux91cdux529bux56deux590dux529bux6bd4ux8f83}

小球横向位移与摆角满足

\[
a=h\sin\theta\simeq h\theta.
\]

重力势能

\[
U_g
=
mgh(1-\cos\theta)
\simeq
\frac{mg}{2h}a^2,
\]

故横向回复力

\[
F_g
=
-\frac{\mathrm dU_g}{\mathrm da}
\simeq
-\frac{mg}{h}a.
\]

总线性力为

\[
F_{\text{总}}
\simeq
-
\left[
\frac{mg}{h}
-
\frac{7\zeta(3)kq^2}{8l^3}
\right]a.
\]

\par\noindent\begin{minipage}{\linewidth}
\begin{center}
\includegraphics[width=\linewidth,height=.40\textheight,keepaspectratio]{\detokenize{assets/module-02/process/M2-177b.pdf}}
\par\smallskip{\small 图 M2-177b：重力与电像力贡献相反的线性刚度}
\end{center}
\end{minipage}\par

图 M2-177b
的红线斜率为正，放大小位移；绿线斜率为负，提供恢复力。中面附近是否稳定取决于两斜率之和的符号，临界相等时必须继续检查更高阶项。

因此中面平衡稳定的条件为

\[
\boxed{
\frac{mg}{h}
>
\frac{7\zeta(3)kq^2}{8l^3}
}.
\]

临界电荷量为

\[
\boxed{
q_c
=
\sqrt{
\frac{8mgl^3}{7\zeta(3)kh}
}
}.
\]

当 \(|q|<q_c\) 时，中面为稳定平衡；当 \(|q|>q_c\)
时，中面失稳。在线性近似中临界点的二次项消失，实际运动由更高阶电像力、板间碰撞和装置尺寸决定。

\subsection{综合题4：偏心圆柱电容器的电容与受力}\label{m2:ux7efcux5408ux98984ux504fux5fc3ux5706ux67f1ux7535ux5bb9ux5668ux7684ux7535ux5bb9ux4e0eux53d7ux529b}

\textbf{题号：M2-178}

\textbf{层级：} 决赛　　\textbf{方法：} 电像法、能量法

真空中的偏心圆柱电容器沿轴向无限延伸。外导体内半径为
\(R\)，内导体外半径为 \(r\)，两轴平行，轴间距离为 \(\Delta\)，满足
\(0<r<R\)、\(0\le\Delta<R-r\)。内外导体保持电势差
\(V\)，忽略端部效应。求单位长度电容
\(C'(\Delta)\)；在电压固定的条件下求内柱所受单位长度横向力
\(F'\)，以增大偏心距为正向，并检验同轴极限。

\par\noindent\begin{minipage}{\linewidth}
\subsubsection{解答}\label{m2:ux89e3ux7b54-3}

\begin{center}
\includegraphics[width=\linewidth,height=.40\textheight,keepaspectratio]{\detokenize{assets/module-02/examples/M2-178.pdf}}
\par\smallskip{\small 图 M2-178：偏心圆柱电容器的电容与受力}
\end{center}
\end{minipage}\par

图 M2-178
的两轴偏距取增大为正，内圆柱靠右使右侧间隙更窄。双极坐标把两圆变成固定坐标边界，尺度因子在电荷积分中抵消；固定电压下使用含电源交换的有效势能，得到继续增大偏距的力。

\subsubsection{1.
双极坐标贴合两圆}\label{m2:ux53ccux6781ux5750ux6807ux8d34ux5408ux4e24ux5706}

偏心圆柱的横截面是两个不相交的偏心圆。可取双极坐标

\[
x
=
\frac{c\sinh\tau}{\cosh\tau-\cos\sigma},
\qquad
y
=
\frac{c\sin\sigma}{\cosh\tau-\cos\sigma}.
\]

\(\tau=\tau_0\) 是圆：

\[
\left(x-c\coth\tau_0\right)^2+y^2
=
c^2\operatorname{csch}^2\tau_0.
\]

设内、外圆对应 \(\tau=\tau_i,\tau_o\)，定义

\[
\chi=|\tau_i-\tau_o|.
\]

消去辅助参数 \(c,\tau_i,\tau_o\)，得到纯几何关系

\[
\boxed{
\cosh\chi
=
\frac{R^2+r^2-\Delta^2}{2Rr}
}.
\]

\par\noindent\begin{minipage}{\linewidth}
\begin{center}
\includegraphics[width=\linewidth,height=.40\textheight,keepaspectratio]{\detokenize{assets/module-02/process/M2-178a.pdf}}
\par\smallskip{\small 图 M2-178a：两圆的几何参数确定双极坐标间距}
\end{center}
\end{minipage}\par

图 M2-178a 的两个圆心相距 \(\Delta\)，双极坐标的两条常坐标边界相隔
\(\chi\)。用半径和偏心距消去辅助参数后，电容只依赖这一坐标间距；同心极限给出独立校验。

当 \(\Delta=0\) 时，

\[
\cosh\chi
=
\frac{R^2+r^2}{2Rr}
=
\cosh\!\left(\ln\frac Rr\right),
\]

所以

\[
\chi=\ln\frac Rr,
\]

可回到同轴圆柱结果。

\subsubsection{2. 电势与电容}\label{m2:ux7535ux52bfux4e0eux7535ux5bb9}

双极坐标的尺度因子为

\[
h_{\tau}=h_{\sigma}
=
\frac{c}{\cosh\tau-\cos\sigma}.
\]

边界电势只随 \(\tau\) 变化，因此拉普拉斯方程允许线性解

\[
\varphi(\tau)
=
V\frac{\tau-\tau_o}{\tau_i-\tau_o}.
\]

法向电场大小为

\[
E_n
=
\frac1{h_{\tau}}
\left|\frac{\partial\varphi}{\partial\tau}\right|
=
\frac{V}{\chi h_{\tau}}.
\]

圆周弧元

\[
\mathrm ds=h_{\sigma}\,\mathrm d\sigma.
\]

内导体单位长度电荷为

\[
\begin{aligned}
\lambda
&=
\varepsilon_0\int_0^{2\pi}E_n\,\mathrm ds\\
&=
\frac{\varepsilon_0V}{\chi}
\int_0^{2\pi}
\frac{h_{\sigma}}{h_{\tau}}\,\mathrm d\sigma\\
&=
\frac{2\pi\varepsilon_0V}{\chi}.
\end{aligned}
\]

因此

\[
\boxed{
C'
=
\frac{\lambda}{V}
=
\frac{2\pi\varepsilon_0}{\chi}
},
\]

其中

\[
\boxed{
\chi
=
\operatorname{arcosh}
\left(
\frac{R^2+r^2-\Delta^2}{2Rr}
\right)
}.
\]

\subsubsection{3.
恒电压条件下的受力}\label{m2:ux6052ux7535ux538bux6761ux4ef6ux4e0bux7684ux53d7ux529b}

当电压由理想电源保持不变时，机械系统的有效势能不是单独的场能

\[
+\frac12C'V^2,
\]

而是计入电源做功后的

\[
\mathcal U_{\text{eff}}
=
-\frac12C'V^2.
\]

所以沿 \(\Delta\) 增大的方向

\[
F'
=
-\frac{\mathrm d\mathcal U_{\text{eff}}}{\mathrm d\Delta}
=
\frac12V^2\frac{\mathrm dC'}{\mathrm d\Delta}.
\]

由

\[
\sinh\chi\frac{\mathrm d\chi}{\mathrm d\Delta}
=
-\frac{\Delta}{Rr}
\]

得

\[
\frac{\mathrm dC'}{\mathrm d\Delta}
=
\frac{2\pi\varepsilon_0\Delta}
{Rr\,\sinh\chi\,\chi^2}.
\]

因此

\[
\boxed{
F'
=
\frac{\pi\varepsilon_0V^2\Delta}
{Rr\,\sinh\chi\,\chi^2}
}.
\]

\par\noindent\begin{minipage}{\linewidth}
\begin{center}
\includegraphics[width=\linewidth,height=.40\textheight,keepaspectratio]{\detokenize{assets/module-02/process/M2-178b.pdf}}
\par\smallskip{\small 图 M2-178b：固定电压时内圆柱趋向较近的外壁}
\end{center}
\end{minipage}\par

图 M2-178b 的红箭头指向较窄间隙。随着偏心距增加，\(\chi\)
减小、电容增加；计入电源后的有效势能随之下降，因此由负能量梯度得到同一受力方向。

因为 \(F'>0\)，内圆柱被拉向更靠近外圆柱的一侧，偏心会继续增大；同轴位置
\(\Delta=0\) 是横向不稳定平衡。

还可定义

\[
l_0
=
\frac{Rr\sinh\chi}{\Delta}
=
\frac{
\sqrt{
\left[(R-r)^2-\Delta^2\right]
\left[(R+r)^2-\Delta^2\right]
}
}{2\Delta},
\]

从而把力写成

\[
\boxed{
F'
=
\frac{\pi\varepsilon_0V^2}{l_0\chi^2}
}.
\]

\subsubsection{4. 两个极限}\label{m2:ux4e24ux4e2aux6781ux9650}

同轴极限 \(\Delta\to0\)：

\[
C'
\to
\frac{2\pi\varepsilon_0}{\ln(R/r)},
\qquad
F'\to0.
\]

接近相切极限 \(\Delta\to R-r\)：

\[
\chi\to0,
\qquad
C'\to\infty,
\]

理想二维模型中的力也发散。真实装置会受到表面粗糙度、击穿、有限长度和机械接触的限制。

\subsection{综合题5：旋转椭球与薄圆盘的孤立电容}\label{m2:ux7efcux5408ux98985ux65cbux8f6cux692dux7403ux4e0eux8584ux5706ux76d8ux7684ux5b64ux7acbux7535ux5bb9}

\textbf{题号：M2-179}

\textbf{层级：} 决赛　　\textbf{方法：} 电像法

在真空中分别考察三个孤立导体，均无外加场，电容定义为总电荷与相对无穷远电势之比：长旋转椭球的极半轴为
\(a\)、赤道半轴为 \(b\)，\(a>b>0\)；扁旋转椭球的极半轴为
\(a\)、赤道半轴为 \(b\)，\(0<a<b\)；理想薄圆盘半径为
\(b\)。求三者自电容，明确长、扁两种几何中焦半距的定义，并通过球形极限和扁椭球趋于薄圆盘的极限互相核对。

\par\noindent\begin{minipage}{\linewidth}
\subsubsection{解答}\label{m2:ux89e3ux7b54-4}

\begin{center}
\includegraphics[width=\linewidth,height=.40\textheight,keepaspectratio]{\detokenize{assets/module-02/examples/M2-179.pdf}}
\par\smallskip{\small 图 M2-179：旋转椭球与薄圆盘的孤立电容}
\end{center}
\end{minipage}\par

图 M2-179
的极半轴与赤道半轴在长、扁两类椭球中角色不同。分别选贴合导体的等势坐标面，再从远区
\(1/r\)
系数读取总电荷；圆盘极限取扁椭球的极半轴趋零，不能把长椭球公式直接代到该极限。

\subsubsection{1. 长旋转椭球}\label{m2:ux957fux65cbux8f6cux692dux7403}

定义焦半距

\[
c=\sqrt{a^2-b^2}.
\]

取长旋转椭球坐标

\[
\rho
=
c\sqrt{(\xi^2-1)(1-\eta^2)},
\qquad
z=c\xi\eta,
\]

其中

\[
\xi\ge1,
\qquad
-1\le\eta\le1.
\]

\(\xi=\xi_0\) 是长旋转椭球，其半轴为

\[
a=c\xi_0,
\qquad
b=c\sqrt{\xi_0^2-1},
\]

所以

\[
\xi_0=\frac ac.
\]

若导体电势为 \(V\)，对称性使外部电势只依赖 \(\xi\)。拉普拉斯方程约化为

\[
\frac{\mathrm d}{\mathrm d\xi}
\left[
(\xi^2-1)\frac{\mathrm d\varphi}{\mathrm d\xi}
\right]
=0.
\]

积分并使用

\[
\varphi(\xi_0)=V,
\qquad
\varphi(\infty)=0,
\]

得到

\[
\varphi(\xi)
=
V
\frac{
\ln\!\left(\dfrac{\xi+1}{\xi-1}\right)
}{
\ln\!\left(\dfrac{\xi_0+1}{\xi_0-1}\right)
}.
\]

当 \(\xi\to\infty\) 时，

\[
\ln\frac{\xi+1}{\xi-1}
=
\frac2\xi+O(\xi^{-3}),
\qquad
r\simeq c\xi.
\]

因此

\[
\varphi
\simeq
\frac{
2cV
}{
\ln\!\left(\dfrac{\xi_0+1}{\xi_0-1}\right)
}
\frac1r.
\]

与

\[
\varphi\simeq\frac{Q}{4\pi\varepsilon_0r}
\]

比较，得

\[
\boxed{
C_{\text{prolate}}
=
\frac{8\pi\varepsilon_0c}
{\ln\!\left(\dfrac{a+c}{a-c}\right)}
}.
\]

\subsubsection{2. 扁旋转椭球}\label{m2:ux6241ux65cbux8f6cux692dux7403}

定义

\[
c=\sqrt{b^2-a^2}.
\]

取扁旋转椭球坐标

\[
\rho
=
c\sqrt{(1+\xi^2)(1-\eta^2)},
\qquad
z=c\xi\eta.
\]

\(\xi=\xi_0\) 的极半轴和赤道半轴分别为

\[
a=c\xi_0,
\qquad
b=c\sqrt{1+\xi_0^2},
\]

故

\[
\xi_0=\frac ac.
\]

只依赖 \(\xi\) 的拉普拉斯方程为

\[
\frac{\mathrm d}{\mathrm d\xi}
\left[
(1+\xi^2)\frac{\mathrm d\varphi}{\mathrm d\xi}
\right]
=0.
\]

满足边界条件的解为

\[
\varphi(\xi)
=
V
\frac{\arctan(1/\xi)}
{\arctan(1/\xi_0)}.
\]

当 \(\xi\to\infty\) 时，

\[
\arctan(1/\xi)
\simeq
\frac1\xi,
\qquad
r\simeq c\xi.
\]

所以

\[
\boxed{
C_{\text{oblate}}
=
\frac{4\pi\varepsilon_0c}
{\arctan(c/a)}
}.
\]

\par\noindent\begin{minipage}{\linewidth}
\begin{center}
\includegraphics[width=\linewidth,height=.40\textheight,keepaspectratio]{\detokenize{assets/module-02/process/M2-179a.pdf}}
\par\smallskip{\small 图 M2-179a：扁椭球极限与球形极限提供交叉检验}
\end{center}
\end{minipage}\par

图 M2-179a 中保持赤道半轴 \(b\)，令极半轴 \(a\) 收缩，得到薄圆盘。另向
\(a=b\) 取极限时须同时令焦半距 \(c\)
趋零，用比值极限回到球电容，不能直接把零代入分子分母。

\subsubsection{3. 薄圆盘极限}\label{m2:ux8584ux5706ux76d8ux6781ux9650}

令扁旋转椭球的极半轴

\[
a\to0,
\]

则

\[
c\to b,
\qquad
\arctan(c/a)\to\frac\pi2.
\]

因此半径为 \(b\) 的薄圆盘电容为

\[
\boxed{
C_{\text{disk}}
=
8\varepsilon_0b
}.
\]

圆盘每一面的面电荷密度为

\[
\boxed{
\sigma(\rho)
=
\frac{2\varepsilon_0V}
{\pi\sqrt{b^2-\rho^2}},
\qquad
0\le\rho<b
}.
\]

两面总电荷为

\[
\begin{aligned}
Q
&=
2\int_0^b
\sigma(\rho)\,2\pi\rho\,\mathrm d\rho\\
&=
8\varepsilon_0bV,
\end{aligned}
\]

与 \(Q=CV\) 一致。由于

\[
\sigma(\rho)
\propto
\frac1{\sqrt{b-\rho}},
\]

电荷在圆盘边缘积聚并出现可积的平方根发散。

\subsubsection{4. 球极限检查}\label{m2:ux7403ux6781ux9650ux68c0ux67e5}

无论从长旋转椭球还是扁旋转椭球令形状趋于半径为 \(R\) 的球，都有
\(c\to0\)。

长旋转椭球中

\[
\ln\frac{a+c}{a-c}
=
\frac{2c}{a}+O(c^3),
\]

故

\[
C_{\text{prolate}}
\to
4\pi\varepsilon_0a.
\]

扁旋转椭球中

\[
\arctan\frac ca
=
\frac ca+O(c^3),
\]

故

\[
C_{\text{oblate}}
\to
4\pi\varepsilon_0a.
\]

两条路径都回到孤立球电容

\[
\boxed{
C_{\text{sphere}}=4\pi\varepsilon_0R
},
\]

说明公式的几何参数和归一化一致。

\section{模块二小结}\label{m2:ux6a21ux5757ux4e8cux5c0fux7ed3}

{\def\LTcaptype{none} % do not increment counter
\begin{longtable}[]{@{}
  >{\raggedright\arraybackslash}p{(\linewidth - 4\tabcolsep) * \real{0.3333}}
  >{\raggedright\arraybackslash}p{(\linewidth - 4\tabcolsep) * \real{0.3333}}
  >{\raggedright\arraybackslash}p{(\linewidth - 4\tabcolsep) * \real{0.3333}}@{}}
\toprule\noalign{}
\begin{minipage}[b]{\linewidth}\raggedright
检验问题
\end{minipage} & \begin{minipage}[b]{\linewidth}\raggedright
结论或判断
\end{minipage} & \begin{minipage}[b]{\linewidth}\raggedright
回看位置
\end{minipage} \\
\midrule\noalign{}
\endhead
\bottomrule\noalign{}
\endlastfoot
面对高对称电荷分布先求场还是先求势？ & 优先使用高斯定理求场 & 专题八 \\
远场展开前最先检查什么？ & 无量纲小参数与净电荷 & 专题十、十一 \\
导体表面切向场为何为零？ & 否则自由电荷不能静止 & 专题十三 \\
固定电荷和固定电势能否使用同一能量导数？ & 必须分别计入外电源功 &
专题十四 \\
唯一性定理替代了什么工作？ & 确认满足全部条件的候选解就是物理解 &
专题十五 \\
分离变量与电像法如何选择？ & 依据几何可分性和边界反射闭合性 &
专题十七至二十 \\
\end{longtable}
}

\section{模块二综合题：接地平面与球形边界}\label{m2:ux6a21ux5757ux4e8cux7efcux5408ux9898ux63a5ux5730ux5e73ux9762ux4e0eux7403ux5f62ux8fb9ux754c}

\subsection{第一题：点电荷、接地平面、能量与远场}\label{m2:ux7b2cux4e00ux9898ux70b9ux7535ux8377ux63a5ux5730ux5e73ux9762ux80fdux91cfux4e0eux8fdcux573a}

\textbf{题号：M2-180}

\textbf{层级：} 复赛　　\textbf{方法：} 多极展开、电像法、能量法

真空半空间 \(z>0\) 中一点电荷 \(q>0\) 固定于
\((0,0,a)\)，\(a>0\)，无限大导体平面 \(z=0\)
接地，要求无穷远势衰减。求真实半空间内电势并检验边界，求平面感应面电荷密度及总电荷、真实点电荷受力和诱导相互作用能（以
\(a\to\infty\) 为零，不含点源自能）。最后在
\(r\gg a\)、观察点仍位于上半空间时求电势最低非零远场项，说明像源延拓与普通局域真实电荷分布的远场解释有何联系和区别。

\par\noindent\begin{minipage}{\linewidth}
\subsubsection{解答}\label{m2:ux89e3ux7b54-5}

\begin{center}
\includegraphics[width=\linewidth,height=.40\textheight,keepaspectratio]{\detokenize{assets/module-02/examples/M2-180.pdf}}
\par\smallskip{\small 图 M2-180：点电荷、接地平面、能量与远场}
\end{center}
\end{minipage}\par

图 M2-180
的真实点源和域外反射像满足上半空间的边值条件。边界法向导数决定感应面电荷，像源场决定力，感应势的一半决定能量；远区虽有偶极形式，完整辅助场只在真实半空间代表原问题。

\[
\boxed{
\varphi(\rho,z)=\frac{q}{4\pi\varepsilon_0}
\left[
\frac{1}{\sqrt{\rho^2+(z-a)^2}}
-\frac{1}{\sqrt{\rho^2+(z+a)^2}}
\right]
}.
\]

它在真实区域内只保留原点电荷的正确奇点，像电荷位于真实区域外；在整张平面上有
\(\varphi(\rho,0)=0\)；无穷远处
\(\varphi\to0\)。因此由第一类边值问题的唯一性，该式就是 \(z>0\)
的物理解。

取导体指向真空的法向为 \(+\mathbf e_z\)，则

\[
\boxed{
\sigma(\rho)=\varepsilon_0E_z(\rho,0^+)
=-\frac{qa}{2\pi(\rho^2+a^2)^{3/2}}
}.
\]

\par\noindent\begin{minipage}{\linewidth}
\begin{center}
\includegraphics[width=\linewidth,height=.40\textheight,keepaspectratio]{\detokenize{assets/module-02/process/M2-180a.pdf}}
\par\smallskip{\small 图 M2-180a：面电荷剖面与同心环积分}
\end{center}
\end{minipage}\par

图 M2-180a
的横坐标是离轴距离，纵坐标是从真空侧法向场读出的面电荷密度。每一段径向宽度应乘同心环周长
\(2\pi\rho\) 后再积分，不能直接把剖面曲线下面积当作总电荷。

积分得到

\[
Q_{\mathrm{ind}}
=\int_0^{\infty}\sigma(\rho)\,2\pi\rho\,d\rho
=-q.
\]

点电荷只受感应电荷产生的场作用。计算力时可用像电荷的场，两者相距
\(2a\)，故

\[
\boxed{
\mathbf F
=-\frac{q^2}{16\pi\varepsilon_0a^2}\,\mathbf e_z
}.
\]

能量不能直接取真实电荷与像电荷相互作用能，而应有因子 \(1/2\)：

\[
\boxed{
W=\frac12q\,\varphi_{\mathrm{ind}}(0,0,a)
=-\frac{q^2}{16\pi\varepsilon_0a}
}.
\]

它满足 \(F_z=-dW/da\)，与上式一致。

令场点球坐标为 \((r,\theta)\)，\(\theta\) 从 \(+z\) 轴量起。对
\(r\gg a\)，

\[
\frac1{|\mathbf r-a\mathbf e_z|}
-\frac1{|\mathbf r+a\mathbf e_z|}
=\frac{2a\cos\theta}{r^2}
+O\!\left(\frac{a^3}{r^4}\right),
\]

所以

\[
\boxed{
\varphi(r,\theta)
=\frac1{4\pi\varepsilon_0}
\frac{(2qa)\cos\theta}{r^2}
+O\!\left(\frac{qa^3}{\varepsilon_0r^4}\right)
}.
\]

上半空间远场在数学上等效于偶极矩
\(\mathbf p_{\mathrm{eff}}=2qa\,\mathbf e_z\) 的像电荷对。这里的
\(\mathbf p_{\mathrm{eff}}\)
是\textbf{带边界问题在上半空间的有效远场参数}，不能不加说明地当成普通全空间局域物理源的偶极矩；感应电荷铺展在无限平面上，通常局域源多极展开所需的紧支撑或绝对矩收敛条件并不成立。

\subsection{第二题：球形边界上的三个轴对称模态}\label{m2:ux7b2cux4e8cux9898ux7403ux5f62ux8fb9ux754cux4e0aux7684ux4e09ux4e2aux8f74ux5bf9ux79f0ux6a21ux6001}

\textbf{题号：M2-181}

\textbf{层级：} 复赛　　\textbf{方法：} 分离变量、边值问题

半径为 \(R>0\)、球心在原点的真空球形区域内部无电荷或其他奇点。采用以
\(+z\) 为极轴的球坐标，球面电势给定为
\(V_0+V_1\cos\theta+V_2P_2(\cos\theta)\)，\(P_2(u)=(3u^2-1)/2\)，各
\(V_i\)
为已知常量。求球心正则的球内电势，求原点电势和电场，最后说明唯一性定理如何保证所得解为全部满足条件的解。

\par\noindent\begin{minipage}{\linewidth}
\subsubsection{解答}\label{m2:ux89e3ux7b54-6}

\begin{center}
\includegraphics[width=\linewidth,height=.40\textheight,keepaspectratio]{\detokenize{assets/module-02/examples/M2-181.pdf}}
\par\smallskip{\small 图 M2-181：球形边界上的三个轴对称模态}
\end{center}
\end{minipage}\par

图 M2-181 的球面数据恰含零、一、二阶三个轴对称模态。球心正则性给它们配
\(1,r,r^2\)，零阶决定中心电势、一次项决定中心匀强场；二次项在中心没有一阶梯度，唯一性也排除额外的零边界调和项。

\[
\varphi(r,\theta)
=\sum_{L=0}^{\infty}A_Lr^LP_L(\cos\theta).
\]

边界数据只含 \(L=0,1,2\) 三个互相正交的模态，逐项匹配得到

\[
\boxed{
\varphi(r,\theta)
=V_0
+V_1\frac rR\cos\theta
+V_2\left(\frac rR\right)^2P_2(\cos\theta)
}.
\]

% 末题两步推导分别完整分页，避免过程图挤走中心场结论。
\newpage
\par\noindent\begin{minipage}{\linewidth}
\begin{center}
\includegraphics[width=\linewidth,height=.40\textheight,keepaspectratio]{\detokenize{assets/module-02/process/M2-181a.pdf}}
\par\smallskip{\small 图 M2-181a：径向正则幂决定球心还留下哪些量}
\end{center}
\end{minipage}\par

图 M2-181a
把三个正则径向因子分开。为看清原点附近的一阶梯度，把刚才的解改写为笛卡尔坐标：

\[
\varphi(x,y,z)=V_0+\frac{V_1z}{R}
+\frac{V_2}{2R^2}(2z^2-x^2-y^2),
\]

\[
\mathbf E(x,y,z)
=\frac{V_2x}{R^2}\mathbf e_x+\frac{V_2y}{R^2}\mathbf e_y
-\left(\frac{V_1}{R}+\frac{2V_2z}{R^2}\right)\mathbf e_z.
\]

二次项的拉普拉斯为
\((V_2/R^2)(2-1-1)=0\)；其电场各分量在原点均为零。常量项只决定中心电势，一次项才决定中心场，这也独立核对了图中的径向结论。

在球心，所有 \(L\ge1\) 的电势项都趋于零，因此

\[
\boxed{\varphi(0)=V_0}.
\]

\(L=1\) 项可写成 \(V_1z/R\)，而 \(L=2\)
项在原点附近是二次函数，所以只有偶极模态对球心的一阶梯度有贡献：

\[
\boxed{
\mathbf E(0)=-\nabla\varphi(0)
=-\frac{V_1}{R}\,\mathbf e_z
}.
\]

该候选式在球内正则，满足拉普拉斯方程，并在整个球面上逐点满足给定电势。第一类边值问题的唯一性保证不能再叠加任何边界值为零而内部非零的调和函数。作为独立检验，球面平均值中
\(P_1\)、\(P_2\) 的平均均为零，故球心电势等于边界平均值
\(V_0\)，与平均值性质一致。
\clearpage
\endgroup
